\documentclass[11pt]{article}
\usepackage[a4paper,margin=1in]{geometry}
\usepackage{amsmath,amssymb,amsthm,mathtools}
\usepackage[T1]{fontenc}
\usepackage{mathptmx}
\renewcommand{\ttdefault}{pcr}
\usepackage{enumitem}
\usepackage[final]{microtype}
\usepackage[hidelinks]{hyperref}
\hypersetup{
  pdfauthor={Tseng},
  pdftitle={Rough Polynomial Fluctuation Theory: Coefficient-Uniform Critical Jumps, Finite Radial Variation, and Endpoint Bounds},
  pdfsubject={Critical jump and finite-radial-variation inequalities for rough polynomial singular integrals},
  pdfkeywords={rough singular integrals, polynomial phases, critical jump inequalities, radial variation, weak type (1,1)}
}
\setlength{\emergencystretch}{3em}
\lefthyphenmin=3
\righthyphenmin=3

\newtheorem{theorem}{Theorem}[section]
\newtheorem{proposition}[theorem]{Proposition}
\newtheorem{lemma}[theorem]{Lemma}
\newtheorem{corollary}[theorem]{Corollary}
\theoremstyle{remark}
\newtheorem{remark}[theorem]{Remark}
\theoremstyle{definition}
\newtheorem{definition}[theorem]{Definition}

\newcommand{\R}{\mathbb R}
\newcommand{\Sph}{\mathbb S}
\newcommand{\supp}{\operatorname{supp}}
\newcommand{\Var}{\operatorname{Var}}
\newcommand{\osc}{\operatorname{osc}}
\newcommand{\Mix}{\operatorname{Mix}}
\newcommand{\Dop}{\mathfrak D}
\newcommand{\cA}{\mathcal A}
\newcommand{\cF}{\mathcal F}
\newcommand{\cL}{\mathcal L}
\newcommand{\cQ}{\mathcal Q}
\newcommand{\cU}{\mathcal U}
\newcommand{\norm}[1]{\lVert #1\rVert}
\newcommand{\abs}[1]{\lvert #1\rvert}
\newcommand{\weak}{L^{1,\infty}}
\newcommand{\BV}{\operatorname{BV}}
\newcommand{\Jop}[1]{\mathcal J_{2,#1}}

\title{Rough Polynomial Fluctuation Theory:\\
Coefficient-Uniform Critical Jumps, Finite Radial Variation,\\
and Endpoint Bounds}
\author{Tseng}
\date{Version 2 --- August 31, 2026}

\begin{document}
\maketitle

\begin{abstract}
We prove coefficient-uniform bounds for the pointwise critical $2$-jump
functional of the full sharp-truncation trajectory
\[
T^{P,b}_{\Omega,\varepsilon}f(x)
 =\int_{|x-y|>\varepsilon}e^{iP(x,y)}b(|x-y|)
 \frac{\Omega((x-y)/|x-y|)}{|x-y|^n}f(y)\,dy,
\]
where $P$ is an arbitrary real polynomial on
$\mathbb R^n\times\mathbb R^n$ of bounded degree,
$\Omega\in L\log L(\mathbb S^{n-1})$ has mean zero, and the radial factor
$b$ has finite $\rho$-variation for some fixed $1\le\rho<\infty$.  For
every $1<p<\infty$ the pointwise functional
\[
 \sup_{\lambda>0}\lambda
 N_\lambda(T^{P,b}_{\Omega,\varepsilon}f:\varepsilon>0)^{1/2}
\]
is bounded on $L^p(\mathbb R^n)$, and it maps $L^1$ to
$L^{1,\infty}$.  The constants depend only on the dimension, the degree
bound, the relevant spatial exponent, and $\rho$; they are independent of
every coefficient of $P$ and of where the radial variation of $b$ is
concentrated.

At $P=0$ and $b\equiv1$, the weak endpoint strictly strengthens, and hence
resolves, the critical-jump problem posed by Jones--Seeger--Wright.  For
general $P$, the principal-value weak type $(1,1)$ consequence answers the
rough-kernel problem posed by Ding and Lai in the mean-zero $L\log L$
class.  The jump theorem implies strong and weak $r$-variation inequalities
for every $r>2$, maximal-truncation bounds, almost-everywhere
principal-value convergence, and a canonical trajectory determined by its
finite-annulus increments and zero-at-infinity anchor.  For radial
$1\le\rho<2$, a direct Young--Stieltjes argument transfers the unweighted
trajectory pointwise.  At and above $\rho=2$, the proof instead uses
operator-specific Fourier, correlation, and packet orthogonality.  The
finiteness of $\rho$ is necessary: a bounded radial Hilbert example rules
out $\rho=\infty$.  An exact Mellin representation of the Poisson semigroup
also shows that true output $V^2$ fails already for sharp Hilbert
truncations.

The proof combines a finite global algebraic flag on polynomial phases
modulo pure input and output phases, a phase-adapted
Calder\'on--Zygmund decomposition with anchored phase--packet induction,
and a selector-free fixed-rank Rademacher--Menshov reduction supported by
concentration-stable equal- and unequal-shell correlation estimates.  The
finite-radial extension is obtained by proving the stability of every
oscillatory, terminal-packet, weak-endpoint, and completion step under a
$V^\rho$ radial amplitude.
\end{abstract}

\medskip\noindent\textit{AI-use disclosure.}
The research direction and target theorem of this paper were formulated
through iterative interaction between the author and an
artificial-intelligence system.  The author directed and coordinated the
mathematical development, including the selection and refinement of the
problems and intermediate objectives pursued during that process.  The AI
system generated the mathematical proof arguments and derivations
establishing the results of the paper; these contributions occur throughout
Sections~2--14 rather than being confined to isolated lemmas or theorems.
The author examined, selected, revised, and organized the resulting
arguments.
\medskip

\section{Introduction and main results}

A maximal-truncation estimate controls only the largest excursion of an
operator trajectory.  An $r$-variation estimate with $r>2$ also quantifies
convergence, but the endpoint $r=2$ is generally inaccessible.  The
critical replacement is the pointwise $2$-jump functional
\begin{equation}\label{eq:intro-pointwise-jump}
 \mathbf J(F)(x)
 :=\sup_{\lambda>0}\lambda N_\lambda(F(x))^{1/2}.
\end{equation}
The spatial norm is taken only after this pointwise supremum; in particular,
the endpoint quantity is
\[
 \norm{\mathbf J(F)}_{L^{1,\infty}}
 =\left\|\sup_{\lambda>0}
 \lambda N_\lambda(F)^{1/2}\right\|_{L^{1,\infty}}.
\]
This is stronger than the conventional mixed quasi-seminorm
$\sup_{\lambda>0}\|\lambda N_\lambda(F)^{1/2}\|_{L^{1,\infty}}$.
The strictness of this comparison, for genuine finite jump paths, is
quantified in Proposition~\ref{prop:mixed-norm-counterexample}.
It measures, simultaneously at every spatial point and over all amplitudes,
the number of disjoint oscillations
of the prescribed size, and it implies every $r$-variation estimate with
$r>2$ through an abstract endpoint principle.  The purpose of this paper is
to prove this critical estimate, including spatial weak type $(1,1)$, for
the complete sharp-truncation trajectory of rough polynomial singular
integrals.

The weak endpoint at $P=0$ strictly strengthens, and therefore resolves,
the problem posed by Jones--Seeger--Wright \cite{JSW}.  We prove it
independently and establish,
in the same framework, the coefficient-uniform extension to arbitrary
real two-variable polynomial phases of bounded degree.  The angular
kernel is assumed only to belong to the endpoint class $L\log L$; the
phase is subject to no nondegeneracy, smallness, or genericity hypothesis.
Thus the theorem simultaneously treats rough angular data, genuinely
nonconvolution phases, critical fluctuations, and the spatial endpoint.

Let $d\sigma$ be the usual surface measure on $\Sph^{n-1}$.  Put
\[
K_\Omega(z)=\frac{\Omega(z/|z|)}{|z|^n},\qquad z\ne0.
\]
For $A=\norm\Omega_1>0$ define the homogeneous gauge
\begin{equation}\label{eq:gauge}
 \mathfrak L(\Omega)
 =A+\int_{\Sph^{n-1}}|\Omega(\theta)|
 \log^+\!\left(\frac{|\Omega(\theta)|}{A}\right)d\sigma(\theta),
\end{equation}
and set $\mathfrak L(0)=0$.  This is equivalent, with dimensional constants,
to the homogeneous Luxemburg $L\log L$ norm.  Its homogeneity and the
boundedness of the odd and even projections
\[
 \Omega_o(\theta)=\frac{\Omega(\theta)-\Omega(-\theta)}2,
 \qquad
 \Omega_e(\theta)=\frac{\Omega(\theta)+\Omega(-\theta)}2
\]
give
\begin{equation}\label{eq:gauge-parity}
 \mathfrak L(\Omega_o)+\mathfrak L(\Omega_e)
 \le C\mathfrak L(\Omega).
\end{equation}
The elementary normalization details are recorded in
Lemma~\ref{lem:gauge-normalization}.

On bounded compactly supported inputs define
\[
 T^P_{\Omega,\varepsilon}f(x)
 =\int_{|x-y|>\varepsilon}e^{iP(x,y)}K_\Omega(x-y)f(y)\,dy.
\]
The final construction for general data is given in
Section~\ref{sec:limits}; it is characterized by the finite-annulus identity
and by convergence to zero as $\varepsilon\to\infty$.

Let $E$ be a linearly ordered set and let $F=(F_t)_{t\in E}$ be a scalar
path.  For $\lambda>0$, its paired strict jump number is
\[
N_{\lambda,E}(F):=
 \sup\bigl(\{0\}\cup\mathcal W_{\lambda,E}(F)\bigr)
 \in\mathbb N_0\cup\{\infty\},
\]
where $\mathcal W_{\lambda,E}(F)$ is the set of all $N\in\mathbb N$ for
which there are $s_i,t_i\in E$, $1\le i\le N$, satisfying
\[
 s_1<t_1\le s_2<t_2\le\cdots\le s_N<t_N,
 \qquad |F_{t_i}-F_{s_i}|>\lambda\quad(1\le i\le N).
\]
Thus a path with no strict $\lambda$-jump has jump number zero, while a path
with arbitrarily large finite witnesses has jump number infinity.
For $r>0$, its homogeneous $r$-variation is
\[
V_E^r(F):=\sup\Biggl(\{0\}\cup
 \Biggl\{\left(\sum_{i=1}^N|F_{t_i}-F_{t_{i-1}}|^r\right)^{1/r}:
 N\in\mathbb N,\quad t_0<\cdots<t_N,\quad t_i\in E\Biggr\}\Biggr).
\]
When the index set is displayed or unambiguous, we omit the subscript $E$;
thus the main positive-radius paths use $N_\lambda(F)$ and
$V^r(F_t:t>0)$.  For a restriction to an ordered subset $D\subset E$, put
\[
 \mathbf J_D(F):=\sup_{\lambda>0}\lambda
 N_{\lambda,D}(F|_D)^{1/2},
\]
and omit $D$ when it is the full parameter set.
For a measurable function $H$, set
$\norm H_{L^{1,\infty}}=\sup_{\tau>0}\tau|\{x:|H(x)|>\tau\}|$.
Call a path coordinatewise measurable if every $F_t$ is measurable in the
spatial variable.  For such a path define $\mathbf J(F)$ pointwise by
\eqref{eq:intro-pointwise-jump}, initially as an extended-valued pointwise
functional without a separate measurability assertion.  We call the path
jump-separable if there are a countable parameter set $D\subset E$ and one common
full-measure spatial set on which, for every $\lambda>0$, its full and
$D$-restricted jump numbers agree.
Whenever $\mathbf J(F)$ is measurable---in particular, for every
jump-separable coordinatewise measurable path---and $1<p<\infty$, set
\begin{equation}\label{eq:inside-jump-norms}
 \Jop{p}(F)=\norm{\mathbf J(F)}_p
 =\left\|\sup_{\lambda>0}
 \lambda N_\lambda(F)^{1/2}\right\|_p,
 \qquad
 \Jop{1,\infty}(F)=\norm{\mathbf J(F)}_{L^{1,\infty}}.
\end{equation}
Because all jumps are strict, the supremum in $\lambda$ may be restricted
to $\mathbb Q_+$ without changing $\mathbf J(F)$.  More precisely, if
$D\subset E$ is countable and the coordinates $F_t$ are measurable
for $t\in D$, then, for every $N\ge1$,
\[
 \{N_\lambda(F_t:t\in D)\ge N\}
 =\bigcup_{(s_i,t_i)_{i=1}^N\in\mathcal O_N(D)}
   \bigcap_{i=1}^N\{|F_{t_i}-F_{s_i}|>\lambda\},
\]
where $\mathcal O_N(D)$ is the countable set of ordered paired
$2N$-tuples in $D$.  Thus the jump count on $D$ is measurable, and every
jump-separable coordinatewise measurable path has a measurable pointwise
jump functional.  If the
path is locally continuous in the truncation parameter, its jump counts on
a dense $D$ agree with those on all positive parameters by
Lemma~\ref{lem:finite-witness}.  These observations and the rational
$\lambda$-supremum prove measurability of $\mathbf J(F)$.

We use the following common-representative convention throughout.  Whenever
one of the continuum-parameter operator trajectories constructed below is
asserted pointwise on the full parameter interval, its construction is
carried out on one spatial set $X_0$ of full measure, independent of the
parameter, on which the parameter path is locally continuous (or locally
absolutely continuous when that stronger property is used) and all stated
finite-annulus identities hold simultaneously.  All coordinates are set
equal to zero on $X_0^c$.  Thus a countable dense parameter set determines
all jumps and variations by Lemma~\ref{lem:finite-witness}, and the preceding
countable-union argument proves their measurability.  This convention does
not impose continuity on the abstract scalar paths used in the jump lemmas.
We do not pass from coordinatewise almost-everywhere identities at
uncountably many parameters to a common operator representative without an
explicit finite-witness or anchored uniform-Cauchy argument.  Joint
measurability in the parameter is stated separately whenever Fubini or a
parameter integral requires it.

Throughout, \(C\) and \(c\) may change from line to line.  Every constant
in a global theorem or corollary depends only on the parameters displayed
there; in particular, it is independent of the support and location of the
input and of every approximation used to construct the operator.  Constants
inside proofs may depend on fixed smooth cutoffs and normalized coordinate
charts chosen once and for all.  Only an explicitly localized statement may
depend on support or chart data named in that statement.  No constant depends
on the coefficients of \(P\), on truncation or dyadic scales, or on auxiliary
finite cutoffs.  Every coefficient norm is a fixed Euclidean norm on the
relevant finite-dimensional polynomial space; changing that norm changes
only the permitted structural constants.

\begin{theorem}[Coefficient-uniform pointwise critical-jump inequalities]
\label{thm:main}
Let $n\in\mathbb N$ and $d\in\mathbb N_0$.  Let $P$ be a real polynomial on
$\R^n\times\R^n$ of total degree at most $d$, and let
$\Omega\in L\log L(\Sph^{n-1})$ satisfy $\int\Omega\,d\sigma=0$.
Then the pointwise critical-jump functional of the sharp-truncation
trajectory satisfies the coefficient-uniform estimates
\begin{align}
 \Jop{p}(T^P_{\Omega,\varepsilon}f:\varepsilon>0)
 &\le C_{n,d,p}\mathfrak L(\Omega)\norm f_p,
 &&1<p<\infty,\label{eq:main-strong-jump}\\
 \Jop{1,\infty}(T^P_{\Omega,\varepsilon}f:\varepsilon>0)
 &\le C_{n,d}\mathfrak L(\Omega)\norm f_1.
 \label{eq:main-weak-jump}
\end{align}
Estimate \eqref{eq:main-strong-jump} holds for every $1<p<\infty$ and
$f\in L^p(\R^n)$; estimate \eqref{eq:main-weak-jump} holds for every
$f\in L^1(\R^n)$.  The constants are independent of every coefficient of
$P$.

For each fixed $\Omega$ as above and each fixed $f$ in the data class of
the asserted estimate, Proposition~\ref{prop:strong-polynomial-jump} in the
strong case and Section~\ref{sec:limits} in the weak case construct the
entire trajectory on one full-measure set, on which all of its finite-annulus
identities hold and
\begin{equation}\label{eq:anchor-main}
 T^P_{\Omega,\varepsilon}f(x)\longrightarrow0
 \qquad\text{as }\varepsilon\to\infty
\end{equation}
holds.  These two properties, rather than the estimates alone, determine
the trajectory uniquely: the finite-annulus identities make the difference
of two such representatives independent of $\varepsilon$ on that common
full-measure set, and the zero anchor makes this difference vanish.  When
$1<p<\infty$, the convergence in \eqref{eq:anchor-main} also holds in
$L^p(\mathbb R^n)$.
\end{theorem}

Put
\begin{equation}\label{eq:kappa-r}
 \kappa_r=\left(\sum_{k=1}^\infty k^{-r/2}\right)^{1/r}
 \le \left(\frac r{r-2}\right)^{1/r},\qquad r>2.
\end{equation}

\begin{corollary}[Variation and convergence consequences]
\label{cor:unweighted-consequences}
Under the hypotheses of Theorem~\ref{thm:main}, for every $r>2$,
\begin{align}
 \norm{V^r(T^P_{\Omega,\varepsilon}f:\varepsilon>0)}_p
 &\le C_{n,d,p}\kappa_r\mathfrak L(\Omega)\norm f_p,
 &&1<p<\infty,\label{eq:main-strong-var}\\
 \norm{V^r(T^P_{\Omega,\varepsilon}f:\varepsilon>0)}_{L^{1,\infty}}
 &\le C_{n,d}\kappa_r\mathfrak L(\Omega)\norm f_1.
 \label{eq:main-weak-var}
\end{align}
Here \eqref{eq:main-strong-var} holds for every $1<p<\infty$ and
$f\in L^p(\R^n)$, and \eqref{eq:main-weak-var} holds for every
$f\in L^1(\R^n)$.  The factor $\kappa_r$ is needed only for the
$V^r$ conclusions.  The zero anchor and
\eqref{eq:anchor-controls-maximal} give directly
\begin{align*}
 \left\|\sup_{\varepsilon>0}|T^P_{\Omega,\varepsilon}f|\right\|_p
 &\le C_{n,d,p}\mathfrak L(\Omega)\norm f_p,
 &&1<p<\infty,\\
 \left\|\sup_{\varepsilon>0}|T^P_{\Omega,\varepsilon}f|\right\|_{L^{1,\infty}}
 &\le C_{n,d}\mathfrak L(\Omega)\norm f_1,
\end{align*}
with no $\kappa_r$ loss.  The finiteness of $V^r$ implies that
$\lim_{\varepsilon\downarrow0}T^P_{\Omega,\varepsilon}f$ exists almost
everywhere in the respective ranges.  Indeed, if a scalar path $F_\varepsilon$
with finite $V^r$ were not Cauchy as $\varepsilon\downarrow0$, there would be
$\delta>0$ and, recursively, pairs
\[
 0<s_{k+1}<t_{k+1}<s_k<t_k,
 \qquad |F_{t_k}-F_{s_k}|>\delta .
\]
For every $N$, inserting the endpoints of the first $N$ pairs in increasing
order into one variation partition would give
$V^r(F)^r\ge N\delta^r$, a contradiction.  Completeness of $\mathbb C$
therefore gives the asserted endpoint limit.  This argument uses only finite
variation; the logically separate zero anchor at infinity is what gives the
maximal estimate above.
\end{corollary}

Consequently, the principal-value limit
$T^P_\Omega f:=\lim_{\varepsilon\downarrow0}T^P_{\Omega,\varepsilon}f$
defines an operator from $L^1$ to $L^{1,\infty}$ with bound
$C_{n,d}\mathfrak L(\Omega)$.  For general $P$, this answers, in the
mean-zero $L\log L$ class, the rough-kernel weak type $(1,1)$ problem posed
by Ding and Lai \cite[\S6.1]{DingLai}; the maximal-truncation and
critical-jump bounds are stronger.

For $1\le\rho<\infty$ and a pointwise bounded radial function
$b:(0,\infty)\to\mathbb C$, put
\begin{equation}\label{eq:radial-variation-norm}
 \|b\|_{\mathcal V^\rho}
 :=\|b\|_\infty+V^\rho(b;(0,\infty)).
\end{equation}
A function of finite $\rho$-variation is regulated.  In
\eqref{eq:radial-variation-norm}, $\|b\|_\infty$ denotes the pointwise
supremum of the chosen representative, consistently with the pointwise
definition of $V^\rho$.  Whenever point values enter a Stieltjes formula,
we use the representative
\[
 b^+(t):=\lim_{s\downarrow t}b(s),\qquad t>0.
\]
Indeed, simultaneous approximation of the points of any finite partition
from the right gives
$V^\rho(b^+;(0,\infty))\le V^\rho(b;(0,\infty))$, and the same
right-limit argument gives $\|b^+\|_\infty\le\|b\|_\infty$.  To quantify
the exceptional point values, fix $\delta>0$ and distinct
$t_1<\cdots<t_N$ such that $|b(t_i)-b^+(t_i)|>\delta$.  Choose, successively
and sufficiently close to the right of each $t_i$, points
$s_i\in(t_i,t_{i+1})$ (with the evident convention for $i=N$) such that
$|b(s_i)-b^+(t_i)|<\delta/2$.  The ordered partition containing every
$t_i,s_i$ then gives
\[
 N(\delta/2)^\rho\le V^\rho(b;(0,\infty))^\rho.
\]
Thus there are only finitely many such points for each $\delta$, and
$b=b^+$ off a countable set.  Polar
coordinates then show that $b(|z|)=b^+(|z|)$ for almost every
$z\in\mathbb R^n$, so replacing $b$ by $b^+$ changes none of the
Lebesgue integral operators below and does not increase
$\|b\|_{\mathcal V^\rho}$.
Initially for bounded compactly supported inputs define
\[
 T^{P,b}_{\Omega,\varepsilon}f(x)
 =\int_{|x-y|>\varepsilon}e^{iP(x,y)}b(|x-y|)
 K_\Omega(x-y)f(y)\,dy.
\]

\begin{theorem}[Coefficient-uniform bounds for every finite radial variation exponent]
\label{thm:finite-radial}
Let the hypotheses on $n,d,P$, and $\Omega$ be those of
Theorem~\ref{thm:main}.  Fix $1\le\rho<\infty$, and let
$b:(0,\infty)\to\mathbb C$ be pointwise bounded with
$\|b\|_{\mathcal V^\rho}<\infty$.  Then
\begin{align}
 \|\mathbf J(T^{P,b}_{\Omega,\varepsilon}f:\varepsilon>0)\|_p
 &\le C_{n,d,p,\rho}\|b\|_{\mathcal V^\rho}
       \mathfrak L(\Omega)\|f\|_p,
 &&1<p<\infty,\label{eq:finite-radial-strong}\\
 \|\mathbf J(T^{P,b}_{\Omega,\varepsilon}f:\varepsilon>0)\|_{L^{1,\infty}}
 &\le C_{n,d,\rho}\|b\|_{\mathcal V^\rho}
       \mathfrak L(\Omega)\|f\|_1.
 \label{eq:finite-radial-weak}
\end{align}
Estimate \eqref{eq:finite-radial-strong} holds for every
$f\in L^p(\mathbb R^n)$, and \eqref{eq:finite-radial-weak} holds for every
$f\in L^1(\mathbb R^n)$.  The constants are independent of every
coefficient of $P$ and of the location of the radial variation of $b$.
The trajectory is uniquely determined by its finite-annulus increments
and the zero-at-infinity anchor.  For every output exponent $r>2$, the
corresponding strong and weak $r$-variation conclusions hold with the
additional factor $\kappa_r$.  The maximal-truncation estimates follow
instead from the zero anchor and \eqref{eq:anchor-controls-maximal}, with
the same right-hand sides as
\eqref{eq:finite-radial-strong}--\eqref{eq:finite-radial-weak} and no
$\kappa_r$ loss.  Almost-everywhere convergence as
$\varepsilon\downarrow0$ follows from the finiteness of $V^r$.
\end{theorem}

The dependence on the fixed finite exponent $\rho$ is not asserted to be
uniform.  Proposition~\ref{prop:radial-infinity-counterexample} shows that
the conclusion is false at $\rho=\infty$.  The range $1\le\rho<2$ also
admits the direct pathwise proof of
Proposition~\ref{prop:young-radial-transfer}; the general proof, including
the critical and supercritical radial ranges, is given in
Section~\ref{sec:finite-radial}.

For a bounded function $b$ of finite total variation, write
\[
 \norm b_{\BV}=\norm b_\infty+\Var_{(0,\infty)}(b).
\]
The case $\rho=1$ of Theorem~\ref{thm:finite-radial} is complemented by the
following direct Stieltjes transfer, whose constant is obtained without the
operator-specific radial argument.

\begin{corollary}[Radial $\BV$ transfer]\label{cor:bv-extension}
Under the hypotheses of Theorem~\ref{thm:main}, let
$b:(0,\infty)\to\mathbb C$ be bounded and of finite total variation.
All conclusions of Theorem~\ref{thm:main} and
Corollary~\ref{cor:unweighted-consequences} remain valid with
$T^P_{\Omega,\varepsilon}$ replaced by $T^{P,b}_{\Omega,\varepsilon}$
and with each right-hand side multiplied by
\[
 \norm b_{\BV}=\norm b_\infty+\Var_{(0,\infty)}(b).
\]
The constants are independent of where the variation of $b$ is
concentrated.
\end{corollary}

The primary result is the weak pointwise critical-jump estimate
\eqref{eq:main-weak-jump}.  Since
$\sup_\lambda\|\lambda N_\lambda^{1/2}\|_{L^{1,\infty}}
\le\|\mathbf J\|_{L^{1,\infty}}$, its $P=0$ case independently resolves
and strictly strengthens the Jones--Seeger--Wright problem, while the full
theorem yields the coefficient-uniform nonconvolution extension for general
$P$.  The strong pointwise critical-jump estimate is proved within the same
architecture.  The variation, maximal, and convergence assertions follow
from the abstract jump calculus and the zero anchor.  Theorem~\ref{thm:finite-radial}
then extends both spatial bounds to every fixed finite radial variation
exponent; Corollary~\ref{cor:bv-extension} records the independent anchored
Stieltjes proof at bounded variation.

\subsection*{Position in the literature}

\paragraph{Static rough polynomial theory.}
For $P=0$, the strong and weak theory of rough homogeneous singular
integrals at the $L\log L$ angular scale was developed in
\cite{CZ56,DRdF86,CRdF88,Seeger96}.  For arbitrary $P(x,y)$,
Jiang--Lu \cite{JiangLu} studied the static principal-value operator with
$\Omega\in L\log L$ and a radial factor in $\BV$, including an $L^p$
boundedness criterion and weighted estimates.  Ojanen \cite{Ojanen}
proved weighted strong estimates for the same polynomial-phase,
angular-$L\log L$, radial-$\BV$ class under the geometric hypotheses
formulated there.  Al-Salman and Grafakos \cite{AlSalmanGrafakos} more
recently obtained static $L^p$ bounds for arbitrary $P(x,y)$, angular
$L\log L$ data, and a rough radial factor satisfying an
$L^\gamma$-type condition.  The corresponding rough-kernel weak type
$(1,1)$ endpoint for the principal-value operator was explicitly listed
as an open problem by Ding and Lai \cite[\S6.1]{DingLai}.
Theorem~\ref{thm:main} and
Corollary~\ref{cor:unweighted-consequences} answer it in the mean-zero
$L\log L$ class, with the stronger critical-jump control of the complete
sharp-truncation trajectory.

\paragraph{Smooth coefficient-uniform weak theory.}
Ricci and Stein \cite{RicciStein} established coefficient-uniform strong
$L^p$ estimates for smooth kernels.  Chanillo and Christ
\cite{ChanilloChrist} proved coefficient-uniform weak type $(1,1)$ for
arbitrary polynomial phases and smooth Calder\'on--Zygmund convolution
kernels.  Their argument removes pure phases, normalizes a leading
coefficient by dilation, and uses a phase-adapted double induction; their
final remarks also yield $L^p$ bounds, $1<p<\infty$, for the maximal
truncation.  Rough $L\log L$ kernels fall outside that theory.  The global
phase flag introduced below is designed to retain coefficient uniformity
while accommodating rough annular decay and phases that degenerate on the
diagonal.

\paragraph{Maximal rough polynomial theory.}
Representative rough polynomial estimates include
\cite{LuZhang,AlSalman}; polynomial maximal-truncation estimates for
smooth Calder\'on--Zygmund kernels include \cite{KrauseLacey}.
Choudhary, Shrivastava, and Shuin \cite{CSS} state coefficient-uniform
sparse and strong maximal estimates for arbitrary $P$ when
$\Omega\in L^\infty$.  For $P=0$, Lai \cite[Theorem~1.1]{Lai} proves the
weak $(1,1)$ maximal estimate at the $L\log L$ angular endpoint.  In the
present paper, maximal truncation is obtained as a consequence of the
stronger critical-jump theorem.

\paragraph{Fluctuation theory.}
For sharply truncated Hilbert transforms,
Campbell, Jones, Reinhold, and Wierdl
\cite[Theorems~1.2--1.4 and Conjecture~7.2]{CJRW} established strong and
weak $r$-variation estimates for $r>2$, the associated noncritical jump
bounds, and short $2$-variation estimates, and isolated the critical jump
exponent.  Their higher-dimensional work \cite{CJRW03} developed
singular-integral variation theory at the rough $L\log L$ scale.
Jones--Seeger--Wright \cite[Theorem~1.6]{JSW} proved the conventional
strong critical-jump inequality for the unphased rough trajectory under
$\Omega\in L^a$, $a>1$, and posed the corresponding weak-type $(1,1)$
estimate for rough homogeneous kernels \cite{JSW}.  Their formulation
places the supremum in the jump amplitude outside the spatial norm.  The
$P=0$ case of \eqref{eq:main-weak-jump} places that supremum pointwise
inside $L^{1,\infty}$, and hence strictly strengthens the problem they
posed.

Chen, Ding, Hong, and Liu \cite[Theorems~1.1 and~1.3]{ChenDingHongLiu}
proved weighted $L^p$ bounds for the pointwise critical-jump functional of
rough homogeneous singular-integral and averaging trajectories when the
angular function belongs to $L^q$, $q>1$.  Their one-dimensional
specializations provide the Hilbert-truncation and one-sided-average input
in Lemma~\ref{lem:one-dimensional-inside-jump}.  Ding, Hong, and Liu
\cite[Theorem~1.2(ii)]{DHL} state, with the angular size tacitly
normalized, the unphased strong critical-jump theorem at the $L\log L$
endpoint.  Bhojak and Shrivastava
\cite[Theorem~1.1]{BhojakShrivastava} subsequently state the unphased weak
critical-jump and weak $r$-variation endpoints.
Remark~\ref{rem:proof-status} explains why neither of the latter two
arguments is used here.

\paragraph{Recent polynomial-modulation variation theory.}
Becker, Jamneshan, and Thiele
\cite[Theorem~1.3 and Corollary~1.4]{BJT} prove strong $L^p$
$r$-variation estimates, $1<p<\infty$ and $r>2$, for sharp truncations of
H\"older Calder\'on--Zygmund kernels satisfying exact two-sided annular
cancellation.  Their estimate takes the supremum over bounded-degree
polynomial input modulations pointwise inside the $L^p$ norm and holds
more generally on homogeneous Lie groups.  In Euclidean space, this
formulation already contains arbitrary two-variable phases: for each
output point $x$, set
\[
 Q_x(y)=P(x,y)-P(x,0),
 \qquad
 T^{P,1}_{\Omega,u}f(x)
   =e^{iP(x,0)}S_u(K_\Omega,Q_x,f)(x).
\]
Because the output factor is independent of $u$, their smooth-kernel
theorem is not restricted to pure-input phases.  Combining their
one-dimensional result with the method of rotations also gives the
fixed-phase, coefficient-uniform strong $r>2$ estimates considered here
for odd rough angular kernels in $L^1$ (and hence for the entire mean-zero
class when $n=1$).  It does not cover the general even rough component at
the $L\log L$ scale, the critical $2$-jump estimate, or spatial weak type
$(1,1)$.  Conversely, the estimates proved here are uniform in the
coefficients of each fixed phase but do not place a pointwise supremum
over phases inside the spatial norm.

\paragraph{Contribution of the present paper.}
The central result is the weak critical-jump estimate
\eqref{eq:main-weak-jump}.  At $P=0$ it independently resolves the
Jones--Seeger--Wright problem at the $L\log L$ angular endpoint; for
general $P$ it gives a coefficient-uniform extension to all
bounded-degree two-variable polynomial phases, with no nondegeneracy
assumption.  The same theorem covers convolution phases, genuinely
nonconvolution phases, and phase geometries that degenerate on the
diagonal.  The strong critical $2$-jump estimate is proved in parallel.
The strong $r>2$ consequence extends the smooth and odd-kernel subclasses
covered by Becker, Jamneshan, and Thiele to the full rough $L\log L$ class,
including its general even component.  The pointwise critical jump, the
spatial weak endpoint, the extension to every finite radial variation
exponent, and the sharpness results of
Section~\ref{sec:sharpness-obstructions} are separate conclusions not
contained in that result.

\begin{remark}[Independence from prior proofs]\label{rem:proof-status}
The proof below does not use the coefficient-detection and annular-decay
arguments in
\cite{AlSalman,CSS,JiangLu,LuWu03,LuWu04,LuZhang} or the endpoint
fluctuation arguments in \cite{DHL,BhojakShrivastava}.  The observations in
this remark concern the cited proofs as written, not the truth of their
theorem statements.  The discussion of \cite{BhojakShrivastava} refers to
version~2, accessed on August~27, 2026.

The obstruction in \cite{CSS} is inherited from the one-ray reduction of
Lu--Zhang \cite[pp.~210--213]{LuZhang}, which \cite[Lemma~3.2]{CSS}
explicitly adapts.  For the top bidegree block
\[
 P_{m,\ell}(x,y)
 =\sum_{\substack{|\alpha|=m\\|\beta|=\ell}}
   a_{\alpha,\beta}x^\alpha y^\beta,
\]
that reduction detects the coefficient tensor only through its diagonal
restriction
\[
 A_{m,\ell}(\theta)=P_{m,\ell}(\theta,\theta)
 =\sum_{\substack{|\alpha|=m\\|\beta|=\ell}}
   a_{\alpha,\beta}\theta^{\alpha+\beta}.
\]
In the genuinely mixed setting $n\ge2$ and $m,\ell\ge1$, the map
$P_{m,\ell}\mapsto A_{m,\ell}$ is in general not injective.  For example,
the alternating bilinear phase $P(x,y)=x_1y_2-x_2y_1$ has a nonzero top
bidegree block, but $A_{1,1}\equiv0$.  Thus the nonzero top-bidegree
hypothesis in Lemma~3.2 does not supply the nonvanishing needed to apply
Corollary~3.4 in the passage to (4.7).
Related coefficient-detection reductions are invoked or reproduced by
Jiang--Lu \cite[p.~142]{JiangLu}, Lu--Wu
\cite[p.~54]{LuWu03} and \cite[p.~308]{LuWu04}, and Al-Salman
\cite[pp.~82--83]{AlSalman}.  The global phase flag used here avoids this
diagonal compression and retains the full mixed coefficient data.

In \cite{DHL}, the squared-multiplier geometric tail preceding (2.16)
carries a factor comparable to $m$ that is not retained in the displayed
estimate; after the Plancherel square root, the operator-norm estimate
(2.16) is therefore missing a factor comparable to $\sqrt m$.
Independently, Lemma~2.11 is false as stated for $0<\theta<1$: as
$t\downarrow0$, the ratio of its left-hand side to its right-hand side is
comparable to $t^{\theta-1}$.  The lemma is used in precisely this range to
reduce a loss of order $m^{2-\theta}$ to the $m$-weighted angular
summability available at the $L\log L$ endpoint.  Consequently, the
published proof does not justify that endpoint summation in its printed
form.

In version~2 of \cite{BhojakShrivastava}, the short-jump packets
$B_{i,\mathrm{in}}^{k-s,m}$ depend on the output point $x$, whereas (2.2)
and Lemma~A.2 control sums whose coefficients are chosen before $x$ is
evaluated; the variational Rademacher--Menshov theorem does not by itself
supply the missing selector estimate.  In the long-jump argument, the
displayed fixed-shift families $\mathcal D_u$ are not nested: for $n=1$ and
$u=1/3$, the intervals $[1/3,4/3)$ and $[2/3,8/3)$ overlap without
either containing
the other.  These points require additional arguments.  For these reasons,
the present paper proves the relevant critical-jump endpoints
independently.
\end{remark}

\subsection*{Endpoints and sharpness}

The theorem reaches both critical output thresholds: jump exponent $2$ and
spatial weak type $(1,1)$.  For $n\ge2$, the angular $L\log L$ assumption
is sharp for coefficient-uniform estimates of this generality in the
following inherited sense.  Al-Qassem, Cheng, and Pan
\cite[Theorem~1.3(b)]{ACP} construct mean-zero angular functions in
$L(\log L)^{1-\varepsilon}$, for every $\varepsilon>0$, whose associated
scalar rough oscillatory integrals fail to be uniformly bounded over
polynomial phases.

The radial and fluctuation endpoints are settled separately in
Section~\ref{sec:sharpness-obstructions}.  Theorem~\ref{thm:finite-radial}
holds for each fixed $1\le\rho<\infty$, but a one-dimensional Hilbert
multiplier example excludes $\rho=\infty$.  An abstract scalar-path example
shows that the direct Young--Stieltjes transfer necessarily stops before
$\rho=2$, even though operator-specific orthogonality recovers every finite
$\rho$.  Finally, an exact Mellin representation reduces Poisson
$2$-variation to sharp Hilbert truncations and, together with the
Jones--Wang counterexample, proves that the pointwise critical jump cannot
be strengthened to true output $V^2$.  The precise optimal growth of the
$r>2$ constants as $r\downarrow2$ is not addressed.

\subsection*{Three obstructions and their resolution}

\paragraph{Phase degeneracy.}
Neither a displayed monomial coefficient nor the diagonal restriction of
the phase provides a stable normalization.  Simultaneous translation mixes
individual coefficients, and the alternating bilinear phase
$x_1y_2-x_2y_1$ vanishes on the diagonal.  Section~\ref{sec:flag}
constructs a finite global flag on polynomial phases modulo pure input and
output phases.  The transverse rows are placed in reduced row-echelon
form, and the final nonterminal rows consist of convolution phases and
alternating bilinear forms.  Dilation normalizes the first nonzero flag
coordinate, while simultaneous translation preserves that coordinate and
creates only later ones.  This triangular structure yields
coefficient-uniform annular decay and closes the local analysis by a finite
backward induction over the flag.

\paragraph{Modulation-sensitive cancellation.}
The parent-dependent input modulations generated by a nonconvolution phase
destroy ordinary mean-zero cancellation.  The phase-adapted
Calder\'on--Zygmund decomposition replaces it by
\[
 \int_Q e^{iP(c_Q,y)}b_Q(y)\,dy=0
\]
and subtracts the same parent anchor from the operator phase.  The anchored
packet proposition then performs a simultaneous phase--packet induction.
A large first coordinate gives scale-covariant decay before any expansion;
a small coordinate is Fourier-expanded with packet coefficients fixed
before the output point is evaluated.  The all-low branch terminates in
the separated-packet theory.  Proposition~\ref{prop:anchoredpacket}
isolates this interaction among phase reduction, modulation-sensitive
cancellation, and separated-packet estimates.

\paragraph{Output-dependent selectors.}
A jump linearization chooses parameter intervals as functions of the
output point.  The endpoint argument first imposes finite cube, shell, and
parameter cutoffs and assigns every relevant increment a fixed rank.  A
variational Rademacher--Menshov estimate then controls every interval sum
through that ranking, without introducing an output-dependent sign choice.
The phase-adapted bad functions enter directly through cancellation and
the average density bound
$\norm{h_Q}_1\lesssim\alpha|Q|$.  On each shell, a descending
homogeneous-degree packet induction combines Fourier expansion of low
blocks with a concentration-stable high-block estimate.  Across high
shells, an exact oriented two-ray calculation retains a coefficient of
size $R^m$ after all projective cancellations and lower-degree
perturbations.  A one-sided polynomial incidence estimate then sums the
small-cube mesh errors geometrically while the larger atom remains fixed.
The resulting fixed-sign all-shell estimate is precisely the input needed
by the rank argument, and it requires no $L^\infty$ decomposition of the
bad data.

The proof uses the following specialized external inputs: critical-jump
normability of Mirek, Stein, and Zorin-Kranich; the pointwise-inside-supremum
one-dimensional and smooth-approximate-identity jump estimates of Chen,
Ding, Hong, and Liu; the long--short mechanisms of Jones, Seeger, and
Wright; the Hilbert short-variation theorem of Campbell, Jones, Reinhold,
and Wierdl; and Seeger's two separated-packet estimates.  The sharpness
argument additionally uses the Jones--Wang Poisson counterexample and the
finite-dimensional restriction principle of de Leeuw.  We also use the
classical M.~Riesz theorem, Young integration, and the Fefferman--Stein
$H^1$--BMO theory.  Each input is stated in the precise form required at
its point of application.  All polynomial, oscillatory, incidence, packet,
angular-summation, radial-stability, and limiting arguments specific to the
present theorems are proved here.

The proof is organized as follows.  Section~\ref{sec:nonosc-strong}
develops the abstract jump calculus and establishes the unphased strong
endpoint.  Sections~\ref{sec:flag} and~\ref{sec:strong-far} construct the
global phase flag, prove coefficient-uniform annular decay, and close the
strong polynomial jump theorem by angular extrapolation and triangular
induction.  Sections~\ref{sec:terminal-packets}
and~\ref{sec:anchored-induction} establish the terminal separated-packet
estimates and the anchored phase--packet induction.
Section~\ref{sec:degree-packets} provides the concentration-stable equal-
and unequal-shell estimates, while Section~\ref{sec:unphased-weak} proves
the unphased weak endpoint.  Sections~\ref{sec:phase-cz}--\ref{sec:weak-local}
combine the phase-adapted Calder\'on--Zygmund decomposition with far-field
and local homogeneous-degree induction to obtain the polynomial weak
endpoint.  Section~\ref{sec:limits} constructs the trajectory for general
data, derives the variation and convergence consequences, and proves the
bounded-variation Stieltjes transfer.  Section~\ref{sec:finite-radial}
extends the complete theory to every fixed finite radial variation exponent
and proves the shellwise Hilbert refinement.  Section~\ref{sec:sharpness-obstructions}
contains the mixed-norm, abstract-transfer, radial-endpoint, and true
$V^2$ obstructions.

\section{Abstract jump inequalities and the unphased strong endpoint}
\label{sec:nonosc-strong}

This section assembles the jump calculus and the one-dimensional
estimates used throughout the paper, then combines them with a
Riesz-transform factorization to prove the unphased strong critical-jump
theorem.

\begin{lemma}[Gauge normalization]\label{lem:gauge-normalization}
For every $\Omega\in L\log L(\Sph^{n-1})$, the quantity $\mathfrak L$
in \eqref{eq:gauge} is equivalent, with
dimensional constants, to the homogeneous Luxemburg norm associated with
$\Phi_0(t)=(1+t)\log(1+t)-t$.  It also satisfies
\eqref{eq:gauge-parity}.
\end{lemma}

\begin{proof}
The assertion is immediate for $\Omega=0$, so assume
$\mathfrak L(\Omega)>0$.
Put $A=\|\Omega\|_1$ and $M=\mathfrak L(\Omega)$.  The elementary bound
\[
 \Phi_0(t)\le C\bigl(t+t\log^+t\bigr),\qquad t\ge0,
\]
and $M\ge A$ give, for a sufficiently large absolute $C_0$,
\begin{align*}
 \int_{\Sph^{n-1}}\Phi_0\!\left(\frac{|\Omega|}{C_0M}\right)d\sigma
 &\le \frac C{C_0M}\left(A+
   \int_{\Sph^{n-1}}|\Omega|
      \log^+\!\left(\frac{|\Omega|}{M}\right)d\sigma\right)\\
 &\le \frac C{C_0M}\left(A+
   \int_{\Sph^{n-1}}|\Omega|
      \log^+\!\left(\frac{|\Omega|}{A}\right)d\sigma\right)
 \le1.
\end{align*}
This proves one direction of the Luxemburg comparison.

Conversely, suppose that
$\int\Phi_0(|\Omega|/L)d\sigma\le1$, and put
$m_n=\sigma(\Sph^{n-1})$.  Jensen's inequality yields
\[
 \Phi_0\!\left(\frac{A}{m_nL}\right)
 \le\frac1{m_n}\int\Phi_0(|\Omega|/L)d\sigma
 \le\frac1{m_n},
\]
and hence $A\le C_nL$.  Since
\[
 \log^+\!\left(\frac{|\Omega|}{A}\right)
 \le \log\!\left(1+\frac{|\Omega|}{L}\right)
      +\log^+\!\left(\frac L A\right),
\]
while $t\log(1+t)\le\Phi_0(t)+t$, we obtain
\begin{align*}
 \int|\Omega|\log^+(|\Omega|/A)d\sigma
 &\le L\int\frac{|\Omega|}{L}
          \log\!\left(1+\frac{|\Omega|}{L}\right)d\sigma
       +A\log^+(L/A)\\
 &\le L\int\Phi_0(|\Omega|/L)d\sigma+A+L/e
 \le C_nL.
\end{align*}
Thus the two quantities are equivalent, with dimensional constants.

If $R\Omega(\theta)=\Omega(-\theta)$, then reflection is an isometry of
the Luxemburg space.  Convexity gives
\[
 \|\Omega_o\|_{L\log L,\mathrm{Lux}}
 +\|\Omega_e\|_{L\log L,\mathrm{Lux}}
 \le2\|\Omega\|_{L\log L,\mathrm{Lux}}.
\]
Homogeneity and the equivalence just proved give
\eqref{eq:gauge-parity}.
\end{proof}

\subsection{Jump calculus and one-dimensional inputs}

For a scalar path indexed by positive radii, write
\begin{equation}\label{eq:dyadic-jump-definition}
 N^{\rm dyad}_\lambda(F):=
 N_\lambda(F_{2^j}:j\in\mathbb Z),
 \qquad
 \mathbf J^{\rm dyad}(F)
 :=\sup_{\lambda>0}\lambda
 N^{\rm dyad}_\lambda(F)^{1/2}.
\end{equation}
For a finite ordered parameter set $I$, $\mathbf J_I(F)$ denotes the same
pointwise functional with the path restricted to $I$.

\begin{lemma}[Pointwise jump and variation calculus]\label{lem:jumpcalc}
Let $E$ be a linearly ordered set.  For scalar $E$-paths $F,G$ and
$\lambda>0$,
\begin{align}
 N_\lambda(F+G)&\le N_{\lambda/2}(F)+N_{\lambda/2}(G),
 \label{eq:jumpadd}\\
 \lambda N_\lambda(F)^{1/2}&\le V^2(F),
 \label{eq:jumpvar}\\
 V^2(F+G)&\le V^2(F)+V^2(G).
 \label{eq:varadd}
\end{align}
Consequently,
\begin{align}
 \mathbf J(F+G)&\le2\mathbf J(F)+2\mathbf J(G),
 \label{eq:pointwise-jump-quasitriangle}\\
 \mathbf J(F)&\le V^2(F).
 \label{eq:pointwise-jump-by-V2}
\end{align}
If $F_t\to0$ at the upper endpoint of its ordered parameter interval
(with value $0$ there when that endpoint belongs to the parameter set), then
\begin{equation}\label{eq:anchor-controls-maximal}
 \sup_t|F_t|\le\mathbf J(F).
\end{equation}
For every $r>2$,
\begin{equation}\label{eq:pointwise-jump-to-variation}
 V^r(F)\le\kappa_r\mathbf J(F),
 \qquad
 \kappa_r=\left(\sum_{k=1}^\infty k^{-r/2}\right)^{1/r}.
\end{equation}
Finally, if $F=(F_t)_{t>0}$ is a positive-radius path decomposed into
dyadic parameter blocks, then
\begin{equation}\label{eq:longshort}
 \lambda N_\lambda(F)^{1/2}
 \le C\left[\left(\sum_{j\in\mathbb Z}
 V^2(F_\varepsilon:2^j\le\varepsilon\le2^{j+1})^2\right)^{1/2}
 +\lambda N^{\rm dyad}_{\lambda/3}(F)^{1/2}\right],
\end{equation}
and hence
\begin{equation}\label{eq:pointwise-longshort}
 \mathbf J(F)
 \le C\left[\left(\sum_{j\in\mathbb Z}
 V^2(F_\varepsilon:2^j\le\varepsilon\le2^{j+1})^2\right)^{1/2}
 +\mathbf J^{\rm dyad}(F)\right].
\end{equation}
Moreover $\lambda N^{\rm dyad}_\lambda(F)^{1/2}\le
V^2(F_{2^j}:j\in\mathbb Z)$.
\end{lemma}

\begin{proof}
For each strict $\lambda$-jump of $F+G$, one of the two corresponding
increments has modulus $>\lambda/2$.  Assign the pair to such a summand;
each assigned subcollection remains ordered and paired, proving
\eqref{eq:jumpadd}.  Inserting the endpoints of a paired collection into a
variation partition gives \eqref{eq:jumpvar}; Minkowski in the finite
$\ell^2$ sum gives \eqref{eq:varadd}.  Taking the supremum in $\lambda$ in
\eqref{eq:jumpadd} proves \eqref{eq:pointwise-jump-quasitriangle}, and
\eqref{eq:pointwise-jump-by-V2} is \eqref{eq:jumpvar} followed by the same
supremum.

For \eqref{eq:anchor-controls-maximal}, first let $t$ lie strictly below the
upper endpoint and let $s>t$ tend to that endpoint.  For every
$0<\lambda<|F_t-F_s|$ there is one strict $\lambda$-jump, so
$\mathbf J(F)\ge|F_t-F_s|$; now let $s$ tend to the endpoint.  If the
upper endpoint belongs to the parameter set, its value is zero by
hypothesis, so the assertion at that endpoint is immediate.

To prove \eqref{eq:pointwise-jump-to-variation}, fix a finite variation
partition and let $a_1^*\ge a_2^*\ge\cdots$ be the decreasing rearrangement
of the moduli of its increments.  For every $0<\lambda<a_k^*$, at least
$k$ partition increments are strict $\lambda$-jumps.  Retain any $k$ of
them and restore their original temporal order.  They are a subfamily of
the consecutive partition intervals, so their interiors are disjoint and
the right endpoint of each retained interval is no later than the left
endpoint of the next.  Hence they form an admissible paired family.  Thus
$\lambda\sqrt{k}\le\mathbf J(F)$, and letting
$\lambda\uparrow a_k^*$ gives $a_k^*\le k^{-1/2}\mathbf J(F)$.  Hence
\[
 \left(\sum_i a_i^r\right)^{1/r}
 \le\mathbf J(F)\left(\sum_{k=1}^\infty k^{-r/2}\right)^{1/r}.
\]
Taking the supremum over partitions proves the claim.  The numerical bound
in \eqref{eq:kappa-r} follows from the integral test.

For \eqref{eq:longshort}, let $j(t)$ be the integer determined by
$2^{j(t)}\le t<2^{j(t)+1}$.  Fix an ordered paired family $(s_i,t_i)$ of
strict $\lambda$-jumps, and denote its finite cardinality by $N$.  If
$j(s_i)=j(t_i)$, assign the pair to that
dyadic block.  Otherwise put
\[
 \widehat s_i=2^{j(s_i)+1},\qquad
 \widehat t_i=2^{j(t_i)}.
\]
Then $s_i\le\widehat s_i\le\widehat t_i\le t_i$, and
\[
 F_{t_i}-F_{s_i}
 =(F_{t_i}-F_{\widehat t_i})
 +(F_{\widehat t_i}-F_{\widehat s_i})
 +(F_{\widehat s_i}-F_{s_i}).
\]
One of the three terms has modulus $>\lambda/3$.  Choose exactly one such
term, using any fixed tie-breaking rule.  The selected middle terms form
an ordered paired family on the dyadic sequence.  The within-block terms
and the selected initial and terminal terms assigned to any fixed dyadic
block have pairwise disjoint interiors, because the original pairs are
ordered.  Insert all their endpoints, together with the two endpoints of
that dyadic block, into a single increasing partition and delete repeated
points.  Every selected short segment is then one increment of this
partition: two such segments may share an endpoint but cannot overlap in
their interiors.  Therefore the sum of their squared increments is bounded
by the $V^2$ square of the path on that block.  Consequently, if
$N_{\rm sh}$ and $N_{\rm long}$ denote the
numbers assigned to short and middle terms, respectively, then
\[
 N=N_{\rm sh}+N_{\rm long},
 \qquad
 \frac{\lambda^2}{9}N_{\rm sh}
 \le \sum_j V^2(F_\varepsilon:2^j\le\varepsilon\le2^{j+1})^2,
 \qquad
 N_{\rm long}\le N^{\rm dyad}_{\lambda/3}(F).
\]
It follows that
\[
 \lambda\sqrt N
 \le3\left(\sum_j
 V^2(F_\varepsilon:2^j\le\varepsilon\le2^{j+1})^2\right)^{1/2}
 +3\frac\lambda3\,N^{\rm dyad}_{\lambda/3}(F)^{1/2}.
\]
Taking the supremum over all finite strict $\lambda$-jump witnesses proves
\eqref{eq:longshort}, including the case in which the jump count is
infinite.  Taking the supremum in $\lambda$ then gives
\eqref{eq:pointwise-longshort}.  The last assertion follows from
\eqref{eq:jumpvar} on the dyadic sequence.
\end{proof}

Normability at the critical jump exponent is essential for the continuum
angular factorizations below.  On a linearly ordered set $E$, define
\[
 N^{\rm ch}_{\lambda,E}(F):=\sup\Bigl(\{0\}\cup
 \bigl\{J\in\mathbb N:\ \exists\ t_0<\cdots<t_J,\quad t_j\in E,\quad
 |F_{t_j}-F_{t_{j-1}}|>\lambda\ (1\le j\le J)\bigr\}\Bigr)
 \in\mathbb N_0\cup\{\infty\}.
\]
As before, the index-set subscript is omitted when it is displayed or
unambiguous.  Then, for every ordered $E$,
\begin{equation}\label{eq:chain-pair}
 N^{\rm ch}_{\lambda,E}(F)\le2N_{\lambda,E}(F),
 \qquad N_{\lambda,E}(F)\le N^{\rm ch}_{\lambda/2,E}(F).
\end{equation}
The first inequality follows by taking alternating edges of a chain.  For
the second, start from $r_0=s_1$ and $r_1=t_1$ in an admissible family of
pairs.  Inductively, after $r_{i-1}$ has been selected from the preceding
pair, at least one of $s_i,t_i$ is more than $\lambda/2$ away from
$F_{r_{i-1}}$; otherwise
\[
 |F_{t_i}-F_{s_i}|
 \le |F_{t_i}-F_{r_{i-1}}|+|F_{s_i}-F_{r_{i-1}}|
 \le\lambda.
\]
Choose such an endpoint as $r_i$.  The temporal ordering of the pairs,
together with the strict increment, gives $r_{i-1}<r_i$.  The resulting
chain has one initial $\lambda$-increment and then strict
$\lambda/2$-increments, hence has at least as many strict
$\lambda/2$-increments as the original paired family has pairs.  For a
linearly ordered set $E$, put
\begin{equation}\label{eq:chain-jump-seminorm-definition}
 \mathbf J_E^{\rm ch}(F)
 :=\sup_{\lambda>0}\lambda
 N^{\rm ch}_{\lambda,E}(F)^{1/2}.
\end{equation}
We omit $E$ when it is the full parameter set.  For a finite ordered set
$I$ and a spatially measurable $I$-path, also put
\begin{equation}\label{eq:finite-chain-jump-spatial-norm}
 \mathcal J^{\rm ch}_{2,p}(F_t:t\in I)
 :=\norm{\mathbf J_I^{\rm ch}(F)}_p.
\end{equation}
Thus \eqref{eq:chain-pair} gives, pointwise and uniformly in every ordered
set $E$,
\begin{equation}\label{eq:chain-pair-functional}
 \mathbf J_E^{\rm ch}(F)\le\sqrt2\,\mathbf J_E(F),
 \qquad
 \mathbf J_E(F)\le2\mathbf J_E^{\rm ch}(F).
\end{equation}

\begin{lemma}[Pathwise critical-jump normability and integration]
\label{lem:jumpnorm}
For every nonempty finite ordered set $I$, the scalar functional
$\mathbf J_I^{\rm ch}$ is equivalent, with absolute constants uniform in
$I$, to a subadditive seminorm on scalar $I$-paths modulo constant paths.
Consequently, if $\mu$ is a complex measure on a measure space $Z$ and
$F^z=(F_t^z)_{t\in I}$ is coordinatewise $|\mu|$-integrable, then
\begin{equation}\label{eq:pathwise-jump-minkowski}
 \mathbf J_I^{\rm ch}\!\left(\int_ZF^z\,d\mu(z)\right)
 \le C\int_Z\mathbf J_I^{\rm ch}(F^z)\,d|\mu|(z).
\end{equation}
If the paths also depend on $x$, assume that, for every $t\in I$, the map
$(z,x)\mapsto F_t^z(x)$ is jointly measurable and that, for almost every
$x$,
\[
 \int_Z|F_t^z(x)|\,d|\mu|(z)<\infty,
 \qquad t\in I.
\]
Then \eqref{eq:pathwise-jump-minkowski} holds for almost every $x$.  In
particular, for $1\le p<\infty$,
\begin{equation}\label{eq:jump-minkowski}
 \mathcal J^{\rm ch}_{2,p}\!\left(\int_ZF^z\,d\mu(z)\right)
 \le C\int_Z\mathcal J^{\rm ch}_{2,p}(F^z)\,d|\mu|(z),
\end{equation}
whenever the right-hand side is finite.
\end{lemma}

\begin{proof}
Apply \cite[Corollary~2.11]{MSZK} with the spatial measure space $X$
consisting of one point of mass one, with scalar Banach space $\mathbb C$,
and with $p=q=\rho=2$ in the notation of that corollary (the spatial,
Lorentz, and jump exponents, respectively).  The jump count in the cited
source is the consecutive-chain count.  On this one-point space the
Lorentz norm is the absolute value, so its functional is exactly the
scalar finite-path functional
$\mathbf J_I^{\rm ch}$.  In particular, on a one-point space taking the
$\lambda$-supremum before or after the spatial Lorentz norm produces the
same scalar quantity.  Thus this specialization supplies the pointwise
inside-supremum functional used here, rather than merely an estimate with
the amplitude supremum outside a nontrivial spatial norm.  The constant in
the cited corollary is uniform in the cardinality of the finite ordered set
$I$.  This gives a
seminorm $\|\cdot\|_{\mathfrak J_I}$ on scalar paths modulo constants such
that
\[
 C^{-1}\mathbf J_I^{\rm ch}(a)
 \le\|a\|_{\mathfrak J_I}
 \le C\mathbf J_I^{\rm ch}(a),
\]
with $C$ independent of $I$.  The cited source counts increments of size
$\ge\lambda$, whereas we use strict jumps; if $N_\lambda^{\ge}$ and
$N_\lambda^>$ denote the two counts, then
$N_\lambda^>\le N_\lambda^{\ge}\le N_{\lambda/2}^>$, so the corresponding
functionals are equivalent with an absolute constant.

Write $d\mu=\omega\,d|\mu|$, $|\omega|=1$, fix $t_*\in I$, and replace
each $F^z$ by $F^z-F_{t_*}^z$.  Coordinatewise integrability and finiteness
of $I$ give
\[
 \int_Z\sum_{t\in I}|F_t^z-F_{t_*}^z|\,d|\mu|(z)<\infty.
\]
Thus the resulting map is Bochner integrable in the finite-dimensional path
quotient.  Subadditivity and homogeneity give
\[
 \left\|\int_Z\omega(z)(F^z-F_{t_*}^z)\,d|\mu|(z)
 \right\|_{\mathfrak J_I}
 \le\int_Z\|F^z-F_{t_*}^z\|_{\mathfrak J_I}\,d|\mu|(z).
\]
The equivalence above proves \eqref{eq:pathwise-jump-minkowski}.  For finite
$I$, the chain functional is a finite maximum of expressions obtained from
coordinate differences, finite minima, and square roots; it is therefore a
Borel function of the finitely many path coordinates.  Under the stated
joint measurability and sectionwise integrability hypotheses, the scalar
inequality consequently applies for almost every $x$ and both sides are
measurable.  Minkowski's integral inequality then proves
\eqref{eq:jump-minkowski}.
\end{proof}

We next record the pointwise critical-jump estimates for the
one-dimensional model paths used in the factorization.  If
$u\in\Sph^{n-1}$, put
\[
 H_{\varepsilon,u}f(x)=\int_{|s|>\varepsilon}f(x-su)\,\frac{ds}{s},
 \qquad
 A^+_{t,u}f(x)=\frac1t\int_0^t f(x-su)\,ds.
\]

\begin{lemma}[One-dimensional pointwise critical jumps]
\label{lem:one-dimensional-inside-jump}
For every $1<p<\infty$, every $f\in L^p(\R^n)$, and every
$u\in\Sph^{n-1}$,
\begin{equation}\label{eq:one-dimensional-inside-jump}
 \norm{\mathbf J(H_{\varepsilon,u}f:\varepsilon>0)}_p
 +\norm{\mathbf J(A^+_{t,u}f:t>0)}_p
 \le C_p\norm f_p,
\end{equation}
with a constant independent of $u$.  The same estimate holds with the chain
functional $\mathbf J^{\rm ch}$ in place of $\mathbf J$.
\end{lemma}

\begin{proof}
Chen, Ding, Hong, and Liu
\cite[Theorems~1.1 and~1.3]{ChenDingHongLiu} prove the stronger weighted
estimates with the supremum in $\lambda$ taken pointwise before the
$L^p(w)$ norm for rough homogeneous singular integrals and rough averaging
families whose angular function belongs to $L^q$, $q>1$.  Specialize to
$n=1$ and $w\equiv1$.  On $\Sph^0=\{-1,1\}$, the choice
\[
 \Omega(1)=1,\qquad \Omega(-1)=-1
\]
gives the sharp Hilbert truncation, while
\[
 \Omega(1)=1,\qquad \Omega(-1)=0
\]
gives $A_t^+$.  Both angular functions belong to every $L^q(\Sph^0)$.
For a given $p>1$, choose $q>\max\{p,p'\}$.  Then $w\equiv1$ satisfies
condition~(ii) in each cited theorem, and the $L^q(\Sph^0)$ norms of the
two fixed angular functions are absolute constants.  With the usual
counting measure on $\Sph^0$, the specializations are exactly
\[
 T_\varepsilon^{(1,-1)}=H_\varepsilon,
 \qquad M_t^{(1,0)}=A_t^+.
\]
Thus the cited estimates have a constant depending only on $p$, and they
prove the one-dimensional estimates.  We spell out the common
continuum-parameter representatives used in this specialization.  Fix a
Borel representative of $f$ and, for general $u$, write
$x=x_\perp+ru$ with $x_\perp\in u^\perp$.  For almost every $x_\perp$, the
function $g(r)=f(x_\perp+ru)$ belongs to $L^p(\mathbb R)$.  On each such
line, for every $r$ and every $\varepsilon>0$, H\"older's inequality gives
\[
 \int_{|s|>1}\frac{|g(r-s)|}{|s|}\,ds
 \le\norm g_p
 \bigl\||s|^{-1}\mathbf1_{\{|s|>1\}}\bigr\|_{p'}<\infty,
\]
while the portion $\varepsilon<|s|\le1$ is finite by local
integrability.  Hence the sharp Hilbert truncation is absolutely defined.
Moreover $\varepsilon\mapsto H_\varepsilon g(r)$ is locally absolutely
continuous on $(0,\infty)$, since on compact $\varepsilon$-intervals its
derivative is represented almost everywhere by
\[
 \frac{g(r+\varepsilon)-g(r-\varepsilon)}{\varepsilon}.
\]
Likewise $t\mapsto A_t^+g(r)$ is continuous because it is a locally
absolutely continuous primitive divided by $t>0$.  Set both paths equal to
zero on the exceptional lines.  These choices give one common full-measure
set on which the parameter paths are continuous, so their pointwise jump
functionals are measurable by the finite-witness principle.  Applying the
one-dimensional bounds on each good line and then Fubini gives
\eqref{eq:one-dimensional-inside-jump}.  The chain form follows from
\eqref{eq:chain-pair-functional}.
\end{proof}

\begin{lemma}[Radial profiles]\label{lem:profiles}
Let $1<p<\infty$, let $f\in L^p(\R^n)$, let
$u\in\Sph^{n-1}$, and let $\varphi$ be represented
by a right-continuous function that is locally of bounded variation on
$(0,\infty)$,
$\varphi(t)\to0$ as $t\to\infty$, and
\[
 B(\varphi)=\int_{(0,\infty)}t\,d|D\varphi|(t)<\infty.
\]
For $P^{\varphi}_{\varepsilon,u}f(x)=
\int_0^\infty\varphi(s)f(x-\varepsilon su)\,ds$,
the pointwise path is represented by the compatible Stieltjes
representative constructed in the proof, and
\begin{equation}\label{eq:profile-jump}
 \norm{\mathbf J(P^{\varphi}_{\varepsilon,u}f:\varepsilon>0)}_p
 \le C_pB(\varphi)\norm f_p.
\end{equation}
Suppose, more generally, that $(\varphi_u)_{u\in\Sph^{n-1}}$ is a family
satisfying the preceding hypotheses, with right-continuous representatives
chosen so that $(u,s)\mapsto\varphi_u(s)$ is jointly measurable on
$\Sph^{n-1}\times(0,\infty)$, and suppose that
$\int_{\Sph^{n-1}}B(\varphi_u)\,d\sigma(u)<\infty$.  Then, for every
$\varepsilon>0$, the map $u\mapsto P^{\varphi_u}_{\varepsilon,u}f$ is
strongly measurable in $L^p$, its Bochner integral is well defined, and
\begin{equation}\label{eq:profile-integrated-jump}
 \left\|\mathbf J\left(
  \int_{\Sph^{n-1}}P^{\varphi_u}_{\varepsilon,u}f\,d\sigma(u):
  \varepsilon>0\right)\right\|_p
 \le C_p\left(\int_{\Sph^{n-1}}B(\varphi_u)d\sigma(u)\right)\norm f_p.
\end{equation}
Both conclusions remain valid with the chain functional.
\end{lemma}

\begin{proof}
Fix a Borel representative of $f$.  Stieltjes integration and the zero
limit at infinity give, for every $s>0$, the following identity with no
endpoint ambiguity.  For $R>s$ our right-continuous convention gives
\[
 D\varphi((s,R])=\varphi(R)-\varphi(s),
 \qquad
 |D\varphi|((s,\infty))\le s^{-1}B(\varphi)<\infty.
\]
Letting $R\to\infty$ and using continuity from below of the finite measure
$|D\varphi|$ gives
\[
 \varphi(s)=-D\varphi((s,\infty)).
\]
Moreover, Tonelli's theorem gives
\[
 \int_0^\infty|\varphi(s)|\,ds
 \le\int_0^\infty\int_{(s,\infty)}d|D\varphi|(t)\,ds
 =\int_{(0,\infty)}\int_0^t ds\,d|D\varphi|(t)
 =B(\varphi).
\]
Thus the defining integral for $P^\varphi_{\varepsilon,u}f$ is a Bochner
integral in $L^p$.  The $L^p$-valued Bochner--Fubini theorem then yields
\begin{equation}\label{eq:profile-stieltjes-representative}
 P^\varphi_{\varepsilon,u}f
 =-\int_{(0,\infty)}tA^+_{\varepsilon t,u}f\,d\varphi(t)
 \qquad\text{in }L^p.
\end{equation}
Indeed, the total variation of the right-hand side is bounded in $L^p$
by $B(\varphi)\norm f_p$.  Below we show that the pointwise Stieltjes
integral on the right is simultaneously defined for every
$\varepsilon>0$ on one full-measure set; that integral is the compatible
representative used here.

First fix a finite ordered parameter set $I\subset(0,\infty)$.  Applying
\eqref{eq:pathwise-jump-minkowski} to
\eqref{eq:profile-stieltjes-representative}, and then applying Minkowski's
integral inequality in $L^p$, gives
\begin{align}
 \bigl\|\mathbf J_I^{\rm ch}
   (P^\varphi_{\varepsilon,u}f:\varepsilon\in I)\bigr\|_p
 &\le C\int_{(0,\infty)}t
  \bigl\|\mathbf J_I^{\rm ch}
   (A^+_{\varepsilon t,u}f:\varepsilon\in I)\bigr\|_p
  \,d|D\varphi|(t)\notag\\
 &\le C_pB(\varphi)\norm f_p.                 \label{eq:profile-finite-parameter-bound}
\end{align}
For the second inequality, multiplication by a fixed $t>0$ preserves the
order of the parameters, and we used
Lemma~\ref{lem:one-dimensional-inside-jump}.  In particular, the bound is
uniform in $I$.

It remains to justify passage from a dense parameter set to the entire
pointwise path.  Define the measurable directional one-sided maximal
function
\begin{equation}\label{eq:directional-one-sided-maximal}
 \mathcal M_u^+f(x)
 :=\sup_{r\in\mathbb Q_+}|A^+_{r,u}f(x)|,
 \qquad
 \norm{\mathcal M_u^+f}_p\le C_p\norm f_p.
\end{equation}
The estimate is the one-dimensional one-sided Hardy--Littlewood maximal
theorem, applied on almost every line parallel to $u$.  On those same
lines the restriction of $f$ is locally integrable and
$r\mapsto A^+_{r,u}f(x)$ is continuous on $(0,\infty)$.  Consequently,
the rational supremum in \eqref{eq:directional-one-sided-maximal} also
dominates $|A^+_{r,u}f(x)|$ for every real $r>0$.

Put $d\nu(t)=t\,d\varphi(t)$.  This is a finite complex measure with
$|\nu|((0,\infty))=B(\varphi)$.  If $\varepsilon_k\to\varepsilon>0$, then
for almost every $x$ the integrand in
\eqref{eq:profile-stieltjes-representative} converges pointwise in $t$,
and
\[
 |A^+_{\varepsilon_k t,u}f(x)-A^+_{\varepsilon t,u}f(x)|
 \le2\mathcal M_u^+f(x).
\]
Dominated convergence with respect to $|\nu|$ proves that
$\varepsilon\mapsto P^\varphi_{\varepsilon,u}f(x)$ is continuous for
almost every $x$.  Let $I_m\uparrow\mathbb Q_+$ be an increasing
exhaustion by finite ordered sets.  Every strict chain witness is finite,
so
\[
 \mathbf J_{I_m}^{\rm ch}(P^\varphi)(x)
 \uparrow\mathbf J_{\mathbb Q_+}^{\rm ch}(P^\varphi)(x).
\]
Monotone convergence applied to the $p$th powers in
\eqref{eq:profile-finite-parameter-bound} yields the same estimate on
$\mathbb Q_+$.  The continuity just proved and the chain form of the
finite-witness principle in Lemma~\ref{lem:finite-witness} then give the
chain form of \eqref{eq:profile-jump} on the full parameter set.

We now treat the parameterized family.  For rational $0<a<b<\infty$, the
identity
\[
 D\varphi_u((a,b])=\varphi_u(b)-\varphi_u(a)
\]
shows that $u\mapsto D\varphi_u((a,b])$ is measurable.  On each
$K_m=[m^{-1},m]$, these rational half-open intervals form a countable
generating algebra and $|D\varphi_u|(K_m)\le mB(\varphi_u)$.  The monotone
class theorem therefore shows that $u\mapsto D\varphi_u(E)$ is measurable
for every Borel $E\subset K_m$.  Approximating $t\mathbf1_E(t)$ by bounded
simple functions on $K_m$, and then letting $m\to\infty$, shows that the
measures
\[
 \nu_u(dt):=t\,D\varphi_u(dt)
\]
form a measurable finite complex-measure kernel.  For completeness, total
variation on a Borel set is the supremum of
$\sum_k|\nu_u(E_k)|$ over finite partitions by members of the preceding
countable algebra; regular approximation then gives the same supremum as
the full Borel definition.  Hence $(u,E)\mapsto|\nu_u|(E)$ is a measurable
positive-measure kernel, and
\[
 |\nu_u|((0,\infty))=B(\varphi_u).
\]
The map
\[
 (u,t)\longmapsto A^+_{\varepsilon t,u}f
 =\int_0^1 f(\,\cdot-\varepsilon tr u)\,dr
\]
is strongly measurable with values in the separable space $L^p$: in fact
strong continuity of translations and approximation of the $r$-integral
by simple functions show that it is continuous in $(u,t)$.  Moreover,
\begin{align*}
 &\int_{\Sph^{n-1}}\int_{(0,\infty)}
   \|A^+_{\varepsilon t,u}f\|_p\,d|\nu_u|(t)d\sigma(u)\\
 &\qquad\le \|f\|_p
   \int_{\Sph^{n-1}}B(\varphi_u)d\sigma(u)<\infty.
\end{align*}
Integration against the preceding
kernel now shows that
\[
 u\longmapsto P^{\varphi_u}_{\varepsilon,u}f
 :=-\int_{(0,\infty)}A^+_{\varepsilon t,u}f\,d\nu_u(t)
\]
is strongly measurable.  It agrees in $L^p$ with the original profile
integral by the same Bochner--Fubini argument as above, and
\[
 \norm{P^{\varphi_u}_{\varepsilon,u}f}_p
 \le B(\varphi_u)\norm f_p.
\]
Hence its integral in $u$ is a well-defined Bochner integral.

For every finite ordered $I$, apply
\eqref{eq:pathwise-jump-minkowski} first to the measure $\nu_u$ and then
to $d\sigma(u)$, and finally use Minkowski's integral inequality and
Lemma~\ref{lem:one-dimensional-inside-jump}.  This gives
\begin{align}
 &\left\|\mathbf J_I^{\rm ch}\left(
   \int_{\Sph^{n-1}}P^{\varphi_u}_{\varepsilon,u}f\,d\sigma(u):
   \varepsilon\in I\right)\right\|_p\notag\\
 &\qquad\le C\int_{\Sph^{n-1}}\int_{(0,\infty)}
   \left\|\mathbf J_I^{\rm ch}
   (A^+_{\varepsilon t,u}f:\varepsilon\in I)\right\|_p
   \,d|\nu_u|(t)\,d\sigma(u)\notag\\
 &\qquad\le C_p\left(\int_{\Sph^{n-1}}B(\varphi_u)\,d\sigma(u)\right)
   \norm f_p.                                  \label{eq:integrated-profile-finite-parameter}
\end{align}

For completeness, we verify that the integrated path also has one
pointwise representative continuous in $\varepsilon$.  For the fixed Borel
representative of $f$, the formula
\[
 A^+_{r,u}f(x)=\int_0^1f(x-rsu)\,ds
\]
is jointly measurable in $(u,r,x)$.  Consequently the rational supremum
$(u,x)\mapsto\mathcal M_u^+f(x)$ is measurable.  The set on which
$r\mapsto A^+_{r,u}f(x)$ is continuous can be expressed through Cauchy
conditions on rational sequences and is therefore measurable.  Moreover,
\begin{equation*}
 G(x):=\int_{\Sph^{n-1}}B(\varphi_u)\mathcal M_u^+f(x)\,d\sigma(u)
 \quad\text{satisfies}\quad
 \norm G_p\le C_p\left(\int_{\Sph^{n-1}}B(\varphi_u)\,d\sigma(u)\right)
 \norm f_p.
\end{equation*}
Fubini's theorem supplies a common full-measure set of $x$ on which
$G(x)<\infty$ and, for almost every $u$, the map
$r\mapsto A^+_{r,u}f(x)$ is continuous.  On this set define
\[
 \mathcal P_\varepsilon f(x)
 :=-\int_{\Sph^{n-1}}\int_{(0,\infty)}
 A^+_{\varepsilon t,u}f(x)\,d\nu_u(t)\,d\sigma(u).
\]
The integral is absolutely convergent because it is bounded by $G(x)$,
and it represents the preceding $L^p$-valued Bochner integral.  If
$\varepsilon_k\to\varepsilon$, its integrand converges for almost every
$(u,t)$ and its difference is dominated by
$2\mathcal M_u^+f(x)$ with respect to
$d\sigma(u)\,d|\nu_u|(t)$.  Dominated convergence therefore proves that
$\varepsilon\mapsto\mathcal P_\varepsilon f(x)$ is continuous on the
same full-measure set.  For an increasing finite exhaustion
$I_m\uparrow\mathbb Q_+$, the finite-parameter chain functionals increase
pointwise to the chain functional on $\mathbb Q_+$.  Monotone convergence
in \eqref{eq:integrated-profile-finite-parameter}, followed by the
continuity just proved and the chain form of
Lemma~\ref{lem:finite-witness}, now proves
\eqref{eq:profile-integrated-jump} for the full parameter set.  Finally,
\eqref{eq:chain-pair-functional} converts both chain estimates into the
stated paired estimates (and gives the asserted chain versions as well).
\end{proof}

We repeatedly use the following elementary compactness principle.

\begin{lemma}[Finite witnesses and dense parameter sets]
\label{lem:finite-witness}
Let $I$ be a finite ordered set, and suppose that
$F^{(m)}_t\to F_t$ for every $t\in I$.  Then, for every $\lambda>0$ and
$r\ge1$,
\[
 N_\lambda(F_t:t\in I)
 \le\liminf_{m\to\infty}N_\lambda(F^{(m)}_t:t\in I),
 \qquad
 N^{\rm ch}_\lambda(F_t:t\in I)
 \le\liminf_{m\to\infty}N^{\rm ch}_\lambda(F^{(m)}_t:t\in I),
\]
\[
 V^r(F_t:t\in I)
 \le\liminf_{m\to\infty}V^r(F^{(m)}_t:t\in I),
\]
and
\begin{equation}\label{eq:finite-J-lsc}
 \mathbf J_I(F)\le\liminf_{m\to\infty}\mathbf J_I(F^{(m)}).
\end{equation}
Moreover,
\begin{equation}\label{eq:finite-chain-J-lsc}
 \mathbf J_I^{\rm ch}(F)
 \le\liminf_{m\to\infty}\mathbf J_I^{\rm ch}(F^{(m)}).
\end{equation}
If a path $F$ is continuous on compact subintervals of an interval $J$
and $D\subset J$ is dense, then its paired and chain strict jump counts,
its paired and chain pointwise critical-jump functionals, and its
$r$-variation on $D$ agree with those on $J$.
\end{lemma}

\begin{proof}
Any finite family of strict $\lambda$-jumps of $F$ has a positive minimum
excess over $\lambda$; all of its witnessing increments therefore persist
for sufficiently large $m$.  This applies verbatim to both paired and
chain witnesses and proves the two jump-count inequalities.  Applying
the same argument to a fixed finite partition and then taking the
supremum proves the variation inequality.  If $a<\mathbf J_I(F)$, choose a jump amplitude
$\lambda$ and a finite strict witness with
$\lambda N_\lambda(F)^{1/2}>a$; that witness persists for all sufficiently
large $m$, proving \eqref{eq:finite-J-lsc}.  The identical persistence
argument for a finite chain witness proves
\eqref{eq:finite-chain-J-lsc}.  Finally, continuity permits the finitely
many points in any paired witness, chain witness, or variation partition
to be perturbed into $D$ while preserving their strict order and changing
every increment by an arbitrarily small amount.  This proves all dense-set
assertions.
\end{proof}

\begin{lemma}[Homogeneous tail counting]\label{lem:tailcount}
Let $\Omega\in L\log L(\Sph^{n-1})$ and put $A=\norm\Omega_1$.
For every $\delta>0$,
\begin{equation}\label{eq:tailcount}
 \sum_{s\ge0}\norm{\Omega\mathbf1_{\{|\Omega|>2^{\delta s}A\}}}_1
 \le C_\delta\mathfrak L(\Omega).
\end{equation}
If $\int\Omega=0$, put $m_n=\sigma(\Sph^{n-1})$ and
$E_s=\{|\Omega|>2^{\delta s}A\}$,
$c_s=m_n^{-1}\int_{E_s}\Omega\,d\sigma$.  Then
\begin{equation}\label{eq:meanzerosplit}
 \Omega_s^\sharp=\Omega\mathbf1_{E_s}-c_s,
 \qquad \omega_s=\Omega\mathbf1_{E_s^c}+c_s
\end{equation}
have mean zero and satisfy
\begin{equation}\label{eq:splitbounds}
 \norm{\omega_s}_\infty\le C_n2^{\delta s}A,
 \quad \norm{\omega_s}_1\le2A,
 \quad \norm{\Omega_s^\sharp}_1\le2\norm{\Omega\mathbf1_{E_s}}_1.
\end{equation}
\end{lemma}

\begin{proof}
If $A=0$, all assertions are immediate.  Hence assume $A>0$.
Tonelli gives
\[
\sum_s\norm{\Omega\mathbf1_{E_s}}_1
 =\int|\Omega(\theta)|\#\{s:2^{\delta s}A<|\Omega(\theta)|\}\,d\sigma(\theta),
\]
and the counting function is bounded by
$C_\delta(1+\log^+(|\Omega|/A))$.  This proves \eqref{eq:tailcount}.
The identities and bounds in \eqref{eq:meanzerosplit}--\eqref{eq:splitbounds}
follow from $\int_{E_s^c}\Omega=-m_nc_s$ and
$m_n|c_s|\le\norm{\Omega\mathbf1_{E_s}}_1\le A$.
\end{proof}

\begin{lemma}[Compact $L\log L$ data in $H^1$]\label{lem:LlogL-H1}
Suppose that $F$ is supported in a fixed ball and $\int F=0$.  Define
$\mathcal E(0)=0$, while for $F\ne0$ put
\[
 \mathcal E(F)=\norm F_1+
 \int |F(x)|\log\!\left(e+\frac{|F(x)|}{\norm F_1}\right)dx.
\]
If $\mathcal E(F)<\infty$, then
$\norm F_{H^1(\R^n)}\le C\mathcal E(F)$, with a constant depending only
on the dimension and the supporting ball.
\end{lemma}

\begin{proof}
If $F=0$ there is nothing to prove.  Put $A=\norm F_1$.
By translation invariance we may suppose that
$\supp F\subset B(0,R)$, where $R$ is fixed.
Let $\phi\in C_c^\infty(B(0,1))$ have integral one and let
$\mathcal M_\phi F=\sup_{t>0}|F*\phi_t|$.  On a fixed enlargement of the
support, $\mathcal M_\phi F\le C MF$.  The refined weak estimate
\[
 |\{MF>2s\}|\le \frac C s\int_{\{|F|>s\}}|F|
\]
follows by decomposing
$F=F\mathbf1_{\{|F|\le s\}}+F\mathbf1_{\{|F|>s\}}$: the maximal
function of the first summand is at most $s$, and the weak $(1,1)$ maximal
theorem applies to the second.  If $B_R$ is a fixed enlargement of
$B(0,R)$, layer cake, with threshold $s_0=A$, gives
\begin{align*}
 \int_{B_R}MF
 &=\int_0^\infty|B_R\cap\{MF>s\}|\,ds\\
 &\le C A+C\int_A^\infty\frac1s
       \int_{\{|F|>s/2\}}|F(x)|\,dx\,ds\\
 &\le C A+C\int|F(x)|
       \log^+\!\left(\frac{2|F(x)|}{A}\right)dx
 \le C\mathcal E(F),
\end{align*}
where Tonelli justifies the interchange.  Outside that ball, cancellation
gives
\[
 |F*\phi_t(x)|
 \le\int|F(y)|\,|\phi_t(x-y)-\phi_t(x)|dy
 \le C_{\phi,R}|x|^{-n-1}\norm F_1.
\]
Indeed, whenever the integrand is nonzero, $t\ge c_R|x|$, and the
mean-value theorem gives
\[
 |\phi_t(x-y)-\phi_t(x)|
 \le |y|\|\nabla\phi_t\|_\infty
 \le C_\phi R t^{-n-1}
 \le C_{\phi,R}|x|^{-n-1}.
\]
The last majorant is integrable off $B_R$.  Thus
$\norm{\mathcal M_\phi F}_1\le C\mathcal E(F)$.  The radial
maximal characterization of real $H^1$ is
\cite[Theorem~11 and the following Remark~2]{FS72}; this proves the
assertion.  The Riesz-transform characterization
$R_j:H^1\to L^1$ is recorded in
\cite[Section~II.2, p.~144]{FS72}.
\end{proof}

\begin{proposition}[Unphased strong critical-jump inequality]
\label{prop:nonosc-strong}
If $\int\Omega=0$ and $\Omega\in L\log L(\Sph^{n-1})$, then, for every
$1<p<\infty$ and every $f\in L^p(\R^n)$,
\begin{equation}\label{eq:nonosc-strong}
 \Jop{p}(T^0_{\Omega,\varepsilon}f:\varepsilon>0)
 \le C_{n,p}\mathfrak L(\Omega)\norm f_p.
\end{equation}
\end{proposition}

\begin{proof}
Fix $1<p<\infty$.  We first take $f\in\mathcal S(\R^n)$, restrict the
truncation parameter to a finite ordered set $I\subset(0,\infty)$, and
use the chain convention.  All constants below are independent of $I$.

We record the measurability convention for the angular integrations.  For
$M>\varepsilon$, put
\[
 H^M_{\varepsilon,u}f
 =\int_{\varepsilon<|s|<M}f(\,\cdot-su)\,\frac{ds}{s}.
\]
Strong continuity of translations makes
$u\mapsto H^M_{\varepsilon,u}f$ continuous as an $L^p$-valued map.  The
one-dimensional fixed-truncation theorem, followed by slicing along lines
parallel to $u$, gives
\[
 \sup_{u\in\Sph^{n-1}}\norm{H_{\varepsilon,u}f}_p
 \le C_p\norm f_p,
\]
and $H^M_{\varepsilon,u}f\to H_{\varepsilon,u}f$ in $L^p$ for every
$u$.  Hence $u\mapsto H_{\varepsilon,u}f$ is strongly measurable and
Bochner integrable against every $L^1$ angular density.  Likewise, for
every $a\in L^1(0,\infty)$, the map
\[
 u\longmapsto\int_0^\infty
 a(s)f(\,\cdot-\varepsilon su)\,ds
\]
is strongly continuous in $L^p$, by dominated convergence.  Thus every
application of \eqref{eq:jump-minkowski} below satisfies its strong
measurability and Bochner-integrability hypotheses.

By density, both conclusions extend to every $h\in L^p(\R^n)$.  Indeed,
the fixed-truncation estimate for $H_{\varepsilon,u}$ is uniform in $u$,
while
\[
 \left\|\int_0^\infty
 a(s)h(\,\cdot-\varepsilon su)\,ds\right\|_p
 \le\norm a_1\norm h_p
\]
uniformly in $u$.  Approximating $h$ in $L^p$ by Schwartz functions
therefore makes each corresponding $L^p$-valued map a uniform $L^p$
limit, in $u$, of strongly measurable maps.  In particular these
conclusions apply to the functions $g_j=-R_jf$ introduced below.

Decompose $\Omega=\Omega_o+\Omega_e$ into its odd and even parts.  The
gauge is stable under this decomposition: this is
\eqref{eq:gauge-parity}.

\smallskip
\noindent\emph{The odd component.}

For odd $\Omega$, polar coordinates give the exact rotation identity
\begin{equation}\label{eq:odd-rotation}
 T^0_{\Omega,\varepsilon}f
 =\frac12\int_{\Sph^{n-1}}\Omega(u)H_{\varepsilon,u}f\,d\sigma(u).
\end{equation}
Lemma~\ref{lem:jumpnorm} and Lemma~\ref{lem:one-dimensional-inside-jump} yield
\begin{equation}\label{eq:odd-jump}
 \mathcal J^{\rm ch}_{2,p}
 (T^0_{\Omega,\varepsilon}f:\varepsilon\in I)
 \le C_p\norm\Omega_1\norm f_p.
\end{equation}

\smallskip
\noindent\emph{The even component and the Riesz-transform factorization.}

Assume now that $\Omega$ is even; when $n=1$ its mean-zero condition makes
it identically zero.  Let $K$ be the homogeneous principal-value
distribution
\[
 K(x)=\Omega(x/|x|)|x|^{-n}.
\]
Its action on a Schwartz function is obtained by integrating over
$\{\delta<|x|<R\}$ and passing to $\delta\downarrow0$,
$R\uparrow\infty$.  More explicitly, angular cancellation gives the
absolutely convergent formula
\[
 \langle K,\phi\rangle
 =\int_{\Sph^{n-1}}\Omega(u)\int_0^\infty
 \bigl[\phi(ru)-\phi(0)\mathbf1_{(0,1)}(r)\bigr]
 \frac{dr}{r}\,d\sigma(u),
 \qquad \phi\in\mathcal S(\R^n).
\]
Near zero the bracket is $O(r)$, while at infinity Schwartz decay applies;
thus this formula also proves that the principal-value limit exists and is
tempered.

Choose a smooth radial function $\Phi$ with $\Phi=0$ on $[0,1/4]$ and
$\Phi=1$ on $[3/4,\infty)$, and put $\psi=1-\Phi$.  We now fix the
distributional normalization of every Riesz transform used below.  Write
\[
 \mathfrak r_j(\xi)=-i\frac{\xi_j}{|\xi|},\qquad \xi\ne0,
\]
so that $\sum_j\mathfrak r_j(\xi)^2=-1$.

The following rough-multiplier estimate fixes all low-frequency choices:
\begin{equation}\label{eq:rough-multiplier-normalization}
 \norm{\widehat K}_{L^\infty(\R^n\setminus\{0\})}
 +\norm{\widehat{K\psi}}_{L^\infty(\R^n)}
 \le C_n\mathfrak L(\Omega).
\end{equation}
For completeness, if $e=\xi/|\xi|$, the polar-coordinate formula for the
even kernel gives, up to a harmless Fourier-normalization constant,
\[
 \widehat K(\xi)
 =c_n\int_{\Sph^{n-1}}\Omega(u)\log|e\cdot u|\,d\sigma(u).
\]
The term containing $\log|\xi|$ is annihilated by $\int\Omega=0$, and
the sine term vanishes because $\Omega$ is even.  Also
\[
 \widehat{K\psi}(\xi)
 =\int_{\Sph^{n-1}}\Omega(u)a_\psi(\xi\cdot u)\,d\sigma(u),
 \qquad
 a_\psi(t)=\int_0^\infty\psi(r)(\cos(rt)-1)\,\frac{dr}{r}.
\]
One has $|a_\psi(t)+\log^+|t||\le C_\psi$ and, uniformly for $R>0$
and $|\tau|\le1$,
\[
 |a_\psi(R\tau)+\log^+R|
 \le C_\psi\bigl(1+|\log|\tau||\bigr).
\]
The functions $u\mapsto1+|\log|e\cdot u||$ have uniformly bounded
exponential-Orlicz norm.  The $L\log L$--exponential-Orlicz H\"older
inequality first gives
\[
 \left|\int\Omega(u)\log|e\cdot u|\,d\sigma(u)\right|
 \le C_n\mathfrak L(\Omega).
\]
Writing $\xi=Re$ and using $\int\Omega=0$ also gives the exact
cancellation
\[
 \widehat{K\psi}(Re)
 =\int_{\Sph^{n-1}}\Omega(u)
   [a_\psi(Re\cdot u)+\log^+R]\,d\sigma(u),
\]
so the preceding uniform estimate with $\tau=e\cdot u$ and the same
Orlicz H\"older inequality prove
\eqref{eq:rough-multiplier-normalization}.

To justify the formulas without a boundedness assumption on $\Omega$, put
$m_n=\sigma(\Sph^{n-1})$ and set
\[
 \Omega_M=\Omega\mathbf1_{\{|\Omega|\le M\}}
 -m_n^{-1}\int_{\Sph^{n-1}}
       \Omega\mathbf1_{\{|\Omega|\le M\}}\,d\sigma.
\]
Then $\Omega_M$ is bounded, even, and mean zero, and
$\mathfrak L(\Omega_M-\Omega)\to0$ by absolute continuity of the
$L\log L$ Luxemburg norm.  The displayed distributional formulas hold
for $\Omega_M$; their uniform Orlicz bounds and convergence in
$\mathcal S'$ allow $M\to\infty$, proving them for the original
$\Omega$.

We define, as tempered distributions,
\[
 \widehat{N_j}=\mathfrak r_j\widehat K,\qquad
 \widehat{R_j(K\psi)}=\mathfrak r_j\widehat{K\psi},
\]
and
\[
 \widehat{R_j(K\Phi)}
 =\mathfrak r_j\bigl(\widehat K-\widehat{K\psi}\bigr).
\]
Values assigned at $\xi=0$ are irrelevant, since all three multipliers
are locally integrable functions.  In particular,
\[
 R_j(K\Phi)=N_j-R_j(K\psi)
\]
holds exactly in $\mathcal S'$, with no term supported at the frequency
origin.  Because $\mathfrak r_j$ is homogeneous of degree zero, these
operators commute with $L^1$-normalized dilations and with
$\mathcal A_0=x\cdot\nabla+n$.  On separated supports the same
normalization agrees with the absolutely convergent Riesz-kernel pairing,
and on $H^1$ pieces it agrees with the usual $H^1\to L^1$ Riesz
transform.  Both assertions follow by smoothing $\mathfrak r_j$ near
$\xi=0$ and passing to the limit.

We repeatedly use the following explicit fixed-annulus consequence
of Lemma~\ref{lem:LlogL-H1}.  If
$\eta\in C_c^\infty((0,\infty))$ is supported in $c\le r\le C$, then
\begin{equation}\label{eq:fixed-annulus-H1}
 F_\eta(x)=\eta(|x|)K(x),\qquad
 \int F_\eta=0,\qquad
 \norm{F_\eta}_{H^1}\le C_{\eta,n}\mathfrak L(\Omega).
\end{equation}
The cancellation follows in polar coordinates.  Put
$c_\eta=\int_0^\infty|\eta(r)|dr/r$.  If
$c_\eta\norm\Omega_1=0$, then $F_\eta=0$ and the claim is immediate.
Otherwise
\[
 \|F_\eta\|_1=c_\eta\|\Omega\|_1.
\]
Since $r$ stays in a fixed compact subinterval of $(0,\infty)$, polar
coordinates also give
\begin{align*}
 &\int|F_\eta|
 \log\!\left(e+\frac{|F_\eta|}{\norm{F_\eta}_1}\right)\\
 &\qquad\le C_{\eta,n}\int_{\Sph^{n-1}}|\Omega(u)|
 \log\!\left(e+C_{\eta,n}
              \frac{|\Omega(u)|}{\|\Omega\|_1}\right)d\sigma(u)
 \le C_{\eta,n}\mathfrak L(\Omega).
\end{align*}
The asserted $H^1$ bound now follows from
Lemma~\ref{lem:LlogL-H1}.

We first establish
\begin{equation}\label{eq:Nj-angular}
 N_j(x)=\frac{\nu_j(x/|x|)}{|x|^n},\qquad
 \nu_j\ \hbox{odd},\qquad
 \norm{\nu_j}_1\le C_n\mathfrak L(\Omega).
\end{equation}
It suffices to estimate $N_j$ on $\{3/4\le|x|\le3/2\}$.  Choose a smooth
radial partition $\chi_<+\chi_0+\chi_>=1$ with
\[
 \supp\chi_<\subset[0,2/3),\qquad
 \supp\chi_0\subset(1/2,3),\qquad
 \supp\chi_>\subset(2,\infty),
\]
so that the supports of the first and last pieces are separated from the
compact $x$-annulus.
If
$\rho_j(z)=c_nz_j|z|^{-n-1}$ is the Riesz kernel, angular cancellation gives
for the first piece
\[
 \int_0^{2/3}\chi_<(r)\int\Omega(u)
 [\rho_j(x-ru)-\rho_j(x)]\,d\sigma(u)\frac{dr}{r},
\]
whose absolute value is at most $C\norm\Omega_1$.  The last piece is
absolutely bounded by the same quantity, because $|x-y|\simeq|y|$ and the
remaining radial integral is $\int_2^\infty r^{-n}dr/r$.  Indeed, on the
first piece the mean-value theorem gives
$|\rho_j(x-ru)-\rho_j(x)|\le Cr$, so its radial integral is
$C\int_0^{2/3}dr$.  The middle piece is
controlled by \eqref{eq:fixed-annulus-H1}; hence its Riesz transform is in
$L^1$ with the same bound.  Integrating over the compact $x$-annulus, and
then using homogeneity and parity of the Riesz multiplier, gives
\[
 \int_{3/4\le|x|\le3/2}|N_j(x)|dx
 =\log2\int_{\Sph^{n-1}}|\nu_j(u)|d\sigma(u)
 \le C_n\mathfrak L(\Omega),
\]
which proves \eqref{eq:Nj-angular}.  The annular calculation gives an actual locally
integrable representative away from the origin.  Two homogeneous
extensions that agree there can differ only by a finite linear combination
of derivatives of $\delta_0$.  Homogeneity of degree $-n$ permits only a
multiple of $\delta_0$, which is even, whereas the Fourier-defined $N_j$
is odd.  Thus there is no origin-supported ambiguity.

Define the commutator correction
\begin{equation}\label{eq:Ej-def}
 E_j=R_j(K\psi)-\psi N_j.
\end{equation}
By the Fourier normalization above and the identity $K\Phi=K-K\psi$
in $\mathcal S'$, one has the exact distributional identity
\begin{equation}\label{eq:RjKPhi}
 R_j(K\Phi)=\Phi N_j-E_j.
\end{equation}
Let $\mathcal A_0=x\cdot\nabla+n$, the generator of $L^1$-normalized
dilations.  We next prove
\begin{equation}\label{eq:Ej-bounds}
 \norm{E_j}_1+\norm{(\mathcal A_0-1)E_j}_1
 \le C_n\mathfrak L(\Omega).
\end{equation}
For $|x|<1/8$, formula \eqref{eq:RjKPhi} reads
$E_j=-R_j(K\Phi)$ and the support separation gives an absolutely convergent
bound.  More precisely, splitting at $|y|=1$ gives
\[
 |R_j(K\Phi)(x)|
 \le C\|\Omega\|_1\left(\int_{1/4}^1\frac{dr}{r}
       +\int_1^\infty r^{-n}\frac{dr}{r}\right).
\]
For $|x|>2$, use $E_j=R_j(K\psi)$ and subtract $\rho_j(x)$ inside the
angular integral.  Here $\psi$ is supported in $[0,3/4]$ and
$\int K(y)\psi(|y|)dy=0$ in the principal-value sense, so the mean-value
theorem applied to
\[
 \int\Omega(u)\int_0^{3/4}\psi(r)
 [\rho_j(x-ru)-\rho_j(x)]\frac{dr}{r}d\sigma(u)
\]
gives $C\norm\Omega_1|x|^{-n-1}$.  On the intervening annulus,
$E_j=\Phi N_j-R_j(K\Phi)$.  The first term is integrable by
\eqref{eq:Nj-angular}.  Choose $\zeta\in C_c^\infty((1/8,8))$ equal to
one on $[1/4,4]$.  The term $\zeta K\Phi$ is controlled by
\eqref{eq:fixed-annulus-H1} and $R_j:H^1\to L^1$.  The remainder
$(1-\zeta)K\Phi$ is separated from the intervening $x$-annulus and is
absolutely integrable after applying the Riesz kernel.  This proves
$\norm{E_j}_1\le C\mathfrak L(\Omega)$.

Since $\mathcal A_0$ commutes with every $R_j$ and annihilates the
homogeneous distributions $K,N_j$,
\[
 \mathcal A_0E_j=R_j(K(r\psi'(r)))-r\psi'(r)N_j.
\]
The factor $r\psi'(r)$ is supported in the fixed annulus
$1/4\le r\le3/4$.  The first term is controlled by
\eqref{eq:fixed-annulus-H1} and $R_j:H^1\to L^1$, and the second by
\eqref{eq:Nj-angular}.  Together with the bound for $E_j$, this proves
\eqref{eq:Ej-bounds}.

We next choose the ray representatives in a jointly measurable way.  In
logarithmic polar coordinates put
\[
 G_j(s,u)=e^{(n-1)s}E_j(e^su),\qquad
 D_j(s,u)=e^{(n-1)s}(\mathcal A_0-1)E_j(e^su).
\]
Both functions are locally integrable in $(s,u)$, and
$\partial_sG_j=D_j$ in the distributional sense, because, with $t=e^s$,
\[
 \partial_s[t^{n-1}E_j(tu)]
 =t^{n-1}(\mathcal A_0-1)E_j(tu).
\]
Applying Fubini first on each compact $s$-interval gives this
one-dimensional distributional identity for almost every $u$.  Choose
$\vartheta\in C_c^\infty(\R)$ with $\int\vartheta=1$, and, for every
$u$ for which the following integrals are defined, set
\[
 c_j(u)=\int_\R\vartheta(s)
 \left(G_j(s,u)-\int_0^sD_j(v,u)\,dv\right)ds,
\]
\[
 \widetilde G_j(s,u)=c_j(u)+\int_0^sD_j(v,u)\,dv.
\]
Let $\mathcal N_j^{(0)}$ be the exceptional null set of directions on
which either these integrals are not defined or the one-dimensional
distributional identity above fails, and set $\widetilde G_j=0$ on
$\mathcal N_j^{(0)}$.
Fubini and the one-dimensional $W^{1,1}_{\rm loc}$ theorem show that
$\widetilde G_j=G_j$ almost everywhere.  This construction makes
$(s,u)\mapsto\widetilde G_j(s,u)$ jointly measurable and makes
$s\mapsto\widetilde G_j(s,u)$ locally absolutely continuous for almost
every $u$.

Define
\[
 \varphi_{j,u}(t)=\widetilde G_j(\log t,u),\qquad t>0.
\]
The map $(u,t)\mapsto\varphi_{j,u}(t)$ is jointly measurable.  For every
$u$, the map $t\mapsto\varphi_{j,u}(t)$ is continuous (locally absolutely
continuous for every nonexceptional direction and identically zero on the
exceptional set), and hence is a right-continuous representative.  Moreover,
\[
 \varphi'_{j,u}(t)
 =t^{n-2}(\mathcal A_0-1)E_j(tu)
\]
for almost every $(t,u)$.  Polar coordinates give
\[
 \int_{\Sph^{n-1}}B(\varphi_{j,u})\,d\sigma(u)
 =\norm{(\mathcal A_0-1)E_j}_1
 \le C\mathfrak L(\Omega),
\]
and also
\[
 \int_{\Sph^{n-1}}\int_0^\infty
 |\varphi_{j,u}(t)|\,dt\,d\sigma(u)=\norm{E_j}_1.
\]
For almost every $u$, the derivative of $\varphi_{j,u}$ is integrable on
$(1,\infty)$, since
\[
 \int_1^\infty|\varphi'_{j,u}(t)|\,dt
 \le\int_1^\infty
 t^{n-1}|(\mathcal A_0-1)E_j(tu)|\,dt.
\]
Integrating this inequality in $u$ gives
\[
 \int_{\Sph^{n-1}}\int_1^\infty
 |\varphi'_{j,u}(t)|\,dt\,d\sigma(u)
 \le\norm{(\mathcal A_0-1)E_j}_1<\infty.
\]
To make the quantifier in the parameterized profile lemma literal, enlarge
the exceptional set to
\[
\begin{split}
 \mathcal N_j:=\mathcal N_j^{(0)}
 &\cup\{u:B(\varphi_{j,u})=\infty\}
 \cup\left\{u:\int_0^\infty|\varphi_{j,u}(t)|\,dt=\infty\right\}\\
 &\cup\left\{u:\int_1^\infty|\varphi'_{j,u}(t)|\,dt=\infty\right\}.
\end{split}
\]
For $u\notin\mathcal N_j^{(0)}$, local absolute continuity and the displayed
derivative identity give
\[
 B(\varphi_{j,u})
 =\int_0^\infty t^{n-1}
   |(\mathcal A_0-1)E_j(tu)|\,dt,
\]
On $\mathcal N_j^{(0)}$ the profile is zero, so its $B$-value is zero;
extend the displayed integral by zero there as well.
Thus $u\mapsto B(\varphi_{j,u})$ is measurable by joint measurability of
the integrand; the other two defining integrals are measurable for the
same reason.  Tonelli and the three displayed integral bounds now show
that every set in the union is null.  Redefine both
$\widetilde G_j(s,u)$ and
$\varphi_{j,u}(t)$ to be zero for $u\in\mathcal N_j$.  This preserves joint
measurability and changes the functions only on a product-null set.  For
$u\notin\mathcal N_j$, the integrable tail derivative gives a limit at
infinity,
and $L^1(0,\infty)$ integrability forces that limit to be zero; for
$u\in\mathcal N_j$ the same conclusion is immediate from the zero definition.
Thus every direction, not merely almost every direction, satisfies the
profile hypotheses.  The product-null modification does not change the
polar-coordinate convolution identity, so all hypotheses of the
parameterized form of Lemma~\ref{lem:profiles} are satisfied.

For a distribution $F$ write $F^{(\varepsilon)}(x)=
\varepsilon^{-n}F(x/\varepsilon)$.  Define
\[
 T^0_\varepsilon f(x)=\int_{|y|<\varepsilon}
 K(y)\Phi(|y|/\varepsilon)f(x-y)dy.
\]
Since $(K\Phi)^{(\varepsilon)}(x)=K(x)\Phi(|x|/\varepsilon)$,
\begin{equation}\label{eq:exact-smooth-split}
 T^0_{\Omega,\varepsilon}f=(K\Phi)^{(\varepsilon)}*f-T^0_\varepsilon f.
\end{equation}
For the current Schwartz function $f$, put $g_j=-R_jf$, so that
$f=\sum_jR_jg_j$.  The multiplier normalization
\eqref{eq:rough-multiplier-normalization} makes the following calculation
an $L^2$ identity.  Since
$\widehat g_j=-\mathfrak r_j\widehat f$ and
$\mathfrak r_j(\varepsilon\xi)=\mathfrak r_j(\xi)$,
\begin{align*}
 \sum_{j=1}^n
 \widehat{(R_j(K\Phi))^{(\varepsilon)}*g_j}(\xi)
 &=-\widehat{K\Phi}(\varepsilon\xi)
   \sum_{j=1}^n\mathfrak r_j(\xi)^2\widehat f(\xi)\\
 &=\widehat{K\Phi}(\varepsilon\xi)\widehat f(\xi).
\end{align*}
Thus
\[
 (K\Phi)^{(\varepsilon)}*f
 =\sum_{j=1}^n(R_j(K\Phi))^{(\varepsilon)}*g_j
\]
in $L^2$ and in $\mathcal S'$.  Substituting
\eqref{eq:RjKPhi} into this exact identity and then using
\eqref{eq:exact-smooth-split} gives
\begin{equation}\label{eq:exact-Riesz-factorization}
 T^0_{\Omega,\varepsilon}f
 =\sum_{j=1}^n(\Phi N_j)^{(\varepsilon)}*g_j
  -\sum_{j=1}^nE_j^{(\varepsilon)}*g_j-T^0_\varepsilon f.
\end{equation}

Each of the three families is now controlled without a dyadic summation.
First, let $q(t)=\Phi(t)t^{-1}\mathbf1_{(0,1)}(t)$, with the
right-continuous value $q(1)=0$.  Since $\Phi$ vanishes near zero,
\[
 B(q)=\int_0^1t\left|\frac d{dt}\frac{\Phi(t)}t\right|dt
      +|\Phi(1)|<\infty.
\]
Moreover,
\[
 ((\Phi N_j)^{(\varepsilon)}*g)(x)
 =\frac12\int_{\Sph^{n-1}}\nu_j(u)
 \int_{\R}\Phi(|s|/\varepsilon)g(x-su)\,\frac{ds}{s}
 \,d\sigma(u),
\]
while
\[
 \int_{\R}\Phi(|s|/\varepsilon)g(x-su)\frac{ds}{s}
 =H_{\varepsilon,u}g+P^q_{\varepsilon,u}g-P^q_{\varepsilon,-u}g.
\]
Thus \eqref{eq:Nj-angular}, Lemma~\ref{lem:jumpnorm},
Lemma~\ref{lem:one-dimensional-inside-jump}, and
Lemma~\ref{lem:profiles} control
the first sum in
\eqref{eq:exact-Riesz-factorization} by
$C\mathfrak L(\Omega)\sum_j\norm{g_j}_p$, uniformly for
$\varepsilon\in I$.  For the second sum, polar coordinates give the
explicit Bochner representation
\[
 (E_j^{(\varepsilon)}*g)(x)
 =\int_{\Sph^{n-1}}\int_0^\infty
   \varphi_{j,u}(t)g(x-\varepsilon tu)\,dt\,d\sigma(u).
\]
The profiles $\varphi_{j,u}$ and Lemma~\ref{lem:profiles} therefore
control it by the same quantity, uniformly on $I$.  Finally,
\[
 T^0_\varepsilon f(x)=\int_{\Sph^{n-1}}\Omega(u)
 \int_0^1q(t)f(x-\varepsilon tu)\,dt\,d\sigma(u),
\]
which is the parameterized profile family
$\varphi_u(t)=\Omega(u)q(t)$ and satisfies
\[
 \int_{\Sph^{n-1}}B(\varphi_u)d\sigma(u)
 =B(q)\|\Omega\|_1.
\]
Thus Lemmas~\ref{lem:jumpnorm} and~\ref{lem:profiles} give
$C\norm\Omega_1\norm f_p$.  The scalar $L^p$ bounds for the Riesz transforms
\cite{SteinSI} prove the required estimate for the even component.  Combining
it with \eqref{eq:odd-jump} by the subadditive jump seminorm in
Lemma~\ref{lem:jumpnorm}, and then using \eqref{eq:gauge-parity}, proves
\eqref{eq:nonosc-strong} for Schwartz functions, uniformly over every
finite ordered parameter set $I$.

\smallskip
\noindent\emph{Extension in the truncation parameter and in the input.}

All estimates above are uniform over finite parameter sets, and the same
factorization with ordinary $L^p$ norms bounds every fixed coordinate.
Fix a countable dense set $D\subset(0,\infty)$, and choose
$f_m\in\mathcal S(\R^n)$ with $f_m\to f$ in $L^p$ so rapidly that
\[
 \sum_m\|f_{m+1}-f_m\|_p<\infty.
\]
For every
$\varepsilon\in D$, the right-hand side of
\eqref{eq:exact-Riesz-factorization} is Cauchy in $L^p$; denote its limit
by $T^0_{\Omega,\varepsilon}f$.

Let $I\subset D$ be finite.  After passing to a subsequence, the
coordinates for $f_m$ converge almost everywhere simultaneously on $I$.
Lemma~\ref{lem:finite-witness}, in the form
\eqref{eq:finite-chain-J-lsc}, and Fatou's lemma give
\[
 \mathcal J^{\rm ch}_{2,p}
 (T^0_{\Omega,\varepsilon}f:\varepsilon\in I)
 \le C_{n,p}\mathfrak L(\Omega)\norm f_p.
\]
The subsequence may depend on $I$, which is harmless because the bound is
uniform in $I$.  If $I_1\subset I_2\subset\cdots$ are finite and exhaust
$D$, then $\mathbf J^{\rm ch}_{I_k}$ increases pointwise to the chain
functional on $D$: every strict jump witness is finite and is therefore
contained in some $I_k$.  Monotone convergence in the spatial variable
gives the same estimate on all of $D$.

It remains to record compatibility of the chosen coordinates.  If
$\varepsilon<\eta$ lie in $D$, then for every $m$
\begin{equation}\label{eq:finite-annulus-identity}
 T^0_{\Omega,\varepsilon}f_m-T^0_{\Omega,\eta}f_m
 =\int_{\varepsilon<|z|\le\eta}K_\Omega(z)f_m(\,\cdot-z)\,dz.
\end{equation}
The finite-annulus kernel is in $L^1$, with norm
$\norm\Omega_1\log(\eta/\varepsilon)$.  Young's inequality passes this
identity in $L^p$ to the limiting coordinates.  Since $D$ is countable,
one null set may be discarded so that the identity holds simultaneously
for every pair in $D$.

Fix $\varepsilon_0\in D$.  For arbitrary $\varepsilon>0$ define
\[
 T^0_{\Omega,\varepsilon}f=
 \begin{cases}
 T^0_{\Omega,\varepsilon_0}f+
 \displaystyle\int_{\varepsilon<|z|\le\varepsilon_0}
 K_\Omega(z)f(\,\cdot-z)dz,&\varepsilon<\varepsilon_0,\\[1.2ex]
 T^0_{\Omega,\varepsilon_0}f-
 \displaystyle\int_{\varepsilon_0<|z|\le\varepsilon}
 K_\Omega(z)f(\,\cdot-z)dz,&\varepsilon>\varepsilon_0,\\[1.2ex]
 T^0_{\Omega,\varepsilon_0}f,&\varepsilon=\varepsilon_0.
 \end{cases}
\]
The preceding identity leaves all coordinates of $D$ unchanged.  On each
$0<a<b<\infty$ and every bounded measurable $U\subset\R^n$, Fubini and
H\"older's inequality give
\begin{align*}
 &\int_U\int_a^b\int_{\Sph^{n-1}}|\Omega(u)|
 |f(x-ru)|\,d\sigma(u)\frac{dr}{r}dx\\
 &\qquad\le |U|^{1/p'}\|f\|_p\|\Omega\|_1\log(b/a)<\infty.
\end{align*}
Thus, for almost every $x$,
\[
 \int_a^b\int_{\Sph^{n-1}}|\Omega(u)|
 |f(x-ru)|\,d\sigma(u)\,\frac{dr}{r}<\infty.
\]
Taking a countable exhaustion by rational $a,b$ and bounded balls $U$
produces one full-measure set on which this holds on every compact
subinterval of $(0,\infty)$.
The resulting path is locally absolutely continuous, with almost-everywhere
derivative
\[
 -\varepsilon^{-1}\int_{\Sph^{n-1}}
   \Omega(u)f(x-\varepsilon u)d\sigma(u).
\]
Its strict jump counts on $D$ equal those on all positive parameters by
Lemma~\ref{lem:finite-witness}.  Finally use \eqref{eq:chain-pair}.

\smallskip
\noindent\emph{The zero-at-infinity anchor.}

For Schwartz data, the defining tail integral tends to zero pointwise as
$\varepsilon\to\infty$ by absolute convergence and Schwartz decay.
Equation~\eqref{eq:anchor-controls-maximal} therefore gives
\[
 \sup_{\varepsilon>0}|T^0_{\Omega,\varepsilon}f|
 \le\mathbf J(T^0_{\Omega,\varepsilon}f:\varepsilon>0).
\]
The right-hand side belongs to $L^p$ by the estimate already proved, so
dominated convergence also gives the $L^p$ anchor on the Schwartz core.

For arbitrary $f\in L^p$, use the rapidly convergent Schwartz sequence
chosen above and put
\[
 M_m(x)=\sup_{\varepsilon>0}
 |T^0_{\Omega,\varepsilon}(f_{m+1}-f_m)(x)|.
\]
Each difference is a Schwartz function and its trajectory is anchored at
infinity; hence \eqref{eq:anchor-controls-maximal} and the pointwise jump
estimate give
\[
 \sum_m\|M_m\|_p
 \le C_{n,p}\mathfrak L(\Omega)
       \sum_m\|f_{m+1}-f_m\|_p<\infty.
\]
Thus $\sum_mM_m<\infty$ almost everywhere, and the core trajectories
converge uniformly in $\varepsilon>0$ outside one null set.  Their limit
agrees at every parameter in $D$ with the coordinates constructed above,
preserves the finite-annulus identities and local continuity, and is
anchored at infinity: first approximate uniformly by one core trajectory
and then send its truncation radius to infinity.  The same estimate shows
convergence in $L^p$ uniformly over $\varepsilon$, so the anchor also
holds in $L^p$.  Finally use \eqref{eq:chain-pair-functional} to return to
the paired convention.  This completes the proof.
\end{proof}

\section{The polynomial phase flag and scale-covariant annular decay}
\label{sec:flag}

This section develops the coefficient-uniform oscillatory input.  Pure
output and input phases do not affect operator norms or pointwise jump
counts, because
\begin{equation}\label{eq:purephase}
 P(x,y)=p(x)+q(y)+\widetilde P(x,y)
 \quad\Longrightarrow\quad
 T^P_{\Omega,\varepsilon}f=e^{ip}
 T^{\widetilde P}_{\Omega,\varepsilon}(e^{iq}f).
\end{equation}

\subsection{Algebraic flag}

Let $\mathcal P_d$ be the space of real polynomials of degree at most $d$.
Set
\[
 \Delta P(y,u,h)=P(y+u,y+h)-P(y+u,y),\qquad
 \Dop P=\Mix\Delta P,
\]
where $\Mix$ retains precisely the monomials having positive degree in both
$y$ and $u$.

\begin{lemma}[Kernel of the transverse map]\label{lem:kernelD}
For every $P\in\mathcal P_d$, one has $\Dop P=0$ if and only if
\begin{equation}\label{eq:kernelD}
 P(x,y)=p(x)+q(y)+Q(x-y)+x^TBy,
 \qquad B^T=-B.
\end{equation}
Equivalently, for every $a\in\R^n$ there are polynomials $p_a,q_a$ such
that, for all $X,Y\in\R^n$,
\[
 P(a+X,a+Y)-P(X,Y)=p_a(X)+q_a(Y).
\]
\end{lemma}

\begin{proof}
The phases in \eqref{eq:kernelD} are annihilated by $\Dop$.  Conversely,
$\Dop P=0$ means that $\Delta P$ has no monomial of positive degree in
both $y$ and $u$.  Equivalently,
\[
 \partial_{y_j}\partial_{u_k}\Delta P(y,u,h)=0
 \qquad(1\le j,k\le n).
\]
Differentiate this identity in $h_i$ at $h=0$ and put $x=y+u$.
Since
\[
 \left.\partial_{h_i}\Delta P(y,u,h)\right|_{h=0}
 =\partial_{y_i}P(y+u,y),
\]
the chain rule shows that the mixed Hessian
$M_{ki}=\partial_{x_k}\partial_{y_i}P$ satisfies
$(\partial_{x_j}+\partial_{y_j})M_{ki}=0$, so $M(x,y)=F(x-y)$.
Commutation of $x$- and $y$-derivatives gives
\[
 \partial_{z_\ell}F_{ki}=\partial_{z_k}F_{\ell i},\qquad
 \partial_{z_\ell}F_{ki}=\partial_{z_i}F_{k\ell}.
\]
Put $B=(F-F^T)/2$.  The two compatibility identities imply, for every
$i,k,\ell$,
\[
 \partial_{z_\ell}F_{ik}
 =\partial_{z_i}F_{\ell k}
 =\partial_{z_k}F_{\ell i}
 =\partial_{z_\ell}F_{ki}.
\]
Thus every entry of $B$ is constant; in particular, $B$ is a constant
skew-symmetric matrix.  Put $S=F-B=(F+F^T)/2$.
The symmetric potential is produced by the following polynomial Poincar\'e
construction.  For each $k$ the polynomial one-form
$\sum_iS_{ki}(z)\,dz_i$ is closed, and therefore
\[
 v_k(z)=\int_0^1\sum_iS_{ki}(tz)z_i\,dt
\]
satisfies $\partial_i v_k=S_{ki}$.  Symmetry of $S$ makes
$\sum_kv_k(z)\,dz_k$ closed, so
\[
 Q(z)=-\int_0^1\sum_kv_k(tz)z_k\,dt
\]
is a polynomial with $-\partial_k\partial_iQ=S_{ki}$.  Subtracting
$Q(x-y)+x^TBy$ leaves zero mixed Hessian, because
\[
 \partial_{x_k}\partial_{y_i}
 \bigl(Q(x-y)+x^TBy\bigr)
 =-\partial_k\partial_iQ(x-y)+B_{ki}
 =S_{ki}(x-y)+B_{ki}=F_{ki}(x-y).
\]
A polynomial with zero mixed Hessian is a sum of a pure $x$-polynomial
and a pure $y$-polynomial, which proves the converse.  The resulting
decomposition plainly has the stated translation property for every $a$.
Conversely, that property implies, after taking the mixed Hessian, that
$M(x+a,y+a)=M(x,y)$ for every $a$; hence $M(x,y)=F(x-y)$, and the
compatibility and polynomial Poincar\'e argument just given recovers
\eqref{eq:kernelD}, and therefore $\Dop P=0$.
\end{proof}

Order the mixed monomials in $(y,u,h)$ first by decreasing $y$-degree,
then by decreasing $u$-degree, decreasing $h$-degree, decreasing total
degree, and finally by a fixed lexicographic order.  Columns are listed in
this order; throughout, the pivot of a nonzero row means its first nonzero
column, normalized to have coefficient one, and rows are listed in the
order of their pivots.  The substitutions defining $\Delta$ are linear and
contain no constant term, and $\Mix$ only deletes monomials.  Consequently
$\Dop$ preserves total homogeneous degree and
$\operatorname{im}\Dop$ is the direct sum of its homogeneous components.
In each such component take reduced row-echelon form with the preceding
column and pivot convention.
Merge the rows from the separate homogeneous components in the global
order of their pivot monomials, and call them $F_1,\ldots,F_M$.
Here homogeneous means homogeneous in the total degree of $(y,u,h)$, and
the global monomial order restricts to each homogeneous component.  To
check the echelon property after merging, let
$\pi_\nu\prec\pi_\mu$ be two pivot monomials.  Every monomial in $F_\mu$
has the same total degree as $\pi_\mu$.  If one of them, say $\tau$, also
satisfied $\tau\prec\pi_\nu$, then $\tau\prec\pi_\mu$ and the reduced
row-echelon construction in the component containing $F_\mu$ would force
its coefficient to vanish.  If the two rows have the same total degree,
reduced row-echelon form also clears the $\pi_\nu$ column from $F_\mu$;
if their total degrees differ, $F_\mu$ cannot contain $\pi_\nu$ at all.
Thus no row later than $F_\nu$ contains a monomial preceding or equal to
the pivot of $F_\nu$.  We record this global echelon property for later
use: if $\pi_\nu$ is the pivot of $F_\nu$, then
\begin{equation}\label{eq:global-echelon-property}
 \mu>\nu
 \quad\Longrightarrow\quad
 \operatorname{supp}F_\mu
 \cap\{\tau:\tau\preceq\pi_\nu\}=\varnothing.
\end{equation}
This includes rows of a different total degree, since a row and all of
its monomials lie in the same homogeneous component.  Homogeneous preimages
may be chosen without an extra assumption: if $F_\mu$ has degree $r$ and
$\Dop H=F_\mu$, write $H=\sum_kH^{[k]}$ as a sum of homogeneous parts.
The graded property gives $\Dop H^{[r]}=F_\mu$ and
$\Dop H^{[k]}=0$ for $k\ne r$.  We therefore take
$H_\mu=H^{[r]}$, so that $\Dop H_\mu=F_\mu$ and $H_\mu$ is homogeneous.
Replacing $H_\mu(x,y)$ by
$H_\mu(x,y)-H_\mu(y,y)$, which changes only a pure input phase, arrange
\begin{equation}\label{eq:Hdiag}
 H_\mu(y,y)=0.
\end{equation}
Append to this list, in decreasing homogeneous degree and fixed monomial
order, the convolution monomials $(x-y)^\alpha$ with
$2\le|\alpha|\le d$, and, when $d\ge2$, a fixed basis
$x_i y_j-x_jy_i$ of skew bilinear forms.  Denote the
full finite list again by $H_1,\ldots,H_{N_d}$, and write
$m_\mu=\deg H_\mu$.
Modulo pure phases every $P\in\mathcal P_d$ has a unique first nonzero flag
coordinate $c_\nu$.

\begin{theorem}[Global polynomial phase flag]\label{thm:phaseflag}
For each $n\in\mathbb N$ and $d\in\mathbb N_0$ there is a finite ordered
family of homogeneous real
polynomials $H_1,\ldots,H_{N_d}$ with the following properties.
\begin{enumerate}[label=\textup{(\roman*)}]
\item Every $P\in\mathcal P_d$ has a unique representation modulo pure
input and output phases,
\begin{equation}\label{eq:flagrepresentation}
 P(x,y)=p(x)+q(y)+\sum_{\mu=1}^{N_d}c_\mu(P)H_\mu(x,y).
\end{equation}
Thus every nonpure phase has a well-defined first nonzero flag coordinate.
\item The coordinates lying in $\ker\Dop$ occur after the transverse
coordinates and, modulo pure phases, span exactly
\[
 Q(x-y)+x^TBy,\qquad B^T=-B.
\]
In particular, the same flag incorporates convolution and
alternating-bilinear phases as its final nonterminal strata.
\item If $\nu$ is the first nonzero coordinate of $P$ and
$m_\nu=\deg H_\nu$, then, for every dilation factor $\varrho>0$, the
simultaneous dilation $(x,y)=(\varrho X,\varrho Y)$ changes that coordinate
to $c_\nu\varrho^{m_\nu}$.  For every $a\in\R^n$,
\begin{equation}\label{eq:triangular}
 P(a+X,a+Y)=p_a(X)+q_a(Y)+c_\nu H_\nu(X,Y)
 +\sum_{\mu>\nu}c_{\mu,a}H_\mu(X,Y).
\end{equation}
Thus translation preserves the first coordinate and creates only later
coordinates.
\item Suppose that $H_\nu$ is transverse and that
$y^\alpha u^\beta h^\gamma$ is the pivot monomial of
$\Dop H_\nu$ in the reduced-row-echelon construction above.  If $P$ has
first coordinate $c_\nu=\pm1$, then the coefficient of
$y^\alpha u^\beta$ in $\Dop P$ is a polynomial $a_P(h)$ of degree
$|\gamma|$.  Its $h^\gamma$ coefficient is exactly $c_\nu=\pm1$, and its
top homogeneous coefficient tensor has norm at least $c_{n,d}>0$.
This conclusion is uniform in every later flag coefficient.
\end{enumerate}
\end{theorem}

\begin{proof}
Lemma~\ref{lem:kernelD} gives spanning; we record explicitly why the
resulting coordinates are unique modulo pure phases.  The reduced
row-echelon rows $F_1,\ldots,F_M$ are linearly independent.  Hence, if a
linear combination of their preimages together with the appended kernel
elements is pure, applying $\Dop$ first forces every transverse coefficient
to vanish.  It remains to consider an identity
\[
 Q(x-y)+x^TBy=p(x)+q(y),\qquad B^T=-B,
\]
where $Q$ is a linear combination of the appended convolution monomials and
therefore has no terms of degree at most one.  Taking the mixed Hessian in
$(x,y)$ gives
\[
 -\nabla^2Q(x-y)+B=0.
\]
The first matrix is symmetric and the second is skew-symmetric, so both
vanish.  Thus $B=0$ and $Q$ is affine; by the chosen degree range for $Q$,
this means $Q=0$.  This proves linear independence of the appended kernel
basis in the quotient and, together with Lemma~\ref{lem:kernelD}, proves the
representation and the description of the kernel strata.  Homogeneity gives
the dilation statement.  For a transverse row,
\[
 \Dop(P(a+\cdot,a+\cdot))=
 \Mix[(\Dop P)(y+a,u,h)].
\]
Translation preserves the highest $y$-degree term and creates only lower
$y$-degree terms, which are later in the chosen order.  Reduced echelon
form, specifically \eqref{eq:global-echelon-property}, prevents a later
row from altering the pivot coefficient.  A
convolution monomial is unchanged by simultaneous translation, and a skew
bilinear form changes only by pure input and output phases.  This proves
(iii).

For (iv), project $\Dop P$ onto the scalar tensor component indexed by
$y^\alpha u^\beta$.  The pivot row contributes $\pm h^\gamma$.  No later
row contains the pivot monomial, by
\eqref{eq:global-echelon-property}, and no later
row contains a monomial with the same $y$- and $u$-component and higher
$h$-order, because such a monomial would precede the pivot in the defining
order.  Thus the $h^\gamma$ coefficient remains $\pm1$, and later
coefficients can change only terms of lower $h$-degree or other
multiindices of the same $h$-degree.  In particular, even in the latter
case the fixed $h^\gamma$ coordinate of the top homogeneous coefficient
tensor remains $\pm1$.  The coordinate
functional selecting $h^\gamma$ has norm at most $C_{n,d}$ on each of the
finitely many fixed coefficient spaces that occur.  Consequently $a_P$
has degree $|\gamma|$, and the norm of its top homogeneous coefficient
tensor is at least $c_{n,d}>0$, uniformly in all later coefficients.
\end{proof}

For packet arguments it is convenient to use anchored representatives
\begin{equation}\label{eq:anchoredbasis}
 \widehat H_\mu(X,Y)=H_\mu(X,Y)-H_\mu(0,Y).
\end{equation}
They have the same flag coordinates modulo pure input phases and satisfy
$\widehat H_\mu(0,Y)=0$.

\subsection{Polynomial oscillation estimates}

We give the elementary estimates needed for rough annular decay.

\begin{lemma}[Sublevel and negative moments]\label{lem:sublevel}
Let $k\ge1$, let $p(t)=a_kt^k+\cdots+a_0$, $a_k\ne0$, on an interval
$I$ of length $L>0$, and let $\eta>0$.  Then
\[
 |\{t\in I:|p(t)|\le\eta\}|\le C_k(\eta/|a_k|)^{1/k}.
\]
Consequently, if $0<\sigma<1/k$, then, uniformly in $c$ and the lower
coefficients,
\begin{equation}\label{eq:neg1d}
 \int_I|p(t)-c|^{-\sigma}dt
 \le C_{k,\sigma}|a_k|^{-\sigma}L^{1-k\sigma}.
\end{equation}
If $a$ is a polynomial of degree $k$ on $\R^n$ whose top homogeneous
coefficient norm is at least one, then, for every $R>0$ and every fixed
$C_0<\infty$,
\begin{equation}\label{eq:negn}
 \int_{|h|\le C_0R}|a(h)|^{-\sigma}dh\le C R^{n-k\sigma},
 \qquad 0<\sigma<1/k.
\end{equation}
Here the last constant may depend on $C_0$.  If $A_k$ is homogeneous of
degree $k$ and $\norm{A_k}_{\rm coeff}=1$, then
\begin{equation}\label{eq:negsphere}
 \int_{\Sph^{n-1}}|A_k(\theta)|^{-\beta}d\sigma(\theta)\le C,
 \qquad 0<\beta<1/k.
\end{equation}
\end{lemma}

\begin{proof}
Let $E=\{t\in I:|p(t)|\le\eta\}$ and suppose $|E|>0$.
Set $\delta=|E|/[2(k+1)]$.  A greedy selection gives $k+1$ distinct
points of $E$ at mutual distance at least $\delta$: after $j\le k$
points have been selected, the union of their open $\delta$-neighborhoods
has measure at most $2j\delta<|E|$ and therefore cannot cover $E$.
After relabeling, we have $t_0<\cdots<t_k$ and
\[
 |t_i-t_j|\ge\delta=c_k|E|
 \qquad(0\le i,j\le k,\ i\ne j).
\]
Lagrange interpolation at these points gives
\[
 |a_k|
 =\left|\sum_{j=0}^k
   \frac{p(t_j)}{\prod_{\ell\ne j}(t_j-t_\ell)}\right|
 \le C_k\eta |E|^{-k}.
\]
This proves the sublevel estimate.

Apply it to $p-c$ and combine it with the trivial bound by $L$:
\[
 |\{t\in I:|p(t)-c|\le\eta\}|
 \le\min\left\{L,C_k(\eta/|a_k|)^{1/k}\right\}.
\]
The layer-cake identity gives
\[
 \int_I|p(t)-c|^{-\sigma}dt
 =\sigma\int_0^\infty
  \eta^{-\sigma-1}
  |\{t\in I:|p(t)-c|<\eta\}|\,d\eta.
\]
Split the last integral where the two bounds in the minimum agree,
at a value $\eta_0\asymp_k |a_k|L^k$.  The part below $\eta_0$ is
bounded by
\[
 C_k|a_k|^{-1/k}
 \int_0^{\eta_0}\eta^{1/k-\sigma-1}\,d\eta,
\]
and the part above it by
$L\int_{\eta_0}^{\infty}\eta^{-\sigma-1}\,d\eta$.
Both are at most
$C_{k,\sigma}|a_k|^{-\sigma}L^{1-k\sigma}$ precisely when
$0<\sigma<1/k$, proving \eqref{eq:neg1d}.

For \eqref{eq:negn}, write $a=A_k+\text{lower-order terms}$.
Finite-dimensional norm equivalence supplies
$v\in\Sph^{n-1}$ such that
\[
 |A_k(v)|\ge c_{n,k}\norm{A_k}_{\rm coeff}\ge c_{n,k}.
\]
Write $h=w+tv$ with $w\in v^\perp$.  For fixed $w$, the polynomial
$t\mapsto a(w+tv)$ has degree $k$ and leading coefficient $A_k(v)$,
independent of $w$.  Its intersection with $|h|\le C_0R$ is an interval
of length at most $2C_0R$.  Hence \eqref{eq:neg1d}, followed by integration
over the $w$-projection of that ball, gives
\[
 \int_{|h|\le C_0R}|a(h)|^{-\sigma}dh
 \le C R^{1-k\sigma}R^{n-1}
 =C R^{n-k\sigma}.
\]
Finally apply this estimate with $a=A_k$ and $R=2$ to the annulus
$1\le|x|\le2$.  Homogeneity and polar coordinates give
\[
 \int_{1\le|x|\le2}|A_k(x)|^{-\beta}dx
 =\left(\int_1^2r^{n-1-k\beta}\,dr\right)
   \int_{\Sph^{n-1}}|A_k(\theta)|^{-\beta}d\sigma(\theta).
\]
The radial factor is positive and finite, so \eqref{eq:negsphere} follows.
\end{proof}

\begin{lemma}[Leading-coefficient van der Corput]\label{lem:vdc}
Let $k\ge1$, let $I$ be an interval of length at most $L_0<\infty$, and
let
\[
 \phi(t)=a_kt^k+\cdots+a_0,\qquad a_k\ne0.
\]
For every bounded-variation amplitude $\psi$ supported in $I$,
\[
 \left|\int e^{i\phi(t)}\psi(t)dt\right|
 \le C_{k,L_0}\min(1,|a_k|^{-1/k})
 \bigl(\norm\psi_\infty+\Var(\psi)\bigr).
\]
The constant is independent of all lower coefficients.
\end{lemma}

\begin{proof}
Extend $\psi$ by zero off $I$.  Its distributional derivative is a finite
measure and
\[
 |D\psi|(\mathbb R)\le \Var_I(\psi)+2\norm\psi_\infty.
\]
For $k=1$, the trivial estimate applies when $|a_1|<1$, while Stieltjes
integration by parts applies when $|a_1|\ge1$ and gives the asserted bound.
Suppose $k\ge2$.  The case $|a_k|<1$ follows from the trivial estimate
$L_0\norm\psi_\infty$, so assume $|a_k|\ge1$ and let $\mu>0$.  Put
$E=\{t\in I:|\phi'(t)|\le\mu\}$.  Lemma~\ref{lem:sublevel}, applied to
$\phi'$, gives
\[
 |E|\le C_k(\mu/|a_k|)^{1/(k-1)}.
\]
The boundary of $E$ is contained in the roots of
$\phi'=\mu$ and $\phi'=-\mu$ together with the endpoints of $I$.
Consequently, after also splitting at the at most $k-2$ roots of $\phi''$,
$I\setminus E$ is the union of $O_k(1)$ intervals $J$ on each of which
$|\phi'|\ge\mu$ and $\phi'$ is monotone.  On every such interval, with
one-sided endpoint values (equivalently, on half-open intervals so that jumps
are counted only once), Stieltjes integration by parts yields
\[
 \int_J e^{i\phi(t)}\psi(t)\,dt
 =\left[\frac{e^{i\phi(t)}\psi(t)}{i\phi'(t)}\right]_{\partial J}
 -\int_J\frac{e^{i\phi(t)}}{i\phi'(t)}\,d\psi(t)
 +\int_J\frac{e^{i\phi(t)}\psi(t)\phi''(t)}
                  {i\phi'(t)^2}\,dt.
\]
Since $\phi'$ is monotone on $J$,
\[
 \int_J\frac{|\phi''(t)|}{|\phi'(t)|^2}\,dt\le \frac{2}{\mu}.
\]
Summing the endpoint, Stieltjes, and absolutely continuous terms over the
$O_k(1)$ intervals, and estimating the integral over $E$ trivially, gives
\[
 \left|\int e^{i\phi(t)}\psi(t)\,dt\right|
 \le C_{k,L_0}\left\{
  (\mu/|a_k|)^{1/(k-1)}\norm\psi_\infty
  +\mu^{-1}\bigl(\norm\psi_\infty+\Var_I(\psi)\bigr)
 \right\}.
\]
Taking $\mu=|a_k|^{1/k}$ proves the result, uniformly in all lower
coefficients.
\end{proof}

\begin{lemma}[Compact polynomial operator]\label{lem:compactop}
Let $R>0$, let $Q,U$ be cubes of diameter at most $C_0R$, and let
$\Phi(y,u)$ be a real
polynomial.  Suppose the largest mixed $y$-degree is $m\ge1$, and among
those terms the largest $u$-degree is $\ell\ge1$.  If the Euclidean norm
$M$ of the coefficient tensor of bidegree $(m,\ell)$ is nonzero, then
\begin{equation}\label{eq:compactop}
 \left\|1_Q(y)\int_Ue^{i\Phi(y,u)}g(u)du\right\|_2
 \le C R^n\min\{1,(R^{m+\ell}M)^{-\epsilon_{m,\ell}}\}\norm g_2
\end{equation}
for every $g\in L^2(U)$, where $\epsilon_{m,\ell}>0$ and the constant
may depend on $C_0$ but is independent of all pure
terms and all lower mixed coefficients.
\end{lemma}

\begin{proof}
Remove pure terms by unitary multiplication.  Let
$\Phi_{m,\ell}(y,u)$ be the component homogeneous of degree $m$ in $y$
and degree $\ell$ in $u$.  On this fixed finite-dimensional space,
equivalence of the coefficient norm with the supremum norm on
$\Sph^{n-1}\times\Sph^{n-1}$ gives unit vectors $v,w$ such that
\[
 |\Phi_{m,\ell}(v,w)|\ge c_{n,m,\ell}M.
\]
Slice $y=y_\perp+sv$, $u=u_\perp+tw$.  Direct expansion shows that
\begin{equation}\label{eq:compact-leading-contraction}
 [s^mt^\ell]\,
 \Phi(y_\perp+sv,u_\perp+tw)=\Phi_{m,\ell}(v,w).
\end{equation}
In particular, this coefficient is independent of the transverse
variables and has modulus comparable to $M$.  We give the fiber
calculation, including the $T^*T$ square root that determines the exponent.

Fix $y_\perp,u_\perp$ and denote the resulting operator from the $t$-line
to the $s$-line by $T_{y_\perp,u_\perp}$.  Cross-sections of $Q$ and $U$
are intervals of length at most $CR$; their characteristic functions have
variation at most two.  Write $p_m(t)$ for the coefficient of $s^m$ in
the sliced phase.  It is a polynomial of degree $\ell$ whose leading
coefficient has modulus at least $cM$.  The kernel of
$T_{y_\perp,u_\perp}^*T_{y_\perp,u_\perp}$ is
\[
 \mathbf1_{I_U}(t)\mathbf1_{I_U}(t')
 \int_{I_Q}e^{i[\Phi(s,t)-\Phi(s,t')]}\,ds.
\]
After $s=R\tau$, Lemma~\ref{lem:vdc}, including its bounded-variation
amplitude form, gives
\begin{equation}\label{eq:compact-fiber-kernel}
 |K(t,t')|\le CR\min\left\{1,
  \bigl(R^m|p_m(t)-p_m(t')|\bigr)^{-1/m}\right\}.
\end{equation}
Choose
\[
 0<\delta<\min\{1/m,1/\ell\}.
\]
Since $\min(1,X^{-1/m})\le X^{-\delta}$, the negative-moment
estimate \eqref{eq:neg1d}, applied to the polynomial in $t'$ for fixed
$t$, yields the required uniform bound: indeed,
$t'\mapsto p_m(t)-p_m(t')$ has degree $\ell$ and leading coefficient
$-\Phi_{m,\ell}(v,w)$, independent of $t,y_\perp,u_\perp$.  Therefore
\begin{align*}
 \sup_t\int|K(t,t')|\,dt'
 &\le C R^{1-m\delta}M^{-\delta}R^{1-\ell\delta}\\
 &=CR^2(R^{m+\ell}M)^{-\delta}.
\end{align*}
The same estimate with $t,t'$ interchanged follows by symmetry.  The
trivial kernel estimate gives $CR^2$ instead, so Schur's test and then the
$T^*T$ identity imply
\begin{equation}\label{eq:compact-slice}
 \norm{T_{y_\perp,u_\perp}}_{2\to2}
 \le CR\min\{1,(R^{m+\ell}M)^{-\delta/2}\}.
\end{equation}

Reassemble the slices, without invoking a direct-integral identification.
Minkowski's inequality in $u_\perp$ and \eqref{eq:compact-slice} give,
for fixed $y_\perp$,
\[
 \norm{Tg(y_\perp,\cdot)}_{L_s^2}
 \le B_R\int_{\pi_{w^\perp}U}
       \norm{g(u_\perp,\cdot)}_{L_t^2}\,du_\perp,
 \quad
 B_R=CR\min\{1,(R^{m+\ell}M)^{-\delta/2}\}.
\]
Cauchy--Schwarz in $u_\perp$ makes the right side at most
\[
 C B_R R^{(n-1)/2}
 \left(\int_{\pi_{w^\perp}U}
   \norm{g(u_\perp,\cdot)}_{L_t^2}^2\,du_\perp\right)^{1/2}
 =C B_RR^{(n-1)/2}\norm g_2.
\]
The $y_\perp$-projection of $Q$ has measure at most $CR^{n-1}$.
Squaring the preceding estimate, integrating in $y_\perp$, and taking a
square root therefore give
\[
 \norm{Tg}_2\le CR^{n-1}B_R\norm g_2,
\]
which is \eqref{eq:compactop} with
$\epsilon_{m,\ell}=\delta/2$.  Every bound used only the displayed
leading coefficient, so all lower mixed coefficients remain unrestricted.
\end{proof}

For $R>0$ and $1/2\le u\le1$, put
\begin{equation}\label{eq:partialannulus}
 A^{P,\Theta}_{R,u}f(x)=
 \int_{uR<|x-y|<R}e^{iP(x,y)}K_\Theta(x-y)f(y)dy.
\end{equation}

\begin{proposition}[Normalized rough annular decay]\label{prop:annulardecay}
Let $P\in\mathcal P_d$.  For every flag position $\mu$ there are
$a_\mu>0$ and $C_{n,d}$ such that,
if the first nonzero flag coordinate of $P$ is $\pm1$ at position $\mu$,
then, for $R\ge1$ and $\Theta\in L^\infty(\Sph^{n-1})$,
\begin{equation}\label{eq:annulardecay}
 \sup_{1/2\le u\le1}\norm{A^{P,\Theta}_{R,u}}_{2\to2}
 \le C_{n,d}R^{-a_\mu}\norm\Theta_\infty.
\end{equation}
Every later coefficient is unrestricted.
\end{proposition}

\begin{proof}
We treat the three parts of the flag.

For a transverse pivot, partition the input into cubes of side comparable
to $R$; the corresponding output enlargements have bounded overlap.  On
one input cube $Q$, write $A_Q=A_{R,u}\mathbf1_Q$.  Its output is supported
in $X_Q=Q+B(0,R)$, and the family $\{X_Q\}$ has bounded overlap.  The
kernel of $A_Q^*A_Q$ is
\[
 L_Q(y,y+h)=\mathbf1_Q(y)\mathbf1_Q(y+h)
 \int e^{i[P(x,y+h)-P(x,y)]}
 \overline{k_{R,u}(x-y)}k_{R,u}(x-y-h)dx,
\]
where $k_{R,u}=K_\Theta\mathbf1_{\{uR<|\cdot|<R\}}$.  Put $x=y+v$ and
remove the pure $y$- and $v$-terms in the phase difference.  Let $\nu$ be
the current transverse flag position.  Since its normalized coefficient is
$\pm1$ and all earlier coordinates vanish,
\begin{equation}\label{eq:transverse-pivot-expansion}
 \Dop P(y,v,h)=\pm F_\nu(y,v,h)
  +\sum_{\nu<\mu\le M}c_\mu F_\mu(y,v,h).
\end{equation}
Let $y^\alpha v^\beta h^\gamma$ be the pivot monomial of $F_\nu$, and put
\[
 m=|\alpha|,\qquad \ell=|\beta|,\qquad k=|\gamma|.
\]
Because $\Mix$ retains only terms of positive $y$- and $v$-degree, while
$\Delta P(y,v,0)=0$, one has $m,\ell,k\ge1$.  Thus every use below of the
compact-operator and negative-moment lemmas lies in their stated
positive-degree range.
Write $\mathsf T_{m,\ell}(h)$ for the coefficient tensor of bidegree
$(m,\ell)$ in $(y,v)$ in the mixed part of
$P(y+v,y+h)-P(y+v,y)$, and let $a(h)$ be its
$(\alpha,\beta)$-component.  By Theorem~\ref{thm:phaseflag}(iv), $a$ has
degree $k$, the coefficient of $h^\gamma$ is exactly $\pm1$, and no later
flag coefficient can change that coefficient or introduce a term of higher
$h$-degree in the same component.  Consequently
\begin{equation}\label{eq:transverse-tensor-lower-bound}
 \norm{\mathsf T_{m,\ell}(h)}_{\mathrm{coeff}}
 \ge c_{n,d}|a(h)|,
\end{equation}
and no mixed bidegree preceding $(m,\ell)$ occurs.  Here is the precise
link between the flag order and the hypothesis of
Lemma~\ref{lem:compactop}.  A mixed monomial with $y$-degree greater than
$m$, or with $y$-degree $m$ and $v$-degree greater than $\ell$, precedes
the pivot regardless of its $h$-degree.  Such a monomial is absent from
$F_\nu$ because its pivot is the first nonzero column, absent from every
later transverse row by \eqref{eq:global-echelon-property}, and absent
from the kernel strata because $\Dop$ vanishes there.  All earlier flag
coordinates are zero by hypothesis.  Thus $(m,\ell)$ is exactly the
leading bidegree to which Lemma~\ref{lem:compactop} applies, uniformly in
all unrestricted later coefficients.
Notice that this argument detects a fixed component of the full
$(m,\ell)$ tensor; it never replaces a nonzero tensor by its diagonal
contraction.  Such a replacement would be false in general (the
alternating-bilinear contraction vanishes identically), which is why the
convolution and skew kernel strata are treated separately below.
For fixed $h$, define
\[
 (E_hG)(y)=\mathbf1_Q(y)\mathbf1_Q(y+h)
 \int_{|v|\le R}e^{i[P(y+v,y+h)-P(y+v,y)]}G(v)\,dv.
\]
To apply Lemma~\ref{lem:compactop}, extend
$G\mathbf1_{\{|v|\le R\}}$ by zero to a cube of diameter $O(R)$; this
leaves the operator unchanged.  Put
\[
 g_h(v)=\overline{k_{R,u}(v)}k_{R,u}(v-h).
\]
Let $E_h^{(0)}$ denote the operator in the preceding display with the
factor $\mathbf1_Q(y+h)$ omitted.  Since
\[
 E_h=M_{\mathbf1_{Q-h}}E_h^{(0)},
 \qquad \norm{E_h}_{2\to2}\le\norm{E_h^{(0)}}_{2\to2},
\]
the additional output cutoff is harmless.  Moreover
$L_Q(y,y+h)=E_hg_h(y)$.  If $a(h)=0$, Cauchy--Schwarz in the input and
output cubes gives the trivial estimate
$\norm{E_h}_{2\to2}\leq CR^n$, which is the desired bound below with the
minimum interpreted as $1$.  If $a(h)\ne0$, then
\eqref{eq:transverse-tensor-lower-bound} makes the leading tensor nonzero,
so Lemma~\ref{lem:compactop} applies to $E_h^{(0)}$.  Absorbing the fixed
factor $c_{n,d}$ into the structural constant, the two cases give
\[
 \norm{E_h}_{2\to2}\le CR^n
 \min\{1,(R^{m+\ell}|a(h)|)^{-\epsilon}\}.
\]
The actual input $g_h$ satisfies, uniformly in $h$ and in $u\in[1/2,1]$,
\begin{equation}\label{eq:gh-L2}
 \norm{g_h}_2\le\norm{k_{R,u}}_\infty\norm{k_{R,u}}_2
 \le CR^{-3n/2}\norm\Theta_\infty^2.
\end{equation}
Consequently
\[
 \int|L_Q(y,y+h)|^2dy
 \le CR^{-n}\norm\Theta_\infty^4
 \min\{1,(R^{m+\ell}|a(h)|)^{-2\epsilon}\}.
\]
Choose $0<\sigma<\min(2\epsilon,1/k)$ and integrate in $|h|\lesssim R$.
The top homogeneous coefficient tensor of $a$ has norm at least
$c_{n,d}$ by Theorem~\ref{thm:phaseflag}(iv); rescaling $a$ by this fixed
constant puts it in the normalization of \eqref{eq:negn}.  Thus, using
$\min(1,X^{-2\epsilon})\le X^{-\sigma}$ and absorbing the resulting
factor $c_{n,d}^{-\sigma}$, \eqref{eq:negn} gives
\[
\begin{split}
 &R^{-n}R^{-\sigma(m+\ell)}
   \int_{|h|\lesssim R}|a(h)|^{-\sigma}\,dh\\
 &\hspace{3em}\le
 C R^{-n}R^{-\sigma(m+\ell)}R^{n-k\sigma}
 =C R^{-\sigma(m+\ell+k)}.
\end{split}
\]
Consequently,
\begin{equation}\label{eq:annular-HS-square}
 \iint|L_Q(y,y+h)|^2\,dy\,dh
 \le C\norm\Theta_\infty^4R^{-\sigma(m+\ell+k)}.
\end{equation}
The change of variables $(y,h)\mapsto(y,y+h)$ has unit Jacobian, so the
left side is exactly
$\norm{A_Q^*A_Q}_{\rm HS}^2$.  Taking its square root first gives
\[
 \norm{A_Q^*A_Q}_{2\to2}
 \le \norm{A_Q^*A_Q}_{\rm HS}
 \le C\norm\Theta_\infty^2R^{-\sigma(m+\ell+k)/2},
\]
and the $A_Q^*A_Q$ identity supplies the second square root:
\[
 \norm{A_Q}_{2\to2}
 \le C\norm\Theta_\infty R^{-\sigma(m+\ell+k)/4}.
\]
Finally, bounded overlap of the $X_Q$ and disjointness of the input cubes
give, with $a=\sigma(m+\ell+k)/4$,
\[
\begin{split}
 \left\|\sum_QA_Q(f\mathbf1_Q)\right\|_2^2
 &\le C\sum_Q\norm{A_Q(f\mathbf1_Q)}_2^2\\
 &\le C R^{-2a}\norm\Theta_\infty^2
       \sum_Q\norm{f\mathbf1_Q}_2^2
 =C R^{-2a}\norm\Theta_\infty^2\norm f_2^2.
\end{split}
\]
This reassembles the cubes and proves
\eqref{eq:annulardecay} with a positive exponent independent of the later
coefficients.

For a convolution pivot, after pure phases are removed write
$P(x,x-z)=Q(z)-x^TBz$.  Let $Q_s$ be the first nonzero homogeneous
convolution component.  Its pivot monomial has coefficient $\pm1$, and
later convolution basis vectors do not change that coordinate; boundedness
of the corresponding coordinate functional therefore gives
$\norm{Q_s}_{\rm coeff}\ge c_{n,d}>0$.  Here $Q_s$ includes every later
convolution coordinate of the same homogeneous degree $s$; those
coordinates may change the other entries of its coefficient tensor but
cannot remove the normalized pivot entry.  Only components of degree
strictly below $s$ become lower coefficients in the one-dimensional
van der Corput argument.  Polar
coordinates give
\[
 A_{R,u}f=\int_{\Sph^{n-1}}\Theta(\theta)
 \int_{uR}^R e^{i[Q(r\theta)-rx^TB\theta]}f(x-r\theta)\frac{dr}{r}
 d\sigma(\theta).
\]
On lines parallel to $\theta$ this is a one-dimensional convolution
multiplier; the skew modulation depends only on the transverse line since
$\theta^TB\theta=0$.  More explicitly, write
$x=x_\perp+t\theta$, with $x_\perp\in\theta^\perp$, and use the Fourier
convention $\widehat g(\xi)=\int e^{-it\xi}g(t)\,dt$ on that line.  The
fiber operator has multiplier
\begin{equation}\label{eq:convolution-fiber-multiplier}
 m_{\theta,x_\perp,u}(\xi)
 =\int_{uR}^{R}
  e^{i[Q(r\theta)-r x_\perp^TB\theta-r\xi]}
  \frac{dr}{r}.
\end{equation}
Indeed, $(x_\perp+t\theta)^TB\theta=x_\perp^TB\theta$ because
$\theta^TB\theta=0$.  After the change $r=R\rho$, both the Fourier
frequency and the transverse skew modulation therefore enter only the
coefficient of the linear term in $\rho$; neither can alter the leading
coefficient.  The
amplitude
\[
 \psi_u(\rho)=\rho^{-1}\mathbf1_{[u,1]}(\rho)
\]
satisfies $\norm{\psi_u}_\infty+\Var(\psi_u)\le C$ uniformly for
$1/2\le u\le1$.  The leading coefficient of the resulting polynomial
phase is $R^sQ_s(\theta)$.  Lemma~\ref{lem:vdc} and Plancherel therefore
give the fiber norm
\[
 C\min\{1,(R^s|Q_s(\theta)|)^{-1/s}\}.
\]
Choose $0<\beta<1/s$.  The last display is bounded by
$CR^{-s\beta}|Q_s(\theta)|^{-\beta}$.  Minkowski in $\theta$ and
\eqref{eq:negsphere} therefore give the required angular integral
explicitly: writing $Q_s=M_sA_s$ with $\norm{A_s}_{\rm coeff}=1$ and
$M_s\ge c_{n,d}$,
\[
 \int_{\Sph^{n-1}}|Q_s(\theta)|^{-\beta}\,d\sigma(\theta)
 =M_s^{-\beta}\int_{\Sph^{n-1}}|A_s(\theta)|^{-\beta}\,d\sigma(\theta)
 \le C_{n,s,\beta}.
\]
Thus Minkowski gives
$CR^{-s\beta}\norm\Theta_\infty$.  All lower homogeneous components of
$Q$ occur only as lower coefficients in van der Corput and are unrestricted.

For a skew pivot, all transverse and convolution coordinates vanish,
because they precede the skew coordinates in the flag.  Thus, modulo pure
phases, the remaining phase is $x^TBy$ with $B^T=-B$.  The normalized
pivot coefficient gives $\norm B\ge c_{n,d}>0$.  Choose an orthogonal
matrix $O$ that puts $B$ in real skew-normal form.  The change of variables
$x=O\widetilde x$, $y=O\widetilde y$ replaces $\Theta$ by
$\Theta_O(\theta)=\Theta(O\theta)$; rotation invariance of surface measure
and of the annulus gives
$\norm{\Theta_O}_2=\norm\Theta_2$ and
$\norm{\Theta_O}_\infty=\norm\Theta_\infty$.  We may therefore assume
that one two-dimensional block is $\beta J$ with
$|\beta|\gtrsim_{n,d}1$.  For the associated twisted convolution
\[
 Sg(u,v)=\int k(a,b)e^{i\beta(ub-va)}g(u-a,v-b)da\,db,
\]
take the Fourier transform
$\widehat g(u,\eta)=\int e^{-iv\eta}g(u,v)\,dv$ in $v$.  A direct
calculation gives
\[
 \widehat{Sg}(u,\eta)
 =\int_{\mathbb R}
  \widehat k^{\,b}\bigl(a,\eta+\beta a-\beta u\bigr)
  \widehat g(u-a,\eta+\beta a)\,da,
\]
where
$\widehat k^{\,b}(a,\xi)=\int e^{-ib\xi}k(a,b)\,db$.
The quantity $q=\eta+\beta u$ is invariant under the passage from
$(u,\eta)$ to the input pair $(u-a,\eta+\beta a)$.  Setting $s=u-a$,
the operator on the fixed-$q$ fiber has kernel
\[
 K_q(u,s)
 =\widehat k^{\,b}(u-s,q-\beta(u+s)).
\]
The change of variables
\[
 (u,s)\longmapsto
 (a=u-s,\rho=q-\beta(u+s))
\]
has Jacobian of absolute value $2|\beta|$.  Therefore Plancherel in the
$b$ variable gives, uniformly in $q$,
\[
\begin{split}
 \norm{K_q}_{L^2_{u,s}}^2
 &=\frac1{2|\beta|}
   \iint|\widehat k^{\,b}(a,\rho)|^2\,da\,d\rho\\
 &\le C|\beta|^{-1}\norm k_2^2.
\end{split}
\]
The operator norm of each fiber is bounded by its Hilbert--Schmidt norm.
Since $(u,\eta)\mapsto(u,q)$ has unit Jacobian, direct integration in $q$
and Plancherel yield
\[
 \norm S_{2\to2}\le C|\beta|^{-1/2}\norm k_2.
\]
For the remaining coordinates, write $B=\beta J\oplus B'$ and
$z=(z_0,z')\in\R^2\times\R^{n-2}$.  For each fixed $z'$, the factor coming
from $B'$ is unimodular and the translation
$x'\mapsto x'-z'$ is an $L^2$ isometry.  Minkowski therefore reduces the
estimate to the preceding two-dimensional bound for the sliced kernel
$k_{z'}(z_0)=k_{R,u}(z_0,z')$.  Cauchy--Schwarz gives
\[
\int_{|z'|\le R}\norm{k_{z'}}_2\,dz'
\le |\{|z'|\le R\}|^{1/2}
 \left(\int\norm{k_{z'}}_2^2\,dz'\right)^{1/2}
\le CR^{(n-2)/2}\norm{k_{R,u}}_2
\le CR^{-1}\norm\Theta_2.
\]
Therefore
$\norm{A_{R,u}}_{2\to2}\le CR^{-1}\norm\Theta_2$, hence
$CR^{-1}\norm\Theta_\infty$.  This proves the proposition in every flag
stratum.  There are only finitely many strata.
\end{proof}

\begin{lemma}[Short variation from partial-annulus decay]\label{lem:shortannular}
Under the hypotheses of Proposition~\ref{prop:annulardecay}, after reducing
the exponent if necessary, every $f\in L^2(\R^n)$ satisfies
\begin{equation}\label{eq:shortannular}
 \norm{V^2(A^{P,\Theta}_{R,u}f:1/2\le u\le1)}_2
 \le C R^{-b_\mu}\norm\Theta_\infty\norm f_2
\end{equation}
for some $b_\mu>0$.
\end{lemma}

\begin{proof}
Let $F(u)=A_{R,u}f$.  Notice that $F(1)=0$.  For test functions,
\[
 F'(u,x)=-u^{-1}\int_{\Sph^{n-1}}\Theta(\theta)
 e^{iP(x,x-uR\theta)}f(x-uR\theta)d\sigma(\theta),
\]
so $\sup_u\norm{F'(u)}_2\le C\norm\Theta_1\norm f_2$.  If
$a=\sup_u\norm{F(u)}_2$ and $b=\sup_u\norm{F'(u)}_2$, the fundamental
theorem of calculus and the anchor at $1$ give pointwise
\[
 \sup_{1/2\le u\le1}|F(u,x)|^2
 \le2\int_{1/2}^1|F(v,x)|\,|F'(v,x)|\,dv.
\]
Integration in $x$ and Cauchy--Schwarz therefore give
$\norm{\sup_u|F(u)|}_2\le C a^{1/2}b^{1/2}$.  For every partition,
\[
 \sum|F(u_i)-F(u_{i-1})|^2
 \le2\sup_u|F(u)|\int_{1/2}^1|F'(u)|du.
\]
To record the exponents, put
\[
 M(x)=\sup_{1/2\le u\le1}|F(u,x)|,
 \qquad A(x)=\int_{1/2}^1|F'(u,x)|\,du.
\]
The preceding partition estimate says $V^2(F)^2\le2MA$.
Minkowski gives $\norm A_2\le b/2$, while the anchored fundamental
theorem of calculus gives $\norm M_2\le Ca^{1/2}b^{1/2}$.  Consequently
\[
 \norm{V^2(F)}_2^2
 \le2\norm M_2\norm A_2
 \le Ca^{1/2}b^{3/2},
\]
and hence $\norm{V^2(F)}_2\le Ca^{1/4}b^{3/4}$.  Insert
\eqref{eq:annulardecay}; since $\norm\Theta_1\le C_n\norm\Theta_\infty$, one
may take $b_\mu=a_\mu/4$.

For completeness, the argument extends from test functions to all
$f\in L^2$.  Choose $f_j\in C_c^\infty(\R^n)$ converging to $f$ in $L^2$ so
rapidly that
$\sum_j\norm{f_{j+1}-f_j}_2<\infty$, and let $F_j$ be the corresponding
paths.  For the path associated with $f_j-f_\ell$,
\[
 V^2\le V^1\le\int_{1/2}^1|F'_{j,\ell}(u)|\,du,
\]
so Minkowski and the derivative bound give
\[
 \norm{V^2(F_j-F_\ell)}_2
 \le C\norm\Theta_1\norm{f_j-f_\ell}_2.
\]
In particular,
\[
 W(x):=\sum_jV^2(F_{j+1}-F_j)(x)\in L^2(\R^n).
\]
All difference paths vanish at $u=1$, and hence, off the single null set
$\{W=\infty\}$,
\[
 \sup_{1/2\le u\le1}|F_{j+1}(u,x)-F_j(u,x)|
 \le V^2(F_{j+1}-F_j)(x).
\]
Thus $F_j(u,x)$ converges uniformly in the entire parameter interval on
one common full-measure set.  Denote its continuous limit by
$\widetilde F(u,x)$ and set every coordinate equal to zero on the
exceptional set.

This limit represents the intended operator at every parameter.  Indeed,
the absolute-value kernel on the fixed annulus has both row and column
integrals bounded by $C\norm\Theta_1$, so Schur's test gives, uniformly in
$u$,
\[
 \norm{F_j(u)-A^{P,\Theta}_{R,u}f}_2
 \le C\norm\Theta_1\norm{f_j-f}_2.
\]
On the other hand, the preceding summable variation bound gives
\[
 \sup_{u}\norm{F_j(u)-\widetilde F(u)}_2
 \le\sum_{k\ge j}\norm{V^2(F_{k+1}-F_k)}_2\longrightarrow0.
\]
Consequently $\widetilde F(u)$ is an $L^2$ representative of
$A^{P,\Theta}_{R,u}f$ for every $u$, with the representative chosen on one
common spatial set.

For clarity, this common representative may also be taken locally
absolutely continuous.  With a fixed Borel representative of $f$, define
in $L^2_x$
\[
 D_f(u,x)=-u^{-1}\int_{\Sph^{n-1}}\Theta(\theta)
 e^{iP(x,x-uR\theta)}f(x-uR\theta)\,d\sigma(\theta).
\]
The map $u\mapsto D_f(u)$ is strongly measurable and
$\sup_u\norm{D_f(u)}_2\le C\norm\Theta_1\norm f_2$.  Approximation by
$f_j$ in the fundamental theorem of calculus gives, in $L^2_x$,
\[
 A^{P,\Theta}_{R,v}f-A^{P,\Theta}_{R,u}f
   =\int_u^vD_f(w)\,dw.
\]
Fubini therefore makes $-\int_u^1D_f(w,x)\,dw$ an absolutely continuous
pointwise representative for almost every $x$.  It agrees with
$\widetilde F$ at all rational $u$ on one full-measure set and hence, by
continuity, at every $u$.

Finally, uniform pointwise convergence implies for each finite partition
$\pi$ that its $\ell^2$ increment sum for $\widetilde F$ is the limit of
the corresponding sums for $F_j$.  Taking the supremum over $\pi$, then
using Fatou's lemma, yields
\[
 \norm{V^2(\widetilde F)}_2
 \le\liminf_{j\to\infty}\norm{V^2(F_j)}_2
 \le C R^{-b_\mu}\norm\Theta_\infty\norm f_2,
\]
which proves \eqref{eq:shortannular} for general $f$ without exchanging an
uncountable supremum with an unsupported coordinatewise limit.
\end{proof}

\begin{corollary}[Scale-covariant form]\label{cor:scalecov}
Let $P$ be a real polynomial of degree at most $d$ whose first nonzero flag
coordinate is $c_\mu\ne0$ of degree $m_\mu$.
Let $\Theta\in L^\infty(\Sph^{n-1})$.  For a shell radius $R>0$ set
\[
 L_\mu=|c_\mu|R^{m_\mu}.
\]
If $L_\mu\ge1$, then
\begin{equation}\label{eq:scalecov}
 \sup_{1/2\le u\le1}\norm{A^{P,\Theta}_{R,u}}_{2\to2}
 +\norm{V^2_u(A^{P,\Theta}_{R,u})}_{2\to2}
 \le C\norm\Theta_\infty L_\mu^{-\beta_\mu}
\end{equation}
for some $\beta_\mu>0$.  Here the variation-operator norm means
\begin{equation}\label{eq:nonlinear-variation-operator-norm}
 \norm{V_u^2(A^{P,\Theta}_{R,u})}_{2\to2}
 :=\sup_{f\ne0}
 \frac{\norm{V^2(A^{P,\Theta}_{R,u}f:1/2\le u\le1)}_2}{\norm f_2}.
\end{equation}
This is the norm of a homogeneous sublinear map; linearity of the
variation map is not asserted.  For the later packet rescaling only (the
following variables are not hypotheses of \eqref{eq:scalecov}), fix a
packet side $q>0$ and a truncation parameter $\varepsilon>0$, and set
$t=\log_2(R/q)$.  Under the normalizing dilation
$x=\rho X$, $y=\rho Y$, $\rho=|c_\mu|^{-1/m_\mu}$,
\begin{equation}\label{eq:scalevariables}
 R\mapsto\widetilde R=R/\rho=L_\mu^{1/m_\mu},\quad
 q\mapsto\widetilde q=q/\rho,\quad
 \varepsilon\mapsto\widetilde\varepsilon=\varepsilon/\rho,
 \quad \log_2(\widetilde R/\widetilde q)=t.
\end{equation}
If $h_\rho(Y)=h(\rho Y)$, then
\begin{equation}\label{eq:scalenorms}
 \norm{h_\rho}_1=\rho^{-n}\norm h_1,
 \qquad \norm{h_\rho}_2=\rho^{-n/2}\norm h_2,
 \qquad \norm{h_\rho}_\infty=\norm h_\infty.
\end{equation}
Thus a normalized bound $\widetilde R^{-a_\mu}$ becomes
$L_\mu^{-a_\mu/m_\mu}$.  For any packet family $h_Q$ supported on
$q$-cubes, put $\widetilde Q=\rho^{-1}Q$ and
$\widetilde h_Q(Y)=h_Q(\rho Y)$.  Then
\[
 \norm{\widetilde h_Q}_1\le\Lambda|\widetilde Q|
 \quad\Longleftrightarrow\quad
 \norm{h_Q}_1\le\Lambda|Q|,
\]
so the density condition is invariant under the same rescaling.
\end{corollary}

\begin{proof}
Homogeneity of $K_\Theta$ gives the exact conjugacy
\[
 A^{P,\Theta}_{R,u}h(\rho X)
 =A^{P_\rho,\Theta}_{R/\rho,u}h_\rho(X).
\]
Here and below
\(P_\rho(X,Y)=P(\rho X,\rho Y)\) and
\(h_\rho(Y)=h(\rho Y)\).
The displayed transformations and norm identities follow by change of
variables.  Apply Proposition~\ref{prop:annulardecay} and
Lemma~\ref{lem:shortannular} at radius $\widetilde R$ and take
$\beta_\mu=\min(a_\mu,b_\mu)/m_\mu$.
\end{proof}

\section{Far-field strong estimates and triangular induction}
\label{sec:strong-far}

We convert the scale-covariant annular decay into an angular
$L\log L$ estimate and close the finite induction over the phase flag.
Proposition~\ref{prop:nonosc-strong} provides the terminal unphased
estimate, and Section~\ref{sec:flag} provides oscillatory decay on every
nonterminal stratum.

Call a phase \emph{normalized nonterminal} if its first nonzero flag
coordinate is $\pm1$.  Since the flag is finite, there is a number
\begin{equation}\label{eq:strong-far-a0}
 a_0=a_0(n,d)>0
\end{equation}
not exceeding the exponent in Proposition~\ref{prop:annulardecay} at any
nonterminal flag position.  If there is no such position, this section is
vacuous and we set $a_0=1$.

\subsection{Angular interpolation and short variation}

\begin{lemma}[$q'$-tracked partial-annulus estimates]
\label{lem:qtracked-annulus}
Let $P$ be a normalized nonterminal phase, let $R\geq1$, and let
$1<q\leq2$.  Suppose that
$\Theta\in L^q(\Sph^{n-1})$.  Uniformly in every later coefficient of $P$,
\begin{equation}\label{eq:qtracked-L2-annulus}
 \sup_{1/2\leq u\leq1}
 \norm{A^{P,\Theta}_{R,u}}_{2\to2}
 \leq C_{n,d}R^{-a_0/q'}\norm\Theta_q,
 \qquad q'=\frac q{q-1}.
\end{equation}
Moreover, for every $1<p<\infty$ put
\begin{equation}\label{eq:gamma-p}
 \gamma_p=\frac{a_0}{2\max(p,p')}>0
\end{equation}
and let $f\in L^p(\R^n)$.  Then
\begin{align}
 \sup_{1/2\leq u\leq1}
 \norm{A^{P,\Theta}_{R,u}f}_p
 &\leq C_{n,d,p}R^{-\gamma_p/q'}
       \norm\Theta_q\norm f_p,
 \label{eq:qtracked-Lp-annulus}\\
 \norm{V^2(A^{P,\Theta}_{R,u}f:1/2\leq u\leq1)}_p
 &\leq C_{n,d,p}R^{-\gamma_p/q'}
       \norm\Theta_q\norm f_p.
 \label{eq:qtracked-short}
\end{align}
All prefactors are uniform for $1<q\leq2$.
\end{lemma}

\begin{proof}
The trivial annular estimate is
\begin{equation}\label{eq:trivial-angular-annulus}
 \sup_u\norm{A^{P,\Theta}_{R,u}}_{2\to2}
 \leq \norm{K_\Theta\mathbf1_{\{R/2<|\cdot|<R\}}}_1
 \leq C\norm\Theta_1.
\end{equation}
For fixed $P,R,u$, the map
$\Theta\mapsto A^{P,\Theta}_{R,u}$ is linear from the angular space into
$\mathcal B(L^2)$.  Complex interpolation of
\eqref{eq:trivial-angular-annulus} with
Proposition~\ref{prop:annulardecay} gives
\eqref{eq:qtracked-L2-annulus}.  Indeed, if $\vartheta$ denotes the weight
of the $L^\infty(\Sph^{n-1})$ endpoint, then
\[
 \frac1q=\frac{1-\vartheta}{1}+\frac{\vartheta}{\infty},
 \qquad\text{so}\qquad \vartheta=\frac1{q'}.
\]
Thus $R^{-a_0}$ becomes $R^{-a_0/q'}$, with no $q$-dependent
interpolation loss.  This interpolation is first performed for each fixed
$u$ in the Banach target $\mathcal B(L^2)$; both endpoint bounds are
uniform in $u$, so the supremum in $u$ may then be taken.  The complex
interpolation constants are uniform in $q$.

For the short variation put $F(u)=A^{P,\Theta}_{R,u}f$.  Initially take
$f$ and $\Theta$ bounded.  The path is absolutely continuous, $F(1)=0$,
and polar coordinates give, for almost every $u\in[1/2,1]$,
\begin{equation}\label{eq:strong-far-derivative}
 F'(u,x)=-\frac1u\int_{\Sph^{n-1}}\Theta(\theta)
 e^{iP(x,x-uR\theta)}f(x-uR\theta)\,d\sigma(\theta).
\end{equation}
Consequently, for every $1\leq s\leq\infty$,
\begin{equation}\label{eq:strong-far-derivative-bound}
 \sup_{1/2\leq u\leq1}\norm{F'(u)}_s
 \leq C\norm\Theta_1\norm f_s.
\end{equation}
Set
\[
 a=\sup_u\norm{F(u)}_2,
 \qquad b=\sup_u\norm{F'(u)}_2.
\]
Since $F(1)=0$, the fundamental theorem of calculus yields pointwise
\[
 \sup_u|F(u,x)|^2
 \leq2\int_{1/2}^1|F(v,x)F'(v,x)|\,dv.
\]
After integration in $x$,
\begin{equation}\label{eq:strong-far-max-ab}
 \norm{\sup_u|F(u)|}_2\leq C a^{1/2}b^{1/2}.
\end{equation}
For every partition $1/2\leq u_0<\cdots<u_N\leq1$,
\[
 \sum_{i=1}^N|F(u_i)-F(u_{i-1})|^2
 \leq2\sup_u|F(u)|\int_{1/2}^1|F'(v)|\,dv.
\]
Taking the supremum over partitions, integrating, and using
\eqref{eq:strong-far-max-ab} gives
\begin{equation}\label{eq:strong-far-a14b34}
 \norm{V^2(F(u):1/2\leq u\leq1)}_2
 \leq C a^{1/4}b^{3/4}.
\end{equation}
Equations~\eqref{eq:qtracked-L2-annulus},
\eqref{eq:strong-far-derivative-bound}, and
$\norm\Theta_1\leq C_n\norm\Theta_q$ therefore imply
\begin{equation}\label{eq:qtracked-short-L2}
 \norm{V^2(F)}_2
 \leq C R^{-a_0/(4q')}\norm\Theta_q\norm f_2.
\end{equation}

The map
\[
 S_Rf:=V^2(A^{P,\Theta}_{R,u}f:1/2\le u\le1)
\]
is nonnegative, absolutely homogeneous, and subadditive.  Moreover, by
\eqref{eq:strong-far-derivative-bound},
\[
 S_Rf\leq V^1(A^{P,\Theta}_{R,u}f:1/2\le u\le1)
 \leq\int_{1/2}^1|F'(u)|\,du,
\]
and hence $S_R$ has nondecaying strong $(1,1)$ and $(\infty,\infty)$
bounds with norm $C\norm\Theta_1$.  We may therefore
apply the Marcinkiewicz interpolation theorem for homogeneous sublinear
operators.  For $1<p<2$, interpolation between the strong $(1,1)$ bound
and \eqref{eq:qtracked-short-L2} puts weight $2/p'$ on the $L^2$
endpoint and gives decay $a_0/(2p'q')$.  For $2<p<\infty$,
interpolation between \eqref{eq:qtracked-short-L2} and the strong
$(\infty,\infty)$ bound puts weight $2/p$ on the $L^2$ endpoint and gives
decay $a_0/(2pq')$.  These exponents are exactly
$\gamma_p/q'$, since
\[
 \frac{a_0}{4q'}\frac{2}{p'}
 =\frac{a_0}{2p'q'}\quad(1<p<2),
 \qquad
 \frac{a_0}{4q'}\frac{2}{p}
 =\frac{a_0}{2pq'}\quad(2<p<\infty).
\]
This proves \eqref{eq:qtracked-short}.  The same spatial
interpolation, now starting from
\eqref{eq:qtracked-L2-annulus}, proves the stronger corresponding bound
for the annular operator; weakening its exponent to $\gamma_p/q'$ gives
\eqref{eq:qtracked-Lp-annulus}.

To remove the boundedness assumptions, choose bounded $f_m\to f$ in
$L^p(\R^n)$ and bounded $\Theta_m\to\Theta$ in
$L^q(\Sph^{n-1})$, and put
$F_m(u)=A^{P,\Theta_m}_{R,u}f_m$.  By
\eqref{eq:strong-far-derivative-bound}, linearity, and $V^2\le V^1$,
\[
 \norm{V^2(F_m-F_\ell)}_p
 \le C\norm{\Theta_m}_1\norm{f_m-f_\ell}_p
 +C\norm{\Theta_m-\Theta_\ell}_1\norm{f_\ell}_p.
\]
The right-hand side tends to zero, since the sphere has finite measure.
Thus the paths are Cauchy in the metric
$\norm{V^2(F_m-F_\ell)}_p$.  We spell out the completion because the
parameter set is uncountable.  Every difference path is anchored at
$u=1$.  Choose a subsequence, still denoted by $F_m$, such that
\[
 \sum_m\norm{V^2(F_{m+1}-F_m:1/2\le u\le1)}_p<\infty.
\]
After deleting one null set, the corresponding pointwise series of
variations is finite.  The anchor then gives
\[
 \sup_{1/2\le u\le1}|F_{m+1}(u,x)-F_m(u,x)|
 \le V^2(F_{m+1}-F_m)(x),
\]
so $F_m(u,x)$ converges uniformly in $u$ on that common full-measure set.
For a fixed rational $u$, Young's inequality on the fixed annulus
identifies the limit with $A^{P,\Theta}_{R,u}f$ in $L^p$.  Passing to a
further diagonal subsequence and deleting one countable union of null
sets makes these identities simultaneous for all
rational $u\in[1/2,1]$.  Polar coordinates and Fubini give the latter
annular integrals a jointly measurable, locally absolutely continuous
representative in $u$: its radial derivative is the right side of
\eqref{eq:strong-far-derivative}, whose $L^1_u(L^p_x)$ norm is bounded by
$C\norm\Theta_1\norm f_p$.  Continuity therefore identifies the uniform
limit at every real $u$.  On each finite partition the increment sum is
continuous under this uniform convergence.  Taking first the limit, then
the supremum over the nested finite rational partitions, and using Fatou
proves \eqref{eq:qtracked-short}.  The same common representative and
coordinatewise convergence prove \eqref{eq:qtracked-Lp-annulus}.
\end{proof}

\subsection{The far path and the \texorpdfstring{$q'$}{q-prime} loss}

For a normalized nonterminal phase define
\begin{equation}\label{eq:strong-far-path}
 \cF^{P,\Theta}_\varepsilon f(x)=
 \int_{|x-y|>\max(\varepsilon,1)}
 e^{iP(x,y)}K_\Theta(x-y)f(y)\,dy
\end{equation}
on bounded compactly supported data.  Equivalently, put
\begin{equation}\label{eq:strong-far-Bj}
 B_j^{P,\Theta}=A^{P,\Theta}_{2^{j+1},1/2},
 \qquad j\geq0.
\end{equation}
If $2^j\leq\varepsilon\leq2^{j+1}$, then
\begin{equation}\label{eq:strong-far-series}
 \cF^{P,\Theta}_\varepsilon f
 =A^{P,\Theta}_{2^{j+1},\varepsilon/2^{j+1}}f
  +\sum_{k=j+1}^\infty B_k^{P,\Theta}f,
\end{equation}
whereas for $0<\varepsilon\leq1$ the path is the constant
$\sum_{k\geq0}B_k^{P,\Theta}f$.

We use the elementary long--short inequality
\begin{equation}\label{eq:V2-long-short}
 V^2(F_\varepsilon:\varepsilon>0)
 \leq C V^2(F_{2^j}:j\in\mathbb Z)
 +C\left(\sum_{j\in\mathbb Z}
 V^2(F_\varepsilon:2^j\leq\varepsilon\leq2^{j+1})^2
 \right)^{1/2}.
\end{equation}
Indeed, insert the adjacent dyadic endpoints into an arbitrary variation
partition.  Each increment is the sum of two within-block increments and
one increment between dyadic endpoints.  The between-block increments
form an ordered subpartition, and every dyadic block contributes at most
two within-block increments.  The inequality follows from
$|z_1+z_2+z_3|^2\leq3\sum|z_i|^2$.

\begin{proposition}[$q'$-tracked far \(V^2\)]
\label{prop:strong-far-q}
Let $P$ be normalized and nonterminal, let $1<p<\infty$, and let
$1<q\leq2$.  Suppose that $\Theta\in L^q(\Sph^{n-1})$; no cancellation
of $\Theta$ is required.  Then, for every $f\in L^p(\R^n)$,
\begin{equation}\label{eq:strong-far-q}
 \norm{\cF^{P,\Theta}_1f}_p
 +\norm{V^2(\cF^{P,\Theta}_\varepsilon f:\varepsilon>0)}_p
 \leq C_{n,d,p}q'\norm\Theta_q\norm f_p.
\end{equation}
The constant is uniform in $q$ and in every unused coefficient of $P$.
\end{proposition}

\begin{proof}
Lemma~\ref{lem:qtracked-annulus} gives
\begin{equation}\label{eq:strong-far-Bj-decay}
 \norm{B_j^{P,\Theta}f}_p
 \leq C_{n,d,p}2^{-\gamma_pj/q'}
 \norm\Theta_q\norm f_p.
\end{equation}
Thus the series in \eqref{eq:strong-far-series} converges in $L^p$ and
\begin{equation}\label{eq:strong-far-geometric}
 \sum_{j\geq0}\norm{B_j^{P,\Theta}f}_p
 \leq C_{n,d,p}q'\norm\Theta_q\norm f_p,
\end{equation}
because, for $q'\geq2$,
\[
 \sum_{j\geq0}2^{-\gamma_pj/q'}
 =\frac1{1-2^{-\gamma_p/q'}}\leq C_{\gamma_p}q'.
\]
This proves the asserted bound for $\cF_1$.
Choose measurable representatives of the countable family $B_jf$.  By
Minkowski and monotone convergence,
\[
 G(x):=\sum_{j\ge0}|B_j^{P,\Theta}f(x)|\in L^p(\R^n),
 \qquad \norm G_p\le\sum_{j\ge0}\norm{B_j^{P,\Theta}f}_p.
\]
Consequently the series is absolutely convergent outside one null set,
and all of its (countably many) dyadic tail identities hold there
simultaneously.  On every block we use the jointly measurable,
locally absolutely continuous partial-annulus representative constructed
in Lemma~\ref{lem:qtracked-annulus}.  Formula
\eqref{eq:strong-far-series} then defines one common representative of the
full far trajectory, with adjacent block definitions agreeing at their
shared dyadic endpoint.  All pointwise variation estimates below refer to
this representative.

At the dyadic parameters,
\[
 \cF^{P,\Theta}_{2^j}f=\sum_{k=j}^\infty B_k^{P,\Theta}f,
 \qquad j\geq0,
\]
and the sequence is constant for $j<0$.  Hence, pointwise,
\begin{equation}\label{eq:strong-far-long}
 V^2(\cF^{P,\Theta}_{2^j}f:j\in\mathbb Z)
 \leq\sum_{j\geq0}|B_j^{P,\Theta}f|.
\end{equation}
Its $L^p$ norm is bounded by \eqref{eq:strong-far-geometric}.

On the block $[2^j,2^{j+1}]$, the tail in
\eqref{eq:strong-far-series} is constant and the nonconstant part is
$A^{P,\Theta}_{2^{j+1},u}$, $1/2\leq u\leq1$.  Therefore
\eqref{eq:qtracked-short} gives
\[
 \norm{V^2(\cF_\varepsilon:2^j\leq\varepsilon\leq2^{j+1})}_p
 \leq C2^{-\gamma_pj/q'}\norm\Theta_q\norm f_p.
\]
Pointwise,
$(\sum_j V_j^2)^{1/2}\leq\sum_jV_j$; hence the $L^p$ norm of the
short term in \eqref{eq:V2-long-short} is again bounded by the geometric
sum in \eqref{eq:strong-far-geometric}.  Combining the long and short
parts proves \eqref{eq:strong-far-q}.
\end{proof}

\subsection{Finite-trajectory extrapolation to
\texorpdfstring{$L\log L$}{L log L}}

Let
\[
 \Phi_0(t)=(1+t)\log(1+t)-t,
 \qquad \Psi_0(s)=e^s-s-1.
\]
These are complementary Young functions.  On the finite measure space
$\Sph^{n-1}$, the Luxemburg norm for $\Phi_0$ is equivalent, with
dimensional constants, to the
homogeneous gauge $\mathfrak L$ in \eqref{eq:gauge}; the Luxemburg norm
for $\Psi_0$ will be denoted by $\norm{\cdot}_{\exp L}$.  The
Orlicz--H\"older inequality therefore gives
\begin{equation}\label{eq:strong-far-Orlicz-Holder}
 \int_{\Sph^{n-1}}|uv|\,d\sigma
 \leq2\norm u_{L\log L}\norm v_{\exp L}.
\end{equation}
For the quantitative direction used below the assertion is trivial when
$\Theta=0$; otherwise let
$A=\norm\Theta_1$ and $B=\mathfrak L(\Theta)$.  Since $B\geq A$ and
$\Phi_0(t)\leq Ct\log(e+t)$,
\[
 \int_{\Sph^{n-1}}\Phi_0\!\left(\frac{|\Theta|}{B}\right)d\sigma
 \leq\frac C B\left[
 A+\int_{\Sph^{n-1}}|\Theta|
 \log^+\!\left(\frac{|\Theta|}{A}\right)d\sigma\right]\leq C.
\]
Convexity and $\Phi_0(0)=0$ imply
$\Phi_0(t/C_1)\leq C_1^{-1}\Phi_0(t)$ for $C_1\geq1$.
Choosing $C_1$ larger than the preceding absolute constant gives
\begin{equation}\label{eq:gauge-controls-Luxemburg}
 \norm\Theta_{L\log L}\leq C_1\mathfrak L(\Theta).
\end{equation}

\begin{proposition}[Angular endpoint for the far path]
\label{prop:strong-far-endpoint}
Let $P$ be normalized and nonterminal, let $1<p<\infty$, and let
$\Theta\in L\log L(\Sph^{n-1})$.  No cancellation of $\Theta$ is
required.  Then, for every $f\in L^p(\R^n)$,
\begin{equation}\label{eq:strong-far-endpoint}
 \norm{\cF^{P,\Theta}_1f}_p
 +\norm{V^2(\cF^{P,\Theta}_\varepsilon f:\varepsilon>0)}_p
 \leq C_{n,d,p}\mathfrak L(\Theta)\norm f_p.
\end{equation}
Consequently,
\begin{equation}\label{eq:strong-far-endpoint-jump}
 \Jop{p}(\cF^{P,\Theta}_\varepsilon f:\varepsilon>0)
 \leq C_{n,d,p}\mathfrak L(\Theta)\norm f_p.
\end{equation}
\end{proposition}

\begin{proof}
Fix $f$ and a finite ordered set $D\subset\mathbb Q_+$ containing $1$.
On $\mathbb C^D$ define
\begin{equation}\label{eq:strong-far-ED}
 \norm z_{E_D}=|z_1|+V_D^2(z_\varepsilon:\varepsilon\in D).
\end{equation}
This is a norm: vanishing variation makes the trajectory constant, and
the value at $1$ then detects that constant.  Since $D$ is finite,
$E_D$ is a finite-dimensional Banach space, hence
$L^p(\R^n;E_D)$ is Banach and its norm is recovered by Hahn--Banach
duality.  For bounded angular data
define, after complexification if necessary, the linear map
\[
 \mathcal A_{f,D}\Theta
 =(\cF^{P,\Theta}_\varepsilon f)_{\varepsilon\in D}
 \quad\hbox{in }L^p(\R^n;E_D).
\]
Proposition~\ref{prop:strong-far-q}, uniformly in $D$, says
\begin{equation}\label{eq:strong-far-ADq}
 \norm{\mathcal A_{f,D}}_{L^q(\Sph^{n-1})\to L^p(E_D)}
 \leq C_{n,d,p}q'\norm f_p,
 \qquad1<q\leq2.
\end{equation}

Let $\Lambda\in[L^p(\R^n;E_D)]^*$ have norm one.  For $s\geq2$ take
$q=s/(s-1)$ in \eqref{eq:strong-far-ADq}.  The scalar functional
$\Theta\mapsto\Lambda(\mathcal A_{f,D}\Theta)$ is represented by a
function $h_s\in L^s(\Sph^{n-1})$ satisfying
\begin{equation}\label{eq:strong-far-hs}
 \norm{h_s}_s\leq Cs\norm f_p.
\end{equation}
The representatives $h_s$ agree almost everywhere.  Indeed, two of them
represent the same functional on the common subspace of bounded angular
functions: all of these extensions agree there with the original
finite-trajectory formula.  Since the sphere has finite measure,
$h_s-h_t\in L^1$ for any $s,t\ge2$.  Its bounded complex sign belongs to
every relevant $L^q$ space, and testing with that sign gives
$\int|h_s-h_t|\,d\sigma=0$.  Thus the representatives agree almost
everywhere; denote their common representative by $h$.

Put $A=\norm f_p$.  If $A=0$ there is nothing to prove.  For every integer $k\geq2$,
$\norm h_k\leq CkA$.  Hence, by Tonelli and the exponential series,
\begin{equation}\label{eq:strong-far-exp-series}
 \int_{\Sph^{n-1}}\Psi_0\!\left(\frac{|h|}{LA}\right)d\sigma
 =\sum_{k=2}^\infty
 \frac{\norm h_k^k}{(LA)^k k!}
 \leq\sum_{k=2}^\infty\frac{(Ck/L)^k}{k!}.
\end{equation}
Since $k^k/k!\leq C_0e^k\sqrt{k}$, the last series is at most one once
$L=L(n,d,p)$ is sufficiently large.  Thus
\begin{equation}\label{eq:strong-far-exp-norm}
 \norm h_{\exp L}\leq C_{n,d,p}\norm f_p.
\end{equation}
Equations~\eqref{eq:strong-far-Orlicz-Holder} and
\eqref{eq:strong-far-exp-norm} give
\[
 |\Lambda(\mathcal A_{f,D}\Theta)|
 \leq C_{n,d,p}\norm\Theta_{L\log L}\norm f_p.
\]
Taking the supremum over $\Lambda$ proves
\begin{equation}\label{eq:strong-far-finite-D}
 \norm{\mathcal A_{f,D}\Theta}_{L^p(E_D)}
 \leq C_{n,d,p}\mathfrak L(\Theta)\norm f_p,
\end{equation}
uniformly in the finite set $D$.

Bounded angular functions are dense in the Luxemburg \(L\log L\) norm.
Indeed, for
$\Theta_N=\Theta\mathbf1_{\{|\Theta|\le N\}}$, the $\Delta_2$ condition
for $\Phi_0$ and dominated convergence for the corresponding modular give
\[
 \norm{\Theta-\Theta_N}_{L\log L}\longrightarrow0.
\]
Thus the finite-trajectory maps extend to all $\Theta\in L\log L$.  To
check consistency, if $D\subset D'$ and both contain $1$, coordinate
restriction satisfies
\[
 |z_1|+V_D^2(z)\le |z_1|+V_{D'}^2(z),
\]
so it is a contraction from $E_{D'}$ to $E_D$.  The two extensions agree
on bounded angular functions and hence, by density, on all of $L\log L$.
Thus the rational trajectory is independent of the exhausting sequence.
Fix once and for all an enumeration of $\mathbb Q_+$ beginning with $1$,
and let $D_N$ be its first $N$ elements.  Choose representatives of the
countable family of rational coordinates simultaneously.  For every $N$,
\eqref{eq:strong-far-finite-D} applies to these same coordinates, and
\[
 |\cF_1f|+V^2_{D_N}(\cF_\varepsilon f)
 \uparrow |\cF_1f|+V^2_{\mathbb Q_+}(\cF_\varepsilon f)
 \quad\hbox{pointwise}.
\]
Monotone convergence (equivalently, Fatou followed by monotonicity) now
passes the uniform finite-$D_N$ estimate to all positive rational
parameters.  The finite-dimensional extension is completed first, and
only then is $N\to\infty$ taken.  For $0<a<b<\infty$, finite-shell
differences satisfy
\begin{equation}\label{eq:strong-far-shell-continuity}
 \norm{(\cF_a-\cF_b)f}_p
 \leq \norm\Theta_1\log(b/a)\norm f_p,
\end{equation}
including when $[a,b]$ meets $1$: the difference is zero if $b\le1$, and
if $a<1<b$ its kernel has $L^1$ norm
$\norm\Theta_1\log b\le\norm\Theta_1\log(b/a)$.

The difference identity and estimate above are first obtained on rational
parameters; approximation by bounded angular data and Young's inequality
preserve that identity.  Put
\[
 \mathcal A_x(r)=\int_{\Sph^{n-1}}
 e^{iP(x,x-r\theta)}\Theta(\theta)f(x-r\theta)\,d\sigma(\theta).
\]
For every $0<a<b<\infty$, polar coordinates give
\[
 \int_a^b|\mathcal A_x(r)|\,\frac{dr}{r}
 \le\int_{a<|z|<b}|K_\Theta(z)|\,|f(x-z)|\,dz.
\]
The right-hand side belongs to $L^p$.  Applying this observation to a
countable family of annuli covering $(0,\infty)$, discard one null set on
whose complement all these radial densities are locally integrable.
Choose one representative of $\cF_1f$, and further remove the countable
union of the null sets on which a rational finite-shell identity fails.
On the resulting common full-measure set define
\[
 \cF_\varepsilon f=
 \begin{cases}
  \cF_1f,&0<\varepsilon\le1,\\[1ex]
  \displaystyle
  \cF_1f-\int_{1<|z|\le\varepsilon}
  e^{iP(x,x-z)}K_\Theta(z)f(x-z)\,dz,&\varepsilon>1.
 \end{cases}
\]
The rational finite-shell identities show that this leaves every rational
coordinate unchanged.  The resulting path is locally absolutely
continuous and satisfies, for almost every $\varepsilon$,
\[
 \frac{d}{d\varepsilon}\cF_\varepsilon f(x)
 =-\mathbf1_{(1,\infty)}(\varepsilon)
   \frac{\mathcal A_x(\varepsilon)}{\varepsilon}.
\]
Lemma~\ref{lem:finite-witness} therefore identifies the rational
variation with the full-parameter variation.  This proves
\eqref{eq:strong-far-endpoint}.
Finally \eqref{eq:jumpvar} gives
\eqref{eq:strong-far-endpoint-jump}.
\end{proof}

\begin{remark}[The far anchor]
\label{rem:strong-far-anchor}
The construction above is anchored at infinity.  For bounded $\Theta$,
take $q=2$ in Lemma~\ref{lem:qtracked-annulus} and interpolate spatially
as in its proof.  More explicitly, for $\varepsilon\ge1$ choose
$j\ge0$ with $2^j\le\varepsilon\le2^{j+1}$.  Then the current partial annulus is bounded
in $L^p$ by
\[
 C_{n,d,p}2^{-\gamma_pj/2}\norm\Theta_2\norm f_p,
\]
and the remaining annular tail is bounded by
\[
 C_{n,d,p}\sum_{k\ge j+1}2^{-\gamma_pk/2}
 \norm\Theta_2\norm f_p.
\]
Both quantities tend to zero as $j\to\infty$, and hence as
$\varepsilon\to\infty$.
For general $\Theta\in L\log L$, choose bounded
$\Theta_m\to\Theta$ in that norm.
Proposition~\ref{prop:strong-far-endpoint} and
\[
 \sup_{\varepsilon>0}|F_\varepsilon|
 \leq |F_1|+V^2(F_\varepsilon:\varepsilon>0)
\]
make the approximation uniform in $L^p$ over $\varepsilon$.  Hence
$\cF^{P,\Theta}_\varepsilon f\to0$ in $L^p$.  Finite \(V^2\) supplies an
almost-everywhere limit at infinity, and an \(L^p\)-convergent subsequence
identifies that limit as zero.
\end{remark}

\subsection{Local triangular reduction}

For a normalized phase put
\begin{equation}\label{eq:strong-local-path}
 \cL^{P,\Omega}_\varepsilon f(x)=
 \begin{cases}
 \displaystyle\int_{\varepsilon<|x-y|<1}
 e^{iP(x,y)}K_\Omega(x-y)f(y)\,dy,&0<\varepsilon<1,\\[1ex]
 0,&\varepsilon\geq1.
 \end{cases}
\end{equation}

We record two elementary endpoint facts.  First, if a path $E$ is zero at
one endpoint and $V^1(E)\leq W$, then
\begin{equation}\label{eq:V1-controls-jump}
 \lambda N_\lambda(E)\leq W,
 \qquad N_\lambda(E)>0\Longrightarrow\lambda\leq W,
 \qquad \lambda N_\lambda(E)^{1/2}\leq W.
\end{equation}
Indeed, the sum of the increments in any ordered family of jump pairs is
at most the total \(1\)-variation.  Second, if
\begin{equation}\label{eq:localized-full-path}
 L_\varepsilon=
 \begin{cases}F_\varepsilon-F_1,&0<\varepsilon<1,\\0,&\varepsilon\geq1,
 \end{cases}
\end{equation}
and $F_\varepsilon\to F_1$ as $\varepsilon\uparrow1$, then
\begin{equation}\label{eq:localized-jump-domination}
 N_\lambda(L)\leq N_\lambda(F).
\end{equation}
To see this, retain every jump pair lying below \(1\); if a pair crosses
\(1\), it is necessarily the last nonzero pair and may be replaced by the
pair with upper endpoint \(1\).

\begin{proposition}[Triangular local reduction in the strong range]
\label{prop:strong-local-reduction}
Let $P$ be a real polynomial of degree at most $d$ whose first nonzero
flag coordinate is $c_\nu=\pm1$.  Let
$\Omega\in L^1(\Sph^{n-1})$ and let $f$ be bounded and compactly
supported.  On each unit output cube, the local path
\eqref{eq:strong-local-path} is, modulo
pure input and output phases, the sum of
\begin{enumerate}[label=\textup{(\roman*)}]
\item a localized path of the form \eqref{eq:localized-full-path} whose
residual phase either has all flag coordinates equal to zero, or has first
nonzero flag coordinate strictly later than \(\nu\);
\item an error path \(E\) satisfying
\begin{equation}\label{eq:strong-local-error}
 V^1(E)(x)\leq C_{n,d}
 \int_{|x-y|<1}
 \frac{|\Omega((x-y)/|x-y|)|}{|x-y|^{n-1}}|f(y)|\,dy.
\end{equation}
\end{enumerate}
The majorant in \eqref{eq:strong-local-error} has \(L^p\) norm at most
$C_{n,d}\norm\Omega_1\norm f_p$ for every \(1\leq p\leq\infty\).
\end{proposition}

\begin{proof}
Let $Q_a$ be the output unit cube centered at \(a\), and let $Q_a^*$ be
the fixed input enlargement relevant to $Q_a$.  Write \(x=a+X\),
\(y=a+Y\).  Theorem~\ref{thm:phaseflag}(iii) gives
\[
 P(a+X,a+Y)=p_a(X)+q_a(Y)+c_\nu H_\nu(X,Y)
 +P_{a,>\nu}(X,Y),
\]
where every flag coordinate of \(P_{a,>\nu}\) is later than \(\nu\).
Set
\[
 f_a(Y)=e^{iq_a(Y)}f(a+Y)\mathbf1_{Q_a^*-a}(Y).
\]
The coordinate translation identifies the original path at $a+X$ with
the translated path at $X$ and preserves Lebesgue measure.  Multiplication
of the output path by $e^{ip_a(X)}$ leaves every pointwise jump unchanged,
while the pure input phase is absorbed into $f_a$.  In particular,
\[
 |f_a(Y)|=|f(a+Y)|\mathbf1_{Q_a^*-a}(Y),
 \qquad
 \norm{f_a}_p=\norm{f\mathbf1_{Q_a^*}}_p.
\]
Thus the local path is the path with phase
$c_\nu H_\nu+P_{a,>\nu}$ applied to $f_a$.  If
\(X\) lies in the reference output cube and \(|X-Y|<1\), both variables
belong to a fixed bounded set.  Since \(H_\nu(Y,Y)=0\),
\begin{equation}\label{eq:strong-local-H-bound}
 |H_\nu(X,Y)|\leq C_{n,d}|X-Y|.
\end{equation}
Thus
\[
 |e^{ic_\nu H_\nu(X,Y)}-1|\leq C_{n,d}|X-Y|.
\]
The total \(1\)-variation of the difference of the two nested truncation
paths is therefore bounded by \eqref{eq:strong-local-error}.  Its kernel
has \(L^1\) norm
\[
 C_{n,d}\int_0^1\int_{\Sph^{n-1}}|\Omega(\theta)|
 \,d\sigma(\theta)\,dr
 =C_{n,d}\norm\Omega_1,
\]
which proves the asserted \(L^p\) bound.  Removing this difference leaves
the local truncation of \(P_{a,>\nu}\) applied to $f_a$.  To identify it
precisely with \eqref{eq:localized-full-path}, let
\[
 F_\varepsilon
 =T^{P_{a,>\nu}}_{\Omega,\varepsilon}f_a
\]
at this bounded compactly supported stage.  The finite-shell identity
gives
\[
 \cL^{P_{a,>\nu},\Omega}_\varepsilon f_a
 =F_\varepsilon-F_1\quad(0<\varepsilon<1),
 \qquad
 \cL^{P_{a,>\nu},\Omega}_\varepsilon f_a=0\quad(\varepsilon\ge1).
\]
Moreover $F_\varepsilon\to F_1$ as $\varepsilon\uparrow1$ by the
$L^1$ norm of the intervening shell kernel.  Thus the remaining path is
exactly the localized full path in \eqref{eq:localized-full-path}, and
\eqref{eq:localized-jump-domination} applies to it.
\end{proof}

\subsection{Coefficient-uniform strong jumps for every phase}

\begin{proposition}[Strong polynomial critical jumps]
\label{prop:strong-polynomial-jump}
Let \(1<p<\infty\), let \(P\) be any real polynomial of degree at most
\(d\), and let \(\Omega\in L\log L(\Sph^{n-1})\) have mean zero.  Then
for every $f\in L^p(\R^n)$,
\begin{equation}\label{eq:strong-polynomial-jump}
 \Jop{p}(T^P_{\Omega,\varepsilon}f:\varepsilon>0)
 \leq C_{n,d,p}\mathfrak L(\Omega)\norm f_p,
\end{equation}
where the constant is independent of every coefficient of \(P\).
Consequently, for every \(r>2\),
\begin{equation}\label{eq:strong-polynomial-variation}
 \norm{V^r(T^P_{\Omega,\varepsilon}f:\varepsilon>0)}_p
 \leq C_{n,d,p,r}\mathfrak L(\Omega)\norm f_p.
\end{equation}
\end{proposition}

\begin{proof}
We first take bounded compactly supported data, for which every finite-shell
identity is classical.  At this stage the induction concerns only these
classical finite-shell trajectories and uses no extension at the flag
position currently being proved.  We use backward induction on the finite
phase flag.  If every flag coordinate vanishes, then \(P\) is a sum of a pure
input phase and a pure output phase.  Equation~\eqref{eq:purephase} and
Proposition~\ref{prop:nonosc-strong} give
\eqref{eq:strong-polynomial-jump}.

Assume the estimate for every phase whose first coordinate is later than
\(\nu\), and suppose that the first coordinate of \(P\) is
\(c_\nu\ne0\), of homogeneous degree \(m_\nu\).  Under
\(x=\rho X\), \(y=\rho Y\), where
\begin{equation}\label{eq:strong-normalizing-rho}
 \rho=|c_\nu|^{-1/m_\nu},
\end{equation}
the first coordinate becomes \(\pm1\).  Work first with this normalized
phase.  Proposition~\ref{prop:strong-far-endpoint} controls its far path.

For the local path partition the output into disjoint unit cubes \(Q_a\).
Only the restriction of \(f\) to a fixed enlargement \(Q_a^*\) contributes
on \(Q_a\).  Proposition~\ref{prop:strong-local-reduction} gives a
residual phase \(P_{a,>\nu}\), the modulated translated input $f_a$ used
in its proof, and an error \(E_a\).  If all flag coordinates of
$P_{a,>\nu}$ vanish, Proposition~\ref{prop:nonosc-strong} together with
the pure input/output factorization \eqref{eq:purephase} supplies the base
case.  Otherwise its first nonzero coordinate is strictly later than
$\nu$, so the induction hypothesis applies.  In either case,
\eqref{eq:localized-jump-domination} and the supremum over the jump
amplitude give
\begin{equation}\label{eq:strong-local-later-bound}
 \norm{\mathbf J(\cL^{P_{a,>\nu},\Omega}_\varepsilon
 f_a:\varepsilon>0)}_p
 \leq C\mathfrak L(\Omega)\norm{f_a}_p
 =C\mathfrak L(\Omega)\norm{f\mathbf1_{Q_a^*}}_p
\end{equation}
uniformly in \(a\) and in all translated residual coefficients.
Equations~\eqref{eq:V1-controls-jump},
\eqref{eq:pointwise-jump-by-V2}, and
\eqref{eq:strong-local-error} give the same estimate for
$\mathbf J(E_a)$, with
\(\norm\Omega_1\leq C\mathfrak L(\Omega)\).

Apply the pointwise quasi-triangle inequality
\eqref{eq:pointwise-jump-quasitriangle} on every output cube, sum the
\(p\)-th powers over the disjoint \(Q_a\), and use bounded overlap of the
\(Q_a^*\).  More explicitly, the overlap multiplicity is bounded by a
dimensional constant and hence
\[
 \sum_a\norm{f\mathbf1_{Q_a^*}}_p^p
 =\int |f(x)|^p\sum_a\mathbf1_{Q_a^*}(x)\,dx
 \le C_n\norm f_p^p.
\]
Combining this identity with \eqref{eq:strong-local-later-bound} and the
corresponding error estimate gives
\begin{equation}\label{eq:strong-local-globalized}
 \Jop{p}(\cL^{P,\Omega}_\varepsilon f:\varepsilon>0)
 \leq C_{n,d,p}\mathfrak L(\Omega)\norm f_p.
\end{equation}
Adding the local and far paths and using
\eqref{eq:pointwise-jump-quasitriangle} proves the normalized case.

It remains to undo the dilation.  If \(f_\rho(Y)=f(\rho Y)\), homogeneity
of the kernel gives the exact conjugacy
\begin{equation}\label{eq:strong-jump-dilation-conjugacy}
 T^P_{\Omega,\varepsilon}f(\rho X)
 =T^{P_\rho}_{\Omega,\varepsilon/\rho}f_\rho(X).
\end{equation}
The map $\varepsilon\mapsto\varepsilon/\rho$ is an order-preserving
bijection of $(0,\infty)$, so it leaves every jump number unchanged.
Changing variables in the input and output norms gives the explicit
identities
\[
 \norm{f_\rho}_{L_X^p}=\rho^{-n/p}\norm f_{L_x^p},
 \qquad
 \norm{\mathbf J(T^{P_\rho}_{\Omega,\varepsilon/\rho}f_\rho)}
       _{L_X^p}
 =\rho^{-n/p}
  \norm{\mathbf J(T^P_{\Omega,\varepsilon}f)}_{L_x^p}.
\]
Thus the input and output acquire the same factor $\rho^{-n/p}$.
Thus the constant does not depend on \(c_\nu\).  The finite backward
induction proves \eqref{eq:strong-polynomial-jump} for every phase.

The variation estimate is now a pointwise consequence rather than a
mixed-norm extrapolation: by \eqref{eq:pointwise-jump-to-variation},
\eqref{eq:kappa-r}, and \eqref{eq:strong-polynomial-jump},
\[
 \norm{V^r(T^P_{\Omega,\varepsilon}f:\varepsilon>0)}_p
 \le \kappa_r\Jop{p}(T^P_{\Omega,\varepsilon}f:\varepsilon>0)
 \le C_{n,d,p}\kappa_r\mathfrak L(\Omega)\norm f_p.
\]
This proves \eqref{eq:strong-polynomial-variation}, with the stated
$r$-dependence absorbed in its constant.

We next construct the anchored trajectory and pass to general data.  For a
normalized nonterminal phase,
Proposition~\ref{prop:strong-far-endpoint} defines the far path on $L^p$,
including
its value at $1$; add the finite-shell local path
\eqref{eq:strong-local-path}.  Their sum has the classical finite-shell
increments, agrees with the far path for $\varepsilon\ge1$, and is
anchored at infinity by Remark~\ref{rem:strong-far-anchor}.  For a
nonnormalized nonterminal phase use the conjugacy
\eqref{eq:strong-jump-dilation-conjugacy}.  If every flag coordinate
vanishes, use the anchored factorization of
Proposition~\ref{prop:nonosc-strong} and the pure modulations
\eqref{eq:purephase}.

Choose bounded compactly supported $f_m\to f$ in $L^p$.  Write
\[
 \widetilde\Omega_m=\Omega\mathbf1_{\{|\Omega|\leq m\}},
 \qquad
 c_m=\frac1{m_n}\int_{\Sph^{n-1}}\widetilde\Omega_m\,d\sigma,
 \qquad m_n=\sigma(\Sph^{n-1}),
\]
and put $\Omega_m=\widetilde\Omega_m-c_m$.  Since the defining Orlicz
function satisfies the $\Delta_2$ condition,
$\widetilde\Omega_m\to\Omega$ in the Luxemburg $L\log L$ norm.
Moreover, using $\int\Omega\,d\sigma=0$ and
$L\log L\hookrightarrow L^1$,
\[
 c_m=-\frac1{m_n}\int_{\Sph^{n-1}}
       \Omega\mathbf1_{\{|\Omega|>m\}}\,d\sigma\longrightarrow0.
\]
On the finite-measure sphere the Luxemburg norm of the constant function
$c_m$ is $C_n|c_m|$.  Hence the bounded mean-zero functions
$\Omega_m$ converge to $\Omega$ in $L\log L$.  By
Lemma~\ref{lem:gauge-normalization}, this convergence is equivalently
\begin{equation}\label{eq:angular-approx-gauge-convergence}
 \mathfrak L(\Omega_m-\Omega)\longrightarrow0.
\end{equation}
The lattice property of
the Luxemburg norm, the embedding $L\log L\hookrightarrow L^1$ on the
sphere, and the triangle inequality also give
\[
 \norm{\widetilde\Omega_m}_{L\log L}
 \leq\norm\Omega_{L\log L},
 \qquad
 \norm{c_m\mathbf1}_{L\log L}
 \leq C_n|c_m|
 \leq C_n\norm\Omega_1
 \leq C_n\norm\Omega_{L\log L}.
\]
Lemma~\ref{lem:gauge-normalization} therefore yields the uniform bound
\begin{equation}\label{eq:angular-approx-uniform-gauge}
 \sup_{m\geq1}\mathfrak L(\Omega_m)
 \leq C_n\mathfrak L(\Omega).
\end{equation}
For a normalized nonterminal
phase, bilinearity gives
\[
 \cF^{P,\Omega_m}_\varepsilon f_m
 -\cF^{P,\Omega}_\varepsilon f
 =
 \cF^{P,\Omega_m}_\varepsilon(f_m-f)
 +\cF^{P,\Omega_m-\Omega}_\varepsilon f.
\]
Proposition~\ref{prop:strong-far-endpoint},
\eqref{eq:angular-approx-uniform-gauge}, and
\eqref{eq:angular-approx-gauge-convergence} show that both terms tend to
zero in the finite-trajectory \(L^p(E_D)\) norm, uniformly over every
fixed finite parameter set $D$.  Thus the far coordinates for
$(\Omega_m,f_m)$ converge in $L^p$ to those for $(\Omega,f)$.
The same bilinear decomposition and Young's inequality give convergence
of every local finite-shell coordinate.
The same conclusion follows after dilation, and in the pure case from the
factorization estimates in Proposition~\ref{prop:nonosc-strong}.  The
finite-shell bound used here is, for \(0<a<b<\infty\),
\[
 \norm{\int_{a<|z|<b}e^{iP(x,x-z)}
 K_\Theta(z)g(x-z)\,dz}_p
 \leq \norm\Theta_1\log(b/a)\norm g_p.
\]
Enumerate $\mathbb Q_+=\{r_1,r_2,\ldots\}$ and put
$D_N=\{r_1,\ldots,r_N\}$.  A single diagonal subsequence of the
coordinatewise $L^p$ convergence is almost everywhere simultaneous at
every rational parameter.  For each fixed $N$, apply
\eqref{eq:finite-J-lsc} on the finite coordinate space indexed by $D_N$
and then Fatou's lemma.  This transfers
\eqref{eq:strong-polynomial-jump} to the limit with a bound independent of
$N$.  Only after this finite-dimensional passage do we let $N\to\infty$.
The pointwise functionals increase monotonically, and every strict rational
jump witness is contained in one $D_N$; monotone convergence therefore
yields the bound on $\mathbb Q_+$.
The finite-shell integrability used in
\eqref{eq:strong-far-shell-continuity} gives local absolute continuity in
the truncation parameter almost everywhere, so
Lemma~\ref{lem:finite-witness} recovers all positive parameters.
The variation estimate then follows pointwise from
\eqref{eq:pointwise-jump-to-variation}.  The construction above supplies
the $L^p$ and almost-everywhere anchor at infinity.
\end{proof}
\medskip

\section{Separated-packet estimates for the terminal phase}
\label{sec:terminal-packets}

All packet families in this and the following packet sections are finite;
we suppress this qualifier in individual statements.  Rankings and
randomizations are always performed at this finite level.  In the endpoint
arguments, cube, shell, and parameter cutoffs are imposed explicitly and
removed only after the estimates have been established, by coordinatewise
convergence on finite annuli and lower semicontinuity of strict jump counts.
Thus no infinite randomization or infinite-rank argument is used.

For $n\ge2$, we use the following two microlocal estimates of Seeger
\cite[Lemmas~2.1 and 2.2]{Seeger96}.
Let $H_j$ be supported where $2^{j-2}\le|x|\le2^{j+2}$, and
assume that $r\mapsto H_j(r\theta)$ is $C^N$ on $(0,\infty)$ for every
$\theta$, with the radial derivatives satisfying
\begin{equation}\label{eq:MN}
 M_N=\sup_{j,\theta,r}\sup_{0\le\ell\le N}
 r^{n+\ell}|\partial_r^\ell H_j(r\theta)|<\infty.
\end{equation}
For an integer $s>3$ and a fixed $0<\kappa<1$, Seeger's construction
\cite[equation~(2.2)]{Seeger96} is
$H_j=\Gamma_j^s+(H_j-\Gamma_j^s)$.  Put $L(Q)=m$ when
$2^{m-1}<\ell(Q)\le2^m$.  If the cubes have disjoint interiors,
$F_m=\sum_{L(Q)=m}f_Q$, and
$\supp f_Q\subset Q$, $\norm{f_Q}_1\le\Lambda|Q|$, then, for
$0<\kappa<1$,
\begin{equation}\label{eq:seeger1}
 \norm{\sum_j\Gamma_j^s*F_{j-s}}_2^2
 \le C M_0^2 2^{-s(1-\kappa)}\Lambda\sum_Q\norm{f_Q}_1.
\end{equation}
If additionally $\supp b_Q\subset Q$, $\int b_Q=0$,
$\ell(Q)=2^{j-s}$, $N\ge n+1$, and
$0\le\epsilon\le1$, then
\begin{equation}\label{eq:seeger2}
 \norm{(H_j-\Gamma_j^s)*b_Q}_1
 \le C_N\left[M_0 2^{-s\epsilon}
 +M_N2^{s(n+(\epsilon-\kappa)N)}\right]\norm{b_Q}_1.
\end{equation}
The constants may depend on the fixed value of $\kappa$ and on Seeger's
fixed auxiliary cutoffs, but not on $j,s$, the cubes, or the functions.
These are precisely the support hypotheses in
\cite[Lemmas~2.1--2.2]{Seeger96}: the first lemma requires no cancellation,
whereas the second uses the displayed mean-zero condition.
In dimension one, the corresponding one-shell estimates below are
proved directly from cancellation, radial derivative bounds, and the same
$L^1$--BMO interpolation mechanism.  This direct branch is what closes
the one-dimensional weak endpoint; no one-dimensional weak
pointwise-jump theorem is imported.

We also need a geometric observation.

\begin{lemma}[Thin radial layers]\label{lem:thinlayer}
Fix $0<c_-<c_+<\infty$.  Let $\Theta\in L^\infty(\Sph^{n-1})$, let
$R>0$, let $t\in\mathbb N_0$, put $q=R2^{-t}$, and let $\Lambda\ge0$.
Suppose $d_U\in L^1$ are supported on pairwise disjoint
$q$-cubes and satisfy $\norm{d_U}_1\le\Lambda|U|$.  If
$E\subset[c_-R,c_+R]$ is a
union of $M_E$ intervals of total length $|E|$, then
\begin{align}
 \left\|K_\Theta\mathbf1_{\{|z|\in E\}}*\sum_Ud_U\right\|_\infty
 &\le C_{n,c_-,c_+}\Lambda\norm\Theta_\infty
 \left(\frac{|E|}{R}+M_E2^{-t}\right),\label{eq:thinLinf}\\
 \left\|K_\Theta\mathbf1_{\{|z|\in E\}}*\sum_Ud_U\right\|_1
 &\le C_{n,c_-,c_+}\norm\Theta_1\frac{|E|}{R}
 \sum_U\norm{d_U}_1.
 \label{eq:thinL1}
\end{align}
Here and in the proof the implicit constants depend only on $n,c_-,c_+$.
\end{lemma}

\begin{proof}
For fixed $x$, every contributing cube lies in a $Cq$-neighborhood of the
translated radial set.  That neighborhood has volume at most
$CR^{n-1}(|E|+M_Eq)$.  Disjointness and the mass-density bound prove
\eqref{eq:thinLinf}.  Equation \eqref{eq:thinL1} is Young's inequality and
$\int_Edr/r\le C|E|/R$.
\end{proof}

\begin{lemma}[Variational Rademacher--Menshov]\label{lem:RM}
For a finite sequence $g_1,\ldots,g_M$ in $L^2$,
\begin{equation}\label{eq:RM}
 \left\|V^2\left(\sum_{\ell\le m}g_\ell:0\le m\le M\right)\right\|_2
 \le C\log(2+M)
 \sup_{\varepsilon_\ell\in\{-1,0,1\}}
 \left\|\sum_\ell\varepsilon_\ell g_\ell\right\|_2.
\end{equation}
\end{lemma}

\begin{proof}
Pad the sequence with zeros to length $2^L$ and let $\mathcal D_r$ be the
fixed partition of $\{1,\ldots,2^L\}$ into dyadic index intervals of
length $2^r$, $0\le r\le L$.  Put $G_I=\sum_{\ell\in I}g_\ell$.  The
canonical dyadic decomposition of an arbitrary index interval uses at most
two members of each $\mathcal D_r$.  If a collection of index intervals is
pairwise disjoint, then all dyadic pieces appearing in their canonical
decompositions are pairwise disjoint.  Consequently, for every fixed
variation partition, Minkowski in the outer finite $\ell^2$ sum gives the
pointwise estimate
\[
 \left(\sum_j\left|\sum_{\ell\in J_j}g_\ell\right|^2\right)^{1/2}
 \le\sqrt2\sum_{r=0}^L
 \left(\sum_{I\in\mathcal D_r}|G_I|^2\right)^{1/2}.
\]
The right side is independent of the variation partition, so taking the
supremum proves the deterministic domination
\begin{equation}\label{eq:RM-deterministic}
 V^2\left(\sum_{\ell\le m}g_\ell:0\le m\le2^L\right)
 \le\sqrt2\sum_{r=0}^L
 \left(\sum_{I\in\mathcal D_r}|G_I|^2\right)^{1/2}.
\end{equation}
For a fixed level, independent Rademacher signs $\epsilon_I$ and Fubini
give
\[
 \left\|\left(\sum_{I\in\mathcal D_r}|G_I|^2\right)^{1/2}\right\|_2^2
 =\mathbb E\left\|\sum_{I\in\mathcal D_r}\epsilon_IG_I\right\|_2^2.
\]
Each randomized block sum is $\sum_\ell\epsilon_\ell g_\ell$ with
$\epsilon_\ell\in\{-1,1\}$ fixed on its dyadic block; hence its norm is
bounded by the supremum in \eqref{eq:RM}.  Take $L^2$ norms in
\eqref{eq:RM-deterministic} and sum the $L+1\le C\log(2+M)$ levels.
\end{proof}

We shall use the following elementary completion fact whenever a finite
Fourier expansion is removed.  Its anchored formulation avoids treating
$V^r$ as a norm modulo constants and also avoids any appeal to an
uncountable-parameter Bochner space.

\begin{lemma}[Anchored variation completion]\label{lem:anchored-variation-completion}
Let $I=[a,b]$, let $1\le p,r<\infty$, and let
$H_N=(H_{N,t})_{t\in I}$ be paths of measurable functions satisfying
$H_{N,b}=0$.  Assume that the variation functions of the paths and their
differences are measurable; this holds, in particular, when their
variation is detected by one fixed countable dense subset of $I$, as in
all applications below.  Suppose
\begin{equation}\label{eq:anchored-variation-cauchy}
 \norm{V^r(H_N-H_M:t\in I)}_p\longrightarrow0
 \qquad(N,M\to\infty).
\end{equation}
Then there is an anchored path $H=(H_t)_{t\in I}$ such that
\begin{equation}\label{eq:anchored-variation-limit}
 \norm{V^r(H_N-H:t\in I)}_p\longrightarrow0.
\end{equation}
After passage to a subsequence, $H_{N,t}(x)\to H_t(x)$ uniformly in
$t\in I$ for almost every $x$.  Moreover,
\begin{equation}\label{eq:anchored-variation-fatou}
 \norm{V^r(H:t\in I)}_p
 \le\liminf_{N\to\infty}\norm{V^r(H_N:t\in I)}_p.
\end{equation}
Consequently, if $P_N=G_N+E_N$, both $(G_N)$ and $(E_N)$ satisfy the
corresponding anchored Cauchy hypotheses (possibly with different
$p,r$), and $P_{N,t}\to P_t$ for every $t$ almost everywhere, then the
limits satisfy $P_t=G_t+E_t$ for every $t$ on one common full-measure
set whenever the convergence of $P_N$ is uniform in $t$ there.
\end{lemma}

\begin{proof}
Choose a subsequence $N_k$ such that
\[
 \sum_k\norm{V^r(H_{N_{k+1}}-H_{N_k}:t\in I)}_p<\infty.
\]
Minkowski's inequality shows that, outside one null set,
\[
 \sum_kV^r(H_{N_{k+1}}-H_{N_k}:t\in I)<\infty.
\]
Since every difference path vanishes at $b$,
\[
 \sup_{t\in I}|H_{N_{k+1},t}-H_{N_k,t}|
 \le V^r(H_{N_{k+1}}-H_{N_k}:t\in I).
\]
  Thus the subsequence converges uniformly in $t$ almost everywhere to an
  anchored path $H$.  We record explicitly the measurability needed below.
  On the same full-measure set put
  \[
   S_k=\sum_{\ell\ge k}
   V^r(H_{N_{\ell+1}}-H_{N_\ell}:t\in I).
  \]
  The functions $S_k$ are measurable, finite almost everywhere, and decrease
  to zero.  Applying the variation triangle inequality first to each finite
  tail and then passing to its uniform limit gives
  \[
   V^r(H_{N_k}-H:t\in I)\le S_k.
  \]
  Consequently, for every fixed $N$,
  \[
   \left|V^r(H_N-H:t\in I)
       -V^r(H_N-H_{N_k}:t\in I)\right|\le S_k\longrightarrow0.
  \]
  Hence $V^r(H_N-H:t\in I)$ is the pointwise limit of measurable
  functions and is measurable.  For every finite parameter partition, its
  increment sum passes to the pointwise limit.  Taking the supremum over
  finite partitions and applying Fatou now gives
\[
 \norm{V^r(H_N-H:t\in I)}_p
 \le\liminf_{k\to\infty}
 \norm{V^r(H_N-H_{N_k}:t\in I)}_p.
\]
Indeed, for every $\delta>0$ the Cauchy hypothesis supplies $N_0$ such
that
\[
 \norm{V^r(H_N-H_M:t\in I)}_p<\delta
 \quad\text{whenever }N,M\ge N_0.
\]
For $N\ge N_0$, the preceding inequality with $M=N_k$ and then
$k\to\infty$ yields
$\norm{V^r(H_N-H:t\in I)}_p\le\delta$.  This proves
\eqref{eq:anchored-variation-limit} for the full sequence.  Pointwise,
the variation triangle inequality gives
\[
 \big|V^r(H_N:t\in I)-V^r(H:t\in I)\big|
 \le V^r(H_N-H:t\in I).
\]
The reverse triangle inequality in $L^p$ and
\eqref{eq:anchored-variation-limit} therefore give convergence of the
two variation norms, which is stronger than
\eqref{eq:anchored-variation-fatou}.  For the last assertion, choose one
diagonal subsequence along which the required summability holds for both
$G_N$ and $E_N$; passing the finite identities to their common uniform
pointwise limits gives the claimed identity.
\end{proof}

Because Seeger's microlocal kernels $\Gamma_j^s$ have noncompact spatial
support, localization and ranking are carried out with the original
compactly supported kernels $H_j$.

\begin{lemma}[Fixed-sign full-kernel estimate]\label{lem:fullkernel}
Let $n\ge1$, $N=4n+1$, and let $H_j$ be supported where
$2^{j-2}\le|x|\le2^{j+2}$, radially $C^N$ as above, and satisfy
\eqref{eq:MN}.  Let $\cQ$ be a
finite family of cubes with disjoint interiors, and suppose
\[
 \supp b_Q\subset Q,\qquad \int b_Q=0,\qquad
 \norm{b_Q}_1\le\Lambda|Q|.
\]
For $s\in\mathbb N$, $s\ge5$, pair $H_j$ with the cubes of side
$2^{j-s}$.  Then, for every fixed choice $|c_Q|\le1$,
\begin{equation}\label{eq:fullkernel}
 \left\|\sum_jH_j*\!\sum_{\ell(Q)=2^{j-s}}c_Qb_Q\right\|_2^2
 \le C_nM_N^2\,2^{-s/24}\Lambda\sum_Q\norm{b_Q}_1.
\end{equation}
\end{lemma}

\begin{proof}
By homogeneity assume $M_N=1$, and put
$\Sigma=\sum_Q\norm{b_Q}_1$.  If $\Sigma=0$ there is nothing to prove.
Write
\[
 F=\sum_jH_j*\!\sum_{\ell(Q)=2^{j-s}}c_Qb_Q.
\]
We first establish an $L^3$ bound valid uniformly at every scale and in
every dimension.  Choose a fixed
smooth radial cutoff equal to one on
$\{2^{-2}\le|x|\le2^2\}$ and dilate a sufficiently large multiple of
$2^{-jn}$ times that cutoff.  The resulting nonnegative radial majorants
$\nu_j$ satisfy
\[
 |H_j|\le\nu_j,\qquad \norm{\nu_j}_1\le C_n,\qquad
 \norm{\nabla\nu_j}_\infty\le C_n2^{-j(n+1)}.
\]
Set
\[
 \mathcal P=\sum_j\nu_j*\!\sum_{\ell(Q)=2^{j-s}}|b_Q|.
\]
Then $|F|\le\mathcal P$ and
$\norm{\mathcal P}_1\le C_n\Sigma$.  We claim that
\begin{equation}\label{eq:P-BMO}
 \norm{\mathcal P}_{\mathrm{BMO}}\le C_n\Lambda.
\end{equation}
Fix a cube $R_0$ with center $x_0$.  For $2^j<\ell(R_0)$, compact support
of $\nu_j$ shows that every $Q$ contributing on $R_0$ lies in a fixed
enlargement of $R_0$.  Since the cubes in $\cQ$ are disjoint across all
scales and $\norm{b_Q}_1\le\Lambda|Q|$, integration over $R_0$ bounds the
total contribution of these scales by $C_n\Lambda|R_0|$.

For $2^j\ge\ell(R_0)$ subtract the value at $x_0$.  Every cube
contributing to either convolution at $x\in R_0$ or at $x_0$ lies in a
ball of radius $C_n2^j$ about $x_0$.  Disjointness and the density bound
therefore give total relevant mass at most $C_n\Lambda2^{jn}$.  The
mean-value theorem and the gradient estimate yield
\[
 \sup_{x\in R_0}\left|
 \nu_j*\!\sum_{\ell(Q)=2^{j-s}}|b_Q|(x)
 -\nu_j*\!\sum_{\ell(Q)=2^{j-s}}|b_Q|(x_0)\right|
 \le C_n\Lambda\ell(R_0)2^{-j}.
\]
  Summing in $j$ proves \eqref{eq:P-BMO}.  We next justify the endpoint
  $L^1$--BMO estimate used below directly.  The Fefferman--Stein
  interpolation identity \cite[Section~III.5, equation~(5.1)]{FS72}
  explicitly includes the endpoint choice $p=\infty$ and yields the same
  estimate; the following distributional argument makes that endpoint
  instance and its constants transparent.  Let
  $h\in L^1(\mathbb R^n)\cap\mathrm{BMO}(\mathbb R^n)$, and put
  $A=\norm h_1$ and $B=\norm h_{\mathrm{BMO}}$.  If $B=0$, then $h$ is
  almost everywhere constant, and its integrability forces $h=0$.  Suppose
  $B>0$ and apply the dyadic Calder\'on--Zygmund decomposition to $|h|$ at
  height $B$.  Its maximal cubes $Q$ satisfy
  \[
   \sum_Q|Q|\le A/B,
   \qquad B<\frac1{|Q|}\int_Q|h|\le2^nB,
  \]
  while $|h|\le B$ almost everywhere off their union.  In particular
  writing $h_Q:=|Q|^{-1}\int_Qh$, we have $|h_Q|\le2^nB$.  The
  John--Nirenberg inequality in
  \cite[Section~II.1]{FS72}, applied on each $Q$, therefore gives constants
  $c_n,C_n>0$ such that, whenever $\lambda\ge2^{n+1}B$,
  \[
   \big|\{x:|h(x)|>\lambda\}\big|
   \le C_n\frac AB\exp(-c_n\lambda/B).
  \]
  For $0<\lambda<2^{n+1}B$, Chebyshev's inequality gives the complementary
  bound $|\{|h|>\lambda\}|\le A/\lambda$.  The layer-cake formula, split at
  $2^{n+1}B$, now yields
  \[
   \int_{\mathbb R^n}|h|^3
   =3\int_0^\infty\lambda^2|\{|h|>\lambda\}|\,d\lambda
   \le C_nAB^2.
  \]
  Thus we have proved the form used below:
\begin{equation}\label{eq:L1BMO}
 \norm h_3^3\le C_n\norm h_1\norm h_{\mathrm{BMO}}^2,
 \qquad h\in L^1\cap\mathrm{BMO}.
\end{equation}
Consequently,
\begin{equation}\label{eq:full-L3}
 \norm F_3^3\le\norm{\mathcal P}_3^3
 \le C_n\Lambda^2\Sigma.
\end{equation}

Suppose first that $n=1$.  Let $x_Q$ be the center of $Q$.  Cancellation
and the $W^{1,1}$ translation estimate give
\begin{align*}
 \norm{H_j*(c_Qb_Q)}_1
 &=\left\|\int_Qc_Qb_Q(y)
   [H_j(\,\cdot-y)-H_j(\,\cdot-x_Q)]\,dy\right\|_1\\
 &\le C\ell(Q)\norm{H_j'}_1\norm{b_Q}_1.
\end{align*}
Since $H_j$ is supported where $|x|\simeq2^j$, the radial derivative
bound in \eqref{eq:MN} gives
$\norm{H_j'}_1\le C2^{-j}$.  For
$\ell(Q)=2^{j-s}$ it follows that
\begin{equation}\label{eq:fullkernel-one-L1}
 \norm F_1\le C2^{-s}\Sigma.
\end{equation}
Interpolation between \eqref{eq:fullkernel-one-L1} and
\eqref{eq:full-L3} now gives
\[
 \norm F_2^2
 \le\norm F_1^{1/2}\norm F_3^{3/2}
 \le C2^{-s/2}\Lambda\Sigma
 \le C2^{-s/24}\Lambda\Sigma.
\]
This proves \eqref{eq:fullkernel} in dimension one.

Assume henceforth that $n\ge2$.  Apply Seeger's decomposition with
$\kappa=1/2$ and write $F=A+R$.  Equation~\eqref{eq:seeger1} gives
\begin{equation}\label{eq:full-A}
 \norm A_2^2\le C_n2^{-s/2}\Lambda\Sigma.
\end{equation}
Summing \eqref{eq:seeger2} with $\epsilon=1/4$ gives
\begin{equation}\label{eq:full-R}
 \norm R_1\le C_n2^{-s/4}\Sigma,
\end{equation}
because $n-(4n+1)/4=-1/4$.  The coefficients $c_Q$ preserve
cancellation and the mass-density estimate.

Put $u=\Lambda2^{-s/8}$ and $E=\{|R|>u\}$.  By
\eqref{eq:full-R},
\[
 |E|\le C_n\Lambda^{-1}2^{-s/8}\Sigma.
\]
On $E^c$, equations \eqref{eq:full-A} and \eqref{eq:full-R} give
\[
 \int_{E^c}|F|^2\le2\norm A_2^2+2u\norm R_1
 \le C_n2^{-3s/8}\Lambda\Sigma.
\]
On $E$, H\"older's inequality and \eqref{eq:full-L3} give
\[
 \int_E|F|^2\le\norm F_3^2|E|^{1/3}
 \le C_n2^{-s/24}\Lambda\Sigma.
\]
This proves \eqref{eq:fullkernel} for $M_N=1$ in every dimension.
Rescaling the kernels restores the factor $M_N^2$.
\end{proof}

\begin{proposition}[Selector-free terminal packet trajectory]
\label{prop:terminalpacket}
There are $t_0\in\mathbb N$ and $\delta_T>0$, depending only on $n$, such
that the following holds.  Let $\Theta\in L^\infty(\Sph^{n-1})$, let
$R>0$, $t\in\mathbb N$, $q=R2^{-t}$, $t\ge t_0$, and let
$\{U\}$ be pairwise disjoint $q$-cubes.  Let $\Lambda\ge0$ and suppose
\[
 \supp d_U\subset U,
 \quad \int d_U=0,
 \quad \norm{d_U}_1\le\Lambda|U|,
 \quad M=\sum_U\norm{d_U}_1.
\]
For fixed complex coefficients $|a_U|\le1$, define
\[
 F_\varepsilon(x)=\int_{\varepsilon<|x-y|\le2R}
 K_\Theta(x-y)\sum_Ua_Ud_U(y)dy,
 \qquad R\le\varepsilon\le2R.
\]
Then
\begin{equation}\label{eq:terminalpacket}
 \norm{V^2(F_\varepsilon:R\le\varepsilon\le2R)}_2^2
 \le C\norm\Theta_\infty^2 2^{-\delta_Tt}\Lambda M.
\end{equation}
The same estimate holds after multiplication of the entire output path by
$e^{i\xi\cdot x}$, for any fixed $\xi\in\R^n$.  It also holds after
replacing each $d_U(y)$ by $e^{i\eta_U\cdot y}d_U(y)$, with arbitrary
$\eta_U\in\R^n$, provided
$\int e^{i\eta_U\cdot y}d_U(y)\,dy=0$ for every $U$.  In the
phase-adapted application below, the input modulation is frozen packet by
packet and the resulting error is estimated separately.
\end{proposition}

\begin{proof}
We first justify the reduction to $R=1$, including the normalization of
the density and mass.  Put $x=RX$, $y=RY$,
$\widetilde U=R^{-1}U$, and $\widetilde d_U(Y)=d_U(RY)$.  Homogeneity gives
$F_\varepsilon(RX)=\widetilde F_{\varepsilon/R}(X)$, while
\[
 \norm{\widetilde d_U}_1=R^{-n}\norm{d_U}_1,
 \qquad
 \norm{\widetilde d_U}_1\le\Lambda|\widetilde U|,
 \qquad
 \widetilde M:=\sum_U\norm{\widetilde d_U}_1=R^{-n}M.
\]
Thus the unit-scale squared $L^2_X$ estimate has right side
$C\norm\Theta_\infty^2 2^{-\delta_Tt}\Lambda R^{-n}M$.  Returning to
$x=RX$ multiplies the squared output norm by $R^n$, exactly cancelling
the input mass factor $R^{-n}$.  It is therefore enough to prove the
estimate at $R=1$.  Throughout the unit-scale proof write
$K_{|\Theta|}(z)=|\Theta(z/|z|)|/|z|^n$.  First consider $n=1$.
For a radially smoothed
one-shell kernel, we use the following fixed convention.  Choose a
nonnegative function $\phi\in C_c^\infty((-1/10,1/10))$ with
$\int_\R\phi=1$, put $\phi_h(r)=h^{-1}\phi(r/h)$, and extend every
radial step profile below by zero to $\R$ before setting $a_h=a*\phi_h$.
Then
\[
 \norm{a_h}_\infty\le\norm a_\infty,\qquad
 \norm{\partial_r^\ell a_h}_\infty
 \le C_\ell h^{-\ell}\norm a_\infty.
\]
If $a$ is constant on $M_t$ intervals, then $a-a_h$ is supported in the
$O(h)$-neighborhood of its jump set, hence in $O(M_t)$ intervals of total
length $O(M_th)$.  Every mollification in this proof is understood in this
sense.

For a radially smoothed
one-shell kernel $H$ at relative scale $h$, cancellation gives
\[
 H*d_U(x)=\int_U[H(x-y)-H(x-c_U)]d_U(y)dy.
\]
Since $\norm{\tau_zH-H}_1\le |z|\norm{H'}_1$ and
$\norm{H'}_1\le C\norm\Theta_\infty h^{-1}$, summing over $U$ yields
\[
 \left\|H*\sum_Ua_Ud_U\right\|_1
 \le C\norm\Theta_\infty2^{-t}h^{-1}M.
\]
Lemma~\ref{lem:thinlayer} gives the simultaneous $L^\infty$ bound
$C\Lambda\norm\Theta_\infty$; hence the squared $L^2$ norm is at most
$C\norm\Theta_\infty^2\Lambda M2^{-t}h^{-1}$.
Take $\eta=1/8$ and $\gamma=1/16$.  Partition $[1,2]$ into
$M_t\simeq2^{\gamma t}$ fixed cells $I_\ell$ of length
$\delta=2^{-\gamma t}$, and take $h=2^{-\eta t}$.  Put
\[
 G_\ell=K_\Theta\mathbf1_{\{|z|\in I_\ell\}}
  *\sum_Ua_Ud_U.
\]
Inserting all cell endpoints into an arbitrary variation partition gives
the pointwise deterministic inequality
\begin{equation}\label{eq:griddecomp-one}
 V^2(F_\varepsilon:1\le\varepsilon\le2)
 \le C V^2\left(\sum_{\ell\le m}G_\ell:0\le m\le M_t\right)
 +C\left(\sum_\ell
     V^2(F_\varepsilon:\varepsilon\in I_\ell)^2\right)^{1/2}.
\end{equation}
Reversing the order of the partial sums, if required by the convention
$\varepsilon<|z|$, does not change their variation.  Thus every
$x$-dependent choice of variation intervals is confined to the second
term; the first term is one fixed finite sequence.  Within a cell
$V^2\le V^1$.  If $W_\ell$ is the positive convolution over that cell,
Lemma~\ref{lem:thinlayer} gives
\[
 \sup_\ell\norm{W_\ell}_\infty
 \le C\Lambda\norm\Theta_\infty(\delta+2^{-t}),
 \qquad
 \sum_\ell\norm{W_\ell}_1\le C\norm\Theta_\infty M.
\]
Consequently,
\[
 \int_{\R}\sum_\ell
 V^2(F_\varepsilon:\varepsilon\in I_\ell)^2\,dx
 \le C\Lambda\norm\Theta_\infty^2M2^{-\gamma t}.
\]

For a fixed choice of signs on the cells, write the resulting radial step
function as
\[
 a(r)=\sum_{\ell=1}^{M_t}\varepsilon_\ell\mathbf1_{I_\ell}(r),
 \qquad \varepsilon_\ell\in\{-1,0,1\},
\]
with $a$ extended by zero outside $[1,2]$, and smooth it at scale $h$.
Let $\mathcal D_a$ consist of the endpoints $1,2$ together with all
interior jump points of $a$, and set
\[
 E_{\rm bdry}=\{r\in\mathbb R:\operatorname{dist}(r,\mathcal D_a)\le Ch\},
\]
where $C$ is fixed large enough in terms of $\supp\phi$.  Then
\begin{equation}\label{eq:terminal-boundary-profile-one}
 \#\text{components}(E_{\rm bdry})\le C M_t,
 \qquad |E_{\rm bdry}|\le C M_th,
 \qquad \supp(a-a_h)\subset E_{\rm bdry},
 \qquad |a-a_h|\le2\mathbf1_{E_{\rm bdry}}.
\end{equation}
After increasing $t_0$ if necessary, $E_{\rm bdry}\subset[1/2,5/2]$,
so Lemma~\ref{lem:thinlayer} applies with fixed radial constants.
Set $H(z)=K_\Theta(z)a_h(|z|)$.  The cancellation estimate above and the
$L^\infty$ part of Lemma~\ref{lem:thinlayer} give
\[
 \norm{H*\textstyle\sum_Ua_Ud_U}_2^2
 \le C\Lambda\norm\Theta_\infty^2M\,2^{-(1-\eta)t}.
\]
Let $G_{\rm bdry}$ be the output obtained by replacing the radial profile
by $a-a_h$.  The last inequality in
\eqref{eq:terminal-boundary-profile-one} gives the pointwise positive-kernel
domination
\begin{equation}\label{eq:terminal-boundary-domination-one}
 |G_{\rm bdry}(x)|
 \le 2\left(K_{|\Theta|}\mathbf1_{\{|z|\in E_{\rm bdry}\}}
       *\sum_U|d_U|\right)(x).
\end{equation}
The functions $|d_U|$ have the same supports and $L^1$ bounds as $d_U$,
so both estimates of Lemma~\ref{lem:thinlayer} apply to the right side of
\eqref{eq:terminal-boundary-domination-one}.  Together with
$\norm G_2^2\le\norm G_\infty\norm G_1$, they give
\[
 \norm{G_{\rm bdry}}_2^2
 \le C\Lambda\norm\Theta_\infty^2M
 \bigl(M_th+M_t2^{-t}\bigr)M_th.
\]
With the stated values of $\eta$ and $\gamma$, the last factor is bounded
by
\[
 C\bigl(2^{-2(\eta-\gamma)t}
       +2^{-(1+\eta-2\gamma)t}\bigr)
 =C\bigl(2^{-t/8}+2^{-t}\bigr).
\]
The smoothed signed sum decays with exponent $1-\eta=7/8$, the boundary
term with exponent at least $1/8$, and the within-cell term with exponent
$\gamma=1/16$.  The signs $\varepsilon_\ell$ above are
fixed before evaluation in $x$ and are independent of the already fixed
packet coefficients $a_U$.  Hence Lemma~\ref{lem:RM} applies to the first
term of \eqref{eq:griddecomp-one}.  Since
\begin{equation}\label{eq:terminal-RM-log-cost}
 \log^2(2+M_t)\lesssim(1+\gamma t)^2\lesssim t^2,
 \qquad
 t^2 2^{-ct}\le C_{c,c'}2^{-c't}\quad(0<c'<c),
\end{equation}
its squared cost is absorbed by reducing the positive decay exponent.  In
particular, $t^2 2^{-t/8}\lesssim2^{-t/16}$; together with the within-cell
estimate this proves \eqref{eq:terminalpacket} in dimension one with the
explicit safe choice $\delta_T=1/32$.

Now assume $n\ge2$ and choose $N=4n+1$, $\eta=(16N)^{-1}$ and
$\gamma=\eta/2$.  Partition
$[1,2]$ into $M_t\simeq2^{\gamma t}$ fixed cells $I_\ell$.  Inserting the
cell endpoints into a variation partition gives
\begin{align}\label{eq:griddecomp}
 V^2(F_\varepsilon:1\le\varepsilon\le2)
 \le C V^2\left(\sum_{\ell\le m}G_\ell:0\le m\le M_t\right)
 +C\left(\sum_\ell V^2(F_\varepsilon:\varepsilon\in I_\ell)^2\right)^{1/2},
\end{align}
where $G_\ell=K_\Theta\mathbf1_{\{|z|\in I_\ell\}}*\sum_Ua_Ud_U$.

Inside one cell, $V^2\le V^1$ and Lemma~\ref{lem:thinlayer} gives
\[
 \int\sum_\ell V^2(F:I_\ell)^2
 \le C\norm\Theta_\infty^2\Lambda M\,2^{-\gamma t}.
\]
Indeed, the supremum of one absolute cell convolution is
$C\Lambda\norm\Theta_\infty(2^{-\gamma t}+2^{-t})$, while the sum of its
$L^1$ norms is $C\norm\Theta_1M$.

It remains to estimate the grid partial sums.  Fix signs
$\varepsilon_\ell\in\{-1,0,1\}$ and let $a(r)=\sum\varepsilon_\ell
\mathbf1_{I_\ell}(r)$.  Mollify $a$ at radial scale
$h=2^{-\eta t}$, using the zero-extension convention above, and set
\[
 a_h=a*\phi_h,\qquad H(z)=K_\Theta(z)a_h(|z|).
\]
The kernel $H$ satisfies
$M_N(H)\le C\norm\Theta_\infty h^{-N}$.  Moreover,
$|a|\le1$ and $\supp a\subset[1,2]$ imply
\[
 |a_h(r)|\le \mathbf1_{[1-Ch,\,2+Ch]}(r).
\]
Thus the possibly signed smooth profile is pointwise dominated by one fixed
enlarged annulus.  In particular, Lemma~\ref{lem:thinlayer}, applied to that
annulus, gives the uniform bound used below.  Apply
\eqref{eq:seeger1}--\eqref{eq:seeger2} with
$\kappa=1/2$, $\epsilon=1/4$, embedding the one-shell kernel into
Seeger's family by taking $H_0=H$ and $H_j=0$ for $j\ne0$.  The
microlocal part $A$ obeys
\[
 \norm A_2^2\le C\norm\Theta_\infty^2 2^{-t/2}\Lambda M,
\]
while the remainder obeys
\[
 \norm R_1\le C\norm\Theta_\infty
 2^{-at}M,
 \qquad a:=\frac14-\eta N=\frac3{16}.
\]
The full smoothed convolution $S=A+R$ therefore has
$\norm S_\infty\le C\Lambda\norm\Theta_\infty$.
$\Lambda\norm\Theta_\infty=0$ only if either $K_\Theta=0$ almost
everywhere or $\Lambda=0$, in which case
$\norm{d_U}_1\le\Lambda|U|$ forces every $d_U=0$ almost everywhere.
Thus all the trajectories in the proposition vanish in this case.  In
what follows assume $\Lambda\norm\Theta_\infty>0$, put
\[
 b:=\frac18-\frac{\eta N}{2}=\frac3{32},
 \qquad u=\Lambda\norm\Theta_\infty2^{-bt}>0,
\]
and let $E=\{|R|>u\}$.  On $E^c$, use
$|S|^2\le2|A|^2+2u|R|$; on $E$, use
$\int_E|S|^2\le\norm S_\infty^2|E|$ and
$|E|\le u^{-1}\norm R_1$.  This yields
\begin{equation}\label{eq:fixedsign}
 \norm S_2^2\le C\norm\Theta_\infty^2
 2^{-3t/32}\Lambda M.
\end{equation}
Indeed, the three contributions have respective decay exponents
$1/2$, $a+b=9/32$, and $a-b=3/32$.  This records both square-root losses
and fixes the exponent in place of an unspecified $c_1$.

Let $\mathcal D_a$ consist of $1,2$ and all interior jump points of $a$,
and put
\[
 E_{\rm bdry}=\{r\in\mathbb R:\operatorname{dist}(r,\mathcal D_a)\le Ch\}.
\]
After increasing the fixed $C$ according to $\supp\phi$, one has
\begin{equation}\label{eq:terminal-boundary-profile-high}
 \#\text{components}(E_{\rm bdry})\le CM_t,
 \quad |E_{\rm bdry}|\le CM_th,
 \quad \supp(a-a_h)\subset E_{\rm bdry},
 \quad |a-a_h|\le2\mathbf1_{E_{\rm bdry}}.
\end{equation}
After increasing $t_0$ if necessary, $E_{\rm bdry}\subset[1/2,5/2]$,
so the radial constants in Lemma~\ref{lem:thinlayer} remain fixed.
Let $G_{\rm bdry}$ denote the output generated by $a-a_h$.  Positivity
and \eqref{eq:terminal-boundary-profile-high} give, pointwise,
\begin{equation}\label{eq:terminal-boundary-domination-high}
 |G_{\rm bdry}(x)|
 \le2\left(K_{|\Theta|}\mathbf1_{\{|z|\in E_{\rm bdry}\}}
       *\sum_U|d_U|\right)(x).
\end{equation}
Both bounds in Lemma~\ref{lem:thinlayer}, applied to the right side of
\eqref{eq:terminal-boundary-domination-high}, and
$\norm{G_{\rm bdry}}_2^2\le
 \norm{G_{\rm bdry}}_\infty\norm{G_{\rm bdry}}_1$ give
\begin{align*}
 \norm{G_{\rm bdry}}_2^2
 &\le C\norm\Theta_\infty^2\Lambda M
       (M_th+M_t2^{-t})M_th\\
 &\le C\norm\Theta_\infty^2\Lambda M
 \left(2^{-2(\eta-\gamma)t}
       +2^{-(1+\eta-2\gamma)t}\right)\\
 &=C\norm\Theta_\infty^2\Lambda M
   \left(2^{-\eta t}+2^{-t}\right),
\end{align*}
where the last identity uses $\gamma=\eta/2$.  Therefore the right side of
\eqref{eq:fixedsign} holds for the unsmoothed signed sum with exponent
$\min(3/32,\eta)=\eta$, uniformly in the signs.  Lemma~\ref{lem:RM}
contributes a
factor $O(t^2)$ after squaring, by
\eqref{eq:terminal-RM-log-cost}; absorb it by reducing the exponent.
Here too the cell signs are fixed globally before evaluation in $x$ and
are independent of the fixed packet coefficients $a_U$.  Together with
the within-cell exponent $\gamma=\eta/2$ and
\eqref{eq:griddecomp}, this proves \eqref{eq:terminalpacket}.  A single
explicit choice valid in this branch is
\[
 \delta_T=\frac{\eta}{4}=\frac1{64(4n+1)}\qquad(n\ge2),
\]
since $t^2 2^{-\eta t}\lesssim2^{-\eta t/2}$ and
$\delta_T<\gamma$.  Together with the choice $\delta_T=1/32$ above, this
also makes the dependence of $\delta_T$ only on $n$ manifest.
\end{proof}

\section{An anchored phase--packet induction}
\label{sec:anchored-induction}

After removal of the parent-dependent input modulation, the
phase-adapted decomposition recovers a mean-zero packet structure.  We
formulate the corresponding estimate in this anchored normalization.

Fix a real polynomial $P\in\mathcal P_d$ and
$\Theta\in L^\infty(\Sph^{n-1})$.  Let $t\in\mathbb N$, and let $R>0$
and $q=R2^{-t}$ be dyadic.  Let
$\mathcal Q$ be a family of
pairwise disjoint dyadic cubes, each contained in a dyadic $R$-cube, and
for $Q\in\mathcal Q$ let $c_Q$ denote its center.
Suppose $d_U$ are mean-zero descendants, supported on pairwise disjoint
$q$-cubes $U\subset Q$, with
\begin{equation}\label{eq:packetdensity}
 \norm{d_U}_\infty\le\Lambda,
 \qquad M=\sum_U\norm{d_U}_1.
\end{equation}
Define the phase-adapted shell trajectory
\begin{equation}\label{eq:anchoredtraj}
 \mathcal P_\varepsilon(x)=
 \sum_{Q\in\mathcal Q}\sum_{U\subset Q}
 \int_{\varepsilon<|x-y|\le2R}
 e^{i[P(x,y)-P(c_Q,y)]}K_\Theta(x-y)d_U(y)dy,
 \quad R\le\varepsilon\le2R.
\end{equation}

We fix the Fourier localization once, rather than separately for each
block.  After the normalization by an interacting $R$-block used below,
$Y$ and every parent center $\xi_Q$ lie in one fixed cube, while $X$ lies
in one of only $O_n(1)$ unit cubes whose difference from that cube meets
$\{1<|Z|\le2\}$.  Choose structural boxes
$\mathfrak X,\mathfrak Y$ containing all these cubes with a fixed margin,
and smooth cutoffs equal to one on the boxes.  The anchored flag contains
only finitely many polynomials (depending on $n,d$); throughout the proof
we use these same boxes and cutoffs and take the maximum of the resulting
Fourier constants over that finite list.

\begin{lemma}[Fourier expansion of a low pivot]\label{lem:fourierpivot}
Let $H$ be one of the fixed anchored flag polynomials, and let $X,Y,\xi$
range in the structural boxes just fixed.  Let
$\chi,\widetilde\chi$ be the corresponding fixed smooth cutoffs.  Then,
for every $L\in\R$,
\[
 \chi(X)\widetilde\chi(Y)
 e^{iL[H(X,Y)-H(\xi,Y)]}
 =\sum_{u,v\in\mathbb Z^n}a_{u,v}(L,\xi)e^{iu\cdot X}e^{iv\cdot Y}.
\]
There exist numbers
$A_{u,v}(L)$ independent of $\xi$ such that
$|a_{u,v}(L,\xi)|\le A_{u,v}(L)$ and
\begin{equation}\label{eq:fourierpivot}
 \sum_{u,v}A_{u,v}(L)(1+|v|)
 \le C(1+|L|)^{C_F}.
\end{equation}
The constants are uniform in $L$ and $\xi$ apart from the displayed
power and, with the preceding convention, depend only on $n,d$.
\end{lemma}

\begin{proof}
Extend the cutoff function periodically on a fixed torus.  Repeated
integration by parts in $X$ and $Y$ gives, uniformly in $\xi$,
$|a_{u,v}|\le C_N(1+|L|)^N(1+|u|+|v|)^{-N}$.  Take the right side as
$A_{u,v}(L)$, choose $N>2n+2$, and sum (increasing $C_F$ if needed).
\end{proof}

\begin{proposition}[Anchored phase--packet principle]\label{prop:anchoredpacket}
There are $c_P,c_E>0$ and $t_0$, depending only on $n,d$, such that every
trajectory \eqref{eq:anchoredtraj} under the finite-packet setup above,
with $P\in\mathcal P_d$, $\Theta\in L^\infty(\Sph^{n-1})$, and
$t\ge t_0$, admits a decomposition
$\mathcal P=G+E$ satisfying
\begin{align}
 \norm{V^2(G_\varepsilon:R\le\varepsilon\le2R)}_2
 &\le C\norm\Theta_\infty 2^{-c_Pt}(\Lambda M)^{1/2},
 \label{eq:anchoredG}\\
 \norm{V^1(E_\varepsilon:R\le\varepsilon\le2R)}_1
 &\le C\norm\Theta_1 2^{-c_Et}M.
 \label{eq:anchoredE}
\end{align}
The constants are independent of $R$, $q$, every coefficient of $P$, the
locations of the cubes, and the parent labels $Q$.
\end{proposition}

Proposition~\ref{prop:anchoredpacket} is included as a self-contained
bounded-density model; it is not used in the weak-endpoint induction.  Its
purpose is to isolate the phase--packet mechanism from the additional
concentration difficulty.  The actual Calder\'on--Zygmund bad functions
satisfy only cancellation and an average mass bound.  They are handled
instead by the concentration-stable estimates of
Section~\ref{sec:degree-packets}.

\begin{proof}
Let $\mathcal D_R$ be the dyadic partition into cubes of side $R$, and let
$\mathcal J_R\subset\mathcal D_R$ be the finite family of cubes containing
at least one parent $Q\in\mathcal Q$.  Define the finite set of interacting
block pairs by
\begin{equation}\label{eq:interacting-R-blocks}
 \mathcal A_R=\left\{(I,J)\in\mathcal D_R\times\mathcal J_R:
   \{x-y:x\in I,\ y\in J\}\cap\{z:R<|z|\le2R\}\ne\varnothing
                   \right\}.
\end{equation}
For $(I,J)\in\mathcal A_R$, set
\begin{equation}\label{eq:anchored-block-path}
 \mathcal P^{I,J}_\varepsilon(x)=\mathbf1_I(x)
 \sum_{\substack{Q\in\mathcal Q,\ Q\subset J\\U\subset Q}}
 \int_{\varepsilon<|x-y|\le2R}
 e^{i[P(x,y)-P(c_Q,y)]}K_\Theta(x-y)d_U(y)\,dy.
\end{equation}
The annular support gives, almost everywhere, the exact block decomposition
$\mathcal P_\varepsilon=\sum_{(I,J)\in\mathcal A_R}
\mathcal P^{I,J}_\varepsilon$.  All block sums in this proof are henceforth
restricted to $\mathcal A_R$.

For a fixed $(I,J)\in\mathcal A_R$, let $a$ be the center of $J$ and scale
$x=a+RX$, $y=a+RY$.  Put
\[
 M_{I,J}=\sum_{\substack{Q\in\mathcal Q,\ Q\subset J\\U\subset Q}}
            \norm{d_U}_1.
\]
Because $I$ and $J$ have side $R$ and their difference meets a fixed
annulus of radii $R$ and $2R$, elementary lattice geometry gives the
two-sided adjacency estimate
\begin{equation}\label{eq:block-adjacency}
 \sup_{I\in\mathcal D_R}\#\{J\in\mathcal J_R:(I,J)\in\mathcal A_R\}
 +\sup_{J\in\mathcal J_R}\#\{I\in\mathcal D_R:(I,J)\in\mathcal A_R\}
 \le C_n.
\end{equation}
In particular, each packet mass is counted for only $O_n(1)$ output
blocks, and hence
\begin{equation}\label{eq:block-mass-overlap}
 \sum_{(I,J)\in\mathcal A_R}M_{I,J}\le C_nM.
\end{equation}
More explicitly, if every path $H_{I,J}$ is supported in $I$, the
variation triangle inequality, Cauchy--Schwarz in the at most $C_n$
adjacent cubes $J$, and disjointness of the output cubes give
\begin{align}
 \left\|V^2\left(\sum_{(I,J)\in\mathcal A_R}H_{I,J}\right)\right\|_2^2
 &=\sum_I\left\|V^2\left(
     \sum_{J:(I,J)\in\mathcal A_R}H_{I,J}\right)\right\|_{L^2(I)}^2
 \notag\\
 &\le C_n\sum_{(I,J)\in\mathcal A_R}
       \norm{V^2(H_{I,J})}_2^2.\label{eq:block-V2-overlap}
\end{align}
Similarly,
\begin{equation}\label{eq:block-V1-overlap}
 \left\|V^1\left(\sum_{(I,J)\in\mathcal A_R}H_{I,J}\right)\right\|_1
 \le\sum_{(I,J)\in\mathcal A_R}\norm{V^1(H_{I,J})}_1.
\end{equation}

For $Q\subset J$ and $U\subset Q$, set
\[
 \widetilde U=R^{-1}(U-a),\qquad
 \widetilde d_U(Y)=d_U(a+RY),\qquad
 \widetilde I=R^{-1}(I-a).
\]
Then $\widetilde U$ has side $2^{-t}$ and
\begin{equation}\label{eq:packet-rescaling-data}
 \norm{\widetilde d_U}_\infty=\norm{d_U}_\infty\le\Lambda,
 \quad
 \norm{\widetilde d_U}_1=R^{-n}\norm{d_U}_1,
 \quad
 \norm{\widetilde d_U}_1\le\Lambda|\widetilde U|,
 \quad
 \sum_{\substack{Q\in\mathcal Q,\ Q\subset J\\U\subset Q}}
       \norm{\widetilde d_U}_1=R^{-n}M_{I,J}.
\end{equation}
Homogeneity, $K_\Theta(RZ)R^n=K_\Theta(Z)$, cancels the Jacobian, so the
scaled block trajectory has the unit-annulus kernel.  Since the output
cutoff is independent of $\varepsilon$, for every $r\ge1$ one has
$V^r(\mathbf1_IH_\varepsilon)=\mathbf1_I V^r(H_\varepsilon)$ pointwise.
Thus restricting a unit-scale output to $\widetilde I$ cannot increase its
variation norm.  If $\widetilde G$ and $\widetilde E$ are paths in the
scaled variables, then
\begin{equation}\label{eq:packet-rescaling-output}
 \norm{\mathbf1_I V^2(G)}_{L_x^2}
 =R^{n/2}\norm{\mathbf1_{\widetilde I}V^2(\widetilde G)}_{L_X^2},
 \qquad
 \norm{\mathbf1_I V^1(E)}_{L_x^1}
 =R^n\norm{\mathbf1_{\widetilde I}V^1(\widetilde E)}_{L_X^1}.
\end{equation}
Thus unit-scale factors
$(\Lambda R^{-n}M_{I,J})^{1/2}$ and $R^{-n}M_{I,J}$ become
$(\Lambda M_{I,J})^{1/2}$ and $M_{I,J}$ in physical variables.
Put
\[
 \Phi(X,Y)=P(a+RX,a+RY)-P(a,a+RY).
\]
Then $\Phi(0,Y)=0$.  If $Q\subset J$ and
$\xi_Q=(c_Q-a)/R$, the phase in \eqref{eq:anchoredtraj} is exactly
\begin{equation}\label{eq:blockphase}
 \Phi(X,Y)-\Phi(\xi_Q,Y).
\end{equation}
Write the anchored flag decomposition, retaining its pure output part,
\begin{equation}\label{eq:blockflag}
 \Phi(X,Y)=p(X)+\sum_{\mu=1}^{N_d} C_\mu\widehat H_\mu(X,Y).
\end{equation}
All $C_\mu$ are dimensionless; they may depend arbitrarily on $a$ and $R$.
The list and its length depend only on $n,d$.  The exact phase difference
is
\begin{equation}\label{eq:block-output-scalar}
 \Phi(X,Y)-\Phi(\xi_Q,Y)
 =p(X)-p(\xi_Q)
  +\sum_{\mu=1}^{N_d}C_\mu
   [\widehat H_\mu(X,Y)-\widehat H_\mu(\xi_Q,Y)].
\end{equation}
Hence $e^{ip(X)}$ is one common output modulation, whereas
$s_Q:=e^{-ip(\xi_Q)}$ is a parent-dependent scalar of modulus one.  We
work with the demodulated block, retain $s_Q$ in every packet coefficient
below, and at the end multiply both constructed component paths by
$e^{ip(X)}$.  This last operation preserves pointwise $V^2$, pointwise
$V^1$, and the identity $\mathcal P=G+E$.

If $N_d=0$, the common output modulation is harmless and the scalars
$s_Q$ are admissible fixed packet coefficients in
Proposition~\ref{prop:terminalpacket}; that proposition gives the
conclusion directly.  We therefore assume $N_d\ge1$ below.

Let $\beta_* =\min_\mu\beta_\mu>0$, where the $\beta_\mu$ are from
Corollary~\ref{cor:scalecov}.  Let $C_F$ be from
Lemma~\ref{lem:fourierpivot}.  Choose $B>1+8C_F/\beta_*$, then choose
$\varepsilon_*>0$ so small that
\begin{equation}\label{eq:kappachoice}
 C_FN_d\varepsilon_*<\min(\delta_T/8,1/8),
 \qquad
 \kappa_\mu=\varepsilon_*B^{\mu-N_d}.
\end{equation}
Thus the $\kappa_\mu$ increase with $\mu$ and
\begin{equation}\label{eq:kappahierarchy}
 C_F\sum_{r<\mu}\kappa_r\le\frac14\beta_*\kappa_\mu.
\end{equation}

Starting at the first flag coordinate, call a coordinate low if
$|C_\mu|\le2^{\kappa_\mu t}$.  Expand every low factor
\[
 e^{iC_\mu[\widehat H_\mu(X,Y)-\widehat H_\mu(\xi_Q,Y)]}
\]
by Lemma~\ref{lem:fourierpivot}.  To record the product estimate explicitly,
let $S$ be any set of low positions already expanded and index a term by
tuples $\mathbf u=(u_r)_{r\in S}$ and
$\mathbf v=(v_r)_{r\in S}$.  Its parent-dependent coefficient and combined
frequencies are
 \[
 a_Q(\mathbf u,\mathbf v)
 =s_Q\prod_{r\in S}a^{(r)}_{u_r,v_r}(C_r,\xi_Q),
 \quad U(\mathbf u)=\sum_{r\in S}u_r,
 \quad V(\mathbf v)=\sum_{r\in S}v_r.
\]
With
$A(\mathbf u,\mathbf v)=\prod_{r\in S}A^{(r)}_{u_r,v_r}(C_r)$,
if $A(\mathbf u,\mathbf v)=0$, then
$|a_Q(\mathbf u,\mathbf v)|\le A(\mathbf u,\mathbf v)$ forces
$a_Q(\mathbf u,\mathbf v)=0$ for every $Q$, and the corresponding term is
omitted.  Thus every division by $A(\mathbf u,\mathbf v)$ below is made
only when $A(\mathbf u,\mathbf v)>0$.  Moreover,
the inequality
$1+|\sum_{r\in S}v_r|\le\prod_{r\in S}(1+|v_r|)$ and
Lemma~\ref{lem:fourierpivot} give
\begin{align}
 \sum_{\mathbf u,\mathbf v}A(\mathbf u,\mathbf v)
     (1+|V(\mathbf v)|)
 &\le\prod_{r\in S}\sum_{u_r,v_r}
 A^{(r)}_{u_r,v_r}(C_r)(1+|v_r|)\notag\\
 &\le C\prod_{r\in S}(1+|C_r|)^{C_F}
 \le C2^{C_Ft\sum_{r\in S}\kappa_r}.
 \label{eq:weighted-fourier-product}
\end{align}
The same estimate without the weight controls the ordinary Fourier loss.
Keeping the tuple indexing and applying the triangle inequality yields the
same bound while preserving uniformity in every $\xi_Q$.  In particular,
the product of the expansions preceding position
$\mu$ has coefficient $\ell^1$ norm at most
\begin{equation}\label{eq:fourierloss}
 C2^{C_Ft\sum_{r<\mu}\kappa_r}.
\end{equation}

We now make the limiting convention needed for the nonlinear variation
norms explicit.  No infinite Fourier series is inserted directly into a
$V^2$ or $V^1$ norm.  For $N\in\mathbb N$, whenever the factor at a low
position $r$ is expanded, first retain only
\[
 \mathfrak F_N=\{(u_r,v_r)\in\mathbb Z^n\times\mathbb Z^n:
                    |u_r|+|v_r|\le N\}.
\]
For a set $S$ of expanded positions, put
\begin{equation}\label{eq:fourier-tail-majorant}
 \omega_N(S)=
 \sum_{\substack{\mathbf u,\mathbf v:\
       \max_{r\in S}(|u_r|+|v_r|)>N}}
 A(\mathbf u,\mathbf v)(1+|V(\mathbf v)|),
 \qquad \omega_N(\varnothing)=0.
\end{equation}
The weighted $\ell^1$ estimate
\eqref{eq:weighted-fourier-product} implies
$\omega_N(S)\to0$.  All component estimates below are first applied to
these finite mode sums; their passage to the full Fourier expansions is
carried out at the end of the proof.

Suppose first that $\mu$ is the first high coordinate.  After fixing a
Fourier mode from the preceding low factors, the remaining common operator
phase is
$\sum_{r\ge\mu}C_r\widehat H_r(X,Y)$; its first flag coordinate is
$C_\mu$.  The factor obtained by evaluating the remaining phase at
$X=\xi_Q$ is a unimodular input multiplier on the descendants of $Q$.
After division by $A(\mathbf u,\mathbf v)$, the parent-dependent
coefficients have modulus at most one.  Absorb them and the preceding
unimodular factor into input multipliers
$\vartheta_{Q,\mathbf u,\mathbf v}$.  Disjointness of the descendants and
\eqref{eq:packet-rescaling-data} give, in scaled variables,
\begin{equation}\label{eq:high-branch-input-L2}
 \left\|\sum_{\substack{Q\in\mathcal Q,\ Q\subset J\\U\subset Q}}
          \vartheta_{Q,\mathbf u,\mathbf v}(Y)
          \widetilde d_U(Y)\right\|_2^2
 \le\sum_{\substack{Q\in\mathcal Q,\ Q\subset J\\U\subset Q}}
             \norm{\widetilde d_U}_\infty
             \norm{\widetilde d_U}_1
 \le\Lambda R^{-n}M_{I,J}.
\end{equation}
Thus Corollary~\ref{cor:scalecov} applies to one common operator phase.
The fixed output modulation $e^{iU(\mathbf u)\cdot X}$ does not change
variation.  Summing the majorants produces exactly the loss in
\eqref{eq:fourierloss}, and \eqref{eq:packet-rescaling-output} restores
physical scale.  The scaled truncation interval is $1\le\varepsilon/R\le2$,
so it is the partial annulus of radius $2$ in
Corollary~\ref{cor:scalecov}.  Its exact scale parameter is therefore
\[
 L_\mu=2^{m_\mu}|C_\mu|\asymp_{n,d}|C_\mu|,
\]
not merely $|C_\mu|$.  Consequently we obtain
\[
 C\norm\Theta_\infty (2^{m_\mu}|C_\mu|)^{-\beta_\mu}
 (\Lambda M_{I,J})^{1/2}
 \le C\norm\Theta_\infty2^{-\beta_*\kappa_\mu t}
 (\Lambda M_{I,J})^{1/2}
\]
for the entire short trajectory.  Multiplication by the loss
\eqref{eq:fourierloss} and \eqref{eq:kappahierarchy} leave decay at least
$2^{-\beta_*\kappa_\mu t/2}$.  Equations
\eqref{eq:block-V2-overlap} and \eqref{eq:block-mass-overlap} prove
\eqref{eq:anchoredG} in every high branch, with $E=0$.

If every coordinate is low, expand all factors.  For a combined tuple and
a descendant $U$, let $\zeta_U$ be the center of $\widetilde U$.  Write
\begin{equation}\label{eq:modsplit}
 e^{iV(\mathbf v)\cdot Y}\widetilde d_U(Y)
 =e^{iV(\mathbf v)\cdot\zeta_U}\widetilde d_U(Y)
   +\widetilde r_{U,\mathbf v}(Y),
 \qquad
 \norm{\widetilde r_{U,\mathbf v}}_1
 \le C|V(\mathbf v)|2^{-t}\norm{\widetilde d_U}_1.
\end{equation}
The first term is still mean zero.  For each fixed mode, its coefficients
are fixed before evaluation in $x$.  More precisely, divide the
parent-dependent coefficient by $A(\mathbf u,\mathbf v)$; the resulting
packet coefficients have modulus at most one, so
Proposition~\ref{prop:terminalpacket} applies to the scaled unit-annulus
packets.  Equation~\eqref{eq:weighted-fourier-product} makes the summation
uniform in every parent label $\xi_Q$.  The total Fourier loss is at most
$2^{C_FN_d\varepsilon_*t}$.  Proposition~\ref{prop:terminalpacket} gives
$2^{-\delta_Tt/2}$ at the (unsquared) $L^2$ level, so
\eqref{eq:kappachoice} leaves at least $2^{-3\delta_Tt/8}$ after the
Fourier summation.  The remainders in \eqref{eq:modsplit}
are controlled by total radial variation and Young's inequality.  Using
\eqref{eq:weighted-fourier-product} and
\eqref{eq:packet-rescaling-data}, their total $L^1$ variation in scaled
variables is
\[
 C\norm\Theta_1 2^{-t}2^{C_FN_d\varepsilon_*t}R^{-n}M_{I,J}.
\]
After \eqref{eq:packet-rescaling-output} and \eqref{eq:kappachoice}, this
is at most $C\norm\Theta_1 2^{-3t/4}M_{I,J}$.  The good term similarly
returns with the factor $(\Lambda M_{I,J})^{1/2}$.  Apply
\eqref{eq:block-V2-overlap} to the squared $L^2$ bounds and
\eqref{eq:block-V1-overlap} to the $L^1$ errors, followed by
\eqref{eq:block-mass-overlap}.  This proves
\eqref{eq:anchoredG}--\eqref{eq:anchoredE}.

The exponents can now be fixed uniformly rather than left implicit.  Take
\begin{equation}\label{eq:anchored-explicit-exponents}
 c_P=\min\left\{
       \min_{1\le\mu\le N_d}\frac{\beta_*\kappa_\mu}{2},
       \frac{3\delta_T}{8}\right\},
 \qquad c_E=\frac34.
\end{equation}
The first term in $c_P$ covers every high branch, the second covers the
all-low branch, and $c_E$ is permitted because
$1-C_FN_d\varepsilon_*>7/8$.  All quantities in
\eqref{eq:anchored-explicit-exponents}, as well as the single system of
Fourier cutoffs fixed above, depend only on $n,d$.  In the already settled
$N_d=0$ terminal case one takes $c_P=3\delta_T/8$ and $c_E=3/4$; its
actual good decay $\delta_T/2$ is stronger.

It remains to remove the finite-mode truncation.  Let $\mathcal P_N$ be
the trajectory in which every low Fourier expansion is restricted to
$\mathfrak F_N$, and let
$\mathcal P_N=G_N+E_N$ be the decomposition just constructed, summed over
$(I,J)\in\mathcal A_R$.  These are genuine finite sums.  If a block has
first high position $\mu$ and $S=\{r:r<\mu\}$, then for $N'>N$ the same
component estimate used above gives
\begin{equation}\label{eq:high-fourier-cauchy}
 \norm{G_{N'}^{I,J}-G_N^{I,J}}_{L^2(V^2)}
 \le C\norm\Theta_\infty
 (2^{m_\mu}|C_\mu|)^{-\beta_\mu}
 (\Lambda M_{I,J})^{1/2}
 \omega_N(S).
\end{equation}
Here and below the superscript $I,J$ includes the fixed output restriction
$\mathbf1_I$.  If every position is low, the terminal-packet estimate and
the remainder estimate in \eqref{eq:modsplit} give instead
\begin{align}
 \norm{G_{N'}^{I,J}-G_N^{I,J}}_{L^2(V^2)}
 &\le C\norm\Theta_\infty 2^{-\delta_Tt/2}
       (\Lambda M_{I,J})^{1/2}\omega_N(S),
 \label{eq:low-good-fourier-cauchy}\\
 \norm{E_{N'}^{I,J}-E_N^{I,J}}_{L^1(V^1)}
 &\le C\norm\Theta_1 2^{-t}M_{I,J}\omega_N(S).
 \label{eq:low-error-fourier-cauchy}
\end{align}
The harmless Fourier losses already displayed in
\eqref{eq:fourierloss} and \eqref{eq:weighted-fourier-product} may be
retained on the right sides; for the fixed block they do not depend on
$N$, while $\omega_N(S)\to0$.  There are only finitely many interacting
blocks.  Therefore \eqref{eq:block-V2-overlap},
\eqref{eq:block-V1-overlap}, and
\eqref{eq:high-fourier-cauchy}--\eqref{eq:low-error-fourier-cauchy} show
that $(G_N)$ is Cauchy in $L^2(V^2)$ and $(E_N)$ is Cauchy in
$L^1(V^1)$.  Every component path is anchored at the upper shell
endpoint: $G_N(2R)=E_N(2R)=0$.  Hence
\[
 \sup_{R\le\varepsilon\le2R}|H_\varepsilon|
 \le V^a(H_\varepsilon:R\le\varepsilon\le2R),
 \qquad a\in\{1,2\},
\]
for every component or difference path $H$ in the corresponding variation
space.  Lemma~\ref{lem:anchored-variation-completion} therefore constructs
anchored paths $G,E$ and gives convergence of the full sequences in the
displayed variation norms, with no uncountable-parameter Bochner-space
assumption.

On the other hand, the Weierstrass majorant in
\eqref{eq:weighted-fourier-product} shows that the Fourier products
defining $\mathcal P_N^{I,J}$ converge absolutely and locally uniformly in
$(X,Y,\xi_Q)$ to the original phase factor.  Since the packets are finite,
bounded, and supported in the fixed block, a constant multiple of the
integrable majorant
\[
 \mathbf1_I(x)
 \left(K_{|\Theta|}\mathbf1_{\{R<|z|\le2R\}}
       *\sum_{\substack{Q\in\mathcal Q,\ Q\subset J\\U\subset Q}}
        |d_U|\right)(x)
\]
is independent of $N$ and of $\varepsilon$, where
$K_{|\Theta|}(z)=|\Theta(z/|z|)|/|z|^n$.  Indeed, disjointness and
\eqref{eq:packetdensity} give $\sum_U|d_U|\le\Lambda$ almost everywhere,
so the displayed function is bounded by
$C\Lambda\norm\Theta_1\mathbf1_I$ and belongs to both $L^1$ and $L^2$.
The same majorant gives convergence uniformly in $\varepsilon$ for almost
every $x$ (and in $L^1+L^2$ on the fixed output block).  The last assertion
of Lemma~\ref{lem:anchored-variation-completion} therefore identifies the
anchored trajectory limit with the original one:
$\mathcal P=G+E$.  Its Fatou conclusion gives
\[
 \norm{V^2(G)}_2\le\liminf_{N\to\infty}\norm{V^2(G_N)}_2,
 \qquad
 \norm{V^1(E)}_1\le\liminf_{N\to\infty}\norm{V^1(E_N)}_1.
\]
Thus the uniform finite-mode bounds pass to $G,E$ with no change in any
decay exponent, completing the proof.
\end{proof}

The proof is a simultaneous phase--packet induction.  At every later phase,
regardless of the sizes of its coefficients, one either reaches a high
coordinate and obtains scale-covariant oscillatory decay, or expands a low
coordinate and proceeds.  The finite hierarchy \eqref{eq:kappahierarchy}
ensures that Fourier losses accumulated before a later high coordinate are
strictly smaller than its oscillatory gain.  Translation-generated
coefficients enter as unrestricted later coordinates in
\eqref{eq:blockflag}, and the flag hierarchy supplies their control.

\section{Concentration-stable degree-block packet estimates}
\label{sec:degree-packets}

The original phase-adapted Calder\'on--Zygmund blocks do not satisfy the
pointwise packet-density hypothesis of
Proposition~\ref{prop:anchoredpacket}.  This section replaces that
hypothesis by estimates using only cancellation and an average $L^1$ mass
bound.  Throughout, coefficient norms are fixed Euclidean norms on the
relevant bounded-degree polynomial spaces.

\subsection{Polynomial incidence and rational finite type}

\begin{lemma}[Multivariable sublevels and discrete smearing]
\label{lem:multi-sublevel}
Let $I\subset\R^N$ be a fixed compact box with nonempty interior, and let
$\mathbf p$ be a nonzero vector of polynomials of degree at most $D$ on
$I$.  Put $A=\norm{\mathbf p}_{\rm coeff}$.  There are $c,C>0$,
depending only on $I,N,D$ and the vector length, such that
\begin{align}
 |\{x\in I:\abs{\mathbf p(x)}\le\tau A\}|&\le C\tau^c,
 \label{eq:multi-sublevel}\\
 \int_I\min\left\{1,
   \left(\frac{L\abs{\mathbf p(x)}}{A}\right)^{-\sigma}\right\}\,dx
 &\le C L^{-\delta_\sigma},
 \qquad \delta_\sigma:=\frac{c\sigma}{c+\sigma}>0,
 \label{eq:multi-negative}
\end{align}
for $0<\tau\le1$, $L\ge1$, and $\sigma>0$.  At a zero of
$\mathbf p$, the integrand in \eqref{eq:multi-negative} is defined to be
$1$.

Let $0<q\le1$, let $\mathcal Q$ be a family of pairwise disjoint
$q$-cubes in $I$, choose $y_Q\in Q$ for every $Q\in\mathcal Q$, and
attach weights $0\le w_Q\le\Lambda q^N$.  If
$M=\sum_{Q\in\mathcal Q}w_Q$, then every polynomial $p$ on
$\R^N\times\R^N$ of degree at most $D$ and coefficient norm one satisfies,
for $0<\tau\le1$,
\begin{equation}\label{eq:weighted-pair-sublevel}
 \sum_{Q,Q'\in\mathcal Q}w_Qw_{Q'}
 \mathbf1_{\{|p(y_Q,y_{Q'})|\le\tau\}}
 \le C\Lambda M(\tau+q)^c.
\end{equation}
There is also the following chart version.  Let
$U_\nu\Subset U_\nu^*\subset\R^N$ be fixed boxes and let
$\chi_\nu:U_\nu^*\to\mathcal M_\nu^*$ be a fixed finite family of rational
charts.  Assume, uniformly in $\nu$, that
\begin{equation}\label{eq:chart-metric-jacobian}
 c|z-z'|\le |\chi_\nu(z)-\chi_\nu(z')|\le C|z-z'|,
 \qquad c\le J_\nu(z)\le C,
\end{equation}
where $J_\nu$ is the density of the induced measure under $\chi_\nu$.
Thus, for a coordinate $q$-cube $Q^\flat\subset U_\nu$ and its image
$Q=\chi_\nu(Q^\flat)$,
\begin{equation}\label{eq:chart-cube-volume}
 c q^N\le |Q|_{\mathcal M_\nu}
 =\int_{Q^\flat}J_\nu(z)\,dz\le Cq^N.
\end{equation}
The chart pair version permits selected points in such chart cubes and
weights $w_Q\le\Lambda|Q|_{\mathcal M_\nu}$; by
\eqref{eq:chart-cube-volume} this is, up to a fixed constant, the same cap
as $w_Q\le\Lambda q^N$.  It also permits physical $q$-cubes after a
bounded proximity coloring: by \eqref{eq:chart-metric-jacobian}, their
inverse images have diameter $O(q)$ and volume comparable to $q^N$, while
the selected coordinate centers in each color admit disjoint smearing
balls of radius $cq$ in the fixed enlarged box.

In a given application let $\mathcal V_\nu$ be a specified fixed
finite-dimensional space of polynomial vectors on $\mathcal M_\nu$, and
let $\mathcal V^{(2)}_{\nu,\nu'}$ be the specified fixed
finite-dimensional pair-polynomial space on
$\mathcal M_\nu\times\mathcal M_{\nu'}$.
Choose a fixed common denominator $\mathcal D_\nu$ which clears every
pullback of an element of $\mathcal V_\nu$ and is bounded above and away
from zero on $U_\nu^*$, and assume explicitly that
\begin{equation}\label{eq:chart-pullback-injective}
 T_\nu:\mathcal V_\nu\longrightarrow
 \mathcal P_{\le D_\nu}(U_\nu),\qquad
 T_\nu\mathbf p
 =\mathcal D_\nu\,\mathbf p\circ\chi_\nu,
 \quad\text{is injective}.
\end{equation}
Then \eqref{eq:multi-sublevel}--\eqref{eq:multi-negative} hold after
pullback to these charts.  The pair estimate
\eqref{eq:weighted-pair-sublevel} also holds provided the analogous
cleared pullback on $\mathcal V^{(2)}_{\nu,\nu'}$ is injective.  The constants may
depend on the finitely many least singular values in
\eqref{eq:chart-pullback-injective}, in addition to the fixed chart data.
No injectivity assertion is made for all ambient polynomials restricted
to a lower-dimensional algebraic chart.
\end{lemma}

\begin{proof}
We first prove \eqref{eq:multi-sublevel}.  A scalar polynomial of degree
at most $D$ on a fixed interval satisfies a uniform sublevel estimate.
Indeed, when a sublevel set $E$ has positive measure, one may choose
$D+1$ of its points with consecutive separations bounded below by a
constant multiple of $|E|$; Lagrange interpolation then bounds the
coefficient norm by $C\tau|E|^{-D}$.  The case $D=0$ is immediate,
so assume henceforth that $D\ge1$.  For several variables, write
$p(x',t)=\sum_{j=0}^D a_j(x')t^j$ and choose $a_j$ with coefficient norm
bounded below.  If the $(N-1)$-variable exponent already obtained is
$c_{N-1}$, split the base according to $|a_j(x')|\le s$.  The induction
hypothesis bounds the first part by $Cs^{c_{N-1}}$, while on the
complement the one-dimensional estimate bounds every fibre by
$C(\tau/s)^{1/D}$.  Taking
$s=\tau^{1/(1+Dc_{N-1})}$ gives a positive exponent depending only on
$N$ and $D$, and proves \eqref{eq:multi-sublevel}.  For a polynomial
vector, one component has coefficient norm at least a fixed multiple of
$A$, and normalization of that component completes the proof.  Choose
$c$ in the statement no larger than this exponent and no larger than one
half of the corresponding scalar sublevel exponent in dimension $2N$;
for the fixed finite chart family, take the minimum of the finitely many
resulting exponents.

For \eqref{eq:multi-negative}, set
$\theta=\sigma/(c+\sigma)$ and
\[
 E_L=\{x\in I:\abs{\mathbf p(x)}\le A L^{-\theta}\}.
\]
Equation~\eqref{eq:multi-sublevel} gives
$|E_L|\le C L^{-c\theta}=C L^{-\delta_\sigma}$.  On
$I\setminus E_L$ the integrand is at most
$L^{-\sigma(1-\theta)}=L^{-\delta_\sigma}$.  This proves
\eqref{eq:multi-negative}, with a constant uniform in $\sigma$.

For \eqref{eq:weighted-pair-sublevel}, join two cubes when their distance
is less than $10q$.  Since the cubes have pairwise disjoint interiors and
equal volume, a packing argument in a fixed enlargement of any one cube
bounds the degree of this proximity graph by a constant depending only on
$N$.  It therefore has a bounded proper coloring.  Within one color the
selected points are $cq$-separated.  Put disjoint $cq$-balls around them
and spread each weight uniformly on its ball.  The resulting density on a
fixed enlargement of $I$ is bounded by $C\Lambda$ and has integral at most
$M$.  A normalized bounded-degree polynomial has bounded gradient on a
fixed enlargement of $I^2$, so every bad selected pair smears into the
$C(\tau+q)$-sublevel.  For any ordered pair of colors, if $F,G$ are the
two smeared densities and $E$ is this sublevel, then
\[
 \iint_EF(x)G(y)\,dx\,dy
 \le\min\{M^2,C\Lambda^2|E|\}
 \le C\Lambda M|E|^{1/2}.
\]
If $\tau+q$ is bounded below by a fixed constant, the desired estimate
follows directly from $M\le C_I\Lambda$.  Otherwise apply the scalar
sublevel estimate just proved, now in $2N$ variables, to $E$ and sum the
finite color pairs.  The choice of
$c$ made above proves \eqref{eq:weighted-pair-sublevel}.

For the chart assertion, fix one of the specified rational charts and
the application-specific space $\mathcal V_\nu$.  The denominator and its
reciprocal are uniformly bounded on the enlarged box.  The required
injectivity is precisely hypothesis
\eqref{eq:chart-pullback-injective}; it is not inferred merely from
vanishing on a possibly lower-dimensional chart image.  The least
singular value of this fixed map and finite-dimensional norm equivalence
give
\[
 c_\nu\norm{\mathbf p}_{\rm coeff}
 \le\norm{T_\nu\mathbf p}_{\rm coeff}
 \le C_\nu\norm{\mathbf p}_{\rm coeff},
 \qquad
 |T_\nu\mathbf p(\zeta)|
 \simeq|\mathbf p(\chi_\nu(\zeta))|.
\]
The change-of-variables formula and \eqref{eq:chart-metric-jacobian} give,
for every measurable $E\subset\mathcal M_\nu$ contained in the chart,
\[
 c|\chi_\nu^{-1}(E)|
 \le |E|_{\mathcal M_\nu}
 =\int_{\chi_\nu^{-1}(E)}J_\nu(z)\,dz
 \le C|\chi_\nu^{-1}(E)|.
\]
Consequently \eqref{eq:multi-sublevel} and
\eqref{eq:multi-negative} transfer to the chart.  For
\eqref{eq:weighted-pair-sublevel}, use the product chart and the explicitly
assumed injectivity of its cleared pullback on
$\mathcal V^{(2)}_{\nu,\nu'}$.  In coordinate variables a chart cube has
mass cap $C\Lambda q^N$ by \eqref{eq:chart-cube-volume}.  Moreover,
\eqref{eq:chart-metric-jacobian} converts a physical separation $cq$ into
a coordinate separation comparable to $q$, and a coordinate ball of
radius $cq$ has both coordinate and induced volume comparable to $q^N$.
Thus the proximity coloring gives disjoint coordinate smearing balls;
spreading $w_Q$ on such a ball produces density at most
\[
 C\frac{\Lambda q^N}{q^N}=C\Lambda.
\]
The coefficient-one cleared polynomial has bounded coordinate gradient on
the fixed enlarged product box, so a bad selected pair and the product of
its two smearing balls lie in the $C(\tau+q)$-sublevel.  The preceding
coordinate proof of \eqref{eq:weighted-pair-sublevel} now applies, with
only the displayed fixed metric and Jacobian constants.  Taking the
minimum over the finite chart family proves all asserted chart versions.
\end{proof}

We first fix the rational quotient and its norm.  This is needed because
all later rational phases have parameters which may approach degenerate
values, and a qualitative assertion that a phase is not separated does not
by itself give a uniform oscillatory constant.

Fix compact intervals $I,J$ and slightly larger intervals $I^*,J^*$.
For fixed $D_0\in\mathbb N$ and $0<\eta<1<M$, call
\begin{equation}\label{eq:admissible-rational-presentation}
 \Psi(u,v)=\frac{N(u,v)}{Q(u,v)}
\end{equation}
an admissible rational presentation if $N,Q$ are real polynomials of degree
at most $D_0$,
\begin{equation}\label{eq:admissible-denominator}
 \norm Q_{\rm coeff}\le M,
 \qquad \eta\le |Q(u,v)|\le M
 \quad ((u,v)\in I^*\times J^*).
\end{equation}
The denominator in \eqref{eq:admissible-rational-presentation} is part of
the presentation.  In every application below it is the displayed common
denominator produced by the coordinate change.

For an admissible presentation define the polynomial
\begin{equation}\label{eq:rational-difference-numerator}
 \mathfrak C_Q\Psi(u,u',v)
 =\{Q(u,v)Q(u',v)\}^2
   \partial_v\{\Psi(u,v)-\Psi(u',v)\}
\end{equation}
and set
\begin{equation}\label{eq:rational-quotient-norm}
 \norm{[\Psi]}_{{\rm rat},Q}
 :=\norm{\mathfrak C_Q\Psi}_{\mathrm{coeff},(u,u',v)}.
\end{equation}
This definition is invariant under addition of $F(u)+G(v)$.  Conversely,
if \eqref{eq:rational-difference-numerator} vanishes, then
$\partial_v\Psi(u,v)$ is independent of $u$, and integration on the
rectangle gives $\Psi(u,v)=F(u)+G(v)$.  Thus
\eqref{eq:rational-quotient-norm} is a genuine norm on the quotient by
separated phases.  If the same rational function is written using two
admissible common denominators of bounded degree, the resulting norms are
uniformly equivalent.  Indeed, if the denominators are $Q_1,Q_2$, then
pointwise on the fixed enlarged box
\begin{equation}\label{eq:rational-denominator-change}
 \mathfrak C_{Q_2}\Psi(u,u',v)
 =\left\{\frac{Q_2(u,v)Q_2(u',v)}
 {Q_1(u,v)Q_1(u',v)}\right\}^{\!2}
 \mathfrak C_{Q_1}\Psi(u,u',v).
\end{equation}
The factor in braces and its reciprocal are bounded by constants depending
only on $\eta,M$.  All difference numerators belong to one fixed
bounded-degree polynomial space, on which coefficient and supremum norms
are uniformly equivalent.  Consequently
\[
 \norm{\mathfrak C_{Q_2}\Psi}_{\rm coeff}
 \gtrsim\norm{\mathfrak C_{Q_2}\Psi}_\infty
 \gtrsim\norm{\mathfrak C_{Q_1}\Psi}_\infty
 \gtrsim\norm{\mathfrak C_{Q_1}\Psi}_{\rm coeff}.
\]
Interchanging $Q_1,Q_2$ gives the reverse inequality and proves the
claimed uniform equivalence.

We will also use the mixed numerator
\begin{equation}\label{eq:rational-mixed-numerator}
 \mathfrak M_Q\Psi(u,v)=Q(u,v)^3\partial_u\partial_v\Psi(u,v).
\end{equation}
It is a polynomial of uniformly bounded degree and satisfies
\begin{equation}\label{eq:mixed-to-quotient}
 \norm{[\Psi]}_{{\rm rat},Q}
 \ge c\norm{\mathfrak M_Q\Psi}_{\rm coeff}.
\end{equation}
To verify this, differentiate \eqref{eq:rational-difference-numerator} in
$u'$ and then put $u'=u$.  Since the phase difference and its $v$
derivative vanish on the diagonal, derivatives falling on the prefactor
make no contribution, and one obtains exactly
\[
 \left.\partial_{u'}\mathfrak C_Q\Psi(u,u',v)\right|_{u'=u}
 =-Q(u,v)^4\partial_u\partial_v\Psi(u,v)
 =-Q(u,v)\mathfrak M_Q\Psi(u,v).
\]
Differentiation and diagonal restriction are bounded maps on coefficient
spaces.  Moreover
$\norm{QH}_{\rm coeff}\ge c\norm H_{\rm coeff}$ for every polynomial $H$
of the relevant bounded degree: this follows from
$\norm{QH}_{L^\infty(I^*\times J^*)}\ge
\eta\norm H_{L^\infty(I^*\times J^*)}$ and finite-dimensional norm
equivalence.  This proves \eqref{eq:mixed-to-quotient}.

\begin{lemma}[Rational van der Corput with a $\BV$ amplitude]
\label{lem:rational-vdc}
Let $\psi=N/q$ on a fixed compact interval, where $N,q$ have degree at
most $D_0$, $\norm q_{\rm coeff}\le M$, and
$\eta\le|q|\le M$ on a fixed enlargement.  Put
\[
 G=\norm{q^2\psi'}_{\rm coeff}
   =\norm{N'q-Nq'}_{\rm coeff}.
\]
If $a$ is a zero-extended bounded-variation function supported in the
interval, then
\begin{equation}\label{eq:rational-vdc}
 \left|\int e^{i\psi(v)}a(v)\,dv\right|
 \le C(\norm a_\infty+\Var(a))\min\{1,G^{-c}\}.
\end{equation}
When \(G=0\), the minimum on the right is understood to be \(1\).
The same assertion holds after a bounded partition into intervals, with
the sum of the sectional $L^\infty+\BV$ norms on the right.
The constants and the positive exponent depend only on
$D_0,\eta,M$ and the fixed interval and its enlargement.
\end{lemma}

\begin{proof}
Only $G>1$ requires proof; the case \(G\le1\) is the trivial \(L^1\)
bound.  Write $p=q^2\psi'$.  By
Lemma~\ref{lem:multi-sublevel}, applied to $p/G$, the set
\[
 E=\{v:|p(v)|\le G^{1/2}\}
\]
has measure at most $CG^{-c_0}$.  Its contribution is bounded by
$C\norm a_\infty G^{-c_0}$.  On the complement,
$|\psi'|\ge cG^{1/2}$.  The numerator of $\psi''$ has bounded degree.
Apart from the two endpoints of the ambient interval, the boundary of
$E$ is contained in the union of the zero sets of the two bounded-degree
polynomials $p-G^{1/2}$ and $p+G^{1/2}$; hence $E$
has only $O_{D_0}(1)$ components.  Partition its complement at those
boundary points and at the zeros of the numerator of $\psi''$; the
number of intervals is therefore bounded in terms of $D_0$.  On each remaining
interval $\psi'$ is monotone and bounded away from zero.  The first
derivative van der Corput estimate, proved for a $\BV$ amplitude by one
Stieltjes integration by parts, gives
$C(\norm a_\infty+\Var(a))G^{-1/2}$.  Summing the boundedly many pieces
and decreasing the exponent proves \eqref{eq:rational-vdc}.  If the
numerator of $\psi''$ vanishes identically, $\psi'$ is already monotone
and the same argument applies without this partition.
\end{proof}

\begin{lemma}[Coefficient-norm polynomial van der Corput]
\label{lem:coefficient-vdc}
Let $I$ be a fixed compact interval, let $D\ge1$, let $p$ be a real
polynomial of degree at most $D$, and put
\[
 \Lambda_p=\norm{p-p(0)}_{\rm coeff}.
\]
If $a$ is supported in $I$ and is of bounded variation, then
\begin{equation}\label{eq:coefficient-vdc}
 \left|\int_I e^{ip(t)}a(t)\,dt\right|
 \le C_D\{\norm a_\infty+\Var(a)\}
       \min\{1,\Lambda_p^{-c_D}\}.
\end{equation}
The value $p(0)$ may be replaced by any constant; equivalently,
$\Lambda_p$ may be any fixed coefficient norm on
the space of real polynomials of degree at most $D$, modulo constants.
When $\Lambda_p=0$, the minimum in
\eqref{eq:coefficient-vdc} is understood to be $1$.
\end{lemma}

\begin{proof}
Apply Lemma~\ref{lem:rational-vdc} with denominator identically one.
Differentiation is injective from the finite-dimensional quotient
of real polynomials of degree at most $D$ by the constants into the space
of polynomials of degree at most $D-1$, and therefore
$\norm{p'}_{\rm coeff}\ge c_D\Lambda_p$ by finite-dimensional norm
equivalence.  Decreasing the exponent in
\eqref{eq:rational-vdc} proves \eqref{eq:coefficient-vdc}.
\end{proof}

\begin{lemma}[Uniform rational oscillatory operator]
\label{lem:rational-operator}
Let $I,J$ be fixed compact intervals, let $A:I\times J\to\mathbb C$ be
measurable, and define $S:L^2(J)\to L^2(I)$ initially on bounded functions
by
\[
 Sf(u)=\int_J e^{i\Psi(u,v)}A(u,v)f(v)\,dv.
\]
Assume that \eqref{eq:admissible-rational-presentation}--
\eqref{eq:admissible-denominator} hold with fixed data.  Suppose $B\ge1$
and
\[
 \sup_{u\in I}\int_J|A(u,v)|\,dv\le B,\qquad
 \sup_{v\in J}\int_I|A(u,v)|\,du\le B.
\]
Assume also that, for every $u\in I$, there is a partition of $J$ into at
most $M_0$ intervals on each of which
\[
 \norm{A(u,\cdot)}_\infty+\Var(A(u,\cdot))\le B.
\]
The endpoints may depend on $u$.  Put
$L=\norm{[\Psi]}_{{\rm rat},Q}$.  Then
\begin{equation}\label{eq:rational-operator}
 \norm S_{2\to2}\le CB^C\min\{1,L^{-c}\}.
\end{equation}
When \(L=0\), the minimum is again understood to be \(1\).
The constants depend only on the fixed intervals, $D_0,\eta$, the
admissibility bound $M$, and $M_0$.
\end{lemma}

\begin{proof}
The trivial part follows from Schur's test, so assume $L\ge1$.  The kernel
of $SS^*$ is
\[
 K(u,u')=\int_J e^{i[\Psi(u,v)-\Psi(u',v)]}
 A(u,v)\overline{A(u',v)}\,dv.
\]
For fixed $(u,u')$, an admissible common denominator of the phase
difference is $q_{u,u'}(v)=Q(u,v)Q(u',v)$; its degree, coefficient norm,
and upper and lower bounds are controlled by the fixed admissibility data.
Its derivative numerator is precisely
\eqref{eq:rational-difference-numerator}.  If
\[
 G(u,u')=\norm{\mathfrak C_Q\Psi(u,u',\cdot)}_{\mathrm{coeff},v},
\]
then refine the two partitions of $J$ belonging to $u$ and $u'$.
There are at most $2M_0$ resulting intervals, and
\[
 \norm{A(u,\cdot)\overline{A(u',\cdot)}}_{L^\infty+\BV}
 \le CB^2
\]
on each of them.  Lemma~\ref{lem:rational-vdc}, summed over this
partition, gives
\begin{equation}\label{eq:rational-K-pointwise}
 |K(u,u')|\le CB^2\min\{1,G(u,u')^{-c_0}\}.
\end{equation}
The entries of the coefficient vector defining $G$ are polynomials in
$(u,u')$ of uniformly bounded degree, and their joint coefficient norm is
exactly $L$.  Consequently \eqref{eq:multi-negative}, applied on
$I\times I$, yields
\[
 \iint_{I\times I}|K(u,u')|^2\,du\,du'
 \le CB^C L^{-c_1}.
\]
There is no deletion of the diagonal or of any algebraically degenerate
branch: all such points are included in this negative-moment integral.
Finally,
\[
 \norm S_{2\to2}^2=\norm{SS^*}_{2\to2}
 \le\norm K_{L^2(I\times I)}.
\]
Taking the two square roots and decreasing $c_1$ proves
\eqref{eq:rational-operator}.
\end{proof}

\subsection{Rough two-ray correlation and a flat high block}

\begin{lemma}[Universal rough polynomial correlation]
\label{lem:rough-correlation}
Let $n\ge1$, let $D\ge1$, let $K=[r_-,r_+]\Subset(0,\infty)$, and let
$a,b$ be radial $\BV$ profiles supported in $K$ with
$L^\infty+\BV$ norm at most $B\ge1$.  Let
$\Theta\in L^\infty(\Sph^{n-1})$ and put
\[
 k_a(z)=|z|^{-n}\Theta(z/|z|)a(|z|),\qquad
 k_b(z)=|z|^{-n}\Theta(z/|z|)b(|z|).
\]
Let
$S$ be a real polynomial of degree at most $D$, and let $\Lambda$ be the
coefficient norm of its nonconstant part.  Then, for every $h\in\R^n$,
if $n\ge2$,
\begin{equation}\label{eq:rough-correlation}
 \left|\int e^{iS(z)}\overline{k_a(z)}k_b(z-h)\,dz\right|
 \le CB^C\norm\Theta_\infty^2
 \min\{1,(\Lambda|h|^D)^{-c}\}.
\end{equation}
For $n\ge2$ the trivial bound is understood when $h=0$ or $\Lambda=0$.
In dimension one, for every $h$, the precise unit-scale analogue is
\[
 \left|\int e^{iS(z)}\overline{k_a(z)}k_b(z-h)\,dz\right|
 \le CB^C\norm\Theta_\infty^2\min\{1,\Lambda^{-c}\}.
\]
\end{lemma}

\begin{lemma}[Quantitative equal-shell pullback]
\label{lem:equal-pullback}
Let $D\ge1$ and let
$\phi(t,r)=\sum_{m=0}^D\phi_m(t,r)$, where $\phi_m$ is homogeneous of
degree $m$, and put
\[
 \Gamma=\norm{\phi-\phi_0}_{\rm coeff}.
\]
Let $0<\alpha\le1$, $0<H\le C_0$, $|U_0|\le C_0$, let
$\Delta\in\R$, and let $q(X,Y)=\Delta+X-Y$ satisfy
$c_0\le|q|\le C_0$ on a fixed enlargement of
a fixed rectangle.  Define
\begin{equation}\label{eq:equal-pullback-phase}
 \Psi(X,Y)
 =\phi\left(\frac{U_0+HX}{q(X,Y)},
             \frac{\alpha}{q(X,Y)}\right)
 =\sum_{m=0}^D\frac{\phi_m(U_0+HX,\alpha)}{q(X,Y)^m}.
\end{equation}
Use the canonical common denominator $Q=q^D$.  Then
\begin{equation}\label{eq:equal-pullback-lower}
 \norm{[\Psi]}_{{\rm rat},Q}
 \ge c(\alpha H)^D\Gamma.
\end{equation}
The constant depends only on $D,c_0,C_0$ and the fixed rectangles.
\end{lemma}

\begin{proof}[Proof of Lemma~\ref{lem:equal-pullback}]
Put $A(X)=U_0+HX$ and use $(X,q)$ rather than $(X,Y)$ as polynomial
coordinates.  Since $X,Y$ range in fixed rectangles and
$c_0\le|q|\le C_0$ on their enlargement, $\Delta$ is uniformly bounded.
Thus $(X,Y)\mapsto(X,q)$ and its inverse have uniformly bounded
coefficient matrices on every bounded-degree polynomial space.
For $m\ge1$, direct differentiation gives
\begin{equation}\label{eq:equal-mixed-formula}
 q^{D+2}\partial_X\partial_Y
 \left\{\phi_m(A(X),\alpha)q^{-m}\right\}
 =q^{D-m}\left\{mHq\,\partial_t\phi_m(A(X),\alpha)
      -m(m+1)\phi_m(A(X),\alpha)\right\}.
\end{equation}
This identity also identifies all possible cancellation between adjacent
homogeneous degrees.  After summing \eqref{eq:equal-mixed-formula}, the
coefficient of $q^{D-m}$ is
\begin{equation}\label{eq:equal-triangular-tag}
 -m(m+1)\phi_m(A,\alpha)
 +(m+1)H\partial_t\phi_{m+1}(A,\alpha),
 \qquad \phi_{D+1}=0.
\end{equation}
Starting with $m=D$ and descending, the triangular system
\eqref{eq:equal-triangular-tag} recovers every
$\phi_m(A,\alpha)$, $1\le m\le D$.  More explicitly, on the finite direct
sum $\bigoplus_{m=1}^D\mathcal P_{\le m}(X)$ the diagonal maps are the
nonzero scalars $-m(m+1)$ and the only off-diagonal maps are
$(m+1)H\partial_t$ from level $m+1$ to level $m$.  Descending
substitution and $0<H\le C_0$ therefore give
\[
 \sum_{m=1}^D\norm{\phi_m(A,\alpha)}_{\mathrm{coeff},X}
 \le C_{D,C_0}
 \norm{q^{D+2}\partial_X\partial_Y\Psi}_{\mathrm{coeff},(X,q)}.
\]
The uniformly invertible affine coefficient change above makes the last
norm equivalent to its $(X,Y)$ coefficient norm.  Hence
\begin{equation}\label{eq:equal-mixed-controls-restriction}
 \norm{q^{D+2}\partial_X\partial_Y\Psi}_{\mathrm{coeff},(X,Y)}
 \ge c\left(\sum_{m=1}^D
       \norm{\phi_m(U_0+HX,\alpha)}_{\mathrm{coeff},X}^2\right)^{1/2}.
\end{equation}
In particular, no choice of a largest homogeneous term is being made:
all homogeneous degrees are recovered simultaneously.

The affine map $t=U_0+HX$ and evaluation at $r=\alpha$ are quantitatively
injective on homogeneous two-variable polynomials.  Indeed, recover
$\phi_m(t,\alpha)$ from its composition with $U_0+HX$ and then recover
the coefficient of $t^jr^{m-j}$ by division by $\alpha^{m-j}$.  Since
$|U_0|\le C_0$, $0<\alpha\le1$, and $0<H\le C_0$, finite-dimensional
norm equivalence gives
\begin{equation}\label{eq:equal-evaluation-lower}
 \norm{\phi_m(U_0+HX,\alpha)}_{\mathrm{coeff},X}
 \ge c(\alpha H)^D\norm{\phi_m}_{\rm coeff}.
\end{equation}
Indeed, the inverse of the affine substitution $t=U_0+HX$ has norm
$O_{m,C_0}(H^{-m})$, while evaluation at $r=\alpha$ multiplies the
coefficient of $t^jr^{m-j}$ by $\alpha^{m-j}$.  The left side is thus
at least $cH^m\alpha^m\norm{\phi_m}_{\rm coeff}$.  If
$\alpha H\le1$, this is at least the right side of
\eqref{eq:equal-evaluation-lower}; if $\alpha H>1$, the bounded upper
ranges of $\alpha$ and $H$ allow the remaining factor to be absorbed
into $c$.
Combining \eqref{eq:equal-mixed-controls-restriction} and
\eqref{eq:equal-evaluation-lower} gives the desired lower bound for the
reduced mixed numerator $q^{D+2}\Psi_{XY}$.

For the presentation with denominator $Q=q^D$, the mixed numerator
$Q^3\Psi_{XY}$ differs from the reduced numerator by multiplication by
$q^{2D-2}$.  This factor is bounded above and away from zero on the
enlarged rectangle, so their coefficient norms are uniformly equivalent.
Equation \eqref{eq:mixed-to-quotient} now proves
\eqref{eq:equal-pullback-lower}.
\end{proof}

\begin{proof}[Proof of Lemma~\ref{lem:rough-correlation}]
Assume $n\ge2$.  Since the radial profiles are supported in a fixed
compact subset of $(0,\infty)$, kernel overlap implies $0<|h|\le C$.
The cases $h=0$ and $\Lambda=0$ are covered by the stated trivial bound.
Put $H=|h|$ and choose $\mathcal O\in O(n)$ such that
$h=H\mathcal Oe_1$.  Under the change of variables
$z=\mathcal O\widetilde z$ replace
\[
 S(z)\quad\hbox{by}\quad
 \widetilde S(\widetilde z)=S(\mathcal O\widetilde z),
 \qquad
 \Theta(\xi)\quad\hbox{by}\quad
 \widetilde\Theta(\xi)=\Theta(\mathcal O\xi).
\]
The radial profiles are unchanged, Lebesgue measure is preserved, and
$\norm{\widetilde\Theta}_\infty=\norm\Theta_\infty$.  On the
finite-dimensional space of nonconstant polynomials of degree at most $D$,
composition with $\mathcal O$ is invertible with inverse composition by
$\mathcal O^T$.  Compactness of $O(n)$ therefore gives, uniformly in
$\mathcal O$,
\[
 c_{n,D}\Lambda\le
 \norm{\widetilde S-\widetilde S(0)}_{\rm coeff}
 \le C_{n,D}\Lambda.
\]
Thus the desired estimate is invariant, up to its allowed structural
constant, under this rotation.  We may rename the transformed quantities,
assume $h=He_1$, and write
\[
 z=(t,\rho\omega),\quad u=t/\rho,\quad v=(t-H)/\rho,
 \quad\omega\in\Sph^{n-2}.
\]
Then exactly
\begin{equation}\label{eq:two-ray-coordinates}
 \rho=\frac{H}{u-v},\qquad
 z=\frac{H(u,\omega)}{u-v},\qquad
 z-h=\frac{H(v,\omega)}{u-v}.
\end{equation}
The two rough factors are respectively functions of $(u,\omega)$ and
$(v,\omega)$, and the Jacobian times the two homogeneous kernels is
\begin{equation}\label{eq:two-ray-jacobian}
 H^{-n}|u-v|^{n-1}(1+u^2)^{-n/2}(1+v^2)^{-n/2}
 \,du\,dv\,d\omega.
\end{equation}
Thus the rough angular factors are input and output $L^2$ functions and
are never differentiated.

Choose once and for all a nonnegative
$\beta\in C_c^\infty((c,C))$ such that, on the relevant positive range,
\[
 \sum_{k\in\mathbb Z}\beta(2^k\rho)=1.
\]
After combining the boundedly many largest pieces, write
$\alpha=2^{-k}\le1$ and
$\beta_\alpha(\rho)=\beta(\rho/\alpha)$.  This is a genuine partition of
unity on the overlap region, its supports have uniformly bounded overlap,
and
\begin{equation}\label{eq:rough-polar-partition-bv}
 \sup_\alpha\left\{\norm{\beta_\alpha}_\infty+
  \Var_\rho(\beta_\alpha)\right\}\le C.
\end{equation}
Thus the transverse radius is decomposed into bands
$\rho\simeq\alpha$, $0<\alpha\le1$.  Put $U=\alpha u,V=\alpha v$.
In these variables \eqref{eq:two-ray-jacobian} becomes
\begin{equation}\label{eq:two-ray-scaled-jacobian}
 H^{-n}\alpha^{n-1}|U-V|^{n-1}
 (\alpha^2+U^2)^{-n/2}(\alpha^2+V^2)^{-n/2}
 \,dU\,dV\,d\omega.
\end{equation}
Annular support gives $|U-V|\simeq H$ in a fixed bounded set and also
$\alpha^2+U^2\simeq\alpha^2+V^2\simeq1$.  Thus, after the factor
$\alpha^{n-1}$ is removed, the kernel in $(U,V)$ has size $O(H^{-1})$
and is supported in a strip of width $O(H)$.

Choose a fixed $\varepsilon_0>0$, smaller than one tenth of the lower
constant in $|U-V|\simeq H$, and partition the $U$- and $V$-axes into
half-open intervals of length $\varepsilon_0H$.  Each input interval
interacts with only boundedly many output intervals and conversely.  On
an interacting pair write
\[
 U=U_0+HX,\qquad V=V_0+HY,\qquad
 \Delta=\frac{U_0-V_0}{H};
\]
then $X,Y$ range over fixed intervals of length $\varepsilon_0$ and
\[
 q(X,Y)=\Delta+X-Y
\]
is bounded above and away from zero on a fixed enlargement of their full
product.  This last assertion concerns the whole product rectangle, not
only the support of the amplitude, and is the reason for the small fixed
choice of $\varepsilon_0$.

The phase on this block is exactly
\begin{equation}\label{eq:equal-block-phase}
 \Psi_{\omega,U_0,V_0,H,\alpha}(X,Y)
 =S\left(\frac{(U_0+HX)e_1+\alpha\omega}{q(X,Y)}\right).
\end{equation}
The factor $dV=H\,dY$ cancels the $H^{-1}$ in the normalized kernel.
The two angular factors are, up to the fixed sign of $q$,
\[
 \overline{\Theta\left(
   \frac{Ue_1+\alpha\omega}{\sqrt{U^2+\alpha^2}}\right)},
 \qquad
 \Theta\left(\frac{Ve_1+\alpha\omega}{\sqrt{V^2+\alpha^2}}\right);
\]
they are respectively functions of the output and input variables and
are never differentiated.  The radial factors are
\[
 \overline{a\left(\frac{\sqrt{U^2+\alpha^2}}{|q|}\right)},\qquad
 b\left(\frac{\sqrt{V^2+\alpha^2}}{|q|}\right).
\]
Here is the uniformity in this assertion.  The sign of $q$ is fixed on
the enlarged block.  Put
$s_U=(U^2+\alpha^2)^{1/2}$ and
$s_V=(V^2+\alpha^2)^{1/2}$.  Up to nonzero factors of fixed sign, the
four sectional derivatives of $s_U/|q|$ and $s_V/|q|$ are therefore
controlled by those of $s_U/q$ and $s_V/q$.  Directly,
\begin{align*}
 \partial_X\frac{s_U}{q}
 &=\frac{HUq-s_U^2}{s_Uq^2},
 &\partial_Y\frac{s_U}{q}&=\frac{s_U}{q^2},\\
 \partial_X\frac{s_V}{q}
 &=-\frac{s_V}{q^2},
 &\partial_Y\frac{s_V}{q}&=\frac{HVq+s_V^2}{s_Vq^2}.
\end{align*}
Thus a derivative can vanish only through
\[
 H Uq-s_U^2=0\quad\hbox{or}\quad H Vq+s_V^2=0;
\]
the other two derivatives never vanish.  Since
$U=U_0+HX$, $V=V_0+HY$, and $q=\Delta+X-Y$, the two displayed
critical-point equations are polynomial equations of degree at most two.
Their zero sets therefore split every section into $O(1)$ monotonicity
intervals (with an identically vanishing derivative requiring no split).
On an interval $I$ on which a profile argument $\gamma$ is monotone, the
one-dimensional BV pullback inequality gives
\[
 \Var_I(a\circ\gamma)
 \le\Var_{\gamma(I)}(a)\le\Var(a),
\]
and the same conclusion is immediate if $\gamma$ is constant.  Summing
over the uniformly bounded number of intervals gives
\[
 \Var(a\circ\gamma)+\Var(b\circ\widetilde\gamma)
 \le C_D\{\Var(a)+\Var(b)\}.
\]
The annular support gives
$\sqrt{U^2+\alpha^2}\simeq\sqrt{V^2+\alpha^2}\simeq1$
and $|q|\simeq1$, so all remaining algebraic factors in
\eqref{eq:two-ray-scaled-jacobian} and their sectional derivatives are
uniformly bounded.  Moreover, on the block
$\rho/\alpha=H/(U-V)=1/q(X,Y)$, so a smooth dyadic cutoff in
$\rho/\alpha$ is the fixed function $\beta(1/q)$ of $q$.  Since $q$ is
bounded above and away from zero on the enlarged rectangle, this cutoff
and all of its sectional derivatives are uniformly bounded, consistently
with \eqref{eq:rough-polar-partition-bv}.
Thus the amplitude in each normalized block
satisfies Lemma~\ref{lem:rational-operator} with bound $CB^C$.

For completeness, the block synthesis has no hidden dependence on $H$.
If $T_{ij}$ is the operator between an interacting $V$-interval and
$U$-interval and
\[
 (\mathcal U_i g)(X)=H^{1/2}g(U_i+HX),\qquad
 (\mathcal V_j f)(Y)=H^{1/2}f(V_j+HY),
\]
then $\mathcal U_iT_{ij}\mathcal V_j^{-1}$ is precisely the normalized
operator above: the factor $dV=H\,dY$ cancels the kernel factor $H^{-1}$.
If $d_{\rm in}$ and $d_{\rm out}$ are the two degrees of the block
interaction graph, Cauchy--Schwarz gives
\[
 \norm T_{2\to2}
 \le(d_{\rm in}d_{\rm out})^{1/2}
     \sup_{i,j}\norm{T_{ij}}_{2\to2}.
\]
Both graph degrees are absolute constants.  The same calculation with
Schur's test proves the trivial band bound $CB^C\alpha^{n-1}$.

For fixed $\omega$, put
\[
 \phi_\omega(t,r)=S(te_1+r\omega),\qquad
 \Gamma(\omega)=\norm{\phi_\omega-\phi_\omega(0,0)}_{\rm coeff}.
\]
Formula \eqref{eq:equal-block-phase} is
\eqref{eq:equal-pullback-phase}.  Lemma~\ref{lem:equal-pullback} gives,
for its canonical denominator $q^D$,
\[
 \norm{[\Psi_{\omega,U_0,V_0,H,\alpha}]}_{{\rm rat},q^D}
 \ge c\Gamma(\omega)(\alpha H)^D.
\]

Lemma~\ref{lem:rational-operator} gives the fixed-$\omega$ block gain,
and the preceding almost-orthogonal synthesis preserves it.  We now spell
out the angular synthesis.  Let $I_0\subset\R$ be a fixed compact interval
containing every possible $U$- and $V$-coordinate on the band, and let
$\beta$ denote the fixed zero-extended cutoff defining this band in
$\rho/\alpha=H/(U-V)$.  After the factor $\alpha^{n-1}$ in
\eqref{eq:two-ray-scaled-jacobian} is removed, put
\begin{align*}
 A_{\alpha,H}(U,V)
 &=\beta\left(\frac{H}{U-V}\right)
 H^{-n}(U-V)^{n-1}
 (\alpha^2+U^2)^{-n/2}(\alpha^2+V^2)^{-n/2}\\
 &\quad\times
 \overline{a\left(\frac{H\sqrt{U^2+\alpha^2}}{U-V}\right)}
 b\left(\frac{H\sqrt{V^2+\alpha^2}}{U-V}\right),
\end{align*}
extended by zero off the band support.  Here $U-V>0$ by
\eqref{eq:two-ray-coordinates}.  Define
\begin{align*}
 (T_{\alpha,\omega}f)(U)
 &=\int_{I_0}e^{iS(H(Ue_1+\alpha\omega)/(U-V))}
 A_{\alpha,H}(U,V)f(V)\,dV,\\
 f_{\alpha,\omega}(V)
 &=\Theta\left(\frac{Ve_1+\alpha\omega}
                    {\sqrt{V^2+\alpha^2}}\right)\mathbf1_{I_0}(V),\\
 g_{\alpha,\omega}(U)
 &=\Theta\left(\frac{Ue_1+\alpha\omega}
                    {\sqrt{U^2+\alpha^2}}\right)\mathbf1_{I_0}(U).
\end{align*}
With the convention $\langle F,G\rangle=\int F\overline G$, the
contribution $\mathcal I_\alpha$ of this transverse band is exactly
\begin{equation}\label{eq:rough-band-direct-integral}
 \mathcal I_\alpha
 =\alpha^{n-1}\int_{\Sph^{n-2}}
 \langle T_{\alpha,\omega}f_{\alpha,\omega},
         g_{\alpha,\omega}\rangle_{L^2(I_0)}\,d\omega.
\end{equation}
The functions in this formula are jointly measurable, and
$\omega\mapsto f_{\alpha,\omega},g_{\alpha,\omega}$ are strongly
measurable $L^2(I_0)$-valued maps.  Moreover
\[
 \norm{f_{\alpha,\omega}}_2+\norm{g_{\alpha,\omega}}_2
 \le C\norm\Theta_\infty.
\]
The block argument above, Lemma~\ref{lem:equal-pullback}, and
Lemma~\ref{lem:rational-operator} give, uniformly in $\omega,\alpha,H$,
\begin{equation}\label{eq:rough-fixed-angular-operator}
 \norm{T_{\alpha,\omega}}_{2\to2}
 \le CB^C\min\{1,[\Gamma(\omega)(\alpha H)^D]^{-c_0}\}.
\end{equation}
The trivial Schur bound supplies an integrable majorant in
\eqref{eq:rough-band-direct-integral}; hence ordinary Fubini and
Cauchy--Schwarz yield
\begin{equation}\label{eq:rough-band-angular-integral}
 |\mathcal I_\alpha|
 \le CB^C\norm\Theta_\infty^2\alpha^{n-1}
 \int_{\Sph^{n-2}}
 \min\{1,[\Gamma(\omega)(\alpha H)^D]^{-c_0}\}\,d\omega.
\end{equation}
No measurable-field or abstract direct-integral assertion is being used.

It remains to estimate the scalar angular integral.  When $n=2$,
$\Sph^0=\{-1,1\}$, and the two restrictions are related by
$r\mapsto-r$; either restriction has coefficient norm comparable to
$\Lambda$, so the required bound follows by summing the two points.
Assume henceforth in this paragraph that $n\ge3$.  Split $\Sph^{n-2}$ by a
fixed measurable partition subordinate to a finite rational atlas
$\omega=\omega_\nu(\zeta)$; the chart Jacobians are bounded above and away
from zero on fixed enlargements.  In the $\nu$-th chart let
$\mathbf g_\nu(\zeta)$ be the vector of all nonconstant $(t,r)$
coefficients of $S(te_1+r\omega_\nu(\zeta))$, after multiplication by one
fixed common chart denominator.  The denominator and its reciprocal are
bounded on the enlargement, so pointwise
\begin{equation}\label{eq:rough-chart-pointwise-comparison}
 |\mathbf g_\nu(\zeta)|\simeq
 \Gamma(\omega_\nu(\zeta)).
\end{equation}
For each nonempty chart, the restriction map
\[
 S-S(0)\longmapsto\mathbf g_\nu
\]
is injective: if its nonconstant $(t,r)$ coefficients vanish throughout
the chart, then $S-S(0)$ vanishes on the corresponding open cone and
therefore identically.  Thus, for the application-specific polynomial
space consisting of these restrictions, this verifies exactly the
injectivity hypothesis \eqref{eq:chart-pullback-injective}; no injectivity
of all ambient polynomials restricted to the sphere is being used.  With
$A_\nu=\norm{\mathbf g_\nu}_{\mathrm{coeff},\zeta}$,
finite-dimensional norm equivalence and the finite atlas give
\begin{equation}\label{eq:rough-chart-coefficient-lower}
 A_\nu\ge c\Lambda
 \qquad\text{for every }\nu.
\end{equation}
If $(\alpha H)^DA_\nu\ge1$, apply
\eqref{eq:multi-negative} to $\mathbf g_\nu/A_\nu$, with
$L=(\alpha H)^DA_\nu$ and $\sigma=c_0$; if
$(\alpha H)^DA_\nu<1$, use the trivial bound.  In view of
\eqref{eq:rough-chart-pointwise-comparison}--
\eqref{eq:rough-chart-coefficient-lower}, summing the finitely many charts
gives
\begin{equation}\label{eq:rough-angular-negative-moment}
 \int_{\Sph^{n-2}}
 \min\{1,[\Gamma(\omega)(\alpha H)^D]^{-c_0}\}\,d\omega
 \le C\min\{1,[\Lambda(\alpha H)^D]^{-c_1}\}.
\end{equation}
All constants here are uniform as $\alpha,H\downarrow0$: the normalized
rectangles, denominator bounds, block-graph degrees, and numbers of
sectional monotonicity intervals are fixed.  Combining
\eqref{eq:rough-band-angular-integral} and
\eqref{eq:rough-angular-negative-moment} proves the band bound
\[
 CB^C\norm\Theta_\infty^2\alpha^{n-1}
 \min\{1,[\Lambda(\alpha H)^D]^{-c}\}.
\]

The original correlation is the sum of these band contributions because
the preceding cutoffs form a partition of unity.  Taking absolute values
only after this identity and using their bounded overlap, write
$A=\Lambda H^D$.  The dyadic band sum is bounded by
\[
 \sum_{\alpha\in2^{-\mathbb N_0}}
 \alpha^{n-1}\min\{1,(A\alpha^D)^{-c}\}.
\]
If $A\le1$, this is $O(1)$.  If $A>1$, choose
$\alpha_0\simeq A^{-1/D}$.  The bands below the threshold satisfy
\[
 \sum_{\alpha\le\alpha_0}\alpha^{n-1}
 \lesssim A^{-(n-1)/D}.
\]
After decreasing the oscillatory exponent so that $Dc<n-1$, the remaining
bands satisfy
\[
 \sum_{\alpha>\alpha_0}
 A^{-c}\alpha^{n-1-Dc}\lesssim A^{-c}.
\]
Thus the complete sum is
$O(\min\{1,A^{-c'}\})$, where
$c'=\min\{c,(n-1)/D\}>0$, which is the right side of
\eqref{eq:rough-correlation}.

When $n=1$, split the support into the finitely many signed intervals on
which the two homogeneous kernel factors and radial profiles are of
bounded variation.  Their product has $L^\infty+\BV$ norm at most
$CB^C\norm\Theta_\infty^2$.
Lemma~\ref{lem:coefficient-vdc} gives the asserted unit-scale bound
uniformly in $h$.
\end{proof}

For a homogeneous polynomial $H(x,y)$, let $\mathcal D H$ denote the
vector of all nonconstant-in-$z$ coefficients of
\begin{equation}\label{eq:same-shell-difference}
 H(y+z,y')-H(y+z,y).
\end{equation}

\begin{lemma}[Same-shell degree-block injectivity]
\label{lem:same-shell-injective}
For every integer $m\ge2$, on homogeneous degree-$m$ polynomials modulo
pure phases, the map
$[H]\mapsto\mathcal D H$ is injective.  Hence
$\norm{\mathcal D H}_{\rm coeff}\ge c_{n,m}\norm{[H]}_{\rm coeff}$.
\end{lemma}

\begin{proof}
If \eqref{eq:same-shell-difference} is independent of $z$, then
$\partial_xH(x,y')-\partial_xH(x,y)=0$.  Thus $\partial_xH$ is
independent of its second variable, and integration makes $H$ pure.
Quantitative injectivity is finite-dimensional norm equivalence.
\end{proof}

\begin{proposition}[Flat-measure high block]
\label{prop:flat-high-block}
Fix a translate $B_0$ of a prescribed bounded unit-scale cube, and assume
that every cube below is contained in $B_0$; the constants are uniform in
the location of $B_0$.
Let $t\ge0$, let $\mathcal Q$ be a finite family of pairwise disjoint
$q$-cubes with $q=2^{-t}$, and, for each $Q\in\mathcal Q$, let $\nu_Q$ be
a complex finite measure supported on $Q$.  Assume $\Lambda\ge0$ and
\begin{equation}\label{eq:flat-measure-data}
 |\nu_Q|(Q)\le\Lambda|Q|,\qquad M=\sum_Q|\nu_Q|(Q).
\end{equation}
Suppose $P$ is a real polynomial of degree at most $d$, its highest
nonpure homogeneous block has coefficient norm $L\ge1$, and
$\Theta\in L^\infty(\Sph^{n-1})$.  For every $|e_Q|\le1$ and every radial
$\BV$ profile $a$, supported in a fixed compact subinterval of
$(0,\infty)$ and of $L^\infty+\BV$ norm at most $B$, where $B\ge1$, write
\[
 A^{P,\Theta}_a\nu(x)=\int e^{iP(x,y)}k_a(x-y)\,d\nu(y).
\]
Then
\begin{equation}\label{eq:flat-high-block}
 \left\|A^{P,\Theta}_a\sum_Qe_Q\nu_Q\right\|_2^2
 \le CB^C\norm\Theta_\infty^2\Lambda M
       (L^{-c}+2^{-ct}).
\end{equation}
\end{proposition}

\begin{proof}
After simultaneously translating the input and output variables by the
center of $B_0$, we may assume that $B_0$ lies in a fixed reference cube.
This translation preserves the kernel geometry and the output $L^2$ norm;
it also preserves the highest nonpure homogeneous block of $P$ and its
coefficient norm, since higher homogeneous blocks are pure and remain pure
under translation.
Put $w_Q=|\nu_Q|(Q)$.  Since the cubes are disjoint and contained in
$B_0$, the data in \eqref{eq:flat-measure-data} give
\[
 M=\sum_Qw_Q\le\Lambda\sum_Q|Q|\le\Lambda|B_0|\le C\Lambda,
 \qquad
 \sum_{Q,Q'}w_Qw_{Q'}=M^2\le C\Lambda M.
\]
Cubes with $w_Q=0$ may be omitted.  Write
$d\nu_Q=h_Q\,d|\nu_Q|$, where $|h_Q|=1$ almost everywhere, and on the
finite product probability space
\[
 \prod_Q\left(Q,\frac{|\nu_Q|}{w_Q}\right)
\]
let $Y_Q$ be the coordinate random variables.  Testing against bounded
Borel functions gives, in the weak sense of complex measures,
\[
 \mathbb E\,[w_Qh_Q(Y_Q)\delta_{Y_Q}]=\nu_Q.
\]
The fixed radial support is separated from zero, so uniformly in $y$
\[
 \norm{A^{P,\Theta}_a\delta_y}_2^2
 \le CB^2\norm\Theta_\infty^2.
\]
Consequently the finite random sum
\[
 \mathcal F(\mathbf Y)=A^{P,\Theta}_a
 \sum_Qe_Qw_Qh_Q(Y_Q)\delta_{Y_Q}
\]
is a strongly measurable, Bochner-integrable $L^2$-valued random
variable, and
\[
 A^{P,\Theta}_a\sum_Qe_Q\nu_Q=\mathbb E\mathcal F
 \quad\hbox{in }L^2.
\]
Hilbert-space Jensen gives
\[
 \left\|A^{P,\Theta}_a\sum_Qe_Q\nu_Q\right\|_2^2
 \le\mathbb E\norm{\mathcal F}_2^2.
\]
It is therefore enough to prove the asserted bound for every fixed
selection $y_Q\in Q$, with atomic coefficient
$b_Q:=e_Qh_Q(y_Q)$ of modulus at most one.  Fix such a selection and colour
the cubes by the bounded proximity coloring used in
Lemma~\ref{lem:multi-sublevel}.

If $1\le L<2$, the trivial correlation bound and
$\sum_{Q,Q'}w_Qw_{Q'}\le C\Lambda M$ on a fixed unit block prove
\eqref{eq:flat-high-block}.  Hence assume $L\ge2$.

After expanding the square, the correlation phase at a selected pair is
$P(y+z,y')-P(y+z,y)$.  To distinguish the two coefficient norms used here,
write
\[
 P(y+z,y')-P(y+z,y)
 =C_0(y,y')+\sum_{1\le|\beta|\le d}C_\beta(y,y')z^\beta,
 \qquad
 \mathbf C(y,y')=(C_\beta(y,y'))_\beta.
\]
For a fixed pair $(y,y')$, the nonconstant-$z$ coefficient norm of the
phase is the value norm $|\mathbf C(y,y')|_{\ell^2_\beta}$.  By contrast,
\begin{equation}\label{eq:flat-double-coefficient-norm}
 N:=\norm{\mathbf C}_{\mathrm{coeff},(y,y');\ell^2_\beta}
 =\left(\sum_{\beta,\gamma,\gamma'}
   |[y^\gamma(y')^{\gamma'}]C_\beta|^2\right)^{1/2}
\end{equation}
is the coefficient norm after the vector-valued function $\mathbf C$ is
again regarded as a polynomial in $(y,y')$.  Lemma~\ref{lem:same-shell-injective}
and projection onto the top total-degree part give $N\ge cL$;
lower-degree homogeneous blocks of $P$ cannot enter that projection.
Put $\mathbf p=\mathbf C/N$, so that this global coefficient norm of
$\mathbf p$ is one, and define
\[
 \mathcal B=\{(Q,Q'):
 |\mathbf p(y_Q,y_{Q'})|_{\ell^2_\beta}\le L^{-1/3}\}.
\]
Choose a scalar component $p_i$ with
$\norm{p_i}_{\rm coeff}\ge c$ and apply
\eqref{eq:weighted-pair-sublevel} to
$p_i/\norm{p_i}_{\rm coeff}$.  The bad vector sublevel is contained
in the resulting scalar $CL^{-1/3}$-sublevel, so that equation bounds its
weighted mass by
$C\Lambda M(L^{-c}+q^c)$.  On the complementary pairs the fixed-pair,
nonconstant-$z$ coefficient norm is at least
$NL^{-1/3}\ge cL^{2/3}$, uniformly even when lower blocks make $N$
arbitrarily large.

For clarity, this reduction enters the atomic square as follows.  If
$\mathcal K(y,y')$ denotes the correlation after the change $x=y+z$,
then, with the convention $\langle F,G\rangle=\int F\overline G$ fixed
above,
\begin{align}
 \left\|A^{P,\Theta}_a\sum_Qb_Qw_Q\delta_{y_Q}\right\|_2^2
 &=\sum_{Q,Q'}\overline{b_Q}b_{Q'}w_Qw_{Q'}
     \mathcal K(y_Q,y_{Q'}),
 \label{eq:flat-atomic-square}\\
 \mathcal K(y,y')
 &=\int e^{i[P(y+z,y')-P(y+z,y)]}
       \overline{k_a(z)}k_a(z-(y'-y))\,dz.
 \label{eq:flat-atomic-correlation}
\end{align}
Thus absolute values are taken only after the exact finite square
expansion, and the contribution of $\mathcal B$ is its weighted mass times
the trivial correlation bound.

If $n=1$, the one-dimensional conclusion of
Lemma~\ref{lem:rough-correlation} has no separation factor.  It therefore
gives $CB^C\norm\Theta_\infty^2L^{-c}$ directly on these complementary
pairs.  Together with the preceding bad-pair mass bound and
$\sum_{Q,Q'}w_Qw_{Q'}\le C\Lambda M$, this proves
\eqref{eq:flat-high-block} in dimension one.  Hence assume $n\ge2$.

For later use, the near-pair estimate is
\begin{equation}\label{eq:flat-near-pairs}
 \sum_{Q,Q'}w_Qw_{Q'}\mathbf1_{\{|y_Q-y_{Q'}|\le r\}}
 \le C\Lambda M(r+q)^n.
\end{equation}
Indeed, use the bounded proximity coloring from
Lemma~\ref{lem:multi-sublevel}.  For one color $\mathfrak c$, spread
$w_Q$ uniformly on a ball $B_Q$ of radius $cq$ centered at $y_Q$ and put
\[
 F_{\mathfrak c}
 =\sum_{Q\in\mathfrak c}\frac{w_Q}{|B_Q|}\mathbf1_{B_Q}.
\]
The balls in one color are disjoint,
$\norm{F_{\mathfrak c}}_\infty\le C\Lambda$, and
$\int F_{\mathfrak c}\le M$.  A selected pair with
$|y_Q-y_{Q'}|\le r$ smears into $|x-x'|\le r+Cq$.  Thus, for every
ordered pair of colors,
\[
 \iint_{\{|x-x'|\le r+Cq\}}
 F_{\mathfrak c}(x)F_{\mathfrak c'}(x')\,dx\,dx'
 \le C\Lambda(r+q)^n\int F_{\mathfrak c'}
 \le C\Lambda M(r+q)^n.
\]
Summing the boundedly many color pairs proves
\eqref{eq:flat-near-pairs}; when $r+q$ is bounded below, the same
conclusion follows directly from $M^2\le C\Lambda M$.
Take $r=L^{-1/(6d_*)}$, where $d_*=\max(d,1)$.  On a pair outside
$\mathcal B$ with $|y_Q-y_{Q'}|>r$, if $D_0\le d$ is the degree of its
nonconstant correlation phase, then
\[
 |\mathbf C(y_Q,y_{Q'})|\,|y_Q-y_{Q'}|^{D_0}
 \ge cL^{2/3}r^{D_0}
 \ge cL^{2/3-d/(6d_*)}\ge cL^{1/2}.
\]
Lemma~\ref{lem:rough-correlation} therefore bounds the remaining
correlations by $CB^C\norm\Theta_\infty^2L^{-c}$.  Splitting
\eqref{eq:flat-atomic-square} into bad pairs, nonbad near pairs, and the
remaining pairs, and then harmlessly enlarging the near sum to all near
pairs, gives, after decreasing $c>0$ once,
\begin{align*}
 \left\|A^{P,\Theta}_a\sum_Qb_Qw_Q\delta_{y_Q}\right\|_2^2
 &\le CB^C\norm\Theta_\infty^2\left[
    \sum_{(Q,Q')\in\mathcal B}w_Qw_{Q'}
   +\sum_{|y_Q-y_{Q'}|\le r}w_Qw_{Q'}
   +L^{-c}\sum_{Q,Q'}w_Qw_{Q'}\right]\\
 &\le CB^C\norm\Theta_\infty^2\Lambda M
   \{L^{-c}+q^c+(r+q)^n\}\\
 &\le CB^C\norm\Theta_\infty^2\Lambda M
   (L^{-c}+2^{-ct}).
\end{align*}
This is the desired bound for every fixed atomic selection; Jensen returns
to the original measures and proves \eqref{eq:flat-high-block}.
\end{proof}

\subsection{A one-shell phase-adapted packet theorem}

\begin{lemma}[Fourier expansion of a low homogeneous block]
\label{lem:block-fourier}
Let $\mathbf H=(H_1,\ldots,H_N)$ be a fixed basis of one homogeneous
degree block modulo pure phases.  Fix smooth cutoffs
$\chi,\widetilde\chi$ supported on fixed bounded cubes and equal to one
on the relevant smaller cubes.  Then, for every
$\mathbf C\in\R^N$ and every $\xi$ in the specified fixed cube,
\[
 \chi(X)\widetilde\chi(Y)
 e^{i\mathbf C\cdot[\mathbf H(X,Y)-\mathbf H(\xi,Y)]}
 =\sum_{u,v\in\mathbb Z^n}a_{u,v}(\mathbf C,\xi)
   e^{iu\cdot X}e^{iv\cdot Y},
\]
and there are numbers $A_{u,v}(\mathbf C)$, independent of $\xi$, such
that $|a_{u,v}(\mathbf C,\xi)|\le A_{u,v}(\mathbf C)$ and
\begin{equation}\label{eq:block-fourier}
 \sum_{u,v}A_{u,v}(\mathbf C)(1+|v|)
 \le C(1+\norm{\mathbf C})^{C_F}.
\end{equation}
\end{lemma}

\begin{proof}
Extend the cutoff expression periodically to a fixed torus and integrate
by parts more than $2n+2$ times.  Every derivative is bounded by a fixed
power of $1+\norm{\mathbf C}$, uniformly on the compact set of $\xi$.
More precisely, for $N>2n+2$ the integration-by-parts estimate is
\[
 |a_{u,v}(\mathbf C,\xi)|
 \le C_N(1+\norm{\mathbf C})^N(1+|u|+|v|)^{-N}.
\]
Taking the right side as $A_{u,v}(\mathbf C)$ and summing proves the
common-majorant assertion.
\end{proof}

\begin{proposition}[Phase-adapted one-shell packet principle]
\label{prop:direct-packet}
There are $c_P,c_E>0$ and $t_0$, depending only on $n,d$, such that the
following holds.  Let $P$ be a real polynomial of degree at most $d$, let
$\Theta\in L^\infty(\Sph^{n-1})$, and let
$R>0$, $t\in\mathbb N$, $q=R2^{-t}$,
$t\ge t_0$, and let $\mathcal Q$ be a finite family of pairwise disjoint
dyadic $q$-cubes (with $R$ and $q$ dyadic scales), with center $c_Q$ for
each $Q\in\mathcal Q$.  Suppose
\begin{equation}\label{eq:direct-packet-data}
 \supp d_Q\subset Q,\qquad \int d_Q=0,\qquad
 \norm{d_Q}_1\le\Lambda|Q|,\qquad
 M=\sum_{Q\in\mathcal Q}\norm{d_Q}_1.
\end{equation}
For $R\le\varepsilon\le2R$ put
\[
 \mathcal P_\varepsilon(x)=\sum_{Q\in\mathcal Q}
 \int_{\varepsilon<|x-y|\le2R}
 e^{i[P(x,y)-P(c_Q,y)]}K_\Theta(x-y)d_Q(y)\,dy.
\]
Then $\mathcal P=G+E$ with
\begin{align}
 \norm{V^2(G_\varepsilon:R\le\varepsilon\le2R)}_2^2
 &\le C\norm\Theta_\infty^2 2^{-c_Pt}\Lambda M,
 \label{eq:direct-packet-G}\\
 \norm{V^1(E_\varepsilon:R\le\varepsilon\le2R)}_1
 &\le C\norm\Theta_1 2^{-c_Et}M.
 \label{eq:direct-packet-E}
\end{align}
No pointwise bound on $d_Q$ is assumed.  Both estimates remain uniform
after replacing every $d_Q$ by $e_Qd_Q$, where $|e_Q|\le1$.
\end{proposition}

\begin{proof}
The optional factors $e_Q$ may be absorbed into the packets, so suppress
them.  Partition input into dyadic $R$-cubes $J$ and output into dyadic
$R$-cubes $I$; only $O_n(1)$ relative positions interact.  For one such
pair put
\[
 M_{I,J}=\sum_{Q\subset J}\norm{d_Q}_1.
\]
Every input block has boundedly many interacting output blocks, and hence
\begin{equation}\label{eq:direct-block-overlap}
 \sum_{I,J}M_{I,J}\le C_nM.
\end{equation}
The same annular support gives
\[
 \sup_I\#\{J:I\sim J\}+
 \sup_J\#\{I:I\sim J\}\le C_n,
\]
where $I\sim J$ denotes an interacting block pair.
Let $a$ be the center of $J$, set $x=a+RX$, $y=a+RY$, and define
\begin{equation}\label{eq:direct-normalized-phase}
 \Phi(X,Y)=P(a+RX,a+RY)-P(a,a+RY),
 \qquad \xi_Q=\frac{c_Q-a}{R}.
\end{equation}
Then $\Phi(0,Y)=0$, and the packet phase is exactly
$\Phi(X,Y)-\Phi(\xi_Q,Y)$.  The rescaled packets
$\widetilde d_Q(Y)=d_Q(a+RY)$ are supported on disjoint dyadic
$2^{-t}$-cubes and satisfy
\begin{equation}\label{eq:direct-rescaled-data}
 \norm{\widetilde d_Q}_1=R^{-n}\norm{d_Q}_1,
 \qquad \norm{\widetilde d_Q}_1\le\Lambda|\widetilde Q|,
 \qquad \widetilde M_{I,J}:=\sum_{Q\subset J}\norm{\widetilde d_Q}_1
 =R^{-n}M_{I,J}.
\end{equation}
All normalized packet cubes lie in the fixed normalized input block
$B_0=R^{-1}(J-a)$, as required in
Proposition~\ref{prop:flat-high-block}.
We prove the unit-scale estimate and suppress tildes on the packets, but
retain $\widetilde M_{I,J}$, until the last paragraph.  Insert fixed
smooth cutoffs which equal one on the normalized interacting input and
output blocks; these are the cutoffs in
Lemma~\ref{lem:block-fourier}.

Write the pure part of $\Phi$ as $p(X)+q(Y)$.  Its mixed representative
vanishes at $X=0$, so $\Phi(0,Y)=0$ forces $q$ to be constant.  In the
anchored difference $p(X)-p(\xi_Q)$ is a common output modulation times a
packet scalar and may be removed.  Modulo these factors, write the common
block phase in descending nonpure homogeneous degrees as
\[
 \Phi(X,Y)=H_1(X,Y)+\cdots+H_N(X,Y),
 \qquad \deg H_1>\cdots>\deg H_N\ge2,
\]
so the anchored phase is $\Phi(X,Y)-\Phi(\xi_Q,Y)$.  Put
$C_r=\norm{H_r}_{\rm coeff}$.

If $N=0$, only the removed output modulation and packet scalars remain;
Proposition~\ref{prop:terminalpacket} gives the required unit-block bound
with its positive terminal decay exponent and $E=0$.  The rescaling and
bounded-overlap summation in the last paragraph then give
\eqref{eq:direct-packet-G}.  Hence assume $N\ge1$ below.

Let $c_{\rm flat},C_B$ be respectively the decay exponent and the $\BV$
power in Proposition~\ref{prop:flat-high-block}, let $\delta_T$ be from
Proposition~\ref{prop:terminalpacket}, and choose
$C_{\rm Four}\ge C_F$ large enough to cover all Fourier losses in
Lemma~\ref{lem:block-fourier}.  Decrease $c_{\rm flat}$ if necessary so that
$0<c_{\rm flat}\le1$.  Set
\[
 \theta=\frac{c_{\rm flat}}{16(C_B+1)}.
\]
Choose $B_*>1$ so large that
\begin{equation}\label{eq:packet-hierarchy-ratio}
 \frac{2C_{\rm Four}}{B_*-1}
 \le\min\{\theta/2,c_{\rm flat}/8\},
\end{equation}
and then choose $0<\kappa_*\le1$ so small that, with
\begin{equation}\label{eq:packet-hierarchy}
 \kappa_r=\kappa_*B_*^{r-N},\qquad
 S_r=\sum_{i<r}\kappa_i,\qquad \gamma_r=\theta\kappa_r,
\end{equation}
one has
\begin{equation}\label{eq:packet-hierarchy-terminal}
 2C_{\rm Four}\sum_{r=1}^N\kappa_r\le\min\{\delta_T/2,1/2\}.
\end{equation}
Then $S_r\le\kappa_r/(B_*-1)$ and
\begin{align}
 \gamma_r-2C_{\rm Four}S_r&\ge\gamma_r/2,
 \label{eq:packet-hierarchy-a}\\
 c_{\rm flat}\kappa_r-C_B\gamma_r-2C_{\rm Four}S_r
 &\ge c_{\rm flat}\kappa_r/2,
 \label{eq:packet-hierarchy-b}\\
 c_{\rm flat}-C_B\gamma_r-2C_{\rm Four}S_r&\ge c_{\rm flat}/2.
 \label{eq:packet-hierarchy-c}
\end{align}

Call block $r$ low if $C_r\le2^{\kappa_rt}$ and expand it by
Lemma~\ref{lem:block-fourier}.  The product of the preceding low
expansions has weighted coefficient $\ell^1$ norm at most
\begin{equation}\label{eq:direct-fourier-loss}
 C2^{C_{\rm Four}S_rt}.
\end{equation}
For the nonlinear variation norms we use only finite Fourier sums and pass to
the limit afterwards.  Write
$a^{(i)}_{u_i,v_i}$ and $A^{(i)}_{u_i,v_i}$ for the coefficient and its
common majorant in the expansion of a low block $i$.  For $K\in\mathbb N$
retain only
\[
 \mathfrak F_K=\{(u_i,v_i)\in\mathbb Z^n\times\mathbb Z^n:
                    |u_i|+|v_i|\le K\}
\]
in each low expansion.  If $S$ is a set of expanded low blocks, define
\begin{equation}\label{eq:direct-fourier-tail}
 \omega_K(S)=
 \sum_{\substack{(u_i,v_i)_{i\in S}:\
       \max_{i\in S}(|u_i|+|v_i|)>K}}
 \left(\prod_{i\in S}A^{(i)}_{u_i,v_i}\right)
 \left(1+\left|\sum_{i\in S}v_i\right|\right),
 \qquad \omega_K(\varnothing)=0.
\end{equation}
Since
\[
 1+\left|\sum_{i\in S}v_i\right|
 \le\prod_{i\in S}(1+|v_i|),
\]
Lemma~\ref{lem:block-fourier} and dominated convergence for the resulting
absolutely summable product series give
\begin{equation}\label{eq:direct-fourier-tail-zero}
 \omega_K(S)\longrightarrow0\qquad(K\to\infty),
\end{equation}
uniformly in all packet centers $\xi_Q$.  Every estimate below is first
applied with this finite-mode convention.  Its constants are independent of
$K$.
Suppose $r$ is the first high block.  After a preceding Fourier tuple is
fixed, denote its parent-dependent Fourier coefficient by
$a_Q(\mathbf u,\mathbf v)$ and the product of the corresponding common
majorants in Lemma~\ref{lem:block-fourier} by
$A(\mathbf u,\mathbf v)$.  Then $|a_Q/A|\le1$, and summation of $A$
gives exactly \eqref{eq:direct-fourier-loss}.  Absorb this normalized
scalar together with the factor
\[
 e^{-i\sum_{i\ge r}H_i(\xi_Q,Y)}e^{iv\cdot Y}
\]
into the complex measure on packet $Q$.  This multiplier need not be
constant; total variation and \eqref{eq:direct-packet-data} are unchanged.
The remaining common operator phase has $H_r$ as its highest nonpure
block, and all lower blocks are unrestricted.

For this fixed Fourier tuple denote the resulting measures by $\mu_Q$,
and write $\Phi_r=\sum_{i\ge r}H_i$ for the common phase.  Partition
$[1,2]$ into $K_r\simeq2^{\gamma_rt}$ radial cells $I_\ell$ and set
\[
 \mathcal G_\ell(X)=
 \sum_Q\int_{\{|X-Y|\in I_\ell\}}
 e^{i\Phi_r(X,Y)}K_\Theta(X-Y)\,d\mu_Q(Y).
\]
Inserting all cell endpoints into an arbitrary variation partition gives
the pointwise deterministic bound
\begin{align}\label{eq:direct-high-grid-decomposition}
 V^2(\mathcal H_\varepsilon:1\le\varepsilon\le2)
 &\le C V^2\left(\sum_{\ell\le j}\mathcal G_\ell:
                  0\le j\le K_r\right)\notag\\
 &\quad+C\left(\sum_\ell
 V^2(\mathcal H_\varepsilon:\varepsilon\in I_\ell)^2\right)^{1/2},
\end{align}
where $\mathcal H$ is this component trajectory.  Reversing the order of
the grid partial sums, if necessary, does not change their variation.
Inside a cell use $V^2\le V^1$.
Lemma~\ref{lem:thinlayer}, whose proof uses only the average mass cap,
gives
\[
 \sup_\ell\norm{W_\ell}_\infty
 \le C\Lambda\norm\Theta_\infty(2^{-\gamma_rt}+2^{-t}),
 \qquad
 \sum_\ell\norm{W_\ell}_1\le C\norm\Theta_1\widetilde M_{I,J}.
\]
Thus the squared $L^2$ norm of the within-cell square function is at most
$C\norm\Theta_\infty^2\Lambda\widetilde M_{I,J}2^{-\gamma_rt}$.

For the grid partial sums, fix
$(\epsilon_\ell)_\ell\in\{-1,0,1\}^{K_r}$ before evaluation in $X$.
The corresponding sharp radial profile
$a_\epsilon=\sum_\ell\epsilon_\ell\mathbf1_{I_\ell}$ satisfies
\[
 \norm{a_\epsilon}_\infty+\Var(a_\epsilon)\le CK_r.
\]
Proposition~\ref{prop:flat-high-block} and
Lemma~\ref{lem:RM} therefore give, uniformly over these fixed signs,
using
$\log^2(2+K_r)\lesssim(1+\gamma_rt)^2\lesssim t^2$,
\begin{equation}\label{eq:direct-high-grid}
 Ct^2\norm\Theta_\infty^2
 2^{C_B\gamma_rt}\Lambda\widetilde M_{I,J}
 \{2^{-c_{\rm flat}\kappa_rt}+2^{-c_{\rm flat}t}\}.
\end{equation}
In particular, no $X$-dependent selector is inserted into the fixed-sign
estimate; all such dependence was already confined to the positive
within-cell term in \eqref{eq:direct-high-grid-decomposition}.  Denote by
$G^{(r),I,J}$ the full component obtained after summing all Fourier tuples
preceding the first high block $r$.  The variation triangle inequality and
the square of the common-majorant sum
\eqref{eq:direct-fourier-loss} give the complete estimate
\begin{align}
 \norm{V^2(G^{(r),I,J}_\varepsilon)}_2^2
 &\le C\norm\Theta_\infty^2\Lambda\widetilde M_{I,J}
  2^{2C_{\rm Four}S_rt}\Bigl[
    2^{-\gamma_rt}
   +t^2 2^{C_B\gamma_rt}
    \{2^{-c_{\rm flat}\kappa_rt}+2^{-c_{\rm flat}t}\}
  \Bigr].
 \label{eq:direct-high-complete-exponents}
\end{align}
There are no suppressed exponential factors in this line: the first term
is the within-cell square function and the other two are respectively the
large-coefficient and mesh terms in \eqref{eq:direct-high-grid}.  The
three hierarchy inequalities say, term by term,
\begin{align*}
 2C_{\rm Four}S_r-\gamma_r&\le-\gamma_r/2,\\
 2C_{\rm Four}S_r+C_B\gamma_r-c_{\rm flat}\kappa_r
   &\le-c_{\rm flat}\kappa_r/2,\\
 2C_{\rm Four}S_r+C_B\gamma_r-c_{\rm flat}
   &\le-c_{\rm flat}/2.
\end{align*}
Since $t^2 2^{-at}\le C_a2^{-at/2}$ for $t\ge t_0$, the right side of
\eqref{eq:direct-high-complete-exponents} is bounded by
\[
 C\norm\Theta_\infty^2\Lambda\widetilde M_{I,J}2^{-c_rt},
 \qquad
 c_r=\frac14\min\{\gamma_r,c_{\rm flat}\kappa_r,c_{\rm flat}\}>0.
\]
This proves every high branch with $E=0$ without changing any exponent
used later.

If every block is low, expand them all.  Under the normalized variables in
\eqref{eq:direct-normalized-phase}, the center of the normalized packet
cube is
\begin{equation}\label{eq:direct-normalized-packet-center}
 \zeta_Q:=\xi_Q=\frac{c_Q-a}{R}.
\end{equation}
For a combined input frequency $v$, write
\[
 e^{iv\cdot Y}d_Q(Y)=e^{iv\cdot\zeta_Q}d_Q(Y)+r_{Q,v}(Y),
 \qquad
 r_{Q,v}(Y)=(e^{iv\cdot Y}-e^{iv\cdot\zeta_Q})d_Q(Y),
\]
and, since $|Y-\zeta_Q|\le C2^{-t}$ on the normalized cube,
\[
 \norm{r_{Q,v}}_1\le C|v|2^{-t}\norm{d_Q}_1.
\]
The frozen term remains mean zero, so
Proposition~\ref{prop:terminalpacket} applies with constant packet
coefficients of modulus at most one: divide the parent-dependent Fourier
coefficient by its product common majorant and absorb the normalized
scalar into the packet.  Summation of the remaining majorants is governed
by \eqref{eq:block-fourier}, and the output Fourier factor is one fixed
modulation.
Equation \eqref{eq:packet-hierarchy-terminal} leaves at least
$2^{-\delta_Tt/2}$ after the squared Fourier loss.  Young's inequality and
the weighted Fourier sum give total $L^1$ variation at most
$C\norm\Theta_1 2^{-t/2}\widetilde M_{I,J}$ for the remainders.

We now remove the finite-mode cutoff.  For a normalized interacting block
pair $(I,J)$, let
\[
 \mathcal P^{I,J,K}=G^{I,J,K}+E^{I,J,K}
\]
be the finite-mode trajectory and decomposition constructed above, with the
output restricted to the normalized cube corresponding to $I$.  If $r$ is
the first high block, put $\mathcal I_r=\{i:i<r\}$ and
\[
 \mathfrak B_{r,t}=2^{-\gamma_rt}
 +t^2 2^{C_B\gamma_rt}
   \{2^{-c_{\rm flat}\kappa_rt}+2^{-c_{\rm flat}t}\}.
\]
For $K'>K$, applying the same within-cell and grid estimates only to the
Fourier tuples omitted at level $K$, and then using the variation triangle
inequality, gives
\begin{equation}\label{eq:direct-high-fourier-cauchy}
 \norm{G^{I,J,K'}-G^{I,J,K}}_{L^2(V^2)}
 \le C\norm\Theta_\infty
       (\Lambda\widetilde M_{I,J})^{1/2}
       \mathfrak B_{r,t}^{1/2}\,\omega_K(\mathcal I_r),
 \qquad E^{I,J,K'}-E^{I,J,K}=0.
\end{equation}
If every block is low, put $S=\{1,\ldots,N\}$.  The terminal-packet
estimate for the frozen terms and Young's inequality for the remainders give
\begin{align}
 \norm{G^{I,J,K'}-G^{I,J,K}}_{L^2(V^2)}
 &\le C\norm\Theta_\infty 2^{-\delta_Tt/2}
       (\Lambda\widetilde M_{I,J})^{1/2}\omega_K(S),
 \label{eq:direct-low-good-fourier-cauchy}\\
 \norm{E^{I,J,K'}-E^{I,J,K}}_{L^1(V^1)}
 &\le C\norm\Theta_1 2^{-t}\widetilde M_{I,J}\omega_K(S).
 \label{eq:direct-low-error-fourier-cauchy}
\end{align}
Indeed, in the first line the unweighted coefficient tail is bounded by
\eqref{eq:direct-fourier-tail}; in the second line the factor
$1+|\sum_i v_i|$ absorbs the freezing error.  These are the same estimates
used above, with the full coefficient sum replaced by its tail, so no decay
exponent is changed.

There are only finitely many interacting block pairs.  Hence
\eqref{eq:direct-fourier-tail-zero} and the bounded-overlap inequalities show
that $(G^K)$ is Cauchy in $L^2(V^2)$ and $(E^K)$ is Cauchy in $L^1(V^1)$.
All finite-mode component paths vanish at the upper endpoint
$\varepsilon=2R$.  Consequently the same is true of every difference
path, and
\[
 \sup_{R\le\varepsilon\le2R}|H_\varepsilon|
 \le V^a(H_\varepsilon:R\le\varepsilon\le2R),
 \qquad a\in\{1,2\},
\]
so convergence in the anchored variation spaces also controls the
supremum of the trajectories.
Lemma~\ref{lem:anchored-variation-completion} constructs anchored limits
$G,E$ for the full finite-mode sequences and gives the corresponding
$L^2(V^2)$ and $L^1(V^1)$ convergence.

It remains to identify this limit.  The common majorants in
Lemma~\ref{lem:block-fourier} imply that the truncated Fourier products
converge absolutely and uniformly on the fixed normalized input, output, and
packet-center boxes.  Since the kernel is restricted to $1<|X-Y|\le2$ and
there are only finitely many packets, for each block pair
\[
 \sup_{1\le\varepsilon\le2}
 |\mathcal P^{I,J,K}_\varepsilon(X)
   -\mathcal P^{I,J}_\varepsilon(X)|
 \le o_K(1)\,C\norm\Theta_\infty
       \sum_{Q\subset J}\norm{d_Q}_1,
 \qquad o_K(1)\longrightarrow0,
\]
on its normalized output cube.  The right side belongs to both $L^1$ and
$L^2$ there, and the convergence is uniform in $\varepsilon$ almost
everywhere.  The last assertion of
Lemma~\ref{lem:anchored-variation-completion} identifies
$\mathcal P^{I,J}=G^{I,J}+E^{I,J}$ (after rescaling, at the physical scale).
The Fatou conclusion of that lemma passes the uniform finite-mode bounds
to these limits without changing an exponent.

Let $G^{I,J}_\varepsilon=\mathbf1_I G_{I,J,\varepsilon}$ and
$E^{I,J}_\varepsilon=\mathbf1_I E_{I,J,\varepsilon}$ denote the resulting
physical paths.  Define
\[
 G_\varepsilon=\sum_I\sum_{J:\,J\sim I}G^{I,J}_\varepsilon,
 \qquad
 E_\varepsilon=\sum_I\sum_{J:\,J\sim I}E^{I,J}_\varepsilon.
\]
The block family is finite, and the limiting blockwise identities are exact
on $I$.  Therefore
$\mathcal P_\varepsilon=G_\varepsilon+E_\varepsilon$ for the original
trajectory.

Take $c_P>0$ to be the minimum of $\delta_T/2$ and the finitely many
high-branch exponents $c_r$, over all possible $N\le d-1$, and take
$c_E=1/2$.  At unit scale
$\widetilde M_{I,J}=R^{-n}M_{I,J}$; on returning to $x$, an
$L^2$-square and an $L^1$ norm both acquire the Jacobian $R^n$.  Thus
\begin{align*}
 \norm{V^2(G^{I,J}_\varepsilon)}_2^2
 &\le C\norm\Theta_\infty^2
       2^{-c_Pt}\Lambda M_{I,J},\\
 \norm{V^1(E^{I,J}_\varepsilon)}_1
 &\le C\norm\Theta_1
       2^{-c_Et}M_{I,J}.
\end{align*}
Since the output cubes $I$ are disjoint and each has only $O_n(1)$ input
neighbors, the variation triangle inequality and Cauchy--Schwarz give
\[
 \norm{V^2(G_\varepsilon)}_2^2
 \le C_n\sum_{I,J}\norm{V^2(G^{I,J}_\varepsilon)}_2^2,
 \qquad
 \norm{V^1(E_\varepsilon)}_1
 \le\sum_{I,J}\norm{V^1(E^{I,J}_\varepsilon)}_1.
\]
Using \eqref{eq:direct-block-overlap} proves
\eqref{eq:direct-packet-G}--\eqref{eq:direct-packet-E}.
\end{proof}

\subsection{Exact oriented algebra and the collective high-shell estimate}

Let
\[
 \mathcal V_m=\{\hbox{homogeneous degree-$m$ polynomials in $(x,y)$}\}/
 \{p(x)+q(y)\}.
\]
Each class has a unique representative consisting only of monomials which
have positive degree in both \(x\) and \(y\).  We use this representative
whenever an element of \(\mathcal V_m\) is evaluated.  Let
\(\mathcal W_m\) be the space of homogeneous degree-\(m\) polynomials in
\((W,X)\).  For such a representative put
\begin{equation}\label{eq:oriented-map}
 \mathcal EH(W,X)=\Mix_{W,X}
 \{H(W+X,0)-H(W+X,W)\}.
\end{equation}

\begin{lemma}[Oriented graph injectivity]\label{lem:oriented-injective}
For every integer $m\ge2$, the map
\(\mathcal E:\mathcal V_m\to\mathcal W_m\) is injective and
\begin{equation}\label{eq:oriented-injective}
 \norm{\mathcal EH}_{\rm coeff}\ge c_{n,m}
 \norm{[H]}_{\rm coeff}.
\end{equation}
Moreover
\begin{equation}\label{eq:oriented-bidegrees}
 \mathcal EH=\sum_{\substack{A+B=m\\ A,B\ge1}}E_{A,B}(W,X),
\end{equation}
where \(E_{A,B}\) is bihomogeneous of bidegree \((A,B)\).
\end{lemma}

\begin{proof}
If the mixed part in \eqref{eq:oriented-map} vanishes, the expression
in braces is \(A(W)+B(X)\).  At \(W=0\) it vanishes, and hence \(B\) is
constant.  Substituting \(x=W+X\) and \(y=W\) shows that
\(H(x,0)-H(x,y)\) is a function of \(y\) alone.  Thus \(H\) is pure.
Quantitative injectivity gives \eqref{eq:oriented-injective} by norm
equivalence on the fixed finite-dimensional spaces.  Homogeneity and the
identities at \(W=0\) and \(X=0\) give
\eqref{eq:oriented-bidegrees}.
\end{proof}

We next record the exact projective calculation used below.  Its purpose is
to identify ordinary coefficients of one common numerator; no formal
asymptotic expansion in the scale ratio is used.

\begin{lemma}[Exact projective numerator tags]\label{lem:projective-tags}
Fix \(D_0\ge1\), a multi-index \(\gamma\) with \(q=|\gamma|\), and a
polynomial \(Q(s)=\sum_{p=0}^Pc_ps^p\), where $P\in\mathbb N_0$ and
\(P+q\le D_0-1\).  Let $\delta>0$, $0<\alpha\le1$, and let
$\omega\in\R^k$, where $k$ is the length of $\gamma$.  On the set
$D:=V-\delta u\ne0$, put
\[
 t=\frac{\delta u}{D},
 \qquad r_\perp=\frac{\delta\alpha\omega}{D},
\]
and put
\[
 F_{q,Q}(u,V,\omega)=r_\perp^\gamma Q(t).
\]
Then
\begin{equation}\label{eq:exact-tag-numerator-new}
 D^{D_0+2}\partial_u\partial_VF_{q,Q}
 =\alpha^q\omega^\gamma
   \sum_{j=0}^{D_0-q}
   \delta^{q+1+j}\tau_{q,j}(Q)
   u^jV^{D_0-q-j},
\end{equation}
where the functionals \(\tau_{q,j}\) depend only on \(D_0,q,j\).  If
\(q=0\) and $P\ge1$, the map
\[
 (c_1,\ldots,c_P)\longmapsto
 (\tau_{0,0}(Q),\ldots,\tau_{0,P-1}(Q))
\]
is invertible.  If \(q\ge1\), the map
\[
 (c_0,\ldots,c_P)\longmapsto
 (\tau_{q,0}(Q),\ldots,\tau_{q,P}(Q))
\]
is invertible.  The inverse norms are bounded in terms of \(D_0\) only.
When $q=P=0$, the first map is the empty map and the assertion is
vacuous.
In particular, every detecting tag in the case $q=0$, $P\ge1$ has order
\(\kappa=q+1+j\le P\), and in the case $q\ge1$ it has order
\(\kappa=q+1+j\le q+P+1\).
After multiplication by $D^{D_0+2}$, the displayed equality is a
polynomial identity.  Although the derivation below uses the dense open
set on which $DV\ne0$, the cleared identity therefore extends across both
$D=0$ and $V=0$.
\end{lemma}

\begin{proof}
Set \(z=\delta u/V\) and
\[
 H_{q,Q}(z)=(1-z)^{-q}Q\left(\frac{z}{1-z}\right).
\]
Apart from the factor \(\alpha^q\omega^\gamma\), one has
\[
 F_{q,Q}=\delta^qV^{-q}H_{q,Q}(z).
\]
Consequently
\begin{equation}\label{eq:projective-mixed-derivative-new}
 \partial_u\partial_VF_{q,Q}
 =-\alpha^q\omega^\gamma\delta^{q+1}V^{-q-2}
 \bigl[(q+1)H_{q,Q}'(z)+zH_{q,Q}''(z)\bigr].
\end{equation}
Multiplication by
\(D^{D_0+2}=V^{D_0+2}(1-z)^{D_0+2}\) gives a polynomial.
The coefficient of \(z^j\) becomes the ordinary coefficient of
\(u^jV^{D_0-q-j}\), with the exact factor
\(\delta^{q+1+j}\).  This proves \eqref{eq:exact-tag-numerator-new}.

Write \(H_{q,Q}(z)=\sum_{s\ge0}h_sz^s\).  The coefficient of
\(z^{s-1}\) in the bracket in
\eqref{eq:projective-mixed-derivative-new} is
\(s(q+s)h_s\).  Multiplication by \((1-z)^{D_0+2}\) is triangular with
diagonal one.  It therefore suffices to inspect
\((h_1,\ldots,h_P)\) when \(q=0\), and
\((h_1,\ldots,h_{P+1})\) when \(q\ge1\).  For \(q=0\),
\(Q(z/(1-z))\) is triangular with diagonal one on
\((c_1,\ldots,c_P)\).  For \(q\ge1\), if the latter \(P+1\)
coefficients vanished, then, after \(s=z/(1-z)\),
\[
 (1+s)^qQ(s)=c_0+O(s^{P+2}).
\]
The first \(P\) equations make \(Q\) equal to \(c_0\) times the
degree-\(P\) Taylor polynomial of \((1+s)^{-q}\), while the coefficient
of \(s^{P+1}\) is a nonzero multiple of
\(c_0\binom{q+P}{P+1}\).  Hence \(c_0=0\), and then all \(c_p=0\).
Finite-dimensional norm equivalence gives the quantitative assertion.
\end{proof}

We next define the unequal-shell phase appearing in the correlation.
Throughout the remainder of this subsection assume that \(P\) has a
highest nonpure homogeneous block \(P_m\); necessarily \(m\ge2\).
Write, modulo pure phases,
\begin{equation}\label{eq:translated-homogeneous-blocks-new}
 P(a+X,a+Y)=p_a(X)+q_a(Y)+\sum_{\ell=2}^mH_{\ell,a}(X,Y),
\end{equation}
where every \(H_{\ell,a}\) is homogeneous of degree \(\ell\), and where
\(H_{m,a}=P_m\).  For \(W,X\in\mathbb R^n\), set
\begin{equation}\label{eq:actual-correlation-phase-new}
 \Phi_{a,S}(W,X)
 =P(a+S(W+X),a)-P(a+S(W+X),a+SW).
\end{equation}
Modulo a function independent of \(X\),
\begin{equation}\label{eq:actual-oriented-decomposition-new}
 \Phi_{a,S}(W,X)
 \equiv\sum_{\ell=2}^mS^\ell\mathcal EH_{\ell,a}(W,X).
\end{equation}
Indeed, the pure output phase cancels, the pure input phase is independent
of \(X\), and a mixed representative has no pure-\(X\) term in the
difference defining \(\mathcal E\).

Fix once and for all annular constants \(0<c_*<C_*<\infty\) containing
all values of \(|W|\) allowed by the two radial supports, and then choose
\(C_0\) sufficiently large, depending only on those fixed supports and
annular constants.  Thus \(R\le S/C_0\) makes
\(\delta=R/(S\rho)\) uniformly small throughout the overlap region.

For completeness we now give a concrete rational atlas and frame.  For
\(1\le j\le n\) and \(s\in\{-1,1\}\) put \(c_{j,s}=se_j\) and
\[
 U_{j,s}=\{\theta\in\Sph^{n-1}:c_{j,s}\cdot\theta>(2\sqrt n)^{-1}\},
 \qquad
 U_{j,s}^*=\{\theta\in\Sph^{n-1}:c_{j,s}\cdot\theta>(4\sqrt n)^{-1}\}.
\]
The smaller sets cover the sphere and
\(U_{j,s}\Subset U_{j,s}^*\).  Choose a fixed signed permutation
\(C_{j,s}\in O(n)\) with \(C_{j,s}e_1=c_{j,s}\).  In stereographic
coordinates \(\vartheta\in\mathbb R^{n-1}\) from the pole \(-c_{j,s}\),
write
\[
 \theta_{j,s}(\vartheta)
 =C_{j,s}\left(
   \frac{1-|\vartheta|^2}{1+|\vartheta|^2},
   \frac{2\vartheta}{1+|\vartheta|^2}
  \right).
\]
On \(U_{j,s}^*\) the coordinate \(\vartheta\) ranges over a fixed bounded
set; the chart Jacobian and its inverse are therefore uniformly bounded.
Let \(a_{j,s}=-c_{j,s}\), choose a fixed signed permutation
\(B_{j,s}\in O(n)\) satisfying \(B_{j,s}e_1=a_{j,s}\), and define the
Householder frame
\begin{equation}\label{eq:rational-householder-frame}
 O_{j,s}(\vartheta)
 =\left(I-2\frac{(a_{j,s}-\theta_{j,s}(\vartheta))
 (a_{j,s}-\theta_{j,s}(\vartheta))^T}
 {|a_{j,s}-\theta_{j,s}(\vartheta)|^2}\right)B_{j,s}.
\end{equation}
This matrix is orthogonal and
\(O_{j,s}(\vartheta)e_1=\theta_{j,s}(\vartheta)\).  Moreover,
\[
 |a_{j,s}-\theta|^2=2(1+c_{j,s}\cdot\theta)
 \ge2\bigl(1+(4\sqrt n)^{-1}\bigr)
 \quad\hbox{on }U_{j,s}^*.
\]
Thus every matrix entry is rational in \(\vartheta\); after taking the
fixed products needed as common denominators, those denominators and their
reciprocals are uniformly bounded on the enlarged chart domains.  This
also remains true for the finitely many fixed powers that occur below.

Order these \(2n\) charts as \(U_1,\ldots,U_{2n}\), relabeling the
corresponding maps as \(\theta_\chi\) and \(O_\chi\), and make the smaller
cover disjoint by putting
\[
 E_1=U_1,\qquad E_\chi=U_\chi\setminus\bigcup_{\nu<\chi}U_\nu.
\]
These are measurable, disjoint, and cover the sphere.  Write \(\chi(W)\)
for the unique active index satisfying \(W/|W|\in E_{\chi(W)}\).
Finally, the annular coordinate map
\((\rho,\vartheta)\mapsto W=\rho\theta_\chi(\vartheta)\) is uniformly
bi-Lipschitz for \(c_*\le\rho\le C_*\) on each enlarged chart.
We use the following precise convention.  A \emph{cleared
chart-polynomial vector} on this annulus is a vector $F_\chi$ for which
\[
 \widetilde F_\chi(\rho,\vartheta)
 :=F_\chi(\rho\theta_\chi(\vartheta))
\]
is an ordinary polynomial vector of degree at most a fixed structural
bound, after the fixed common denominator has been incorporated into
$F_\chi$.  Its chart coefficient norm is
\begin{equation}\label{eq:chart-polynomial-norm}
 \norm{F_\chi}_{\mathrm{ch},\chi}
 :=\norm{\widetilde F_\chi}_{\mathrm{coeff},(\rho,\vartheta)}.
\end{equation}
On the fixed enlarged coordinate box, finite-dimensional norm equivalence
gives
\begin{equation}\label{eq:chart-polynomial-gradient}
 \norm{\widetilde F_\chi}_\infty+
 \norm{\nabla_{\rho,\vartheta}\widetilde F_\chi}_\infty
 \le C\norm{F_\chi}_{\mathrm{ch},\chi}.
\end{equation}
Any two fixed clearing conventions give equivalent norms: after comparing
supremum norms their bounded, nowhere-vanishing denominator ratio is
absorbed, and coefficient/supremum norm equivalence is used on the fixed
polynomial spaces.  In particular, all later polynomial sublevel
arguments are made in $(\rho,\vartheta)$ coordinates with
\eqref{eq:chart-polynomial-norm}; they do not regard $\rho=|W|$ as an
ambient polynomial in $W$.

In frame coordinates define
\begin{equation}\label{eq:oriented-chart-polynomial-new}
 G_{\ell,a,\chi}(\vartheta;T)
 =\rho^{-\ell}\mathcal EH_{\ell,a}
 \bigl(\rho O_\chi e_1,\rho O_\chi T\bigr).
\end{equation}
This is independent of \(\rho\), is a polynomial in \(T\) with no
constant term, and has degree at most \(\ell-1\).

\begin{proposition}[Oriented unequal-shell coefficient]
\label{prop:oriented-coefficient}
Assume \(n\ge2\).  Let \(P_m\) be the highest nonpure homogeneous block
of a real polynomial \(P\) of degree at most $d$, put
\[
 L_R=R^m\norm{P_m}_{\rm coeff},
\]
and suppose \(0<R\le S/C_0\).  For \(y=a+SW\), assume
\(c_*\le|W|\le C_*\), where $a,W\in\R^n$; put $\rho=|W|$.
On every chart \(\chi\) there is a bounded-degree cleared
chart-polynomial vector
\(\mathbf p_{a,R,S,\chi}(W)\), independent of the polar band parameter
\(\alpha\), such that
\begin{equation}\label{eq:oriented-p-norm}
 \norm{\mathbf p_{a,R,S,\chi}}_{\rm coeff}\ge1.
\end{equation}
On the disjoint active-chart partition define
\begin{equation}\label{eq:active-oriented-p-new}
 \mathbf p_{a,R,S}(W)
 :=\mathbf p_{a,R,S,\chi(W)}(W).
\end{equation}
Thus \(\mathbf p\) is piecewise chart-polynomial, with a unique value at
every \(W\).  Every later incidence argument is applied after splitting
the atoms into the finitely many active sets \(E_\chi\).
Here $\norm{\mathbf p}_{\rm coeff}$ means the chart norm
\eqref{eq:chart-polynomial-norm}, whereas
$|\mathbf p(W)|$ denotes the Euclidean norm in its fixed finite-dimensional
value space.  The coefficient lower bound in
\eqref{eq:oriented-p-norm} is not asserted to be a pointwise lower bound;
small values of $|\mathbf p(W)|$ are handled by the incidence estimate.

More precisely, put \(\delta=R/(S\rho)\).  For
\(0<\alpha\le1\), \(\omega\in\Sph^{n-2}\), and on each member of the
fixed finite family of \((u,V)\)-rectangles arising from the radial
supports, use variables
\begin{equation}\label{eq:unequal-projective-map-new}
 D=V-\delta u,
 \qquad
 T=T_{\delta,\alpha,\omega}(u,V)
 =\frac{(\delta u,\delta\alpha\omega)}{D}.
\end{equation}
Every such rectangle has a fixed enlargement on which
\(D=V-\delta u\) is uniformly bounded above and away from zero.  The
actual unequal-shell phase is
\begin{align}
 \Psi_{a,R,S,W,\alpha,\omega}(u,V)
 &=\Phi_{a,S}\bigl(W,\rho O_\chi
 T_{\delta,\alpha,\omega}(u,V)\bigr) \notag\\
 &=P\left(a+S\rho O_\chi
      \frac{(V,\delta\alpha\omega)}{D},a\right) \notag\\
 &\quad-P\left(a+S\rho O_\chi
      \frac{(V,\delta\alpha\omega)}{D},
      a+S\rho O_\chi e_1\right).
 \label{eq:actual-unequal-phase-new}
\end{align}
Since \(\deg P\le d\), the
phase in \eqref{eq:actual-unequal-phase-new} has the canonical common
denominator
\[
 Q(u,V)=D^d=(V-\delta u)^d.
\]
Define its literal mixed numerator by
\begin{equation}\label{eq:actual-mixed-numerator-new}
 \mathcal N_\Psi(u,V)
 :=D^{d+2}\partial_u\partial_V\Psi(u,V).
\end{equation}
This is a polynomial of bounded degree.  For the canonical presentation,
\[
 \mathfrak M_{D^d}\Psi
 =D^{3d}\partial_u\partial_V\Psi
 =D^{2d-2}\mathcal N_\Psi.
\]
Multiplication by \(D^{2d-2}\) is uniformly bounded below in coefficient
norm on these fixed polynomial spaces: this follows from the equivalence
of coefficient and supremum norms and the lower bound for \(|D|\) on the enlarged
rectangle.  Equation~\eqref{eq:mixed-to-quotient} therefore gives
\begin{equation}\label{eq:actual-mixed-to-quotient-new}
 \norm{[\Psi]}_{{\rm rat},D^d}
 \ge c\norm{\mathcal N_\Psi}_{\rm coeff(u,V)}.
\end{equation}

For every \(\alpha\) there is a bounded-degree vector
\(\mathbf q_{\alpha,W}(\omega)\), made from selected ordinary
coefficients of \(\mathcal N_\Psi\), such that
\begin{align}
 \norm{[\Psi_{a,R,S,W,\alpha,\omega}]}_{{\rm rat},D^d}
 &\ge c\abs{\mathbf q_{\alpha,W}(\omega)},
 \label{eq:oriented-band-lower}\\
 \norm{\mathbf q_{\alpha,W}}_{{\rm coeff},\omega}
 &\ge c\alpha^dL_R
 \abs{\mathbf p_{a,R,S,\chi}(W)}.
 \label{eq:oriented-band-lower-new-b}
\end{align}
The constants are independent of \(a\) and of all coefficients of total
degree less than \(m\).
\end{proposition}

\begin{proof}
By \eqref{eq:actual-oriented-decomposition-new} and
\eqref{eq:oriented-chart-polynomial-new}, modulo
\(F(u)+G(V)\) the phase in \eqref{eq:actual-unequal-phase-new} is
\begin{equation}\label{eq:actual-phase-block-sum-new}
 \sum_{\ell=2}^mS^\ell\rho^\ell
 G_{\ell,a,\chi}\left(\vartheta;
 \frac{(\delta u,\delta\alpha\omega)}{D}\right).
\end{equation}
Expand
\[
 G_{\ell,a,\chi}(\vartheta;t,r_\perp)
 =\sum_{q,\,|\gamma|=q}r_\perp^\gamma
 Q_{\ell,a,\chi,\gamma}(\vartheta;t).
\]
Here
\(\deg Q_{\ell,a,\chi,\gamma}+q\le\ell-1\), and the constant
coefficient of \(Q_{\ell,a,\chi,0}\) is zero because
\(G_{\ell,a,\chi}(\vartheta;0)=0\).
Lemma~\ref{lem:projective-tags}, used with \(D_0=d\), shows that the
coefficient of
\(\omega^\gamma u^jV^{d-q-j}\) in \(\mathcal N_\Psi\) is exactly
\begin{equation}\label{eq:actual-tag-coefficient-new}
 C_{q,\gamma,j,\alpha}(\rho,\vartheta)
 =\alpha^q\sum_{\ell=2}^m
 R^\kappa S^{\ell-\kappa}\rho^{\ell-\kappa}
 \tau_{q,j}(Q_{\ell,a,\chi,\gamma}),
 \qquad \kappa=q+1+j.
\end{equation}
Thus different tag orders are different ordinary numerator monomials;
there is no appeal to a formal expansion in \(\delta\).

We make the finite-dimensional coefficient step explicit.  Fix the mixed
monomial basis of $\mathcal V_m$, and use ordinary monomial coefficients in
the chart variable $\vartheta$.  Choose once and for all a nonvanishing
polynomial $\Delta_\chi(\vartheta)$ which clears every denominator
introduced by the frame $O_\chi$ in degrees at most $d$; both
$\Delta_\chi$ and its reciprocal are bounded on $U_\chi^*$.  For
$T=(t,r_\perp)\in\R\times\R^{n-1}$ define the cleared restriction map
\begin{equation}\label{eq:cleared-oriented-restriction-map}
 (\mathsf R_\chi H)(\vartheta;T)
 =\Delta_\chi(\vartheta)
   \mathcal EH\bigl(O_\chi(\vartheta)e_1,
                    O_\chi(\vartheta)T\bigr).
\end{equation}
By homogeneity this is precisely
$\Delta_\chi G_{m,a,\chi}$ when $H=P_m$; in particular it is independent
of $a$ and $\rho$.  Write
\[
 \mathsf R_\chi H
 =\sum_{\gamma}r_\perp^\gamma
   \widetilde Q_{\chi,\gamma}(\vartheta;t),
 \qquad q=|\gamma|.
\]
The $q=0$ polynomial has zero constant coefficient, and
$\deg_t\widetilde Q_{\chi,\gamma}+q\le m-1$.

Let
\begin{equation}\label{eq:detecting-tag-index-set}
 \mathfrak I_m
 =\{(0,j):0\le j\le m-2\}
  \cup\{(\gamma,j):1\le|\gamma|\le m-1,
                     \ 0\le j\le m-1-|\gamma|\},
\end{equation}
where $0$ denotes the zero multi-index, and put
$\kappa(\gamma,j)=|\gamma|+1+j$.  Let $D_\chi$ be a fixed degree bound for
the cleared chart polynomials and set
\[
 \mathcal Z_{m,\chi}
 =\bigoplus_{(\gamma,j)\in\mathfrak I_m}
   \mathcal P_{\le D_\chi}(\vartheta),
\]
with the Euclidean norm of all ordinary monomial coefficients.  All tag
functionals below are those of Lemma~\ref{lem:projective-tags} with
$D_0=d$, extended coefficientwise in $\vartheta$.  Define
\begin{equation}\label{eq:explicit-detecting-tag-map}
 \mathsf T_\chi:\mathcal V_m\longrightarrow\mathcal Z_{m,\chi},
 \qquad
 \mathsf T_\chi H
 =\bigl(
   \tau_{|\gamma|,j}
   (\widetilde Q_{\chi,\gamma}(\vartheta;\cdot))
   \bigr)_{(\gamma,j)\in\mathfrak I_m}.
\end{equation}
Multiplication by $\Delta_\chi$ commutes with the tag functionals because
it is independent of $t$ and $r_\perp$.  Lemma~\ref{lem:projective-tags}
therefore recovers, separately for each $\gamma$, all coefficients of
$\widetilde Q_{\chi,\gamma}$: for $q=0$ it recovers the nonconstant
$t$-coefficients, and for $q\ge1$ it recovers all of them.  Thus vanishing
of $\mathsf T_\chi H$ implies $\mathsf R_\chi H=0$.

The restriction map $\mathsf R_\chi$ is injective.  Indeed, its vanishing
implies $\mathcal EH(\theta,X)=0$ for every $X$ and every direction
$\theta$ in the open chart.  In the bidegree decomposition
\eqref{eq:oriented-bidegrees}, distinct $X$-degrees separate the terms
$E_{A,B}$; each is homogeneous in $W$, so vanishing on an open patch of
the unit sphere extends to the corresponding open cone.  Hence
$\mathcal EH=0$, and Lemma~\ref{lem:oriented-injective} gives $H=0$ in
$\mathcal V_m$.  Consequently $\mathsf T_\chi$ is injective.  Its uniform
injectivity verifies the application-specific hypothesis
\eqref{eq:chart-pullback-injective} for the oriented tag space; no blanket
injectivity statement for ambient polynomials on the sphere is used.  Its
uniform constant is the genuine least singular value
\begin{equation}\label{eq:detecting-tag-singular-value}
 \sigma_\chi
 =\min_{\norm H_{\rm coeff}=1}
   \norm{\mathsf T_\chi H}_{\rm coeff}>0,
 \qquad
 \sigma_*:=\min_\chi\sigma_\chi>0.
\end{equation}
The minimum is taken on the unit sphere of the fixed finite-dimensional
space $\mathcal V_m$; the maps depend only on $n,m,d$ and the fixed atlas,
not on $a,R,S,\alpha$ or any coefficient of a lower homogeneous block.

Now retain exactly the coordinates in \eqref{eq:detecting-tag-index-set}
and multiply \eqref{eq:actual-tag-coefficient-new} by $\rho^\kappa$.
This is a uniformly invertible diagonal change on
$c_*\le\rho\le C_*$.  For $\alpha=1$, after the fixed denominator is
cleared, the selected polynomial vector is
\begin{equation}\label{eq:protected-tag-vector-new}
 \widehat C_{q,\gamma,j}(\rho,\vartheta)
 =\Delta_\chi(\vartheta)
  \sum_{\ell=2}^mR^\kappa S^{\ell-\kappa}\rho^\ell
  \tau_{q,j}(Q_{\ell,a,\chi,\gamma}).
\end{equation}
Let $\pi_m$ be ordinary coefficient projection onto radial degree
$\rho^m$.  Since the chart denominator and the tag selection are
independent of $\rho$, both commute with $\pi_m$, and therefore
\begin{equation}\label{eq:protected-radial-projection}
 \pi_m\widehat{\mathbf C}_{a,R,S,\chi}
 =D_{R,S}\,\mathsf T_\chi P_m,
 \qquad
 (D_{R,S}v)_{\gamma,j}
 =R^{\kappa(\gamma,j)}S^{m-\kappa(\gamma,j)}v_{\gamma,j}.
\end{equation}
This identity also shows directly why no translated lower-degree term can
enter a protected coordinate: the block of total degree $\ell$ carries
exactly the radial monomial $\rho^\ell$, while $H_{m,a}=P_m$.
For every retained coordinate,
$1\le\kappa(\gamma,j)\le m$; hence, since $R\le S$,
\begin{equation}\label{eq:protected-scale-diagonal-lower}
 \norm{D_{R,S}v}_{\rm coeff}
 \ge R^m\norm v_{\rm coeff}.
\end{equation}
Combining \eqref{eq:detecting-tag-singular-value}--
\eqref{eq:protected-scale-diagonal-lower} with the bounded coefficient
projection gives
\begin{equation}\label{eq:protected-vector-norm-new}
 \norm{\widehat{\mathbf C}_{a,R,S,\chi}}_{\rm coeff}
 \ge c_0R^m\norm{P_m}_{\rm coeff}=c_0L_R,
\end{equation}
uniformly in all lower blocks.  For orientation, the retained orders are
$\{1,2\}$ when $m=2$ and $\{1,2,3\}$ when $m=3$, giving respectively the
scale factors $RS,R^2$ and $RS^2,R^2S,R^3$.
Define
\[
 \widetilde{\mathbf p}_{a,R,S,\chi}(\rho,\vartheta)
 =\frac{\widehat{\mathbf C}_{a,R,S,\chi}}{c_0L_R}.
\]
This is an ordinary polynomial vector in $(\rho,\vartheta)$ because the
fixed factor $\Delta_\chi$ has already cleared all frame denominators.
Define
\[
 \mathbf p_{a,R,S,\chi}(\rho\theta_\chi(\vartheta))
 :=\widetilde{\mathbf p}_{a,R,S,\chi}(\rho,\vartheta).
\]
This proves \eqref{eq:oriented-p-norm}.

For fixed \(W\), group the selected coefficients in
\eqref{eq:actual-tag-coefficient-new} according to \((q,j)\), leaving
the homogeneous degree-\(q\) polynomial in \(\omega\) unexpanded; after
the harmless factors \(\rho^\kappa\Delta_\chi\) are inserted, call the
resulting vector \(\mathbf q_{\alpha,W}(\omega)\).  It consists of
selected coefficients of the literal numerator \(\mathcal N_\Psi\), up to a uniformly
invertible diagonal change, and therefore
\eqref{eq:oriented-band-lower} follows from
\eqref{eq:actual-mixed-to-quotient-new}.  The dependence on
\(\alpha\) is exactly diagonal:
\[
 \widehat C_{q,\gamma,j,\alpha}
 =\alpha^q\widehat C_{q,\gamma,j,1}.
\]
Since \(0<\alpha\le1\) and \(0\le q\le d\), the least singular value of
this diagonal map is at least \(\alpha^d\).  Consequently
\[
 \norm{\mathbf q_{\alpha,W}}_{{\rm coeff},\omega}
 \ge c\alpha^d\abs{\widehat{\mathbf C}_{a,R,S,\chi}(W)}
 \ge c\alpha^dL_R\abs{\mathbf p_{a,R,S,\chi}(W)}.
\]
This proves \eqref{eq:oriented-band-lower-new-b} and also proves that the
same \(\mathbf p\) works on every polar band.
\end{proof}

\begin{lemma}[Homogeneous angular restriction]\label{lem:homogeneous-angular-restriction}
Let $k\ge1$ and
\[
 \mathbf Q=\bigoplus_{q=0}^d\mathbf Q^{[q]},
\]
where the summands occupy disjoint value coordinates and every entry of
$\mathbf Q^{[q]}$ is a homogeneous degree-$q$ polynomial on $\R^k$.
If $k\ge2$, fix a finite rational atlas
$\omega=\omega_\nu(\zeta)$ on $\Sph^{k-1}$ and fixed denominators
$\Delta_\nu$ clearing all pullbacks through degree $d$.  Then, on every
nonempty chart,
\begin{equation}\label{eq:homogeneous-angular-restriction}
 \norm{\bigoplus_{q=0}^d
 \Delta_\nu(\zeta)\mathbf Q^{[q]}(\omega_\nu(\zeta))}
      _{\mathrm{coeff},\zeta}
 \ge c_{k,d,\nu}\norm{\mathbf Q}_{\mathrm{coeff}}.
\end{equation}
For $k=1$, evaluation on $\Sph^0=\{-1,1\}$ has the same property.
\end{lemma}

\begin{proof}
The displayed pullback is a linear map between fixed finite-dimensional
coefficient spaces.  If it vanishes on a nonempty chart, the disjoint
value coordinates show separately that every homogeneous block
$\mathbf Q^{[q]}$ vanishes on an open patch of the sphere.  Homogeneity
extends this vanishing to the corresponding open cone in $\R^k$, so every
block is the zero polynomial.  Thus the pullback is injective, and its
least singular value gives \eqref{eq:homogeneous-angular-restriction}.
When $k=1$, an entry is $c\omega^q$, whose value at either point of
$\Sph^0$ has modulus $|c|$.  Notice that keeping distinct degrees in
disjoint value coordinates is essential; no quotient by the sphere
relation $|\omega|^2=1$ is being taken.
\end{proof}

\begin{corollary}[Rough unequal-shell alternative]
\label{cor:unequal-correlation}
Let $P$ be a real polynomial of degree at most $d$ with highest nonpure
homogeneous block $P_m$, let
$\Theta\in L^\infty(\Sph^{n-1})$, and let \(0<R\le S/C_0\).  Put
$L_R=R^m\norm{P_m}_{\rm coeff}$, and let \(b_R,b_S\) be supported in a fixed compact
subinterval of \((0,\infty)\), with
\(L^\infty+\BV\) norm at most \(B\), where \(B\ge1\).  Put
\[
 k_R(z)=|z|^{-n}\Theta(z/|z|)b_R(|z|/R),\qquad
 k_S(z)=|z|^{-n}\Theta(z/|z|)b_S(|z|/S).
\]
Fix $a,W\in\R^n$, put \(y=a+SW\), and assume
\(c_*\le|W|\le C_*\).  Use the active-chart vector
\(\mathbf p(W)=\mathbf p_{a,R,S,\chi(W)}(W)\) from
\eqref{eq:active-oriented-p-new}.  Then, for \(n\ge2\),
\begin{equation}\label{eq:unequal-correlation}
 \left|\int e^{i[P(y+z,a)-P(y+z,y)]}
 \overline{k_R(z)}k_S(z+y-a)\,dz\right|
 \le CB^C\norm\Theta_\infty^2S^{-n}
 \min\{1,(L_R\abs{\mathbf p(W)})^{-c}\}.
\end{equation}
If \(\mathbf p(W)=0\), the last factor is understood to be \(1\).
For \(n=1\) the same conclusion holds, with \(S^{-1}\), for a
bounded-degree polynomial vector \(\mathbf p(W)\) satisfying
\(\norm{\mathbf p}_{\rm coeff}\ge1\).
\end{corollary}

\begin{proof}
If the two radial supports have no overlap after the translation by
\(SW\), the integral is zero in every dimension.  Assume that the overlap
is nonempty, and first let \(n\ge2\).  Then \(|z|\simeq R\),
\(|z+SW|\simeq S\), and the contributing
set has measure \(O(R^n)\); hence
\begin{equation}\label{eq:unequal-trivial-new}
 \int |k_R(z)k_S(z+SW)|\,dz
 \le CB^2\norm\Theta_\infty^2S^{-n}.
\end{equation}
This proves \eqref{eq:unequal-correlation} when \(\mathbf p(W)=0\).
Henceforth assume that \(\mathbf p(W)\ne0\).  Write
\(W=\rho\theta\), choose the active chart so that
\(\theta=O_\chi(\vartheta)e_1\), and let
\(\omega\in\Sph^{n-2}\subset\mathbb R^{n-1}\) denote a transverse
frame coordinate.  Put
\[
 \omega_\chi=O_\chi(\vartheta)(0,\omega)\in\theta^\perp
\]
and decompose \(z=t\theta+r\omega_\chi\), where \(r>0\).  Put
\[
 \xi=\frac{t}{r},\qquad
 \eta=\frac{t+S\rho}{r}.
\]
Thus the two rays are \(z=r(\xi\theta+\omega_\chi)\) and
\(z+SW=r(\eta\theta+\omega_\chi)\).  Choose a zero-extended bounded-variation
partition $\{\beta_k\}_{k\ge0}$ in the dimensionless transverse radius
such that, on the relevant range,
\[
 \sum_{k\ge0}\beta_k(r/R)=1,\qquad
 \supp\beta_k\subset[c2^{-k},C2^{-k}],
 \qquad
 \sup_k\bigl(\norm{\beta_k}_\infty+\Var(\beta_k)\bigr)\le C,
\]
and the supports have uniformly bounded overlap.  Absorb the boundedly
many largest transverse pieces into $k=0$, put $\alpha=2^{-k}$, and on
the $k$-th band write $r\simeq\alpha R$.  Set
\begin{equation}\label{eq:unequal-ray-normalization-new}
 \delta=\frac{R}{S\rho},\qquad
 u=\alpha\xi,\qquad V=\alpha\delta\eta,\qquad
 D=V-\delta u.
\end{equation}
Then, exactly,
\begin{equation}\label{eq:unequal-two-ray-exact-new}
 t=\frac{Ru}{D},\qquad r=\frac{\alpha R}{D},\qquad
 z=\frac{R(u\theta+\alpha\omega_\chi)}{D},\qquad
 z+SW=\frac{S\rho(V\theta+\alpha\delta\omega_\chi)}{D}.
\end{equation}
Here $D=\alpha R/r>0$.  Moreover
\(z/S=\rho O_\chi T_{\delta,\alpha,\omega}(u,V)\), so
the phase in these variables is exactly
\eqref{eq:actual-unequal-phase-new}.

Indeed, cylindrical measure is \(r^{n-2}\,dt\,dr\,d\omega\), and
\[
 \left|\frac{\partial(t,r)}{\partial(u,V)}\right|
 =\frac{\alpha R^2}{D^3}.
\]
Since
\[
 |z|=\frac{R\sqrt{u^2+\alpha^2}}{D},\qquad
 |z+SW|=\frac{S\rho\sqrt{V^2+\alpha^2\delta^2}}{D},
\]
the cylindrical Jacobian and the two homogeneous kernels give the exact
factor
\begin{align}
 &(S\rho)^{-n}\alpha^{n-1}D^{n-1}
 (\alpha^2+u^2)^{-n/2}
 (V^2+\alpha^2\delta^2)^{-n/2}\,du\,dV\,d\omega.
 \label{eq:unequal-jacobian-new}
\end{align}
The radial profiles become
\begin{equation}\label{eq:unequal-profile-pullback-new}
 \overline{b_R\left(\frac{\sqrt{u^2+\alpha^2}}{D}\right)},
 \qquad
 b_S\left(\frac{\rho\sqrt{V^2+\alpha^2\delta^2}}{D}\right),
\end{equation}
while the two rough factors are, respectively,
\begin{equation}\label{eq:unequal-angular-separation-new}
 \overline{\Theta\left(\frac{u\theta+\alpha\omega_\chi}
                    {\sqrt{u^2+\alpha^2}}\right)},
 \qquad
 \Theta\left(\frac{V\theta+\alpha\delta\omega_\chi}
                   {\sqrt{V^2+\alpha^2\delta^2}}\right).
\end{equation}
Thus neither rough angular factor is differentiated.

On the support of one transverse band and of the two radial profiles,
\(u,V\) lie in fixed compact sets and \(D\) is bounded above and away
from zero.  The profile support also gives
\(\sqrt{u^2+\alpha^2}\gtrsim1\) and
\(\sqrt{V^2+\alpha^2\delta^2}\gtrsim1\); hence the two possibly singular
weights in \eqref{eq:unequal-jacobian-new} and all of their sectional
derivatives are uniformly bounded.  After a bounded number of sign and
monotonicity pieces, the remaining amplitude in
\eqref{eq:unequal-jacobian-new}--
\eqref{eq:unequal-profile-pullback-new}, including the transverse band
cutoff, has uniform row and column \(L^1\) bounds and sectional
\(L^\infty+\BV\) norm at most \(CB^C\).  This follows from the
one-dimensional BV chain rule on each monotonicity piece.  Explicitly,
the two profile arguments are
\[
 \gamma_1(u,V)=\frac{\sqrt{u^2+\alpha^2}}{V-\delta u},\qquad
 \gamma_2(u,V)=
 \frac{\rho\sqrt{V^2+\alpha^2\delta^2}}{V-\delta u}.
\]
Writing $s_1=(u^2+\alpha^2)^{1/2}$,
$s_2=(V^2+\alpha^2\delta^2)^{1/2}$, and $D=V-\delta u$, direct
differentiation gives
\[
 \partial_u\gamma_1=\frac{uV+\delta\alpha^2}{s_1D^2},\qquad
 \partial_V\gamma_1=-\frac{s_1}{D^2},\qquad
 \partial_u\gamma_2=\frac{\rho\delta s_2}{D^2},\qquad
 \partial_V\gamma_2=-\frac{\rho\delta(uV+\delta\alpha^2)}{s_2D^2}.
\]
Thus the only nontrivial sectional critical-point equation is the linear
equation $uV+\delta\alpha^2=0$ in the varying variable.  Each section
therefore has a uniformly bounded number of monotonicity pieces; on each
piece
$\Var(b_R\circ\gamma_1)\le\Var(b_R)$ and
$\Var(b_S\circ\gamma_2)\le\Var(b_S)$.  More explicitly, on an interval
$I$ on which $\gamma$ is monotone,
\[
 \Var_I(b\circ\gamma)
 \le\Var_{\gamma(I)}(b)\le\Var(b),
\]
and the conclusion is trivial if $\gamma$ is constant.  The numerator
$uV+\delta\alpha^2$ is affine in either varying coordinate, so each
section requires at most one additional cut, independently of
$\alpha,\delta,\rho$.  Products are controlled by
\[
 \Var(fg)\le\norm f_\infty\Var(g)+\norm g_\infty\Var(f).
\]
After extracting the explicit factor $\alpha^{n-1}$ in
\eqref{eq:unequal-jacobian-new}, the displayed lower bounds for $D$ and
the two square roots give uniformly bounded sectional
$W^{1,\infty}$ norms for all remaining algebraic factors.  Hence the
sectional $L^\infty+\BV$ bounds and the Schur row and column bounds remain
uniform as $\alpha,\delta\downarrow0$.

Fix a compact rectangle containing all the $(u,V)$-supports above, and
partition it by a rectangular grid of sufficiently small fixed mesh.  For
each $\alpha,\delta,\rho$, retain only the grid rectangles meeting the
corresponding support.  Since on that support
\[
 D,\qquad \sqrt{u^2+\alpha^2},\qquad
 \sqrt{V^2+\alpha^2\delta^2}
\]
are bounded below uniformly, and these functions are uniformly Lipschitz
in $(u,V)$ for $0<\alpha\le1$ and $0<\delta\le\delta_0$, the mesh may
be chosen so that the same quantities remain bounded below on a fixed
enlargement of every retained rectangle.  A fixed tensor-product
partition subordinate to the grid has uniformly bounded sectional
$L^\infty+\BV$ norms.  Refining each section at the at most one zero of
$uV+\delta\alpha^2$ gives the monotonicity partitions required by
Lemma~\ref{lem:rational-operator}.  The number and interaction degree of
all rectangles are uniform in $\alpha,\delta,\rho$, so their $L^2$
bounds synthesize by Cauchy--Schwarz with only a fixed loss.

Let $\mathcal R_{\alpha,\delta,\rho}^{\rm rect}$ denote the uniformly finite family
of retained rectangles, including the tensor-product partition just
constructed.  For a rectangle
$\mathfrak b\in\mathcal R_{\alpha,\delta,\rho}^{\rm rect}$ and
fixed $\omega$, let $T^{\mathfrak b}_{\alpha,W,\omega}$ be the corresponding operator
from the $V$-interval to the $u$-interval, after the explicit factor
$(S\rho)^{-n}\alpha^{n-1}$ in
\eqref{eq:unequal-jacobian-new} has been removed.  Put, with zero extension
to the fixed ambient intervals,
\begin{align*}
 f_{\alpha,W,\omega}(V)
 &=\Theta\left(\frac{V\theta+\alpha\delta\omega_\chi}
                 {\sqrt{V^2+\alpha^2\delta^2}}\right),\\
 g_{\alpha,W,\omega}(u)
 &=\Theta\left(\frac{u\theta+\alpha\omega_\chi}
                 {\sqrt{u^2+\alpha^2}}\right).
\end{align*}
If the tensor cutoff of $\mathfrak b$ is
$\eta_{\mathfrak b}(u)\widetilde\eta_{\mathfrak b}(V)$, set
$f^{\mathfrak b}_{\alpha,W,\omega}
=\widetilde\eta_{\mathfrak b}f_{\alpha,W,\omega}$ and
$g^{\mathfrak b}_{\alpha,W,\omega}
=\eta_{\mathfrak b}g_{\alpha,W,\omega}$; their $L^2$ norms on the fixed
intervals are at most $C\norm\Theta_\infty$.  Lemma~\ref{lem:rational-operator},
\eqref{eq:oriented-band-lower}, and the finite interaction degree give
\begin{equation}\label{eq:unequal-fixed-angular-operator}
 \sup_{\mathfrak b\in\mathcal R_{\alpha,\delta,\rho}^{\rm rect}}
 \norm{T^{\mathfrak b}_{\alpha,W,\omega}}_{2\to2}
 \le CB^C\min\{1,|\mathbf q_{\alpha,W}(\omega)|^{-c_0}\}.
\end{equation}
Here and below the value at a zero of $\mathbf q_{\alpha,W}$ is defined to
be one.  If $\mathcal I_\alpha(W)$ is the contribution of the current
transverse band, the changes of variables above and the tensor partition
give the exact synthesis identity
\begin{equation}\label{eq:unequal-angular-synthesis-identity}
 \mathcal I_\alpha(W)
 =(S\rho)^{-n}\alpha^{n-1}
 \int_{\Sph^{n-2}}
  \sum_{\mathfrak b\in\mathcal R_{\alpha,\delta,\rho}^{\rm rect}}
  \left\langle T^{\mathfrak b}_{\alpha,W,\omega}
       f^{\mathfrak b}_{\alpha,W,\omega},
       g^{\mathfrak b}_{\alpha,W,\omega}\right\rangle
 \,d\omega.
\end{equation}
All functions here are jointly measurable.  The trivial Schur bound is a
uniform integrable majorant, so ordinary Fubini justifies
\eqref{eq:unequal-angular-synthesis-identity}.  Cauchy--Schwarz,
\eqref{eq:unequal-fixed-angular-operator}, the uniformly finite rectangle
sum, and $c_*\le\rho\le C_*$ now yield the explicit first half of the
angular synthesis:
\begin{equation}\label{eq:unequal-angular-synthesis-first}
 |\mathcal I_\alpha(W)|
 \le CB^C\norm\Theta_\infty^2S^{-n}\alpha^{n-1}
 \int_{\Sph^{n-2}}
  \min\{1,|\mathbf q_{\alpha,W}(\omega)|^{-c_0}\}\,d\omega.
\end{equation}

Write
\[
 \mathbf q_{\alpha,W}
 =\bigoplus_{q=0}^d\mathbf q_{\alpha,W}^{[q]},
\]
where different homogeneous degrees $q$ occupy disjoint value
coordinates, and every entry of $\mathbf q_{\alpha,W}^{[q]}$ is
homogeneous of degree $q$ in $\omega$.
Set
\[
 A_\alpha=\norm{\mathbf q_{\alpha,W}}_{{\rm coeff},\omega}.
\]
Lemma~\ref{lem:homogeneous-angular-restriction}, applied before any
degrees are combined, shows that on every nonempty rational chart the
cleared chart-polynomial vector has coefficient norm at least $cA_\alpha$.
Let $A_*\ge1$ be fixed so large that the chart lower bound
$cA_\alpha$ is at least one whenever $A_\alpha\ge A_*$.  If
$A_\alpha<A_*$, the integral in
\eqref{eq:unequal-angular-synthesis-first} is bounded by the surface
measure.  If $A_\alpha\ge A_*$, normalize the cleared vector on each chart by
its own coefficient norm and apply \eqref{eq:multi-negative} with
$\sigma=c_0$.  The chart Jacobians and clearing denominators are bounded
above and away from zero, so summing the finite atlas gives
\begin{align}
 \int_{\Sph^{n-2}}
  \min\{1,|\mathbf q_{\alpha,W}(\omega)|^{-c_0}\}\,d\omega
 &\le C\min\{1,A_\alpha^{-c_1}\}\notag\\
 &\le C\min\{1,
    (\alpha^dL_R|\mathbf p(W)|)^{-c_1}\},
 \label{eq:unequal-angular-negative-moment}
\end{align}
where the last inequality is
\eqref{eq:oriented-band-lower-new-b}.  When $n=2$, the same calculation is
the finite sum over $\Sph^0$: the $k=1$ clause of
Lemma~\ref{lem:homogeneous-angular-restriction} gives the required lower
bound at either point.  Combining
\eqref{eq:unequal-angular-synthesis-first} and
\eqref{eq:unequal-angular-negative-moment} gives the band bound
\begin{equation}\label{eq:unequal-band-bound-new}
 CB^C\norm\Theta_\infty^2S^{-n}\alpha^{n-1}
 \min\{1,(L_R\alpha^d\abs{\mathbf p(W)})^{-c}\}.
\end{equation}
Here \(\rho\) is confined to a fixed annulus and has been absorbed into
the constant.  If
\(A=L_R|\mathbf p(W)|\le1\), summability of \(\alpha^{n-1}\) gives the
trivial alternative.  If \(A>1\), split the dyadic sum at
\(\alpha\simeq A^{-1/d}\); the two resulting geometric series are
\(O(A^{-c'})\) after decreasing the exponent.  This proves
\eqref{eq:unequal-correlation}.

For \(n=1\), work on one of the finitely many signed overlap intervals
and write \(z=Rt\).  Let \(\mathbf C(W)\) be the vector of all
nonconstant-in-\(t\) coefficients of the complete translated correlation
phase.  Its degree-\(m\) part is
\begin{equation}\label{eq:one-dimensional-oriented}
 \sum_{\substack{A+B=m\\ A,B\ge1}}c_{A,B}S^AR^BW^At^B.
\end{equation}
For a fixed \(t^B\) coefficient, every degree \(\ell<m\) contribution
has \(W\)-degree at most \(\ell-B<m-B\), and cannot alter the coefficient
of \(W^{m-B}\).  Since \(S^AR^B\ge R^m\),
Lemma~\ref{lem:oriented-injective} gives
\[
 \norm{\mathbf C}_{\mathrm{coeff},W}
 \ge cR^m\norm{P_m}_{\rm coeff}=cL_R.
\]
Put \(\mathbf p=\mathbf C/(cL_R)\).  The product of the two kernels and
\(dz\) is \(S^{-1}b(t)dt\), on finitely many intervals, with
\(\norm b_{L^\infty+\BV}\le CB^C\norm\Theta_\infty^2\).
Lemma~\ref{lem:coefficient-vdc} gives the asserted estimate when
\(\mathbf p(W)\ne0\); when \(\mathbf p(W)=0\), the same conclusion is the
trivial \(S^{-1}\) amplitude bound.
\end{proof}

\begin{lemma}[One-sided variable-mesh incidence]
\label{lem:one-sided-incidence}
Let $0<\varepsilon\le1$, let $q_R=\varepsilon R$, let the scale-$R$
packet cubes form a finite disjoint family $\mathcal Q_R$ of dyadic cubes
in one aligned grid, where $R>0$.  For each $Q\in\mathcal Q_R$ choose a
point $y_Q\in Q$, and attach a weight
$0\le w_Q\le\Lambda q_R^n$, where $\Lambda\ge0$.  Fix $a\in\R^n$ and a
large shell scale $S\ge C_0R$.  If $\mathbf p$ is as in
Proposition~\ref{prop:oriented-coefficient} for $n\ge2$, or as constructed
in Corollary~\ref{cor:unequal-correlation} for $n=1$, then, for the fixed
annular-overlap constants $0<c_*<C_*<\infty$ and every $0<\tau\le1$,
\begin{equation}\label{eq:one-sided-incidence}
 \sum_{\substack{Q:\ c_*S\le|y_Q-a|\le C_*S\\
          |\mathbf p((y_Q-a)/S)|\le\tau}}w_Q
 \le C\Lambda S^n\{\tau^c+(\varepsilon R/S)^c\}.
\end{equation}
\end{lemma}

\begin{proof}
For \(n\ge2\), first split the sum into the finitely many disjoint active
sets \(E_\chi\) from \eqref{eq:active-oriented-p-new}; on the \(\chi\)-th
piece use \(\mathbf p_{a,R,S,\chi}\).  (For \(n=1\) no chart split is
needed.)  There is no second denominator clearing at this stage.  Indeed,
for $n\ge2$, equations \eqref{eq:protected-tag-vector-new}--
\eqref{eq:active-oriented-p-new} define
\[
 \widetilde{\mathbf p}_{a,R,S,\chi}(\rho,\vartheta)
 =\mathbf p_{a,R,S,\chi}(\rho\theta_\chi(\vartheta))
\]
as an ordinary coordinate-polynomial vector after the fixed denominator
has already been cleared.  For $n=1$ the vector constructed in the proof
of Corollary~\ref{cor:unequal-correlation} is already an ordinary
polynomial vector in $W$.  In either case its vector coefficient norm is
at least one.  Hence one fixed scalar component on the current chart,
called $p_0$, satisfies
$A_0=\norm{p_0}_{\rm coeff}\ge c>0$, where $c$ depends only on the fixed
vector length.  Put
\(\widetilde p_0=p_0/A_0\).  This last polynomial
has coefficient norm one.  At a bad point it satisfies
\(|\widetilde p_0|\le\tau/A_0\le C\tau\), regardless of how large the
translated lower blocks make \(A_0\).  After division by $S$, color the
$q_R$ grid and put disjoint balls of radius
$c\varepsilon R/S$ around the selected points.  Work in the fixed enlarged
radial annulus $c_*/2<|W|<2C_*$.  If this radius exceeds a fixed fraction
of either the radial collar or, when $n\ge2$, the angular chart collar
$\operatorname{dist}(U_\chi,\partial U_\chi^*)$,
the mesh term on the right of \eqref{eq:one-sided-incidence} is already
bounded below by a fixed constant.  Disjointness and the weight cap give
total weight at most $C\Lambda S^n$ in the fixed annulus, so the assertion
is then trivial.  Otherwise
the balls stay in the enlarged radial annulus and, for $n\ge2$, in the
enlarged chart $U_\chi^*$.  Uniform bi-Lipschitz
equivalence between physical \(W\)-balls and chart-coordinate balls
(ordinary \(W\)-intervals when $n=1$) changes their radii and volumes by
fixed factors only.  The normalized, already-cleared polynomial
\(\widetilde p_0\) has uniformly bounded gradient on the enlarged
coordinate domain by \eqref{eq:chart-polynomial-gradient},
so a ball centered at a bad point lies in
the $C(\tau+\varepsilon R/S)$-sublevel.  To record the density calculation, put
\[
 h=\frac{\varepsilon R}{S},\qquad W_Q=\frac{y_Q-a}{S}.
\]
For $n\ge2$ write $z_Q=(\rho_Q,\vartheta_Q)$ for the active chart
coordinate of $W_Q$; for $n=1$ put $z_Q=W_Q$.  For one color
$\mathfrak c$ the coordinate balls $B_Q=B(z_Q,ch)$ are disjoint, after a
fixed decrease of $c$, and
\[
 F_{\mathfrak c}(z)
 =\sum_{Q\in\mathfrak c}\frac{w_Q}{|B_Q|}\mathbf1_{B_Q}(z)
\]
satisfies
\begin{equation}\label{eq:one-sided-smeared-density}
 \norm{F_{\mathfrak c}}_\infty
 \le C\frac{\Lambda(\varepsilon R)^n}{h^n}
 =C\Lambda S^n.
\end{equation}
Every ball centered at a bad point is contained in the
$C_1(\tau+h)$-sublevel of $\widetilde p_0$.  If
$C_1(\tau+h)\ge1$, then
$\tau^c+h^c\ge(\tau+h)^c\ge C_1^{-c}$.  The disjoint packet cubes whose
centers lie in the fixed $S$-annulus have total volume $O(S^n)$, and hence
\[
 \sum_{Q\ {\rm bad}}w_Q\le\sum_Qw_Q\le C\Lambda S^n
 \le C\Lambda S^n(\tau^c+h^c).
\]
This is the required trivial branch for the current color and chart,
without applying a sublevel estimate outside its stated threshold range.
For the remaining color--chart pieces we may therefore assume
$C_1(\tau+h)<1$.  Lemma~\ref{lem:multi-sublevel}, now legitimately
applied in the coordinate box, gives
\[
 \sum_{\substack{Q\in\mathfrak c\\Q\ {\rm bad}}}w_Q
 \le\int_{\{|\widetilde p_0|\le C_1(\tau+h)\}}
       F_{\mathfrak c}(z)\,dz
 \le C\Lambda S^n(\tau+h)^{c_0}.
\]
Sum the boundedly many colors and active charts, and decrease the exponent
so that $(\tau+h)^{c_0}\le C(\tau^c+h^c)$.  This proves
\eqref{eq:one-sided-incidence}.  The large atom was fixed before
smearing, so $h=\varepsilon R/S$ and no mesh term of size
$\varepsilon S/S$ occurs.
\end{proof}

\begin{proposition}[Collective fixed-sign high-shell estimate]
\label{prop:collective-high}
Let $P$ be a real polynomial of degree at most $d$, let $P_m$ be its
highest nonpure homogeneous block, let
$\Theta\in L^\infty(\Sph^{n-1})$, fix
$\kappa>0$, let $s\in\mathbb N$, let $\varepsilon=2^{-s}$, and suppose
that $\mathcal R$ is a finite set of positive dyadic radii such that every
$R\in\mathcal R$ satisfies
\begin{equation}\label{eq:collective-high-threshold}
 L_R=R^m\norm{P_m}_{\rm coeff}\ge\varepsilon^{-\kappa}.
\end{equation}
At radius $R\in\mathcal R$, let $\mathcal Q_R$ be a finite family of
disjoint dyadic cubes of side $\varepsilon R$ in one aligned grid; assume
the cubes over all radii are globally disjoint and
each complex measure $\nu_Q$ is supported on $Q$ and satisfies
$|\nu_Q|(Q)\le\Lambda|Q|$, where $\Lambda\ge0$.  Let
\[
 T_R\nu(x)=\int e^{iP(x,y)}K_\Theta(x-y)
 b_R(|x-y|/R)\,d\nu(y),
\]
where $B\ge1$ and all normalized radial profiles are supported in one
fixed compact subinterval of $(0,\infty)$ with $L^\infty+\BV$ norm at
most $B$.
Then, uniformly for arbitrary complex coefficients $|e_Q|\le1$,
\begin{equation}\label{eq:collective-high}
 \left\|\sum_{R\in\mathcal R} T_R
       \sum_{Q\in\mathcal Q_R}e_Q\nu_Q\right\|_2^2
 \le CB^C\norm\Theta_\infty^2\Lambda M 2^{-c_\kappa s},
 \qquad M=\sum_{R\in\mathcal R}\sum_{Q\in\mathcal Q_R}|\nu_Q|(Q).
\end{equation}
Here $c_\kappa>0$, and the implicit constant may depend on $\kappa$
(in addition to the fixed structural data).
\end{proposition}

\begin{proof}
If $\mathcal R=\varnothing$, the assertion is immediate; hence assume
$\mathcal R\ne\varnothing$.
Let $c_{\rm flat},c_{\rm corr},c_{\rm inc}>0$ be fixed admissible
exponents in Proposition~\ref{prop:flat-high-block},
Corollary~\ref{cor:unequal-correlation}, and
Lemma~\ref{lem:one-sided-incidence}, respectively.  Decrease them once
and for all to the common exponent
\begin{equation}\label{eq:collective-common-exponent}
 c_0=\min\{c_{\rm flat},c_{\rm corr}/2,c_{\rm inc}/2\}>0,
 \qquad
 c_\kappa=\tfrac12c_0\min\{1,\kappa\}>0.
\end{equation}
All occurrences of a decay exponent $c$ in this proof are henceforth
understood to have been lowered to $c_0$; the halves above account for
the choice $\tau_R=L_R^{-1/2}$.  Likewise, enlarge the harmless power of
$B$ to the maximum of the finitely many cited powers.
Put $w_Q=|\nu_Q|(Q)$ and write the polar decomposition
$d\nu_Q=h_Q\,d|\nu_Q|$, with $|h_Q|=1$ $|\nu_Q|$-almost everywhere; cubes
with $w_Q=0$ may be omitted.  On the finite product probability space
\[
 (\mathfrak X,\mathbb P)
 =\prod_{R\in\mathcal R}\prod_{Q\in\mathcal Q_R}
   \left(Q,\frac{|\nu_Q|}{w_Q}\right),
\]
let $Y_Q$ be the coordinate random variables.  Then, in the weak sense of
complex measures,
\begin{equation}\label{eq:collective-polar-sampling}
 \mathbb E\bigl[w_Qh_Q(Y_Q)\delta_{Y_Q}\bigr]=\nu_Q;
\end{equation}
indeed, testing against a bounded Borel function gives the defining polar
integral.  Set
\[
 c_Q(\mathbf Y)=e_Qh_Q(Y_Q),\qquad
 \mathcal F(\mathbf Y)
 =\sum_{R\in\mathcal R}T_R
   \sum_{Q\in\mathcal Q_R}c_Q(\mathbf Y)w_Q\delta_{Y_Q}.
\]
The radial support is separated from zero, and hence, uniformly in $y$,
\begin{equation}\label{eq:collective-atom-L2}
 \norm{T_R\delta_y}_2^2
 \le CB^2\norm\Theta_\infty^2R^{-n}.
\end{equation}
Since all index sets are finite, this makes $\mathcal F$ a strongly
measurable Bochner-integrable $L^2$-valued random variable and justifies
all exchanges of expectation, finite sums, and integration.  By
\eqref{eq:collective-polar-sampling},
\[
 \mathbb E\mathcal F
 =\sum_R T_R\sum_Qe_Q\nu_Q
 \qquad\text{in }L^2.
\]
Hilbert-space Jensen therefore gives
\begin{equation}\label{eq:collective-jensen-reduction}
 \left\|\sum_R T_R\sum_Qe_Q\nu_Q\right\|_2^2
 =\norm{\mathbb E\mathcal F}_2^2
 \le\mathbb E\norm{\mathcal F}_2^2.
\end{equation}
It therefore suffices to prove the desired bound for every fixed selection
$y_Q\in Q$, with weights $w_Q\le\Lambda|Q|$ and arbitrary packet
coefficients $c_Q$ satisfying $|c_Q|\le1$.  Fix such a selection and put
\[
 F_R=T_R\sum_{Q\in\mathcal Q_R}c_Qw_Q\delta_{y_Q}.
\]
Partition physical input and output space into $R$-cubes; annular support
leaves only boundedly many relative input--output positions.  Write
$J\sim I$ when the pair can interact.  Then
\[
 \sup_I\#\{J:J\sim I\}+
 \sup_J\#\{I:J\sim I\}\le C_n.
\]
To make the scale normalization explicit, fix one interacting
output--input pair $(I,J)$, let $x_I$ be the chosen center of $I$, and let
$\mathcal Q_{R,I,J}$ be the packet cubes in $J$ which contribute to this
pair.  Put
\[
 M_{R,I,J}=\sum_{Q\in\mathcal Q_{R,I,J}}w_Q,
 \qquad \sigma_{x_I,R}(y)=\frac{y-x_I}{R},
\]
and, for each such $Q$, set
\[
 \widetilde\nu_Q
 =R^{-n}(\sigma_{x_I,R})_\#(w_Q\delta_{y_Q}),
 \qquad \widetilde Q=\sigma_{x_I,R}(Q).
\]
Then $\widetilde Q$ has side length $\varepsilon$ and, exactly,
\[
 |\widetilde\nu_Q|(\widetilde Q)
 =R^{-n}w_Q
 \le R^{-n}\Lambda|Q|=\Lambda|\widetilde Q|,
 \qquad
 \sum_{Q\in\mathcal Q_{R,I,J}}
 |\widetilde\nu_Q|(\widetilde Q)=R^{-n}M_{R,I,J}.
\]
All cubes $\widetilde Q$ lie in the fixed normalized input block
$B_0=\sigma_{x_I,R}(J)$.  Homogeneity of the kernel gives
\[
 T_R\sum_{Q\in\mathcal Q_{R,I,J}}c_Qw_Q\delta_{y_Q}(x_I+RX)
 =\widetilde T_1\sum_{Q\in\mathcal Q_{R,I,J}}c_Q\widetilde\nu_Q(X),
\]
where $\widetilde T_1$ has phase
$P(x_I+RX,x_I+RY)$ and normalized radial profile $b_R(|X-Y|)$.  Modulo
pure phases, the highest nonpure homogeneous block of the normalized phase
is $R^mP_m$, whose coefficient norm is $L_R$.
Proposition~\ref{prop:flat-high-block} therefore produces a normalized
bound proportional to $\Lambda R^{-n}M_{R,I,J}$.  Since
$dx=R^n\,dX$, the output $L^2$-square Jacobian cancels this $R^{-n}$
factor exactly.  For an interacting pair define
\[
 F_{R,I,J}=\mathbf1_I T_R
   \sum_{Q\in\mathcal Q_{R,I,J}}c_Qw_Q\delta_{y_Q}.
\]
Then $F_R=\sum_I\sum_{J\sim I}F_{R,I,J}$, and disjointness of the output
cubes together with the preceding adjacency bound gives
\[
 \norm{F_R}_2^2
 =\sum_I\left\|\sum_{J\sim I}F_{R,I,J}\right\|_2^2
 \le C_n\sum_I\sum_{J\sim I}\norm{F_{R,I,J}}_2^2.
\]
Moreover, every packet cube occurs for only $O_n(1)$ output blocks, so
\[
 \sum_I\sum_{J\sim I}M_{R,I,J}\le C_nM_R,
\]
where $M_R$ is the total packet mass at radius $R$.  Applying
Proposition~\ref{prop:flat-high-block} to every block pair and using these
two inequalities gives
\[
 \norm{F_R}_2^2\le CB^C\norm\Theta_\infty^2\Lambda M_R
 (L_R^{-c}+\varepsilon^c).
\]
The high threshold and \eqref{eq:collective-common-exponent} make the
parenthesis at most
$C(\varepsilon^{c_0\kappa}+\varepsilon^{c_0})
\le C2^{-c_\kappa s}$.
The finite squared-norm expansion is
\begin{equation}\label{eq:collective-pair-expansion}
 \left\|\sum_{R\in\mathcal R}F_R\right\|_2^2
 =\sum_R\norm{F_R}_2^2
  +2\operatorname{Re}\sum_{\substack{R,S\in\mathcal R\\R<S}}
    \langle F_R,F_S\rangle.
\end{equation}
For $R<S$, split the pairs into the comparable range $S<C_0R$ and the
separated range $R\le S/C_0$; these ranges are disjoint and exhaustive.
If $R$ and $S$ are comparable, then
$2|\langle F_R,F_S\rangle|\le\norm{F_R}_2^2+\norm{F_S}_2^2$ reduces the
cross term to those two diagonal bounds.  Each radius has only boundedly
many comparable dyadic neighbours, so their total is bounded by the right
side of \eqref{eq:collective-high}.

For $R\le S/C_0$, fix a large-shell atom $y_A$ of weight $w_A$ and put
\begin{equation}\label{eq:collective-tau}
 \tau_R=L_R^{-1/2}.
\end{equation}
For this fixed atomic selection, the scale-$R$ contribution to the
inner product with the fixed scale-$S$ atom has the form, up to packet
coefficients of modulus at most one,
as follows.  Indeed, write
$F_{S,A}=c_Aw_A T_S\delta_{y_A}$, set
$a=y_A$, $y=y_Q$, and $W=(y_Q-y_A)/S$, and change variables
$x=y_Q+z$.  Irrespective of the convention chosen for which argument of
the inner product is complex-linear, taking the absolute value gives the
exact identity
\begin{align*}
 &\left|\left\langle T_R\delta_{y_Q},
                         T_S\delta_{y_A}\right\rangle\right|\\
 &\quad=
 \left|\int e^{i[P(y_Q+z,y_A)-P(y_Q+z,y_Q)]}
       \overline{k_R(z)}k_S(z+y_Q-y_A)\,dz\right|.
\end{align*}
On the support of this integral the annular overlap gives
$c_*\le|W|\le C_*$, so this is exactly the left side of
Corollary~\ref{cor:unequal-correlation} with $y=a+SW$.  Expanding the
small-shell sum and using $|c_A|,|c_Q|\le1$ now gives
\[
 |\langle F_R,F_{S,A}\rangle|
 \le w_A\sum_{Q\in\mathcal Q_R}w_Q
 |\mathcal C_{R,S}(y_Q,y_A)|,
\]
where $w_Q$ is the fixed weight attached to $Q$ and
$\mathcal C_{R,S}$ is precisely the correlation estimated in
Corollary~\ref{cor:unequal-correlation}.
For this fixed triple $(S,A,R)$, define the bad small-shell packets by
\[
 \operatorname{Bad}_R(A)
 =\left\{Q:\ c_*S\le|y_Q-y_A|\le C_*S,\quad
   \left|\mathbf p\left(\frac{y_Q-y_A}{S}\right)\right|
   \le\tau_R\right\},
\]
with the finite active-chart split from
Lemma~\ref{lem:one-sided-incidence} understood when $n\ge2$.  On this bad
set, that lemma and the trivial correlation bound give, before
multiplication by $w_A$,
\[
 CB^C\norm\Theta_\infty^2\Lambda
 \{\tau_R^{c_{\rm inc}}+(\varepsilon R/S)^{c_{\rm inc}}\}
 =CB^C\norm\Theta_\infty^2\Lambda
 \{L_R^{-c_{\rm inc}/2}+(\varepsilon R/S)^{c_{\rm inc}}\}.
\]
On its complement, Corollary~\ref{cor:unequal-correlation} gives
\[
 CB^C\norm\Theta_\infty^2S^{-n}
 (L_R\tau_R)^{-c_{\rm corr}}
 =CB^C\norm\Theta_\infty^2S^{-n}L_R^{-c_{\rm corr}/2}.
\]
After lowering all exponents to the common $c_0$ from
\eqref{eq:collective-common-exponent}, these are the displayed $L_R^{-c}$
and mesh terms used below.  At one fixed scale the total
small-atom weight in the interacting $S$-ball is at most
$C\Lambda S^n$, by packing the disjoint $\varepsilon R$-cubes.  Thus the
good row has the same bound
$CB^C\norm\Theta_\infty^2\Lambda L_R^{-c}$.

Let $R_0$ be the smallest retained dyadic radius.  Since
$L_R=L_{R_0}(R/R_0)^m$ and
$L_{R_0}\ge\varepsilon^{-\kappa}$,
\[
 \sum_{\substack{R\in\mathcal R\\R\ge R_0}}L_R^{-c}
 \le\varepsilon^{c\kappa}
 \sum_{\ell\ge0}2^{-mc\ell}
 \lesssim\varepsilon^{c\kappa}.
\]
Likewise,
\[
 \sum_{R\le S/C_0}(\varepsilon R/S)^c
 \lesssim\varepsilon^c\sum_{\ell\ge0}2^{-c\ell}
 \lesssim\varepsilon^c.
\]
By \eqref{eq:collective-common-exponent} and $\varepsilon=2^{-s}$,
\[
 \varepsilon^{c_0\kappa}+\varepsilon^{c_0}
 \le2\,2^{-c_0\min\{1,\kappa\}s}
 \le C2^{-c_\kappa s}.
\]
Thus the complete small-shell row is at most
$CB^C\norm\Theta_\infty^2\Lambda2^{-c_\kappa s}$.  Multiply by $w_A$,
sum the large atoms and then $S$.  Equivalently, if $F_{S,A}$ denotes the
contribution of the fixed large-shell atom, the preceding row estimate says
\[
 \sum_{R\le S/C_0}|\langle F_R,F_{S,A}\rangle|
 \le CB^C\norm\Theta_\infty^2\Lambda
       2^{-c_\kappa s}w_A.
\]
Summing first in the atoms at scale $S$ and then in $S$ gives
$\sum_{S,A}w_A=M$, so no additional large-shell counting factor occurs.
The result is
$CB^C\norm\Theta_\infty^2\Lambda M2^{-c_\kappa s}$ for every fixed atomic
selection.  Together with \eqref{eq:collective-pair-expansion}, this proves
the deterministic atomic estimate; taking expectation in
\eqref{eq:collective-jensen-reduction} proves \eqref{eq:collective-high}
for the original complex measures.  Every separated unordered pair has
been oriented by its unique smaller and larger radius and counted once;
the conjugate orientation is exactly the factor two in
\eqref{eq:collective-pair-expansion}.  All estimates used absolute values,
so the coefficients may be fixed independently for every packet, as
required by the later rank argument.
\end{proof}

\section{The unphased weak endpoint}
\label{sec:unphased-weak}

We prove the unphased weak endpoint simultaneously in every dimension
$n\ge1$.  The terminal one-shell estimate already has a direct
one-dimensional branch in Proposition~\ref{prop:terminalpacket}, and the
long-variation fixed-sign estimate has the direct one-dimensional branch
of Lemma~\ref{lem:fullkernel}.  Thus the selector-free rank argument below
requires no imported weak endpoint for Hilbert truncations.  The good part
of the Calder\'on--Zygmund decomposition is controlled by the $p=2$ case
of Proposition~\ref{prop:nonosc-strong}:
\begin{equation}\label{eq:nonosc-L2}
 \norm{\mathbf J(T^0_{\Omega,\varepsilon}h:\varepsilon>0)}_2
 \le C_n\mathfrak L(\Omega)\norm h_2,
 \qquad \int\Omega=0.
\end{equation}

\begin{proposition}[Unphased weak critical-jump inequality on core data]
\label{prop:nonoscweak}
For mean-zero $\Omega\in L\log L(\Sph^{n-1})$ and every
bounded compactly supported $f$,
\begin{equation}\label{eq:nonoscweak}
 \norm{\mathbf J(T^0_{\Omega,\varepsilon}f:\varepsilon>0)}_{L^{1,\infty}}
 \le C_n\mathfrak L(\Omega)\norm f_1.
\end{equation}
\end{proposition}

\begin{proof}
\smallskip
\noindent\emph{Finite reduction and Calder\'on--Zygmund decomposition.}

Throughout this proposition $f$ is bounded and compactly supported.  We
impose the finite restrictions needed for every later ranking argument at
the outset.  After forming the full Calder\'on--Zygmund decomposition, fix
a finite subfamily $\mathcal Q_0$ of its bad cubes, a finite interval of
shell indices $j_-\le j\le j_+$, and a finite ordered set of positive
rational truncation parameters, enlarged by the finitely many dyadic and
cell endpoints used below.  The bad sums, localized families, and
trajectories are initially interpreted with these restrictions.  In
particular, every family $\mathcal R_{s,c}$ ranked below is finite
when Lemma~\ref{lem:RM} is invoked.  All bounds are uniform in the three
cutoffs, which are suppressed from the notation.

If \(\mathfrak L(\Omega)=0\), then \(\Omega=0\) and the assertion is
immediate.  By homogeneity normalize $\mathfrak L(\Omega)=1$.  Fix a
Calder\'on--Zygmund height $\alpha>0$ and perform the ordinary dyadic
Calder\'on--Zygmund decomposition at height $\alpha$ as follows.  Take the
maximal dyadic cubes $Q$ for which
$|Q|^{-1}\int_Q|f|>\alpha$.  They are disjoint,
$\sum_Q|Q|\le\alpha^{-1}\norm f_1$, and dyadic differentiation gives
$|f|\le\alpha$ off their union.  Maximality makes the average over the
dyadic parent at most $\alpha$, hence
$|Q|^{-1}\int_Q|f|\le2^n\alpha$.  Put
\[
 g=f\mathbf1_{(\cup Q)^c}+\sum_Q f_Q\mathbf1_Q,\qquad
 b_Q=(f-f_Q)\mathbf1_Q,\qquad f_Q=|Q|^{-1}\int_Qf.
\]
The preceding facts, followed by
$\norm g_2^2\le\norm g_\infty\norm g_1$, give
\begin{equation}\label{eq:CZterminal}
 \norm g_2^2\le C\alpha\norm f_1,
 \quad \int b_Q=0,
 \quad \norm{b_Q}_1\le C\alpha|Q|,
 \quad \Sigma:=\sum_Q\norm{b_Q}_1\le C\norm f_1.
\end{equation}
By \eqref{eq:nonosc-L2}, \eqref{eq:CZterminal}, and Chebyshev,
\begin{equation}\label{eq:nonosc-good-inside}
 \left|\left\{x:\mathbf J(T^0_{\Omega,\varepsilon}g:
 \varepsilon>0)>\alpha\right\}\right|
 \le \alpha^{-2}\norm{\mathbf J(T^0_{\Omega,\varepsilon}g)}_2^2
 \le C\alpha^{-1}\norm f_1.
\end{equation}
Choose $s_0\ge\max\{5,t_0^{\rm term}\}$, where $t_0^{\rm term}$ denotes
the constant $t_0$ in Proposition~\ref{prop:terminalpacket}, and enlarge
every bad cube by a fixed factor large enough that, outside the resulting
set $E_0$, an annulus of radius $2^j$ can meet a cube of side $2^m$ only
if $s=j-m\ge s_0$.  This factor depends only on $n$, and
$|E_0|\le C\alpha^{-1}\norm f_1$.
Put $b=\sum_Q b_Q$ and write
$B_m=\sum_{\ell(Q)=2^m}b_Q$.

\smallskip
\noindent\emph{Angular decomposition and short variation.}

For every $s\ge s_0$, make the mean-zero angular split
\eqref{eq:meanzerosplit} with a parameter $\delta_a>0$.  With the finite
cutoffs still in force, define the bounded-angular and tail trajectories by
\begin{align}
 U_{s,\varepsilon}
 &=\sum_{j=j_-}^{j_+}
 K_{\omega_s}
 \mathbf1_{\{\max\{\varepsilon,2^j\}<|z|\le2^{j+1}\}}*B_{j-s},
 \label{eq:nonosc-bounded-path}\\
 \mathrm{Tail}_{s,\varepsilon}
 &=\sum_{j=j_-}^{j_+}
 K_{\Omega_s^\sharp}
 \mathbf1_{\{\max\{\varepsilon,2^j\}<|z|\le2^{j+1}\}}*B_{j-s}.
 \label{eq:nonosc-angular-tail}
\end{align}
Outside $E_0$, the finite-cutoff bad trajectory is exactly
$\sum_{s\ge s_0}(U_{s,\varepsilon}+\mathrm{Tail}_{s,\varepsilon})$.
The change of variables $(j,m)\mapsto(j,s=j-m)$ is a bijective reindexing
of the shell--cube-scale pairs.  For fixed $s$, each bad cube can therefore
occur in at most one summand in $j$.  We spell out the variation estimate,
since neither positivity nor monotonicity of the complex kernel is being
used.  For fixed $j,Q$ put
\[
 A_{j,Q}(\varepsilon,x)
 =\int_{\max\{\varepsilon,2^j\}<|z|\le2^{j+1}}
   K_{\Omega_s^\sharp}(z)b_Q(x-z)\,dz.
\]
If $0<\varepsilon_0<\cdots<\varepsilon_N$, the radial regions occurring in
the increments $A_{j,Q}(\varepsilon_k)-A_{j,Q}(\varepsilon_{k-1})$ are
pairwise disjoint.  Hence
\[
 \sum_{k=1}^N
 |A_{j,Q}(\varepsilon_k,x)-A_{j,Q}(\varepsilon_{k-1},x)|
 \le
 \bigl(|K_{\Omega_s^\sharp}|\mathbf1_{\{2^j<|z|\le2^{j+1}\}}
       *|b_Q|\bigr)(x).
\]
Taking the supremum over the partition, summing in $j,Q$, and then using
Tonelli gives the factor
$\int_{2^j<|z|\le2^{j+1}}|K_{\Omega_s^\sharp}(z)|\,dz
=\log2\,\|\Omega_s^\sharp\|_1$.  Young's inequality consequently gives
\begin{align}
 \norm{V^1(\mathrm{Tail}_{s,\varepsilon}:\varepsilon>0)}_1
 &\le C\norm{\Omega_s^\sharp}_1
   \sum_j\sum_{\ell(Q)=2^{j-s}}\norm{b_Q}_1\notag\\
 &\le C\norm{\Omega_s^\sharp}_1\Sigma.
 \label{eq:nonosc-angular-tail-V1}
\end{align}
There is no additional shell-counting factor.  Summing in $s$ and using
Lemma~\ref{lem:tailcount} yields
\[
 \sum_{s\ge s_0}
 \norm{V^1(\mathrm{Tail}_{s,\varepsilon}:\varepsilon>0)}_1
 \le C\mathfrak L(\Omega)\Sigma.
\]
Under the present normalization this is at most $C\Sigma$.

For the bounded pieces $\omega_s$, the short square function is controlled
by Proposition~\ref{prop:terminalpacket}:
\begin{equation}\label{eq:nonoscshort}
 \left\|\left(\sum_j
 V^2(K_{\omega_s}\mathbf1_{\{\varepsilon<|z|\le2^{j+1}\}}*B_{j-s}:
 2^j\le\varepsilon\le2^{j+1})^2\right)^{1/2}\right\|_2^2
 \le C2^{2\delta_as}2^{-\delta_Ts}\alpha\Sigma.
\end{equation}
Choose $\delta_a<\delta_T/8$ and sum the square roots over $s$.

\smallskip
\noindent\emph{Fixed-rank control of the long variation.}

It remains to control the dyadic long variation.  Smooth each sharp annular
kernel at relative radial scale $h_s=2^{-\eta_0s}$, where
$N=4n+1$, $\eta_0=1/(192N)$, and choose
\begin{equation}\label{eq:nonosc-parameters}
 0<\delta_a<\min(\delta_T/8,1/192).
\end{equation}
To make the smoothing precise, fix a nonnegative
$\phi\in C_c^\infty((-1/10,1/10))$ with $\int\phi=1$, put
$\phi_h(t)=h^{-1}\phi(t/h)$, and set
\[
 a_s=\mathbf1_{(1,2]}*\phi_{h_s},\qquad
 H_j^s(z)=K_{\omega_s}(z)a_s(2^{-j}|z|).
\]
The difference between $a_s$ and $\mathbf1_{(1,2]}$ is supported in two
endpoint collars of total length $O(h_s)$ and is bounded by one.  Hence,
if $E_j^s$ is the difference between the sharp and smoothed scale-$j$
kernels, polar coordinates give
\[
 \norm{E_j^s}_1\le C\norm{\omega_s}_1h_s.
\]
The long paths are suffix sums of the shell contributions, so their
$V^1$ distance is bounded by the sum of the shell errors.  More precisely,
at the finite-cutoff stage put $e_j=E_j^s*B_{j-s}$ and
$D_k=\sum_{j\ge k}e_j$.  Then $D_{k+1}-D_k=-e_k$, and insertion of the
zero and total endpoints gives pointwise
\[
 V^1(D_k:k)=\sum_k|e_k|.
\]
Young's inequality and the disjoint assignment of the bad cubes to the
functions $B_{j-s}$ therefore give
\[
 \norm{V^1(\text{sharp long path}-\text{smoothed long path})}_1
 \le\sum_j\norm{E_j^s}_1\norm{B_{j-s}}_1
 \le C\norm{\omega_s}_1h_s\Sigma.
\]
This is summable in $s$, since
$\|\omega_s\|_1\le2\|\Omega\|_1$ and $h_s=2^{-\eta_0s}$.  The kernels
$H_j^s$ are smooth in the radial
variable and supported in a fixed annular enlargement.  Leibniz' rule,
$\norm{\partial_r^\ell a_s}_\infty\le C_\ell h_s^{-\ell}$, and
$r\simeq2^j$ on their support give
\begin{equation}\label{eq:nonosc-MN}
 M_N(H^s)\le C\norm{\omega_s}_\infty h_s^{-N}
 \le C2^{(\delta_a+N\eta_0)s}.
\end{equation}

We organize the long variation by a fixed global rank.  This
selector-free construction converts each point-dependent shell tail into a
suffix of one finite ranked sequence.  For a bad cube $Q$ with
$\ell(Q)=2^{j-s}$, let $A_j(Q)$ be its ancestor of side $2^j$ in
the original dyadic grid.  Since $H_j^s$ is supported in a fixed annular
enlargement of radius $C2^j$, the function $H_j^s*b_Q$ is supported in a
fixed enlargement of $A_j(Q)$.  We use the following nested adjacent
dyadic systems.  For
$a\in\{0,1/3,2/3\}^n$, put
\begin{equation}\label{eq:adjacent-grids}
 \mathcal D^a=\left\{2^{-k}\bigl([0,1)^n+m+(-1)^ka\bigr):
 k\in\mathbb Z,\ m\in\mathbb Z^n\right\}.
\end{equation}
These systems are nested.  In one coordinate, if
$s=(-1)^k$, the left endpoint of a child differs from a level-
$(k-1)$ endpoint, in units of $2^{-k}$, by
$m-2m'+3sa$.  Since $N=m+3sa$ is integral, take
$m'=\lfloor N/2\rfloor$; then $N-2m'\in\{0,1\}$.  Taking products proves
nesting.
They also have the adjacent covering property: choose a dyadic length
$L$ with $3\ell(I)<L\le6\ell(I)$ for a given interval $I$.  At scale
$L$, the three boundary classes are separated by $L/3$, so at least one
class has no boundary in $I$ and its length-$L$ interval contains $I$.
Choose this class in every coordinate.  Thus every cube is contained in a
member of some $\mathcal D^a$ whose side is at most six times larger.

Applying this property to a closed cube containing the fixed enlargement,
and using the strict inequality $3\ell(I)<L$ in the covering argument,
gives one of the $3^n$ systems and a cube $R(Q)$ in that system such that
the closed enlargement is contained in the interior of $R(Q)$ and
$\ell(R(Q))\le C_n2^j$.  Thus, after choosing the usual almost-everywhere
representatives of the convolutions,
\begin{equation}\label{eq:localized-essential-support}
 \operatorname{ess\,supp}(H_j^s*b_Q)\subset R(Q).
\end{equation}
All pointwise statements below are made off the countable union of the
boundaries of the adjacent dyadic systems, a null set independent of the
finite family.  For the fixed $s$, assign
each pair $(j,Q)$ with $\ell(Q)=2^{j-s}$ to exactly one of the finitely
many classes
\[
 c=(a,\iota_{\rm gen})\in\mathfrak C,
\]
determined by the adjacent system and the generation offset, and make one
fixed choice of $R(Q)$ in that class.  The set $\mathfrak C$ depends only on
the dimension.  Write $\mathcal Q_{s,c}$ for the cubes assigned to $c$ and
\[
 \Sigma_{s,c}=\sum_{Q\in\mathcal Q_{s,c}}\norm{b_Q}_1,
 \qquad \sum_{c\in\mathfrak C}\Sigma_{s,c}\le\Sigma.
\]

Fix a class $c$.  All its output cubes belong to one nested dyadic system
and satisfy $\ell(R(Q))=2^{j+\iota_c}$ for a fixed integer $\iota_c$.  For
such an output cube define
\begin{equation}\label{eq:localized-FR}
 F_{R,c}=H_j^s*\!\sum_{\substack{Q\in\mathcal Q_{s,c}:R(Q)=R\\
                  \ell(Q)=2^{j-s}}}b_Q,
 \qquad \ell(R)=2^{j+\iota_c}.
\end{equation}
By \eqref{eq:localized-essential-support},
$\operatorname{ess\,supp}F_{R,c}\subset R$.  Let $\mathcal R_{s,c}$ be
the finite family of nonzero $R$.  Each assigned $Q$ occurs once, and each
nonempty $R$ contains at least one assigned $Q$ of side $2^{j-s}$.
Choosing one such cube for each $R$ proves
\begin{equation}\label{eq:Rpacking}
 \sum_{R\in\mathcal R_{s,c}}|R|
 \le C2^{ns}\sum_{Q\in\mathcal Q_{s,c}}|Q|.
\end{equation}
Indeed $|R|\le C2^{ns}|Q|$, and the selected $Q$ are distinct.  Set
\[
 \mathcal N_{s,c}(x)=\sum_{R\in\mathcal R_{s,c}}\mathbf1_R(x),\qquad
 D_s=\lceil2^{(n+1)s}\rceil,\qquad
 X_{s,c}=\{\mathcal N_{s,c}>D_s\},
\]
and put $X_s=\bigcup_{c\in\mathfrak C}X_{s,c}$.  Chebyshev,
\eqref{eq:Rpacking}, the Calder\'on--Zygmund packing bound, and finiteness of
$\mathfrak C$ give
\begin{equation}\label{eq:Xs}
 |X_s|\le C2^{-s}\alpha^{-1}\norm f_1.
\end{equation}
Because $F_{R,c}$ is essentially supported in $R$, cubes
$R\subset X_{s,c}$ contribute nothing almost everywhere on
$X_{s,c}^c$ and may be discarded.  Put
\[
 \mathcal R_{s,c}^{\rm ret}
 =\{R\in\mathcal R_{s,c}:R\not\subset X_{s,c}\}.
\]
Every chain in $\mathcal R_{s,c}^{\rm ret}$ has length at most $D_s$: if
$R_1\supsetneq\cdots\supsetneq R_L$, choose
$x\in R_L\setminus X_{s,c}$; then
$L\le\mathcal N_{s,c}(x)\le D_s$.

Give a remaining cube the recursive classwise rank
\begin{equation}\label{eq:rank-definition}
 \operatorname{rk}_c(R)=1+\max\{\operatorname{rk}_c(R'):
 R'\in\mathcal R_{s,c}^{\rm ret},\ R'\subsetneq R\},
\end{equation}
with the maximum of the empty set equal to zero.  Thus
$1\le\operatorname{rk}_c(R)\le D_s$, equal-rank cubes form an antichain,
and strict inclusion strictly increases rank.  Put
\[
 g_{r,c}=\sum_{\substack{R\in\mathcal R_{s,c}^{\rm ret}\\
                      \operatorname{rk}_c(R)=r}}F_{R,c},
 \qquad1\le r\le D_s.
\]
Fix $x\notin X_{s,c}$ outside the null boundary set just specified.  At
each generation there is at most one cube of the current adjacent system
containing $x$, and the active cubes over all generations form a chain.
Write them in increasing shell order as
\[
 R_{j_1}(x)\subsetneq\cdots\subsetneq R_{j_L}(x),
 \qquad j_1<\cdots<j_L.
\]
Their ranks strictly increase in this order.  Every term arising from a
cube not containing $x$ vanishes by
\eqref{eq:localized-essential-support}.  Therefore, for every shell
threshold $k$, there is an integer $r=r(k,x,c)$ such that
\begin{equation}\label{eq:rank-shell-pointwise}
 \sum_{j\ge k}\sum_{\ell(R)=2^{j+\iota_c}}F_{R,c}(x)
 =\sum_{q\ge r}g_{q,c}(x).
\end{equation}
Conversely every distinct suffix of the active ranks is obtained by placing
$k$ between the corresponding consecutive generations.  A rank not
assumed by an active cube contributes zero at $x$ and merely inserts a
repeated suffix value.  Repetitions do not affect variation, and inclusion
of the zero and total endpoints shows that the classwise shell-tail
$2$-variation at $x$ equals the $2$-variation of the rank suffixes.
Subtracting each suffix from the fixed total sum identifies suffix
variation with prefix variation exactly.  Notice also that the rank is a
deterministic function of the finite cube family: once the
Rademacher--Menshov signs $\epsilon_r$ are fixed, the induced packet
coefficient $\epsilon_{\operatorname{rk}_c(R(Q))}$ is fixed before $x$ is
evaluated.

Apply Lemma~\ref{lem:RM} to $g_{1,c},\ldots,g_{D_s,c}$.  For a fixed choice
of rank signs, \eqref{eq:localized-FR} becomes
\[
 \sum_r\epsilon_rg_{r,c}
 =\sum_jH_j^s*\!\sum_{\ell(Q)=2^{j-s}}c_Qb_Q,
\]
where
\[
 c_Q=
 \begin{cases}
 \epsilon_{\operatorname{rk}_c(R(Q))},
 &Q\in\mathcal Q_{s,c},\ R(Q)\in\mathcal R_{s,c}^{\rm ret},\\
 0,&\text{otherwise}.
 \end{cases}
\]
Thus discarded cubes and cubes in other adjacent-grid classes have
coefficient zero; in particular every displayed rank is defined.  These
coefficients are fixed before $x$ is evaluated and satisfy the fixed-sign
hypothesis of Lemma~\ref{lem:fullkernel}; the $x$-dependent threshold acts
only through \eqref{eq:rank-shell-pointwise}.  Define
\[
 B_{j-s}^{s,c}
 =\sum_{\substack{Q\in\mathcal Q_{s,c}\\\ell(Q)=2^{j-s}}}b_Q.
\]
Lemma~\ref{lem:fullkernel}, \eqref{eq:nonosc-MN}, and
$\log(2+D_s)\lesssim s$ give
\[
 \norm{\mathbf1_{X_{s,c}^c}
 V^2(\sum_{j\ge k}H_j^s*B_{j-s}^{s,c}:k)}_2^2
 \le Cs^2 2^{-[1/24-2\delta_a-2N\eta_0]s}\alpha\Sigma_{s,c}.
\]
Finally, $B_{j-s}=\sum_{c\in\mathfrak C}B_{j-s}^{s,c}$ and
$X_s^c\subset X_{s,c}^c$ for every $c$.  The $V^2$ triangle inequality,
Cauchy--Schwarz over the finite class set, and
$\sum_c\Sigma_{s,c}\le\Sigma$ therefore yield
\begin{equation}\label{eq:nonosclong}
 \norm{\mathbf1_{X_s^c}V^2(\sum_{j\ge k}H_j^s*B_{j-s}:k)}_2^2
 \le Cs^2 2^{-[1/24-2\delta_a-2N\eta_0]s}\alpha\Sigma.
\end{equation}
The exponent in brackets is positive by \eqref{eq:nonosc-parameters}.
Indeed, $2N\eta_0=1/96$, and $\delta_a<1/192$ gives
\[
 \frac1{24}-2\delta_a-2N\eta_0
 >\frac1{24}-\frac1{96}-\frac1{96}=\frac1{48}.
\]
Equations \eqref{eq:Xs}--\eqref{eq:nonosclong} are summable in $s$.

Let $X=\bigcup_sX_s$.  We now record explicitly the scalar majorants
that make the passage to the pointwise supremum in the jump amplitude
legitimate.  Put
\begin{align*}
 \mathcal S_s
 &=\left(\sum_j
 V^2(U_{s,\varepsilon}:2^j\le\varepsilon\le2^{j+1})^2
 \right)^{1/2},\\
 \widetilde U_s(k)
 &=\sum_{j\ge k}H_j^s*B_{j-s},\\
 \mathcal L_s
 &=\mathbf1_{X_s^c}V^2(\widetilde U_s(k):k),\\
 \mathcal E_s
 &=V^1(U_{s,2^k}-\widetilde U_s(k):k),\\
 \mathcal T_s
 &=V^1(\mathrm{Tail}_{s,\varepsilon}:\varepsilon>0).
\end{align*}
The smoothing calculation preceding \eqref{eq:nonosc-MN} gives
$\norm{\mathcal E_s}_1\le
 C\norm{\omega_s}_1h_s\Sigma$.  Equations
\eqref{eq:nonoscshort} and \eqref{eq:nonosclong} give the corresponding
$L^2$ estimates for $\mathcal S_s$ and $\mathcal L_s$, while
\eqref{eq:nonosc-angular-tail-V1} controls $\mathcal T_s$ in $L^1$.
Consequently, with
\[
 \mathcal G=\sum_{s\ge s_0}(\mathcal S_s+\mathcal L_s),
 \qquad
 \mathcal W=\sum_{s\ge s_0}(\mathcal E_s+\mathcal T_s),
\]
Minkowski's inequality and the geometric decay in $s$ yield
\begin{equation}\label{eq:nonosc-global-majorants}
 \norm{\mathcal G}_2^2\le C\alpha\Sigma,
 \qquad
 \norm{\mathcal W}_1\le C\Sigma,
 \qquad
 |X|\le C\alpha^{-1}\norm f_1.
\end{equation}

Fix the finite truncation set $I$ imposed at the start of the proof.  The
pointwise long--short inequality \eqref{eq:pointwise-longshort}, followed
by $\mathbf J^{\rm dyad}\le V^2$ and the $V^2$ triangle inequality, gives
on $X_s^c$
\[
 \mathbf J_I(U_s)
 \le C(\mathcal S_s+\mathcal L_s+\mathcal E_s),
 \qquad
 \mathbf J_I(\mathrm{Tail}_s)\le\mathcal T_s.
\]
Apply Lemma~\ref{lem:jumpnorm} to counting measure on the finite set of
$s$-indices, and then use \eqref{eq:chain-pair-functional}.  Its constant
is independent of the number of retained shells.  Since the decomposition
in \eqref{eq:nonosc-bounded-path}--\eqref{eq:nonosc-angular-tail} is exact,
we obtain the pointwise, jump-amplitude-independent domination
\begin{equation}\label{eq:nonosc-bad-inside-majorant}
 \mathbf1_{(E_0\cup X)^c}
 \mathbf J_I(T^0_{\Omega,\varepsilon}b:\varepsilon\in I)
 \le C(\mathcal G+\mathcal W).
\end{equation}
This is the precise step that permits the supremum in $\lambda$ to be
placed before the spatial weak norm; no union over jump amplitudes is
taken.  Chebyshev and \eqref{eq:nonosc-global-majorants} now yield
\[
 \left|\left\{x\notin E_0\cup X:
 \mathbf J_I(T^0_{\Omega,\varepsilon}b:\varepsilon\in I)>C\alpha
 \right\}\right|
 \le C\alpha^{-2}\norm{\mathcal G}_2^2
     +C\alpha^{-1}\norm{\mathcal W}_1
 \le C\alpha^{-1}\Sigma.
\]
Add \eqref{eq:nonosc-good-inside}, the measure bounds for $E_0$ and $X$,
and use \eqref{eq:pointwise-jump-quasitriangle} for $f=g+b$.  Thus there
is an absolute $C_0$ such that, uniformly in all finite cutoffs,
\[
 \left|\left\{x:
  \mathbf J_I(T^0_{\Omega,\varepsilon}f:\varepsilon\in I)
       >C_0\alpha\right\}\right|
 \le C\alpha^{-1}\|f\|_1.
\]
This is an estimate for the pointwise supremum in the jump amplitude, not
for the weaker mixed quasi-seminorm.  After the cutoffs have been removed,
putting $\alpha=\tau/C_0$ and taking the supremum over $\tau>0$ gives the
claimed $L^{1,\infty}$ quasi-norm bound.

\smallskip
\noindent\emph{Removal of the auxiliary cutoffs.}\par\nobreak\smallskip
We remove the cutoffs in their declared order.  With shells and parameters fixed, increase
$\mathcal Q_0$ to the full bad family.  Writing
\[
 f_{\mathcal Q_0}=g+\sum_{Q\in\mathcal Q_0}b_Q,
\]
we have $f_{\mathcal Q_0}\to f$ in $L^1$, since
$\sum_Q\norm{b_Q}_1<\infty$.  Every coordinate has a finite-annulus
$L^1$ kernel, so Young's inequality gives the corresponding coordinatewise
$L^1$ convergence.  After a subsequence, convergence is pointwise
simultaneously on the finite parameter set; write $F^{(\nu)}$ and $F$ for
the corresponding approximating and limiting finite paths.  For finite
$I$ the functional
$\mathbf J_I$ is a Borel function of finitely many measurable coordinates.
The lower semicontinuity \eqref{eq:finite-J-lsc} therefore implies, for
every $\tau>0$,
\[
 \mathbf1_{\{\mathbf J_I(F)>\tau\}}
 \le\liminf_{\nu\to\infty}
   \mathbf1_{\{\mathbf J_I(F^{(\nu)})>\tau\}}
 \quad\text{almost everywhere}.
\]
Indeed, if \(\mathbf J_I(F)>\tau\), the definition of the supremum supplies
one jump amplitude and one finite strict witness whose value exceeds
\(\tau\) by a positive amount.  Coordinatewise convergence preserves all
of its strict increments and its excess for every sufficiently large
\(\nu\).  Fatou's lemma now preserves the distribution bound for the
finite-parameter pointwise functional, with a constant independent of
\(I\).

Next send $j_-\to-\infty$ and $j_+\to\infty$.  For core data, each fixed
spatial point $x$ and positive truncation $q$ sees only finitely many
shells: if $2^{j+1}\le q$, the shell lies below the truncation, while for
all sufficiently large $j$ the annulus centered at $x$ no longer meets the
compact support of the input.  Thus every coordinate is eventually
constant in the shell cutoff.  The truncated coordinates converge
pointwise on each fixed finite \(I\), and the same strict-witness
lower-semicontinuity and level-set Fatou argument applies again.  Finally
enumerate \(\mathbb Q_+=\{q_1,q_2,\ldots\}\) and put
\(I_N=\{q_1,\ldots,q_N\}\).  These sets are nested, and finiteness of every
strict witness gives, for each \(\tau>0\),
\[
 \{\mathbf J(T^0_{\Omega,q}f:q\in\mathbb Q_+)>\tau\}
 =\bigcup_{N\ge1}
   \{\mathbf J_{I_N}(T^0_{\Omega,q}f:q\in I_N)>\tau\}.
\]
Continuity from below transfers the uniform distribution bound.  For the
present bounded compactly supported data,
Fubini on every finite annulus gives local absolute continuity of the
truncation path (as recorded again in Section~\ref{sec:limits}); hence
Lemma~\ref{lem:finite-witness} identifies the rational and full-parameter
pointwise functionals: every finite full-parameter witness lies in one
compact parameter interval, where its points can be perturbed to positive
rationals while preserving all strict inequalities.  Every
Rademacher--Menshov ranking is therefore
finite, and the cutoff limits recover the full trajectory from these
finite fixed-rank sums.  The resulting distribution bound is
$C\norm f_1/\alpha$ for $\mathbf J$, and hence gives
\eqref{eq:nonoscweak} after restoring the homogeneous angular scale.
\end{proof}

\section{Phase-adapted Calder\'on--Zygmund decomposition}
\label{sec:phase-cz}

We modify the Calder\'on--Zygmund decomposition so that each bad block
cancels against the input modulation generated by its parent cube.  This
is exactly the cancellation condition required by the packet estimates in
the next section.

Fix a polynomial $P$ and take the maximal dyadic family $\cQ$ constructed
in the proof of Proposition~\ref{prop:nonoscweak}, at height $\alpha$ for
$f$.  For each $Q\in\cQ$, let $c_Q$ be its center and
put
\begin{equation}\label{eq:phaseCZ}
 a_Q=|Q|^{-1}\int_Qe^{iP(c_Q,y)}f(y)dy,
 \quad
 g=f\mathbf1_{(\cup Q)^c}+\sum_Qa_Qe^{-iP(c_Q,\cdot)}\mathbf1_Q,
 \quad b_Q=(f-a_Qe^{-iP(c_Q,\cdot)})\mathbf1_Q.
\end{equation}
Then
\begin{equation}\label{eq:phaseCZbounds}
 \norm g_\infty\le C\alpha,
 \quad \norm g_2^2\le C\alpha\norm f_1,
 \quad \sum_Q\norm{b_Q}_1\le C\norm f_1,
 \quad \norm{b_Q}_1\le C\alpha |Q|.
\end{equation}
Indeed, maximality gives
$|Q|^{-1}\int_Q|f|\le 2^n\alpha$, and hence
$|a_Q|\le2^n\alpha$ and
$\norm{b_Q}_1\le2\int_Q|f|$.  The cubes are disjoint.  Off their union
$|f|\le\alpha$ almost everywhere.  Moreover,
\[
 |a_Q||Q|
 =\left|\int_Qe^{iP(c_Q,y)}f(y)\,dy\right|
 \le\int_Q|f(y)|\,dy
\]
and, on $Q$,
$|g|^2=|a_Q|^2\le2^n\alpha|a_Q|$.  Consequently
\[
 \norm g_2^2
 \le \alpha\int_{(\cup Q)^c}|f|
      +2^n\alpha\sum_Q |a_Q||Q|
 \le C\alpha\norm f_1.
\]
These observations give all four estimates in \eqref{eq:phaseCZbounds}.
Most importantly, phase adaptation gives the exact cancellation
\begin{equation}\label{eq:adaptcancel}
 h_Q(y):=e^{iP(c_Q,y)}b_Q(y),
 \qquad \int_Qh_Q(y)\,dy=0,
 \qquad |h_Q(y)|=|b_Q(y)|.
\end{equation}
Indeed,
\[
 \int_Qh_Q(y)\,dy
 =\int_Qe^{iP(c_Q,y)}f(y)\,dy-a_Q|Q|=0,
 \qquad \|h_Q\|_1=\|b_Q\|_1.
\]
No pointwise bound for $b_Q$ or $h_Q$ is asserted.  The packet estimates
below are deliberately formulated with only the density
$\norm{h_Q}_1/|Q|$ and are applied directly to the original blocks; thus
no charge-neutralized stopping-tree decomposition is used.  In the
one-shell estimate below we take $d_Q=h_Q$ and use its cancellation.  In
the collective high-shell estimate we instead take
$d\nu_Q=b_Q(y)\,dy$; that proposition requires the same mass-density
bound but no cancellation.

\section{The far-field weak endpoint}
\label{sec:weak-far}

This section proves the weak critical-jump estimate for the normalized
far-field trajectory.  Assume that the highest nonpure homogeneous block
$P_m$ has been
normalized by
\begin{equation}\label{eq:top-block-normalized}
 \norm{P_m}_{\rm coeff}=1.
\end{equation}
Define the far family by
\[
 \cF^P_\varepsilon f(x)=
 \int_{\substack{|x-y|>\varepsilon\\|x-y|>1}}
 e^{iP(x,y)}K_\Omega(x-y)f(y)\,dy,
\]
extended constantly for $0<\varepsilon\le1$.  Equivalently,
\begin{equation}\label{eq:far-as-full-restriction}
 \cF^P_\varepsilon f
 =T^P_{\Omega,\max\{\varepsilon,1\}}f.
\end{equation}
Thus the far trajectory is a restriction of the full trajectory, with
repeated values for $0<\varepsilon\le1$, and hence
\begin{equation}\label{eq:far-jump-domination}
 N_\lambda(\cF^P_\varepsilon f:\varepsilon>0)
 \le N_\lambda(T^P_{\Omega,\varepsilon}f:\varepsilon>0),
 \qquad
 \mathbf J(\cF^P_\varepsilon f:\varepsilon>0)
 \le\mathbf J(T^P_{\Omega,\varepsilon}f:\varepsilon>0).
\end{equation}
The strong $p=2$ theorem gives
\begin{equation}\label{eq:L2jump}
 \norm{\mathbf J(T^P_{\Omega,\varepsilon}h:\varepsilon>0)}_2
 \le C_{n,d}\mathfrak L(\Omega)\norm h_2.
\end{equation}

\begin{proposition}[Far weak endpoint on core data]
\label{prop:farweak}
Let $P$ be a real polynomial of degree at most $d$ with highest nonpure
homogeneous block $P_m$, and assume
$\norm{P_m}_{\rm coeff}=1$.  Let
$\Omega\in L\log L(\Sph^{n-1})$ have mean zero.  Then, for every bounded
compactly supported $f$,
\begin{equation}\label{eq:farweak}
 \norm{\mathbf J(\cF^P_\varepsilon f:\varepsilon>0)}_{L^{1,\infty}}
 \le C_{n,d}\mathfrak L(\Omega)\norm f_1.
\end{equation}
\end{proposition}

\begin{proof}
We first take bounded compactly supported $f$ and impose finite cube,
shell, and rational-parameter cutoffs, exactly as in
Proposition~\ref{prop:nonoscweak}.  Every estimate below is uniform in
those cutoffs.

If \(\mathfrak L(\Omega)=0\), the assertion is immediate.  By angular
homogeneity normalize $\mathfrak L(\Omega)=1$, fix a
Calder\'on--Zygmund height
$\alpha>0$, and use \eqref{eq:phaseCZ}.  By \eqref{eq:far-jump-domination}, \eqref{eq:L2jump},
\eqref{eq:phaseCZbounds}, and Chebyshev,
\begin{equation}\label{eq:far-good-inside}
 \left|\left\{x:\mathbf J(\cF^P_\varepsilon g:\varepsilon>0)
 >\alpha\right\}\right|
 \le\alpha^{-2}\norm{\mathbf J(\cF^P_\varepsilon g)}_2^2
 \le C\alpha^{-2}\norm g_2^2
 \le C\alpha^{-1}\norm f_1.
\end{equation}
Thus the good part is controlled outside a set of the required measure.
Enlarge the bad cubes
by a fixed factor chosen as follows.  Let
$s_0\ge\max\{5,t_0\}$, where $t_0$ is the threshold in
Proposition~\ref{prop:direct-packet}, and use dilation factor
$C_n2^{s_0}$.  This produces another exceptional set $E_0$ of measure
$C\alpha^{-1}\norm f_1$.  Outside $E_0$, a cube meeting a shell
$R=2^j\ge1$ has
\begin{equation}\label{eq:direct-shell-cube}
 \ell(Q)=R2^{-s},\qquad s\ge s_0.
\end{equation}
Put $b=\sum_Qb_Q$ and
$\Sigma=\sum_Q\norm{b_Q}_1$.  For every fixed $s$, the cubes in
\eqref{eq:direct-shell-cube}, as $j$ varies, form a subfamily of the
original disjoint Calder\'on--Zygmund cubes and
\begin{equation}\label{eq:direct-shell-mass}
 \sum_j\sum_{\ell(Q)=2^{j-s}}\norm{h_Q}_1\le\Sigma,
 \qquad h_Q=e^{iP(c_Q,\cdot)}b_Q.
\end{equation}

Split $\Omega=\omega_s+\Omega_s^\sharp$ by
\eqref{eq:meanzerosplit}.  With the finite shell cutoff understood, define
the bounded-angular and angular-tail trajectories
\begin{align}
 U^P_{s,\varepsilon}(x)
 &=\sum_{j\ge0}
 \int_{\max\{\varepsilon,2^j\}<|x-y|\le2^{j+1}}
 e^{iP(x,y)}K_{\omega_s}(x-y)
 \sum_{\ell(Q)=2^{j-s}}b_Q(y)\,dy,
 \label{eq:far-bounded-trajectory}\\
 \mathrm{Tail}^{P}_{s,\varepsilon}(x)
 &=\sum_{j\ge0}
 \int_{\max\{\varepsilon,2^j\}<|x-y|\le2^{j+1}}
 e^{iP(x,y)}K_{\Omega_s^\sharp}(x-y)
 \sum_{\ell(Q)=2^{j-s}}b_Q(y)\,dy.
 \label{eq:far-angular-tail}
\end{align}
Since $|e^{iP(x,y)}|=1$, the same Fubini--Young calculation as in
\eqref{eq:nonosc-angular-tail-V1}, together with
\eqref{eq:direct-shell-mass}, gives
\begin{equation}\label{eq:far-angular-tail-V1}
 \norm{V^1(\mathrm{Tail}^{P}_{s,\varepsilon}:\varepsilon>0)}_1
 \le C\norm{\Omega_s^\sharp}_1\Sigma.
\end{equation}
Here again the reindexing $(j,m)\mapsto(j,s=j-m)$ assigns each bad cube to
at most one $j$ for fixed $s$, so no shell-counting factor occurs.  Summing
\eqref{eq:far-angular-tail-V1} in $s$ and using
Lemma~\ref{lem:tailcount} proves that the angular tails have total
$V^1L^1$ norm at most $C\mathfrak L(\Omega)\Sigma$.  For the bounded pieces,
\begin{equation}\label{eq:far-angular-bounded}
 \norm{\omega_s}_\infty\le C2^{\delta_as}\norm\Omega_1,
 \qquad \norm{\omega_s}_1\le2\norm\Omega_1.
\end{equation}

\smallskip
\noindent\emph{Short variation.}

We first control the short variation inside every dyadic annulus.  Apply
Proposition~\ref{prop:direct-packet} directly to the $h_Q$ at each
$(j,s)$.  The required phase matching is the exact identity
\begin{equation}\label{eq:far-direct-phase-matching}
 e^{iP(x,y)}b_Q(y)
 =e^{i[P(x,y)-P(c_Q,y)]}h_Q(y).
\end{equation}
Thus the packet datum in that proposition is $d_Q=h_Q$; by
\eqref{eq:adaptcancel} and \eqref{eq:phaseCZbounds} it is supported in
$Q$, has integral zero, satisfies
$\|d_Q\|_1\le C\alpha|Q|$, and the packet cubes are globally disjoint.
We may therefore take $\Lambda=C\alpha$ and $t=s$.  Square-summing in
$j$ and using \eqref{eq:direct-shell-mass} gives
\begin{equation}\label{eq:far-short-direct}
 \left\|\left(\sum_jV^2(G_{j,s})^2\right)^{1/2}\right\|_2^2
 \le C\norm{\omega_s}_\infty^2 2^{-c_Ps}\alpha\Sigma,
\end{equation}
and the packet errors satisfy
\begin{equation}\label{eq:far-short-error}
 \sum_j\norm{V^1(E_{j,s})}_1
 \le C\norm{\omega_s}_1 2^{-c_Es}\Sigma.
\end{equation}

\smallskip
\noindent\emph{Long variation: low shells.}

It remains to control the dyadic long variation.  Fix a small
$\kappa>0$ and split the shells by the single global top block:
\begin{equation}\label{eq:global-low-high}
 \text{low: }1\le R^m\le2^{\kappa s},
 \qquad
 \text{high: }R^m>2^{\kappa s}.
\end{equation}
Since $R=2^j\ge1$, the low condition implies
$0\le j\le\kappa s/m$; hence there are at most
$1+\kappa s/m\le C_d(1+s)$ low shells.  If $F_{j,s}$ is their full
annular output, then
\[
 V^2\left(\sum_{j\ge k,\ j\ {\rm low}}F_{j,s}:k\right)
 \le2\sum_{j\ {\rm low}}|F_{j,s}|.
\]
The endpoint difference of a one-shell trajectory is dominated by its
full variation.  More explicitly write
$F_{j,s}=F^G_{j,s}+F^E_{j,s}$ using the endpoint differences of the
decomposition in Proposition~\ref{prop:direct-packet}.  Thus, with
$R=2^j$,
\[
 F_{j,s}=\mathcal P_R-\mathcal P_{2R}
 =(G_R-G_{2R})+(E_R-E_{2R}),
\]
and consequently
\[
 |F^G_{j,s}|\le V^2(G_{j,s}),\qquad
 |F^E_{j,s}|\le V^1(E_{j,s}).
\]
Cauchy--Schwarz in the $O(1+s)$ low shells and
\eqref{eq:direct-packet-G}--\eqref{eq:direct-packet-E} therefore give
\begin{equation}\label{eq:far-low-long}
 \left\|2\sum_{j\ \mathrm{low}}|F^G_{j,s}|\right\|_2^2
 \le C(1+s)\sum_{j\ \mathrm{low}}
       \norm{V^2(G_{j,s})}_2^2
 \le C\norm{\omega_s}_\infty^2 2^{-cs}\alpha\Sigma,
 \qquad c=c_P/2,
\end{equation}
where $(1+s)2^{-c_Ps}\lesssim2^{-c_Ps/2}$ has been absorbed, together
with a summable $V^1L^1$ error.  This is the only place where translated
lower blocks are Fourier-expanded across low shells.

\smallskip
\noindent\emph{Long variation: high shells.}

On the high shells no phase block is expanded.  Return to the original
form $e^{iP(x,y)}b_Q(y)$; every packet and shell now has the same global
phase $P$, and \eqref{eq:global-low-high} is precisely the threshold
\eqref{eq:collective-high-threshold}.  Their short variation has already
been included in the square sum \eqref{eq:far-short-direct}; the
collective estimate is needed only for the dyadic long variation.

For the high dyadic long variation use the output localization and fixed
rank from the proof of Proposition~\ref{prop:nonoscweak}.  For every $Q$
at a high shell $2^j$, assign its natural annular output to a cube $R(Q)$
in one of the finitely many nested adjacent systems
\eqref{eq:adjacent-grids}, with $\ell(R(Q))\simeq2^j$.  As in the
closed-enlargement construction preceding
\eqref{eq:localized-essential-support}, choose $R(Q)$ so that the
essential support of this sharp-annular output is contained in its
interior.  Let
$\mathfrak C$ be the resulting finite set of adjacent-system and generation-
offset classes; it depends only on the dimension.  Make one fixed class
assignment for each pair $(j,Q)$ and write $\mathcal Q_{s,c}^{\rm high}$
for the packets assigned to $c$.  Within a fixed class, all $R(Q)$ lie in
one nested system and $\ell(R(Q))=2^{j+\iota_c}$ for one fixed integer
$\iota_c$.  Define the current bounded-angular annular operator by
\[
 T_j^{P,\omega_s}g(x)=
 \int_{2^j<|x-y|\le2^{j+1}}e^{iP(x,y)}
 K_{\omega_s}(x-y)g(y)\,dy,
\]
and, for the classwise output cubes, put
\[
 F_{R,c}=T_j^{P,\omega_s}
 \sum_{\substack{Q\in\mathcal Q_{s,c}^{\rm high}:\ R(Q)=R\\
                  \ell(Q)=2^{j-s}}}b_Q,
 \qquad \ell(R)=2^{j+\iota_c}.
\]
Let $\mathcal R_{s,c}$ be the finite family of nonzero classwise output
cubes.  Then $\operatorname{ess\,supp}F_{R,c}\subset R$ and
\[
 \sum_{R\in\mathcal R_{s,c}}|R|
 \le C2^{ns}\sum_{Q\in\mathcal Q_{s,c}^{\rm high}}|Q|.
\]
Let $D_s=\lceil2^{(n+1)s}\rceil$ and define
\[
 X_{s,c}=\left\{\sum_{R\in\mathcal R_{s,c}}\mathbf1_R>D_s\right\},
 \qquad X_s=\bigcup_{c\in\mathfrak C}X_{s,c}.
\]
Since $\mathfrak C$ is finite, the packing estimate gives
\begin{equation}\label{eq:far-Xs}
 |X_s|\le C2^{-s}\alpha^{-1}\norm f_1.
\end{equation}
For each class discard the cubes $R\subset X_{s,c}$ and rank the retained
finite family by strict inclusion, as in \eqref{eq:rank-definition}.  Put
\[
 \mathcal Q_{s,c}^{\rm high,ret}
 =\{Q\in\mathcal Q_{s,c}^{\rm high}:R(Q)\not\subset X_{s,c}\}.
\]
At every $x\notin X_s$ outside the fixed null set of adjacent-grid
boundaries, the high-shell tail in each class is a suffix of its rank
sequence by the active-chain argument proving
\eqref{eq:rank-shell-pointwise}; gaps caused by low shells only repeat a
suffix value.
A fixed rank sign induces a fixed coefficient in $\{-1,0,1\}$ on each
original packet, with coefficient zero on discarded cubes and on packets
outside the current class.  Thus no rank is ever invoked for a discarded
or differently classified packet.

We verify explicitly the hypotheses of
Proposition~\ref{prop:collective-high} class by class.  For fixed $s$ and
$c$, the retained set of high radii is finite because the shell cutoff is
finite.  Across all of those radii, the packet cubes remain a subfamily of
the original globally disjoint Calder\'on--Zygmund cubes.  Take
$d\nu_Q=b_Q(y)\,dy$; unlike the short estimate, this application uses no
zero-mean hypothesis.  Then
\[
 \supp\nu_Q\subset Q,\qquad
 |\nu_Q|(Q)=\norm{b_Q}_1=\norm{h_Q}_1\le C\alpha|Q|.
\]
The normalized profile is the sharp function $\mathbf1_{(1,2]}$, whose
$L^\infty+\BV$ norm is bounded by an absolute constant.  Moreover,
normalization of the top block and the high-shell condition give
$R^m\|P_m\|_{\rm coeff}=R^m>2^{\kappa s}$, exactly
\eqref{eq:collective-high-threshold}.  The Rademacher--Menshov signs are
fixed packet by packet before the output point $x$ is evaluated.  If
\[
 \Sigma_{s,c}^{\rm high}
 =\sum_{Q\in\mathcal Q_{s,c}^{\rm high,ret}}\norm{b_Q}_1,
\]
then Lemma~\ref{lem:RM} and Proposition~\ref{prop:collective-high} give
\[
 \norm{\mathbf1_{X_{s,c}^c}
 V^2(\text{high dyadic tails in class }c)}_2^2
 \le Cs^2\norm{\omega_s}_\infty^2
 2^{-c_\kappa s}\alpha\Sigma_{s,c}^{\rm high}.
\]
The $V^2$ triangle inequality and Cauchy--Schwarz over the finite class set,
together with $\sum_c\Sigma_{s,c}^{\rm high}\le\Sigma$, yield the full
high-shell estimate
\begin{equation}\label{eq:far-high-long}
 \norm{\mathbf1_{X_s^c}V^2(\text{high dyadic tails})}_2^2
 \le Cs^2\norm{\omega_s}_\infty^2
 2^{-c_\kappa s}\alpha\Sigma.
\end{equation}
This is where the estimate must allow packetwise, rather than merely
shellwise, signs.

Choose
$0<\delta_a<\frac14\min(c_P,c_E,c_\kappa)$.  Equations
\eqref{eq:far-short-direct}--\eqref{eq:far-high-long}, polynomial factors
in $s$, and \eqref{eq:far-angular-bounded} imply the following precise
majorant bounds.  Put
\begin{align*}
 \mathcal S_s&=\left(\sum_jV^2(G_{j,s})^2\right)^{1/2},\\
 \mathcal L_s&=2\sum_{j\ \mathrm{low}}|F^G_{j,s}|,\\
 \mathcal H_s&=\mathbf1_{X_s^c}
 V^2(\text{high dyadic tails}).
\end{align*}
For later use retain explicitly the error produced by the one-shell
decompositions and put
\begin{equation}\label{eq:far-shell-error-majorant}
 \mathcal E_s^*=\sum_jV^1(E_{j,s}).
\end{equation}
At the present finite-cutoff stage this is a finite pointwise sum.  In
particular, no compatibility of the choices of \(G_{j,s}\) and \(E_{j,s}\)
from one shell to another is being assumed.  Let \(\mathcal W\) be the
total \(1\)-variation majorant of all angular
tail paths plus
$C\sum_s\mathcal E_s^*$; the constant also covers the low-shell
endpoint errors $2\sum_{j\,\mathrm{low}}|F^E_{j,s}|$.  Set
\[
 \mathcal G=\sum_s(\mathcal S_s+\mathcal L_s+\mathcal H_s),
 \qquad X=\bigcup_sX_s.
\]
Minkowski's inequality, applied to the square roots of the fixed-$s$
estimates, now gives
\begin{equation}\label{eq:far-global-GW}
 \norm{\mathcal G}_2^2\le C\alpha\Sigma,
 \qquad \norm{\mathcal W}_1\le C\Sigma,
 \qquad |X|
 \le C\alpha^{-1}\norm f_1.
\end{equation}
The geometric decay makes all three sums finite.  More precisely, at the
finite-cutoff stage these are finite sums, while summability of their
$L^2$, respectively $L^1$, norms makes the partial sums of $\mathcal G$,
respectively $\mathcal W$, converge in those spaces with the displayed
bounds; monotone subsequences may be used for the nonnegative majorants.
We now spell out the passage from the shellwise decompositions to the full
trajectory.  Denote the moving bounded-angular one-shell path on
\([2^j,2^{j+1}]\) by \(\mathcal P_{j,s}=G_{j,s}+E_{j,s}\), and denote its
endpoint difference by \(F_{j,s}=F^G_{j,s}+F^E_{j,s}\), as above.  The
outer-shell terms in \eqref{eq:far-bounded-trajectory} are constant while
\(\varepsilon\in[2^j,2^{j+1}]\).  Hence the variation triangle inequality,
\(V^2(E_{j,s})\le V^1(E_{j,s})\), and \(\ell^2\)-Minkowski give the
pointwise estimate
\begin{align}
 &\left(\sum_j
 V^2(U^P_{s,\varepsilon}:2^j\le\varepsilon\le2^{j+1})^2\right)^{1/2}
 =\left(\sum_jV^2(\mathcal P_{j,s})^2\right)^{1/2}\notag\\
 &\le \left(\sum_j
       \bigl(V^2(G_{j,s})+V^1(E_{j,s})\bigr)^2\right)^{1/2}\notag\\
 &\le \mathcal S_s+
       \left(\sum_jV^1(E_{j,s})^2\right)^{1/2}
 \le \mathcal S_s+\mathcal E_s^*.
 \label{eq:far-full-short-gluing}
\end{align}
For the dyadic endpoint path, split the full suffix sum into its high- and
low-shell parts.  The exact endpoint identity is
\[
 U^P_{s,2^k}=\sum_{j\ge k}F_{j,s}.
\]
Here and below \(F_{j,s}=0\) for \(j<0\), so the identity also records the
constant part of the far trajectory for \(2^k\le1\).
On \(X_s^c\), the high part is bounded by
\(\mathcal H_s\).  The elementary suffix bound used in
\eqref{eq:far-low-long}, together with
\(|F^E_{j,s}|\le V^1(E_{j,s})\), gives
\begin{align}
 \mathbf1_{X_s^c}
 V^2\left(\sum_{j\ge k}F_{j,s}:k\right)
 &\le \mathcal H_s+
       V^2\left(\sum_{\substack{j\ge k\\j\ {\rm low}}}F_{j,s}:k\right)
       \notag\\
 &\le \mathcal H_s+2\sum_{j\ {\rm low}}|F_{j,s}|\notag\\
 &\le \mathcal H_s+\mathcal L_s+2\mathcal E_s^*.
 \label{eq:far-full-long-gluing}
\end{align}
Thus the deterministic full-variation long--short inequality
\eqref{eq:V2-long-short}, applied to the \emph{original} bounded-angular
path rather than to separately chosen global \(G/E\) paths, yields on
\(X_s^c\)
\begin{equation}\label{eq:far-full-shell-gluing}
 V^2(U^P_{s,\varepsilon}:\varepsilon>0)
 \le C(\mathcal S_s+\mathcal L_s+\mathcal H_s+\mathcal E_s^*).
\end{equation}
Summing \eqref{eq:far-full-shell-gluing} over \(s\) with
\eqref{eq:varadd}, and adding the angular-tail \(V^1\) paths, gives on
\(E_0^c\cap X^c\) a \(2\)-variation majorant bounded by
\(C(\mathcal G+\mathcal W)\).  This construction uses only the displayed
shellwise decompositions and therefore makes no unstated cross-shell
selection.  Since the majorants do not depend on the jump amplitude,
equations
\eqref{eq:pointwise-jump-quasitriangle} and
\eqref{eq:pointwise-jump-by-V2} give directly
\begin{equation}\label{eq:far-pointwise-majorant}
 \mathbf1_{E_0^c\cap X^c}
 \mathbf J(\cF^P_\varepsilon b:\varepsilon>0)
 \le C(\mathcal G+\mathcal W).
\end{equation}

Chebyshev applied to \eqref{eq:far-pointwise-majorant}, together with
\eqref{eq:far-good-inside}, the measures of $E_0$ and $X$,
\eqref{eq:far-global-GW}, and the pointwise jump quasi-triangle inequality,
gives
\[
 \left|\left\{x:\mathbf J(\cF^P_\varepsilon f:\varepsilon>0)
 >C_0\alpha\right\}\right|
 \le C\alpha^{-1}\norm f_1.
\]
Here $C_0$ is fixed and independent of every cutoff.  Once the cutoffs
are removed, setting $\alpha=\tau/C_0$ and taking the supremum over
$\tau>0$ gives precisely the $L^{1,\infty}$ quasi-norm of the pointwise
critical-jump functional.
The proof in dimension one is identical after using the one-dimensional
parts of Corollary~\ref{cor:unequal-correlation} and
Proposition~\ref{prop:collective-high}; mean zero on $\Sph^0$ makes the
angular function a Hilbert-kernel multiple, so no angular-chart issue
arises.

\smallskip
\noindent\emph{Removal of cutoffs.}

Finally remove the cutoffs in the same order as in the proof of
Proposition~\ref{prop:nonoscweak}.  With the shell interval and the finite
rational parameter set \(I\) fixed, exhaust the phase-adapted bad-cube
family.
Since $\sum_Q\norm{b_Q}_1<\infty$, the corresponding inputs converge in
$L^1$; Young's inequality for each finite-annulus coordinate, followed by
a subsequence, gives pointwise convergence simultaneously on the finite
parameter set; denote the approximating and limiting finite paths by
$F^{(\nu)}$ and $F$.  The finite functional is Borel in its coordinates.
If \(\mathbf J_I(F)>\tau\), one amplitude and one finite strict witness
exceed \(\tau\) by a positive amount, and coordinatewise convergence
preserves that witness for all sufficiently large \(\nu\).  Hence the
level-set form of \eqref{eq:finite-J-lsc},
\[
 \mathbf1_{\{\mathbf J_I(F)>\tau\}}
 \le\liminf_{\nu\to\infty}
   \mathbf1_{\{\mathbf J_I(F^{(\nu)})>\tau\}},
\]
followed by Fatou preserves the distribution estimate for every
$\tau>0$.

Next exhaust the dyadic shell interval, still with \(I\) fixed.  For a
fixed output point $x$ and truncation $q>0$, every shell with
$2^{j+1}\le q$ lies below the
truncation, while for all sufficiently large $j$ the annulus centered at
$x$ misses the compact support of the input.  Thus each coordinate is
eventually a finite shell sum, and the shell-cutoff limit is pointwise on
the finite set \(I\).  Apply the same strict-witness lower-semicontinuity
and level-set Fatou argument once more.  Finally enumerate
\(\mathbb Q_+=\{q_1,q_2,\ldots\}\), put
\(I_N=\{q_1,\ldots,q_N\}\), and use
\[
 \{\mathbf J(\cF^P_qf:q\in\mathbb Q_+)>\tau\}
 =\bigcup_{N\ge1}\{\mathbf J_{I_N}(\cF^P_qf:q\in I_N)>\tau\}.
\]
The equality follows because every strict witness is finite; continuity
from below therefore preserves the uniform distribution bound.
For bounded compactly supported data, Fubini on finite annuli gives local
continuity in the truncation parameter.  Every full-parameter witness is
contained in a compact interval, and Lemma~\ref{lem:finite-witness}
therefore identifies the rational and full-parameter functionals.  The
ranks, signs, and
exceptional sets are auxiliary devices used only to prove estimates
uniform in the cutoffs; none of them is passed to the limit.  Only the
original finite-cutoff trajectories are passed to the limit.  In
particular, no convergence of the ranks, signs, or exceptional sets is
asserted or needed.  This proves \eqref{eq:farweak} for core data.
Angular homogeneity restores $\mathfrak L(\Omega)$.
\end{proof}

\section{Local weak estimates and homogeneous-degree induction}
\label{sec:weak-local}

We close the weak endpoint with a coarser triangular order than in the
strong argument: each step removes the entire highest nonpure homogeneous
class rather than a single flag coordinate.

\begin{proposition}[Local homogeneous-degree reduction]
\label{prop:degree-local-reduction}
Let $P$ be a real polynomial of degree at most $d$ whose highest nonpure
homogeneous block $P_m$ satisfies $\norm{P_m}_{\rm coeff}=1$.  Let
$\Omega\in L^1(\Sph^{n-1})$ and let $f$ be bounded and compactly
supported.  On each unit output cube, the local path
\eqref{eq:strong-local-path} is, modulo pure input and output phases, the
sum of
\begin{enumerate}[label=\textup{(\roman*)}]
\item a localized path of the form \eqref{eq:localized-full-path} whose
phase has nonpure total degree at most $m-1$;
\item an error $E$ satisfying
\begin{equation}\label{eq:degree-local-error}
 V^1(E)(x)\le C_{n,d}
 \int_{|x-y|<1}\frac{|\Omega((x-y)/|x-y|)|}{|x-y|^{n-1}}|f(y)|\,dy.
\end{equation}
\end{enumerate}
The majorant has $L^1$ norm at most
$C_{n,d}\norm\Omega_1\norm f_1$.
\end{proposition}

\begin{proof}
Let $Q_a$ be the output unit cube centered at $a$, and let $Q_a^*$ be the
fixed input enlargement relevant to $Q_a$.  Write $x=a+X,y=a+Y$.
Simultaneous translation preserves the degree-$m$ class and creates only
lower degrees.  After separating pure terms, write
\[
 P(a+X,a+Y)=p_a(X)+q_a(Y)+H_m(X,Y)+R_{a,<m}(X,Y),
\]
where $H_m$ is the mixed representative of the normalized top class and
$R_{a,<m}$ has nonpure degree at most $m-1$.  Replace $H_m$ by the anchored
representative
\[
 \widehat H_m(X,Y)=H_m(X,Y)-H_m(Y,Y),
\]
and absorb the pure input polynomial $H_m(Y,Y)$ into $q_a$.  Set
\[
 f_a(Y)=e^{iq_a(Y)}f(a+Y)\mathbf1_{Q_a^*-a}(Y).
\]
The coordinate translation identifies the path at $a+X$ with the path at
$X$ and preserves Lebesgue measure.  Multiplication of the output path by
$e^{ip_a(X)}$ leaves pointwise jump counts unchanged, whereas the pure
input factor is now part of $f_a$.  In particular,
\[
 |f_a(Y)|=|f(a+Y)|\mathbf1_{Q_a^*-a}(Y),
 \qquad
 \norm{f_a}_1=\norm{f\mathbf1_{Q_a^*}}_1.
\]
The anchored representative satisfies $\widehat H_m(Y,Y)=0$.  On the
fixed bounded set relevant to
$|X-Y|<1$,
\[
 |\widehat H_m(X,Y)|
 \le |X-Y|\int_0^1
 |\nabla_XH_m(Y+t(X-Y),Y)|\,dt
 \le C_{n,d}|X-Y|.
\]
The last constant is uniform because the coefficient space
$\mathcal V_m$ is finite dimensional, its normalized unit sphere is
compact, and $X,Y$ range over one fixed bounded set.  Suppressing the
harmless pure output factor, define, for $0<\varepsilon<1$,
\[
 E_{a,\varepsilon}(X)
 =\int_{\varepsilon<|X-Y|<1}
   \bigl(e^{i\widehat H_m(X,Y)}-1\bigr)
   e^{iR_{a,<m}(X,Y)}
   K_\Omega(X-Y)f_a(Y)\,dY,
\]
and put $E_{a,\varepsilon}=0$ for $\varepsilon\ge1$.  For an arbitrary
finite increasing truncation partition, the radial layers occurring in
the successive increments are disjoint.  Using
$|e^{iu}-1|\le|u|$ and the preceding bound therefore gives
\[
 V^1(E_a)(X)
 \le C_{n,d}\int_{|X-Y|<1}
   \frac{|\Omega((X-Y)/|X-Y|)|}{|X-Y|^{n-1}}
   |f_a(Y)|\,dY.
\]
After translating back and discarding the input restriction this is
\eqref{eq:degree-local-error}.  Since the output cubes are disjoint, the
piecewise error satisfies the same pointwise majorant, and polar
coordinates give
\[
 \int_{|z|<1}
  \frac{|\Omega(z/|z|)|}{|z|^{n-1}}\,dz
 =\|\Omega\|_1.
\]
Tonelli thus proves the asserted global $L^1$ estimate.  What remains
after subtracting $E_a$ is exactly the localized path with phase
$R_{a,<m}$ applied to $f_a$, whose nonpure degree is at most $m-1$.
\end{proof}

\begin{proposition}[Weak polynomial critical jumps on core data]
\label{prop:weak-polynomial-jump}
For every real polynomial $P$ of degree at most $d$, every mean-zero
$\Omega\in L\log L(\Sph^{n-1})$, and every bounded compactly supported
$f$,
\[
 \norm{\mathbf J(T^P_{\Omega,\varepsilon}f:\varepsilon>0)}_{L^{1,\infty}}
 \le C_{n,d}\mathfrak L(\Omega)\norm f_1,
\]
with a constant independent of the coefficients of $P$.
\end{proposition}

\begin{proof}
All weak estimates invoked in this induction are the core-data versions
proved above.  More precisely, the base case uses the bounded compactly
supported conclusion of Proposition~\ref{prop:nonoscweak}, and the far
part uses the bounded compactly supported conclusion of
Proposition~\ref{prop:farweak}.  The modulated translated inputs $f_a$
arising in Proposition~\ref{prop:degree-local-reduction} remain bounded
and compactly supported, and satisfy
$\norm{f_a}_1=\norm{f\mathbf1_{Q_a^*}}_1$.  Thus no
operator value for a general $L^1$ input, and no result from
Section~\ref{sec:limits}, is used in the induction.

Induct on the largest total degree $m$ carrying a nonpure homogeneous
class.  If there is no such class, the phase is pure and
Proposition~\ref{prop:nonoscweak} applies after input and output
conjugation.  Assume the assertion for nonpure degree at most $m-1$ and
let
\[
 C_m=\norm{P_m}_{\rm coeff}>0,
 \qquad \rho=C_m^{-1/m}.
\]
Under $x=\rho X,y=\rho Y$, the entire degree-$m$ class, rather than one
scalar coordinate, has norm one.  Proposition~\ref{prop:farweak} controls
the normalized far path.

For the local path partition the output into disjoint unit cubes $Q_a$,
and let $Q_a^*$ denote the fixed input enlargement relevant to $Q_a$.
Proposition~\ref{prop:degree-local-reduction} gives on $Q_a$ a residual
localized path $L_a$ with nonpure degree at most $m-1$, applied to the
corresponding modulated translated input $f_a$, and an error $E_a$.  The
piecewise error $E=E_a$ on $Q_a$ satisfies, by
\eqref{eq:degree-local-error},
\[
 \norm{V^1(E)}_1\le C_{n,d}\norm\Omega_1\norm f_1
 \le C_{n,d}\mathfrak L(\Omega)\norm f_1.
\]
For the residual phase, subtracting the fixed truncation at one does not
change increments below $1$, and adjoining zero for $\varepsilon\ge1$
creates no additional jump by \eqref{eq:localized-jump-domination}.
Taking the supremum in the jump amplitude in that pointwise domination,
the induction hypothesis gives, for every $\alpha>0$,
\[
 \left|\left\{x\in Q_a:\mathbf J(L_a:\varepsilon>0)>\alpha\right\}\right|
 \le C_{n,d}\alpha^{-1}\mathfrak L(\Omega)
       \norm{f_a}_1
  = C_{n,d}\alpha^{-1}\mathfrak L(\Omega)
       \norm{f\mathbf1_{Q_a^*}}_1.
\]
Let $L_{\rm res}=L_a$ on $Q_a$.  Since the output cubes are disjoint and
the input enlargements have bounded overlap,
\begin{align*}
 \left|\left\{x:\mathbf J(L_{\rm res}:\varepsilon>0)>\alpha\right\}\right|
 &=\sum_a
 \left|\left\{x\in Q_a:\mathbf J(L_a:\varepsilon>0)>\alpha\right\}\right|\\
 &\le C_{n,d}\alpha^{-1}\mathfrak L(\Omega)
       \sum_a\norm{f\mathbf1_{Q_a^*}}_1\\
 &\le C_{n,d}\alpha^{-1}\mathfrak L(\Omega)\norm f_1.
\end{align*}
For the error, \eqref{eq:pointwise-jump-by-V2} and $V^2\le V^1$ give
$\mathbf J(E)\le V^1(E)$; hence its $L^1$ majorant yields the required
weak estimate.  Thus the local residual and error paths have the desired
pointwise critical-jump bound.

Combine local and far paths with
\eqref{eq:pointwise-jump-quasitriangle}.  More explicitly, the two weak
distribution estimates and the $L^1$ error majorant give a fixed $C_0$
such that, with $P_\rho(X,Y)=P(\rho X,\rho Y)$ and
$f_\rho(X)=f(\rho X)$,
\[
 \left|\left\{X:
  \mathbf J(T^{P_\rho}_{\Omega,\varepsilon}f_\rho:
            \varepsilon>0)(X)>C_0\alpha\right\}\right|
 \le C\alpha^{-1}\mathfrak L(\Omega)\|f_\rho\|_1.
\]
Putting $\alpha=\tau/C_0$ gives the normalized weak quasi-norm estimate.
The degree-$m$ block of $P_\rho$ has
coefficient norm $\rho^mC_m=1$, and the exact conjugacy is
\[
 T^P_{\Omega,\varepsilon}f(\rho X)
 =T^{P_\rho}_{\Omega,\varepsilon/\rho}f_\rho(X),
 \qquad f_\rho(X)=f(\rho X),
\]
where the reparameterization
$\varepsilon\mapsto\varepsilon/\rho$ preserves every jump count.  The
two scaling identities are, with $H$ denoting any measurable scalar
functional on $\mathbb R^n$ (in the application, the pointwise jump
functional),
\[
 \|f_\rho\|_{L^1_X}=\rho^{-n}\|f\|_{L^1_x},
 \qquad
 \|H(\rho\,\cdot)\|_{L^{1,\infty}_X}
 =\rho^{-n}\|H\|_{L^{1,\infty}_x}.
\]
Thus the factors $\rho^{-n}$ on the normalized input and output sides
cancel exactly, and the bound is independent of $C_m$.  The finite degree
induction proves the proposition for every polynomial phase.
\end{proof}

\section{Completion, consequences, and the direct radial
\texorpdfstring{$\BV$}{BV} transfer}\label{sec:limits}

\subsection{Completion of the critical-jump trajectory}

For core data, the weak endpoint propositions already remove the
auxiliary cube, shell, and rational-parameter cutoffs.  We pass to
general $L^1$ data and construct the complete trajectory over all positive
truncation parameters.

For bounded compactly supported $f$, put, for $N\ge2$,
\[
 J_N(x)=\int_{N^{-1}<|z|<N}|K_\Omega(z)|\,|f(x-z)|\,dz.
\]
Fubini gives
\[
 \int_{\R^n}J_N(x)\,dx
 =\norm f_1
  \norm{K_\Omega\mathbf1_{\{N^{-1}<|\cdot|<N\}}}_1<\infty.
\]
After intersecting the countably many sets $\{J_N<\infty\}$, we obtain
one full-measure set on which every finite-annulus integral is finite.
Every compact interval $[a,b]\Subset(0,\infty)$ is contained in
$[N^{-1},N]$ for some $N$; polar coordinates then show that, on this
single set, the truncation path is locally absolutely continuous, hence
continuous, simultaneously on all compact parameter intervals.  Therefore
the paired jump counts, and hence $\mathbf J$, computed on positive
rationals agree there with those computed on all positive parameters by
Lemma~\ref{lem:finite-witness}.

For such core data the path is zero for sufficiently large truncation at
each fixed $x$.  Equation~\eqref{eq:anchor-controls-maximal} therefore
gives pointwise
\[
 \sup_{\varepsilon>0}|T^P_{\Omega,\varepsilon}h|
 \le\mathbf J(T^P_{\Omega,\varepsilon}h:\varepsilon>0).
\]
The weak pointwise critical-jump estimate implies the maximal estimate
\begin{equation}\label{eq:maximalfromjump}
 \norm{\sup_{q\in\mathbb Q_+}|T^P_{\Omega,q}h|}_{L^{1,\infty}}
 \le C_{n,d}\mathfrak L(\Omega)\norm h_1.
\end{equation}

Start with any bounded compactly supported approximating sequence and
pass to a subsequence, still denoted by $f_k$, such that
\[
 f_k\to f\quad\hbox{in }L^1,\qquad
 \norm{f_{k+1}-f_k}_1\le2^{-3k}.
\]
In particular $\|f_k\|_1\to\|f\|_1$.  By
\eqref{eq:maximalfromjump},
\[
 \left|\left\{
 \sup_{q\in\mathbb Q_+}|T^P_{\Omega,q}(f_{k+1}-f_k)|>2^{-k}
 \right\}\right|
 \le C\mathfrak L(\Omega)2^k\|f_{k+1}-f_k\|_1
 \le C\mathfrak L(\Omega)2^{-2k}.
\]
The right side is summable in $k$.
Borel--Cantelli gives almost-everywhere uniform convergence on
$\mathbb Q_+$.  More precisely, let
\[
 E_{\rm BC}=\limsup_{k\to\infty}\left\{x:
  \sup_{q\in\mathbb Q_+}
  |T^P_{\Omega,q}(f_{k+1}-f_k)(x)|>2^{-k}\right\}.
\]
This is a measurable null set, since the supremum is countable.  For every
rational \(q>0\), define on all of \(\mathbb R^n\)
\begin{equation}\label{eq:rational-coordinate-definition}
 T^P_{\Omega,q}f(x)=
 \begin{cases}
  \displaystyle\lim_{k\to\infty}T^P_{\Omega,q}f_k(x),&x\notin E_{\rm BC},\\
  0,&x\in E_{\rm BC}.
 \end{cases}
\end{equation}
Thus all rational coordinates are defined on the same set, every one is
measurable, and every countable supremum formed from them is measurable.
Outside \(E_{\rm BC}\) there is $k_0(x)$ such that, for every
$k\ge k_0(x)$,
\[
 \sup_{q\in\mathbb Q_+}
 |T^P_{\Omega,q}(f_{k+1}-f_k)(x)|\le2^{-k}.
\]
Consequently the series
\[
 \mathcal R_m(q,x)
 :=\sum_{k\ge m}T^P_{\Omega,q}(f_{k+1}-f_k)(x)
\]
converges uniformly in $q\in\mathbb Q_+$ outside that exceptional set,
and, for every $m\ge k_0(x)$,
\begin{equation}\label{eq:rational-uniform-tail}
 \sup_{q\in\mathbb Q_+}
 |\mathcal R_m(q,x)|
 \le\sum_{k\ge m}2^{-k}.
\end{equation}
By \eqref{eq:rational-coordinate-definition}, outside \(E_{\rm BC}\),
$T^P_{\Omega,q}f=T^P_{\Omega,q}f_m+\mathcal R_m(q,\cdot)$; thus no
operator value on $f-f_m$ is being invoked before the completion is
constructed.
For the core function $f_m$, the trajectory is zero at $x$ for all
sufficiently large $q$.  Combining this fact with
\eqref{eq:rational-uniform-tail} gives the stronger uniform rational anchor
\begin{equation}\label{eq:uniform-rational-anchor}
 \lim_{R\to\infty}
 \sup_{\substack{q\in\mathbb Q_+\\q\ge R}}
 |T^P_{\Omega,q}f(x)|=0
\end{equation}
for almost every $x$: given $m$, take $R$ beyond the support-dependent
vanishing threshold for $f_m$, and then let $m\to\infty$.  Equivalently,
for every $m\ge k_0(x)$,
\[
 \limsup_{R\to\infty}
 \sup_{\substack{q\in\mathbb Q_+\\q\ge R}}
 |T^P_{\Omega,q}f(x)|
 \le\sum_{k\ge m}2^{-k};
\]
letting $m\to\infty$ proves \eqref{eq:uniform-rational-anchor} with the
quantifiers in the asserted order.

Select a further subsequence solely to obtain the pointwise finite-annulus
identities below, and do not relabel it.  This selection does not alter any
of the rational limits already obtained, because the convergence on
$\mathbb Q_+$ was uniform along the preceding sequence.
For every rational pair $0<a<b<\infty$, the absolute-value Young bound for
the finite-annulus kernel,
\[
 \|K_\Omega\mathbf1_{\{a<|\cdot|\le b\}}\|_1
 =\|\Omega\|_1\log(b/a),
\]
gives $L^1$ convergence of
\[
 T^P_{\Omega,a}f_k-T^P_{\Omega,b}f_k
 \quad\hbox{to}\quad
 \int_{a<|z|\le b}e^{iP(x,x-z)}K_\Omega(z)f(x-z)\,dz.
\]
A diagonal choice over the countable set of rational pairs gives pointwise
convergence of all these identities on one full-measure set.  Separately,
Fubini applied to the countable annuli
$\{N^{-1}<|z|<N\}$, $N\ge2$, gives a common full-measure set on which every
finite-annulus integral is finite and depends locally absolutely
continuously on its radial endpoints.  Intersecting these sets with the
complement of \(E_{\rm BC}\) preserves full measure; denote the resulting
measurable set by \(\mathcal X\).

On \(\mathcal X\) define the value at arbitrary $\varepsilon$ from the
base point $1$ by
\[
 T^P_{\Omega,\varepsilon}f=T^P_{\Omega,1}f+
 \begin{cases}
 \displaystyle\int_{\varepsilon<|z|\le1}e^{iP(x,x-z)}K_\Omega(z)f(x-z)dz,
 &\varepsilon<1,\\[1ex]
 \displaystyle-\int_{1<|z|\le\varepsilon}e^{iP(x,x-z)}K_\Omega(z)f(x-z)dz,
 &\varepsilon>1.
 \end{cases}
\]
On \(\mathcal X^c\) define \(T^P_{\Omega,\varepsilon}f=0\) simultaneously
for every \(\varepsilon>0\).  For each fixed \(\varepsilon\), this is a
measurable representative: the finite-annulus integral is measurable by
the parameter integral theorem and \(T^P_{\Omega,1}f\) is measurable by
\eqref{eq:rational-coordinate-definition}.  The modification on
\(\mathcal X^c\) merely replaces all rational representatives on one
additional common null set.
The resulting path is locally absolutely continuous, hence continuous,
in $\varepsilon$ and agrees with the already constructed rational values
on \(\mathcal X\) by the rational finite-annulus identities.

Fix an enumeration \(\mathbb Q_+=\{q_1,q_2,\ldots\}\) and set
\(D_N=\{q_1,\ldots,q_N\}\).  For every \(N\),
\eqref{eq:finite-J-lsc} gives on \(\mathcal X\), hence almost everywhere,
\[
 \mathbf J_{D_N}(T^P_{\Omega,\varepsilon}f:\varepsilon\in D_N)
 \le\liminf_{k\to\infty}
 \mathbf J_{D_N}(T^P_{\Omega,\varepsilon}f_k:\varepsilon\in D_N).
\]
The finite functional is Borel in its finitely many coordinates.  The
lower-semicontinuity statement is exactly the strict-witness statement:
if the left side is \(>\tau\), one may choose a jump amplitude and a finite
strict witness with a positive excess over \(\tau\); coordinatewise
convergence preserves that same witness for all sufficiently large \(k\).
Consequently, for every $\tau>0$,
\[
 \mathbf1_{\{\mathbf J_{D_N}(T^P_{\Omega,\varepsilon}f:
                    \varepsilon\in D_N)>\tau\}}
 \le\liminf_{k\to\infty}
   \mathbf1_{\{\mathbf J_{D_N}(T^P_{\Omega,\varepsilon}f_k:
                    \varepsilon\in D_N)>\tau\}}.
\]
Fatou's lemma and $\|f_k\|_1\to\|f\|_1$ therefore transfer the core
distribution estimate with the same right-hand side
$C\tau^{-1}\mathfrak L(\Omega)\|f\|_1$, uniformly in \(N\).  Because the
sets \(D_N\) are nested and every strict rational jump witness is finite,
\begin{equation}\label{eq:rational-finite-exhaustion}
 \{\mathbf J(T^P_{\Omega,q}f:q\in\mathbb Q_+)>\tau\}
 =\bigcup_{N=1}^\infty
   \{\mathbf J_{D_N}(T^P_{\Omega,q}f:q\in D_N)>\tau\}.
\end{equation}
Continuity from below therefore transfers the same distribution estimate
to the measurable rational pointwise functional.  Every finite
full-parameter witness lies in a compact interval; local continuity and
Lemma~\ref{lem:finite-witness} then identify it almost everywhere with
the full positive-parameter functional.  Finally, for every $R>0$,
continuity and density of $\mathbb Q_+$ imply
\[
 \sup_{\varepsilon\ge R}|T^P_{\Omega,\varepsilon}f(x)|
 \le \sup_{\substack{q\in\mathbb Q_+\\q\ge R}}
       |T^P_{\Omega,q}f(x)|.
\]
Letting $R\to\infty$ in \eqref{eq:uniform-rational-anchor} proves the
zero-at-infinity anchor for all real radii.  In particular the construction
is independent of the chosen core approximation: two resulting paths
have the same finite-annulus increments, so on the intersection of their
full-measure defining sets their difference is independent of
$\varepsilon$; their common zero anchor forces that difference to vanish.
This also proves uniqueness among all paths with the finite-annulus
identities and the zero anchor.  This completes the unweighted weak critical-jump
construction and, together with the strong theorem already proved,
establishes Theorem~\ref{thm:main}.

\subsection{Abstract consequences}

The pointwise inequality \eqref{eq:pointwise-jump-to-variation} applied to
the completed trajectory gives, for every $r>2$,
\[
 V^r(T^P_{\Omega,\varepsilon}f:\varepsilon>0)
 \le\kappa_r\mathbf J(T^P_{\Omega,\varepsilon}f:\varepsilon>0)
 \qquad\text{almost everywhere}.
\]
Taking the strong or weak spatial norm and using
Theorem~\ref{thm:main} proves \eqref{eq:main-strong-var} and
\begin{equation}\label{eq:unweighted-weak-var}
 \norm{V^r(T^P_{\Omega,\varepsilon}f:\varepsilon>0)}_{L^{1,\infty}}
 \le C_{n,d}\kappa_r\mathfrak L(\Omega)\norm f_1.
\end{equation}
No interchange of a spatial norm with the jump-amplitude supremum is used
here.  The zero anchor and \eqref{eq:anchor-controls-maximal} give the
maximal estimate, while finite $r$-variation gives the limit at zero.
This proves Corollary~\ref{cor:unweighted-consequences}.

\subsection{Radial \texorpdfstring{$\BV$}{BV} transfer}

Write
\[
 F_\varepsilon=T^P_{\Omega,\varepsilon}f,
 \qquad G_\varepsilon=T^{P,b}_{\Omega,\varepsilon}f.
\]
Fix throughout the right-continuous representative of $b$ specified
after \eqref{eq:radial-variation-norm}.  A right-continuous BV function
determines a finite complex Borel measure by
\[
 db((s,t])=b(t)-b(s),\qquad
 |db|(0,\infty)\le\Var_{(0,\infty)}(b).
\]
For test data the path $F$ is locally absolutely continuous.  The
Lebesgue--Stieltjes product formula on
$0<\varepsilon<R<\infty$ reads
\[
 b(R)F_R-b(\varepsilon)F_\varepsilon
 =\int_{(\varepsilon,R]}F_s\,db(s)
  +\int_{(\varepsilon,R]}b(s-)\,dF_s.
\]
Since $dF\ll ds$ has no atoms, the last integrand may equivalently be
written $b(s)$.  Moreover
$G_\varepsilon-G_R=-\int_\varepsilon^Rb(s)\,dF_s$.  Rearranging gives
\begin{equation}\label{eq:stieltjes-finite}
 G_\varepsilon-G_R
 =b(\varepsilon)F_\varepsilon-b(R)F_R
   +\int_{(\varepsilon,R]}F_s\,db(s).
\end{equation}
Altering $b$ on a countable set of radii does not change the original
Lebesgue shell integral, whereas the Stieltjes formula depends on the
chosen representative.  We next justify the passage to general data on
one common full-measure set.

Use the representatives of $F_\varepsilon$ constructed in the strong
completion and in the preceding weak endpoint completion.  In both cases
the base coordinate $F_1$ is the almost-everywhere limit of core-data
coordinates, and every other coordinate is fixed by its signed
finite-annulus increment from $1$.  Choose the corresponding core
approximation $f_k$ as follows.  In the strong $L^p$ case, pass to a
subsequence such that
\[
 \sum_k\norm{f_{k+1}-f_k}_p<\infty.
\]
Put
\[
 M_k=\sup_{q\in\mathbb Q_+}
 |T^P_{\Omega,q}(f_{k+1}-f_k)|.
\]
The strong maximal estimate and Minkowski's inequality give
\[
 \sum_k\norm{M_k}_p
 \le C\mathfrak L(\Omega)
      \sum_k\norm{f_{k+1}-f_k}_p<\infty.
\]
Applying this bound to finite partial sums and then using monotone
convergence shows that $\sum_kM_k\in L^p$, and hence is finite almost
everywhere.  In the weak $L^1$ case, choose instead
$\norm{f_{k+1}-f_k}_1\le2^{-3k}$.  By
\eqref{eq:maximalfromjump},
\[
 \sum_k\left|\left\{
  \sup_{q\in\mathbb Q_+}
  |T^P_{\Omega,q}(f_{k+1}-f_k)|>2^{-k}
 \right\}\right|<\infty,
\]
so Borel--Cantelli gives the same uniform Cauchy conclusion outside one
null set.  Thus in either case $F^{(k)}_q$ converges uniformly in
$q\in\mathbb Q_+$ on a common full-measure set.  Since every core-data
difference is locally continuous in the truncation parameter, the same
sequence is uniformly Cauchy on each compact parameter interval: for a
continuous difference, the supremum on such an interval is its supremum
on the rational parameters there.

The locally uniform limit of the continuous core trajectories is locally
continuous.  Fix at this point Borel representatives of $f$ and
$\Omega$.  Also choose a Borel function $F_1^{\rm B}$ which agrees
almost everywhere with the already completed base coordinate $F_1$.
All finite-annulus integrals below are formed with these fixed Borel
representatives.  For every rational pair $0<a<c<\infty$, Young's
inequality uses
\begin{equation}\label{eq:finite-annulus-L1}
 \norm{K_\Omega\mathbf1_{\{a<|\cdot|\le c\}}}_1
 \le\norm\Omega_1\log(c/a)
\end{equation}
to pass the core-data finite-annulus identity to the limit in $L^p$ in
the strong case and in $L^1$ in the weak case.  Pass to one further
diagonal subsequence.  The locally uniform convergence is unaffected, and
outside one null set the identity holds simultaneously for every rational
pair.

For $j\ge1$, put
\[
 J_j(x)=\int_{2^{-j}<|z|<2^j}
 |K_\Omega(z)|\,|f(x-z)|\,dz.
\]
The parameter integral theorem makes every $J_j$ Borel measurable, and
Young's inequality shows that it belongs to $L^p$ in the strong case and
to $L^1$ in the weak case.  Form one exceptional set from the following
null sets: the failure set for the locally uniform core limit; the set on
which $F_1^{\rm B}\ne F_1$; the simultaneous failure set for the rational
finite-annulus identities just obtained; the sets where some $J_j$ is
infinite; and the failure sets for the zero anchor and for finiteness of
the pointwise critical-jump functional.  This is a countable union of
Lebesgue null sets.  Contain it in a Borel null set
\(\mathcal N_{\rm BV}\), and put
\[
 \mathcal X_{\rm BV}=\mathbb R^n\setminus\mathcal N_{\rm BV}.
\]
Replace $F_1^{\rm B}$ by zero on \(\mathcal N_{\rm BV}\), and henceforth
denote this Borel base coordinate simply by $F_1$.  Thus, at every point
of \(\mathcal X_{\rm BV}\), it is the previously completed base coordinate,
all rational finite-annulus identities hold for the fixed Borel versions
of $f$ and \(\Omega\), the core trajectories converge locally uniformly,
and the zero anchor and finite critical-jump bound hold.

The function
\[
 \Phi(x,r,\theta)=e^{iP(x,x-r\theta)}
       \Omega(\theta)f(x-r\theta),
 \qquad (x,r,\theta)\in\mathbb R^n\times(0,\infty)\times\Sph^{n-1},
\]
is jointly Borel measurable.  By the parameter integral theorem, both
\((x,r)\mapsto\int_{\Sph^{n-1}}|\Phi(x,r,\theta)|\,d\sigma(\theta)\) and
the integrals of the real and imaginary parts of \(\Phi\), wherever
absolutely convergent, are Borel measurable.  Define on the entire product
space
\begin{equation}\label{eq:joint-radial-density-definition}
 \mathcal A(x,r)=
 \begin{cases}
  \displaystyle\int_{\Sph^{n-1}}\Phi(x,r,\theta)\,d\sigma(\theta),
  &x\in\mathcal X_{\rm BV}\ \text{and }\
   \displaystyle\int_{\Sph^{n-1}}|\Phi(x,r,\theta)|\,d\sigma(\theta)<\infty,\\
  0,&\text{otherwise}.
 \end{cases}
\end{equation}
Thus \((x,r)\mapsto\mathcal A(x,r)\) is a jointly Borel version, not merely
a separately defined family of radial sections.  Polar coordinates and
the finiteness of the $J_j$ give, for \(x\in\mathcal X_{\rm BV}\) and on
every compact $0<a<c<\infty$,
\[
 \int_a^c|\mathcal A(x,r)|\,\frac{dr}{r}
 \le\int_{a<|z|<c}|K_\Omega(z)|\,|f(x-z)|\,dz<\infty.
\]
Thus the sphere integral is absolutely convergent for $dr/r$-almost every
$r$, which is exactly the representative needed below; no assertion is
made about every individual spherical restriction of an $L^p$ function.
For rational $a<c$, the finite-annulus identity is
\[
 F_a(x)-F_c(x)=\int_a^c\mathcal A(x,r)\,\frac{dr}{r}.
\]
The left-hand side is continuous in the endpoints, while the indefinite
integral on the right is locally absolutely continuous.  Density of the
rational endpoints therefore extends the identity to all real
$0<a<c<\infty$.  Thus, on this same full-measure set, $F$ is locally
absolutely continuous and
\begin{equation}\label{eq:radial-density}
 \int_a^c|\mathcal A(x,r)|\,\frac{dr}{r}<\infty,
 \qquad
 F_a(x)-F_c(x)=\int_a^c\mathcal A(x,r)\,\frac{dr}{r}.
\end{equation}
Using the Borel base coordinate $F_1$ fixed above, define on the entire
product space
\begin{equation}\label{eq:joint-F-representative}
 F(x,s)=F_1(x)
 +\mathbf1_{\{s<1\}}\int_s^1\mathcal A(x,r)\,\frac{dr}{r}
 -\mathbf1_{\{s>1\}}\int_1^s\mathcal A(x,r)\,\frac{dr}{r}.
\end{equation}
The parameter integral theorem makes \((x,s)\mapsto F(x,s)\) jointly
Borel measurable.  On \(\mathcal X_{\rm BV}\) it agrees with every
previously constructed coordinate by \eqref{eq:radial-density}; on its
complement it is identically zero.  We henceforth use this common
representative.
Finally, the zero anchor and
\eqref{eq:anchor-controls-maximal} give
\[
 F_R\longrightarrow0\quad(R\to\infty),
 \qquad M_F:=\sup_{s>0}|F_s|\le\mathbf J(F)<\infty.
\]

On the entire product space define
\begin{equation}\label{eq:stieltjes}
 G(x,\varepsilon)=b(\varepsilon)F(x,\varepsilon)
 +\int_{(\varepsilon,\infty)}F(x,s)\,db(s).
\end{equation}
The Stieltjes integral is absolutely convergent because
$|db|(0,\infty)<\infty$ and $M_F<\infty$ on
\(\mathcal X_{\rm BV}\), and both \(F\) and \(G\) are zero on its
complement.  Moreover, the integrand
\((x,\varepsilon,s)\mapsto\mathbf1_{\{s>\varepsilon\}}F(x,s)\) is jointly
Borel.  Applying the parameter integral theorem to the total-variation
measure \(|db|\), and then to the real and imaginary parts of the complex
measure \(db\), proves that \((x,\varepsilon)\mapsto G(x,\varepsilon)\) is
jointly Borel measurable.  For
$0<\varepsilon<R<\infty$, integration by parts for the locally absolutely
continuous path $F$ and the fixed right-continuous representative of $b$
uses the product formula
\begin{equation}\label{eq:stieltjes-product-rule}
 b(R)F_R-b(\varepsilon)F_\varepsilon
 =\int_{(\varepsilon,R]}F_s\,db(s)
  +\int_\varepsilon^R b(s)\,dF_s.
\end{equation}
There is no ambiguity between $b(s)$ and $b(s-)$ in the second integral,
because $dF$ has no atoms.  By \eqref{eq:radial-density},
\[
 dF_s=-\mathcal A(x,s)\,\frac{ds}{s}.
\]
This equality is an equality of locally finite absolutely continuous
measures and uses the preceding $L^1_{\rm loc}(ds/s)$ representative of
$\mathcal A$; in particular it is only an almost-everywhere density
assertion.
Consequently, \eqref{eq:stieltjes-product-rule} gives
\begin{align}
 G_\varepsilon-G_R
 &=b(\varepsilon)F_\varepsilon-b(R)F_R
   +\int_{(\varepsilon,R]}F_s\,db(s)\notag\\
 &=-\int_{(\varepsilon,R]}b(s)\,dF_s\notag\\
 &=\int_{\varepsilon<|z|\le R}
   e^{iP(x,x-z)}K_\Omega(z)b(|z|)f(x-z)\,dz.
 \label{eq:general-bv-finite-shell}
\end{align}
Indeed,
$-\int_\varepsilon^R b(s)\,dF_s
=\int_\varepsilon^R b(s)\mathcal A(x,s)\,ds/s$, and polar coordinates
give exactly the last line, including its sign and half-open endpoint
convention.
The last integral is absolutely meaningful on the same set; its kernel has
$L^1$ norm at most
$\norm b_\infty\norm\Omega_1\log(R/\varepsilon)$.  Hence
\eqref{eq:stieltjes} has exactly the required finite-annulus increments for
general data.  For clarity, it also agrees with the original operator on
core data, rather than merely defining an abstract extension.  Indeed, if
\(f\) is bounded and compactly supported, then
\[
 \widetilde G_\varepsilon(x)
 =\int_{|z|>\varepsilon}e^{iP(x,x-z)}
   K_\Omega(z)b(|z|)f(x-z)\,dz
\]
is absolutely convergent for each \(\varepsilon>0\), has the same
finite-annulus increments as \(G\) by
\eqref{eq:general-bv-finite-shell}, and satisfies
\(\widetilde G_R(x)=0\) for all sufficiently large \(R\).  Since
\(G_R(x)\to0\), letting \(R\to\infty\) in the common increment identity
gives \(G_\varepsilon(x)=\widetilde G_\varepsilon(x)\) on the common
full-measure set.  For general data, the finite-annulus increments and
zero anchor uniquely determine the extension.  Right continuity fixes the
Stieltjes representative, while altering $b$ at one radius does not change
the original Lebesgue shell integral.  Finally,
\begin{equation}\label{eq:bv-anchor-explicit}
 |G_\varepsilon|
 \le\norm b_\infty|F_\varepsilon|
   +M_F|db|((\varepsilon,\infty))\longrightarrow0
 \qquad(\varepsilon\to\infty),
\end{equation}
so the zero anchor uniquely determines the extension.

Put $B=\norm b_\infty$, $V=\Var(b)$, and
$M_F=\sup_{s>0}|F_s|$.  The zero anchor gives
$M_F\le\mathbf J(F)$ by \eqref{eq:anchor-controls-maximal}.  For every
$r\ge1$, the same anchor also gives
\begin{equation}\label{eq:anchor-controls-variation}
 M_F\le V^r(F).
\end{equation}
Indeed, for fixed $s>0$ and $R>s$, the two-point partition gives
$|F_s-F_R|\le V^r(F)$; letting $R\to\infty$ and then taking the supremum
in $s$ proves \eqref{eq:anchor-controls-variation}.  For a variation
partition $t_0<\cdots<t_N$, use the exact product increment
\[
 b(t_i)F_{t_i}-b(t_{i-1})F_{t_{i-1}}
 =b(t_i)(F_{t_i}-F_{t_{i-1}})
  +(b(t_i)-b(t_{i-1}))F_{t_{i-1}}.
\]
Minkowski's inequality and
$(\sum_i|\Delta_i b|^r)^{1/r}\le\sum_i|\Delta_i b|\le V$ now give the
first line below.  Together with \eqref{eq:stieltjes} this yields,
pointwise,
\begin{align}
 V^r(bF)&\le B V^r(F)+VM_F,\notag\\
 V^r\left(\int_{(\varepsilon,\infty)}F_s\,db(s):\varepsilon>0\right)
 &\le VM_F,\notag\\
 V^r(G)&\le(B+2V)V^r(F).\label{eq:bv-variation}
\end{align}
Indeed, an increment of the Stieltjes-tail path over $(t_{i-1},t_i]$ is
the integral of $F$ over that half-open interval.  These intervals are
pairwise disjoint, so its $\ell^r$ variation is bounded by the sum of
their $|db|$-masses times $M_F$, hence by $VM_F$.

We record separately the critical-jump transfer.  If $M_F=0$, then
$F_s=0$ for every $s>0$, and \eqref{eq:stieltjes} gives $G_\varepsilon=0$;
all assertions are then immediate.  Hence assume $M_F>0$.  Suppose first
that $B>0$.  For an ordered family of $\lambda$-jumps of $bF$, write
\[
 b(t)F_t-b(s)F_s
 =b(t)(F_t-F_s)+(b(t)-b(s))F_s.
\]
Assign the pair to the first term if its modulus is $>\lambda/2$, and
otherwise to the second.  The first subfamily contributes at most
$N_{\lambda/(2B)}(F)$.  The parameter intervals in the second subfamily
are disjoint, and hence
\[
 \frac{\lambda}{2M_F}N_{\rm second}
 <\sum_{\rm second}|b(t)-b(s)|\le V.
\]
This proves
\begin{equation}\label{eq:bF-count}
 N_\lambda(bF)
 \le N_{\lambda/(2B)}(F)+\frac{2M_FV}{\lambda}.
\end{equation}
Use $\sqrt{u+v}\le\sqrt u+\sqrt v$.  On the set where the left side is
nonzero, $\lambda<2BM_F$.  With $\mu=\lambda/(2B)$, the contribution of
the second term after taking square roots is bounded by
\[
 \sqrt{2\lambda M_FV}
 \le2M_F\sqrt{BV}\le(B+V)M_F.
\]
It follows that
\begin{equation}\label{eq:bF-jump}
 \lambda N_\lambda(bF)^{1/2}
 \le2B\mu N_\mu(F)^{1/2}+C(B+V)M_F.
\end{equation}
For $B=0$ the path is zero.  Define the Stieltjes-tail path by
\[
H_\varepsilon=\int_{(\varepsilon,\infty)}F_s\,db(s).
\]
It satisfies $V^1(H)\le VM_F$.  The disjoint jump intervals give
$\lambda N_\lambda(H)<V^1(H)\le VM_F$, whereas a nonzero jump also forces
$\lambda<2\sup_\varepsilon|H_\varepsilon|\le2VM_F$.  Hence
\begin{equation}\label{eq:tail-jump}
 \lambda N_\lambda(H)^{1/2}\le C VM_F.
\end{equation}
Taking the supremum in $\lambda$ in
\eqref{eq:bF-jump}--\eqref{eq:tail-jump} gives
\[
 \mathbf J(bF)\le2B\mathbf J(F)+C(B+V)M_F,
 \qquad
 \mathbf J(H)\le CVM_F.
\]
Since $G=bF+H$, equations
\eqref{eq:pointwise-jump-quasitriangle} and
\eqref{eq:anchor-controls-maximal} imply the pointwise transfer
\begin{equation}\label{eq:bv-pointwise-jump-transfer}
 \mathbf J(G)\le C(B+V)\mathbf J(F).
\end{equation}
The strong and weak pointwise critical-jump estimates for $G$ now follow
directly from Theorem~\ref{thm:main}.  Formula
\eqref{eq:bv-variation} gives the asserted strong and weak $r$-variation
bounds.  This proves the estimates in Corollary~\ref{cor:bv-extension}.

The almost-everywhere zero anchor was proved in
\eqref{eq:bv-anchor-explicit}.  For $f\in L^p$, the strong zero anchor for
$F$, the strong maximal estimate, and continuity from above for the finite
measure $|db|$ give the explicit bound
\[
 \|G_\varepsilon\|_p
 \le B\|F_\varepsilon\|_p
   +|db|((\varepsilon,\infty))\|M_F\|_p
 \longrightarrow0.
\]
Thus both terms also tend to zero in $L^p$.  Since a scalar path of finite $r$-variation has
one-sided limits at both endpoints, $G_\varepsilon$ has an almost-everywhere
limit as $\varepsilon\downarrow0$.  The zero anchor gives
$\sup_{\varepsilon>0}|G_\varepsilon|\le\mathbf J(G)$ by
\eqref{eq:anchor-controls-maximal}, and hence the maximal truncation
assertion.  This completes the proof of
Corollary~\ref{cor:bv-extension}.


\section{Finite radial variation beyond bounded variation}
\label{sec:finite-radial}

For $1\le\rho<\infty$ recall the radial norm
\[
 \|b\|_{\mathcal V^\rho}=\|b\|_\infty+V^\rho(b;(0,\infty)).
\]
The purpose of this section is twofold.  First, we record a direct pathwise
Young--Stieltjes transfer when $\rho<2$.  Second, we prove
Theorem~\ref{thm:finite-radial} for every fixed finite $\rho$ by showing that
each oscillatory, packet, endpoint, and completion component of the
unweighted proof is stable under a radial $V^\rho$ amplitude.  The latter
argument is genuinely operator-specific at and above $\rho=2$.


\subsection{A pathwise Young--Stieltjes transfer below the radial exponent two}

For radial variation below two, the full operator-specific proof is not
needed.  The following abstract statement isolates a direct consequence of
the pointwise critical-jump theorem.  It is also useful for identifying the
precise obstruction at the exponent two.

Throughout this subsection, for a right-continuous integrator $b$ we use
the Young increment notation
\[
 db((s,t])=b(t)-b(s).
\]
This is only a half-open interval-increment convention; when
$\rho>1$, it does not assert that $db$ is a countably additive Stieltjes
measure.  Such a measure is available in the BV case, whereas the argument
below uses Young sums and the Young--Loeve estimate.
Young sums on a finite annulus use the corresponding half-open interval
$(a,c]$: a tagged sum over
$a=t_0<\cdots<t_N=c$ is
\[
 \sum_{j=1}^N F(\xi_j)\{b(t_j)-b(t_{j-1})\},
 \qquad \xi_j\in[t_{j-1},t_j].
\]
Thus a jump of the right-continuous integrator at the right endpoint is
included and one at the left endpoint is not.  For $F=F_1+iF_2$ and
$b=b_1+ib_2$, the complex Young integral is defined by
\[
 \int F\,db
 :=\int F_1\,db_1-\int F_2\,db_2
   +i\left(\int F_1\,db_2+\int F_2\,db_1\right).
\]
Thus it is the complex-bilinear extension of the four real integrals; the real
Young--Loeve estimate therefore gives the same bound up to an absolute
constant.  Since the path $F$ below is locally absolutely continuous and
hence continuous, there is no common-jump correction in integration by
parts.  The integral against $dF$ is the ordinary Lebesgue--Stieltjes
integral; since local absolute continuity makes $dF$ atomless, its value is
unchanged if $b(s)$ is replaced by $b(s-)$.

\begin{proposition}[Young--Stieltjes critical-jump transfer]
\label{prop:young-radial-transfer}
Let $F=(F_\varepsilon)_{\varepsilon>0}$ be a locally absolutely continuous
scalar path such that
\[
 F_\varepsilon\longrightarrow0\quad(\varepsilon\to\infty),
 \qquad \mathbf J(F)<\infty.
\]
Fix $1\le\rho<2$, and let $b:(0,\infty)\to\mathbb C$ be bounded and of
finite $\rho$-variation.  Choose its right-continuous representative.  Then
the improper Young integral
\begin{equation}\label{eq:young-tail-definition}
 H_\varepsilon:=\int_{(\varepsilon,\infty)}F_s\,db(s)
\end{equation}
exists, and the anchored path
\begin{equation}\label{eq:young-weighted-path}
 G_\varepsilon:=b(\varepsilon)F_\varepsilon+H_\varepsilon
\end{equation}
satisfies
\begin{equation}\label{eq:young-pointwise-transfer}
 \mathbf J(G)
 \le C_\rho\bigl(\|b\|_\infty+V^\rho(b)\bigr)\mathbf J(F).
\end{equation}
Moreover $G_\varepsilon\to0$ as $\varepsilon\to\infty$, and for every
$0<a<c<\infty$,
\begin{equation}\label{eq:young-finite-annulus}
 G_a-G_c=-\int_a^c b(s)\,dF_s.
\end{equation}
\end{proposition}

\begin{proof}
Use the norm-nonincreasing right-continuous representative constructed
after \eqref{eq:radial-variation-norm}.  To keep every constant quantitative,
fix once and for all
\[
 u=u_\rho:=
 \begin{cases}
  3,&\rho=1,\\[1mm]
  \frac12(2+\rho'),&1<\rho<2,
 \end{cases}
 \qquad \rho'=\frac{\rho}{\rho-1}.
\]
Then
\begin{equation}\label{eq:young-exponents}
 2<u<\rho',
 \qquad \frac1u+\frac1\rho>1,
\end{equation}
where $\rho'=\infty$ when $\rho=1$.  Put
\[
 J=\mathbf J(F),\quad R=V^u(F),\quad Q=V^\rho(b),
 \quad B=\|b\|_\infty,\quad M=\sup_{s>0}|F_s|.
\]
Equations~\eqref{eq:pointwise-jump-to-variation} and
\eqref{eq:anchor-controls-maximal} give
\begin{equation}\label{eq:young-RM-control}
 R\le\kappa_uJ,\qquad M\le J.
\end{equation}

Because $F$ is continuous and \eqref{eq:young-exponents} holds, the
real-valued Young theorem \cite{Young1936}, applied componentwise as fixed
above, applies on every compact interval.  Its Young--Loeve estimate gives
\begin{equation}\label{eq:young-loeve}
 \left|\int_{(a,c]}(F_s-F_a)\,db(s)\right|
 \le C_{u,\rho}V^u(F;[a,c])V^\rho(b;[a,c]).
\end{equation}
Consequently,
\begin{equation}\label{eq:young-local-tail-bound}
 \left|\int_{(a,c]}F_s\,db(s)\right|
 \le \sup_{a\le s\le c}|F_s|\,|b(c)-b(a)|
 +C_{u,\rho}V^u(F;[a,c])V^\rho(b;[a,c]).
\end{equation}
Finite $\rho$-variation gives a limit of $b$ at infinity, while finite
$u$-variation and the zero anchor imply that both
$\sup_{s\ge R}|F_s|$ and $V^u(F;[R,\infty))$ tend to zero as
$R\to\infty$.  For completeness, the latter assertion follows directly
from finite variation.  If it failed, there would be an $\eta>0$ and,
recursively, disjoint finite intervals $[R_k,S_k]$, tending to infinity,
with
\[
 V^u(F;[R_k,S_k])>\eta.
\]
Choose on each such interval a finite partition whose $u$-power sum is
larger than $\eta^u$.  Concatenating the first $N$ partitions gives
$V^u(F)^u\ge N\eta^u$, a contradiction.  Thus the two terms on the
right of \eqref{eq:young-local-tail-bound}, with $a,c\ge R$, tend to zero
uniformly as $R\to\infty$.  Hence that estimate makes the compact
Young integrals Cauchy and proves \eqref{eq:young-tail-definition}.  It also
yields
\begin{equation}\label{eq:young-tail-supremum}
 \sup_{\varepsilon>0}|H_\varepsilon|
 \le C_{u,\rho}Q(M+R)\le C_\rho QJ.
\end{equation}
With the half-open convention fixed above, Young integration by parts
(valid because $F$ is continuous) gives for
$0<\varepsilon<R_0<\infty$,
\[
 b(R_0)F_{R_0}-b(\varepsilon)F_\varepsilon
 =\int_{(\varepsilon,R_0]}F_s\,db(s)
  +\int_\varepsilon^{R_0}b(s)\,dF_s.
\]
As $R_0\to\infty$, the first integral on the right converges to
$H_\varepsilon$ by its Cauchy construction, while
$|b(R_0)F_{R_0}|\le B|F_{R_0}|\to0$.  The identity therefore also proves
the existence of the improper integral against $dF$ and gives
\begin{equation}\label{eq:young-tail-integration-by-parts}
 G_\varepsilon
 =b(\varepsilon)F_\varepsilon+H_\varepsilon
 =-\int_\varepsilon^\infty b(s)\,dF_s.
\end{equation}
Separately, as $\varepsilon\to\infty$,
\[
 |b(\varepsilon)F_\varepsilon|
 \le B\sup_{s\ge\varepsilon}|F_s|\longrightarrow0,
 \qquad
 |H_\varepsilon|
 \le Q\sup_{s\ge\varepsilon}|F_s|
   +C_{u,\rho}V^u(F;[\varepsilon,\infty))Q\longrightarrow0.
\]
Thus $G_\varepsilon\to0$ at infinity.  For two finite endpoints, the
tail definition and finite-interval integration by parts give explicitly
\begin{align*}
 G_a-G_c
 &=b(a)F_a-b(c)F_c+\int_{(a,c]}F_s\,db(s)\\
 &=-\int_a^c b(s)\,dF_s.
\end{align*}
This proves \eqref{eq:young-finite-annulus}, including its sign and
half-open endpoint convention.

It remains to estimate jumps.  Consider an ordered family of strict
$\lambda$-jumps of the product path $bF$.  On each jump interval $(s,t)$,
\[
 b(t)F_t-b(s)F_s
 =b(t)(F_t-F_s)+(b(t)-b(s))F_s.
\]
Assign the interval to the first term when its modulus exceeds
$\lambda/2$, and otherwise to the second.  For $B>0$ this gives
\begin{equation}\label{eq:young-product-count}
 N_\lambda(bF)
 \le N_{\lambda/(2B)}(F)
 +\left(\frac{2MQ}{\lambda}\right)^\rho.
\end{equation}
A nonzero jump also forces $\lambda\le2BM$.  Since $\rho<2$,
\[
 \lambda\left(\frac{2MQ}{\lambda}\right)^{\rho/2}
 \le C_\rho M B^{1-\rho/2}Q^{\rho/2}
 \le C_\rho M(B+Q).
\]
Moreover,
\[
 \lambda N_{\lambda/(2B)}(F)^{1/2}
 =2B\frac{\lambda}{2B}N_{\lambda/(2B)}(F)^{1/2}
 \le2BJ.
\]
Using $\sqrt{x+y}\le\sqrt x+\sqrt y$ in
\eqref{eq:young-product-count}, and then $M\le J$, therefore gives
\begin{equation}\label{eq:young-product-jump}
 \mathbf J(bF)\le C_\rho(B+Q)J.
\end{equation}
The case $B=0$ is immediate.

Now let $(s_i,t_i)$ be ordered $\lambda$-jump intervals for $H$, and put
\[
 A_i=V^u(F;[s_i,t_i]),\qquad C_i=V^\rho(b;[s_i,t_i]).
\]
By \eqref{eq:young-local-tail-bound},
\[
 |H_{s_i}-H_{t_i}|
 \le M|b(t_i)-b(s_i)|+C_{u,\rho}A_iC_i.
\]
If $Q=0$, the fixed right-continuous representative of $b$ is constant,
the Young tail $H$ vanishes, and there is nothing to prove.  Hence assume
$Q>0$ below.
Assign each interval to a family $\mathcal I_1$ or $\mathcal I_2$
according as the first or the second term on the right exceeds a fixed
positive multiple of $\lambda$, breaking ties arbitrarily.  If
$N_\nu=\#\mathcal I_\nu$, then
\[
 N_1\le C\left(\frac{MQ}{\lambda}\right)^\rho.
\]
Whenever a jump exists, \eqref{eq:young-tail-supremum} gives
$\lambda\le C_\rho QJ$.  Since $\rho<2$ and $M\le J$, the precise
interpolation is
\[
 \lambda N_1^{1/2}
 \le C(MQ)^{\rho/2}\lambda^{1-\rho/2}
 \le C_\rho(MQ)^{\rho/2}(QJ)^{1-\rho/2}
 \le C_\rho QJ.
\]
For the second family define $0<\theta<1$ by
\[
 \frac1\theta=\frac1u+\frac1\rho.
\]
Disjointness of the jump intervals and H\"older's inequality give
\[
 N_2(c\lambda)^\theta
 \le\sum_{i\in\mathcal I_2}(A_iC_i)^\theta
 \le
 \left(\sum_{i\in\mathcal I_2}A_i^u\right)^{\theta/u}
 \left(\sum_{i\in\mathcal I_2}C_i^\rho\right)^{\theta/\rho}
 \le R^\theta Q^\theta,
 \qquad
 N_2\le C_{u,\rho}\left(\frac{RQ}{\lambda}\right)^\theta.
\]
Using again $\lambda\le C_\rho QJ$ and $R\le\kappa_uJ$ gives
\[
 \lambda N_2^{1/2}
 \le C_{u,\rho}(RQ)^{\theta/2}\lambda^{1-\theta/2}
 \le C_\rho(RQ)^{\theta/2}(QJ)^{1-\theta/2}
 \le C_\rho QJ.
\]
For the original witness $N=N_1+N_2$, and hence
$\lambda\sqrt N\le\lambda\sqrt{N_1}+\lambda\sqrt{N_2}\le C_\rho QJ$.
Taking the supremum over all finite strict witnesses and then over
$\lambda>0$ gives
\begin{equation}\label{eq:young-tail-jump}
 \mathbf J(H)\le C_\rho QJ.
\end{equation}
Finally apply the pointwise quasi-triangle inequality
\eqref{eq:pointwise-jump-quasitriangle} to $G=bF+H$ and combine
\eqref{eq:young-product-jump} with \eqref{eq:young-tail-jump}.
\end{proof}


\subsection{The operator-specific theorem for every finite radial exponent}

We now prove Theorem~\ref{thm:finite-radial}.  The argument retains the
notation of the preceding sections so that every dependency can be checked
directly.  By homogeneity in $b$, it is enough to impose
\begin{equation}\label{eq:allq-normalize-b}
 \|b\|_{\mathcal V^\rho}\le1.
\end{equation}
The factor $\|b\|_{\mathcal V^\rho}$ is restored only at the end.  Notice
also the exact dilation invariance
\begin{equation}\label{eq:allq-dilation-b}
 \|b(R\,\cdot)\|_\infty=\|b\|_\infty,
 \qquad V^\rho(b(R\,\cdot))=V^\rho(b),
 \qquad R>0.
\end{equation}


\paragraph{Quantitative convention.}
Fix $n,d$, and $1\le\rho<\infty$ for the remainder of this section.
Every polynomial space, rational chart, compact interval, cutoff, and
finite-dimensional norm below belongs to one of the finitely many structural
families already fixed in Sections~\ref{sec:flag}--\ref{sec:degree-packets}.
We therefore take the minimum of the finitely many positive sublevel
exponents in Lemma~\ref{lem:multi-sublevel} and denote it by
$c_{\rm sub}=c_{\rm sub}(n,d)>0$. Constants denoted by $C_\rho$ may increase from line to line and
depend only on $n,d,\rho$ and the fixed structural data.  Polynomial powers
of an amplitude norm are structural and finite; whenever such a power enters
a later parameter choice, it is fixed once and for all at that point.
Positive exponents are always named before they are used. Replacing an
exponent by a smaller positive number is allowed without further comment. This convention makes the
parameter choices below finite and noncircular: first the oscillatory
exponents are fixed, then the packet hierarchy, and only afterwards the
angular splitting parameters.

\subsection{Finite \texorpdfstring{$\rho$}{rho}-variation calculus}

We first collect the elementary facts used everywhere below.
Throughout this subsection a compactly supported radial profile is first
replaced by its fixed right-continuous representative and then extended by
zero to the real line.  Variations without an indicated interval refer to
this zero extension.  The endpoint convention for consecutive cells is
specified in each displayed partition, and no endpoint value is counted
twice.  Thus every ordinary supremum and every endpoint jump refers to one
fixed representative.

\begin{lemma}[Disjoint intervals, pullbacks, products, and step profiles]
\label{lem:allq-calculus}
Let $1\le\rho<\infty$.

\begin{enumerate}[label=\textup{(\roman*)}]
\item If $a$ is a scalar function with $V^\rho(a)<\infty$ and
$I_1,\ldots,I_M$ are pairwise disjoint intervals arranged from left to
right, then
\begin{equation}\label{eq:allq-disjoint-var}
 \sum_{m=1}^M V^\rho(a;I_m)^\rho\le V^\rho(a)^\rho.
\end{equation}

\item If $a:I\to\mathbb C$ has finite $\rho$-variation and
$\gamma:J\to I$ is monotone, then
\begin{equation}\label{eq:allq-pullback}
 V^\rho(a\circ\gamma;J)\le V^\rho(a;I).
\end{equation}

\item For bounded $f,g$ with $V^\rho(f)+V^\rho(g)<\infty$,
\begin{equation}\label{eq:allq-product}
 V^\rho(fg)
 \le \|f\|_\infty V^\rho(g)+\|g\|_\infty V^\rho(f).
\end{equation}

\item Let $x_0<\cdots<x_K$ and
$I_k=[x_{k-1},x_k)$ for $1\le k\le K$.  Extend $\sigma$ by zero off
$[x_0,x_K)$, and suppose that $\sigma$ is constant on each $I_k$ and
$|\sigma|\le1$.  Then
\begin{equation}\label{eq:allq-step}
 V^\rho(\sigma)\le 2K^{1/\rho}.
\end{equation}
Consequently, if additionally $0<x_0<\cdots<x_K<\infty$ and
\eqref{eq:allq-normalize-b} holds, then, with $b\sigma$ extended by zero
off $[x_0,x_K)$,
\begin{equation}\label{eq:allq-b-step}
 \|b\sigma\|_\infty+V^\rho(b\sigma)
 \le C_\rho(1+K^{1/\rho}).
\end{equation}
\end{enumerate}
\end{lemma}

\begin{proof}
For (i), choose on every $I_m$ a partition which is within an arbitrary
$\epsilon$ of its variation and concatenate these partitions in increasing
order.  The within-$I_m$ increments are a subcollection of the increments in
the resulting global partition.  Letting $\epsilon\downarrow0$ proves
\eqref{eq:allq-disjoint-var}.  Assertion (ii) follows by pulling an ordered
partition of $J$ through $\gamma$ and deleting repetitions; if $\gamma$ is
decreasing, reverse the resulting list, which does not change the moduli of
its successive increments.  For (iii), on
every increment
\[
 f(t)g(t)-f(s)g(s)
 =f(t)(g(t)-g(s))+g(s)(f(t)-f(s)),
\]
and Minkowski in the finite $\ell^\rho$ norm gives
\eqref{eq:allq-product}.  Finally, fix an arbitrary variation partition
for the zero-extended step profile.  Delete zero increments and all but one
point from each constant cell that the partition visits.  The remaining
ordered list has at most $K-1$ increments between two cells, each of modulus
at most two, and at most two increments between the exterior zero value and
an interior cell, each of modulus at most one.  Thus every such partition
satisfies
\[
 \sum_j|\sigma(t_j)-\sigma(t_{j-1})|^\rho
 \le 2+(K-1)2^\rho\le K2^\rho.
\]
Taking the supremum proves \eqref{eq:allq-step}.  More explicitly, the zero extension
creates at most two endpoint increments, and hence
\[
 V^\rho(b\sigma;\mathbb R)
 \le \|b\|_{L^\infty(\supp\sigma)}V^\rho(\sigma;\mathbb R)
 +\|\sigma\|_\infty
  \{V^\rho(b;\supp\sigma)+2\|b\|_\infty\}.
\]
This proves \eqref{eq:allq-b-step} without extending $b$ itself outside
$(0,\infty)$.
\end{proof}

\begin{lemma}[Translation modulus and mollification]
\label{lem:allq-translation}
Let $I\Subset\R$ be fixed and let $a$ be bounded, supported in the
interior of $I$, and extended by zero to $\R$.  All variation seminorms in
this lemma are taken for that zero extension.  Then for $|h|\le1$,
\begin{equation}\label{eq:allq-translation}
 \|a(\cdot+h)-a\|_{L^1(\R)}
 \le C_{I,\rho}|h|^{1/\rho}
 \bigl(\|a\|_\infty+V^\rho(a;I)\bigr).
\end{equation}
If $0<h\le1$, $\phi$ is a standard compactly supported mollifier,
$\phi_h(t):=h^{-1}\phi(t/h)$, and $a_h:=a*\phi_h$, then, for every
integer $N\ge0$,
\begin{align}
 \|a-a_h\|_1
 &\le C_{I,\rho}h^{1/\rho}
  \bigl(\|a\|_\infty+V^\rho(a;I)\bigr),
 \label{eq:allq-moll-L1}\\
 V^\rho(a_h)&\le V^\rho(a),
 \label{eq:allq-moll-var}\\
 \|\partial^Na_h\|_\infty&\le C_Nh^{-N}\|a\|_\infty.
 \label{eq:allq-moll-derivative}
\end{align}
\end{lemma}

\begin{proof}
It is enough to treat $h>0$.  For $x\in[0,h)$ consider the ordered chain
$x+kh$.  Only $O_I(1+h^{-1})$ increments meet the support of the zero
extension.  H\"older's inequality and the definition of $V^\rho$ give
\[
 \sum_k|a(x+(k+1)h)-a(x+kh)|
 \le C_Ih^{-(1-1/\rho)}
 \bigl(\|a\|_\infty+V^\rho(a;I)\bigr).
\]
Integrating over $x\in[0,h)$ proves \eqref{eq:allq-translation}.  Averaging
this inequality against the mollifier proves \eqref{eq:allq-moll-L1}.
For any finite partition $(t_i)$, Minkowski in $\ell^\rho$ gives
\[
 \left(\sum_i|a_h(t_i)-a_h(t_{i-1})|^\rho\right)^{1/\rho}
 \le \int\phi_h(u)
 \left(\sum_i|a(t_i-u)-a(t_{i-1}-u)|^\rho\right)^{1/\rho}du,
\]
which yields \eqref{eq:allq-moll-var}; the derivative estimate is standard.
\end{proof}

The next elementary upper-sum estimate is the missing ingredient in the
terminal packet argument.  Small $L^1$ error alone is not enough there,
because packet inputs are controlled by average mass rather than pointwise
density.

\begin{lemma}[Upper Riemann sums for finite variation]
\label{lem:allq-upper-riemann}
Let $I\Subset\R$ be a fixed bounded interval, and let
$a:\R\to\mathbb C$ be bounded, supported in $I$, and satisfy
$V^\rho(a;\R)<\infty$.  Partition $I$ into consecutive intervals $J_m$
whose lengths are comparable to $\delta\in(0,1]$.  Then
\begin{equation}\label{eq:allq-upper-riemann}
 \delta\sum_m\sup_{J_m}|a|
 \le C_I\left[
 \|a\|_{L^1(\R)}+
 \delta^{1/\rho}\bigl(\|a\|_\infty+V^\rho(a;\R)\bigr)
 \right].
\end{equation}
If $J_m^*$ is obtained from $J_m$ by an enlargement about its centre by a
fixed factor, then the same estimate holds with $\sup_{J_m^*}|a|$ on the
left; the constant may additionally depend on that fixed factor.
\end{lemma}

\begin{proof}
For an un-enlarged interval,
\[
 \sup_{J_m}|a|
 \le |J_m|^{-1}\int_{J_m}|a|+\operatorname{osc}_{J_m}(a).
\]
For each $m$, choose two ordered points whose increment approaches
$\operatorname{osc}_{J_m}(a)$.  Concatenating these pairs in increasing
order and then taking limits shows
\[
 \sum_m\operatorname{osc}_{J_m}(a)^\rho\le V^\rho(a;\R)^\rho.
\]
H\"older's inequality and the fact that there are
$O_I(\delta^{-1})$ intervals give
\[
 \delta\sum_m\operatorname{osc}_{J_m}(a)
 \le C_I\delta^{1/\rho}V^\rho(a;\R),
\]
which proves \eqref{eq:allq-upper-riemann}.  For fixed enlargements, the
interval graph of $(J_m^*)_m$ has uniformly bounded degree and hence a
bounded proper coloring.  Intervals in one color have disjoint interiors.
Apply the preceding oscillation argument to each color; after using the
same average-plus-oscillation inequality on $J_m^*$, the average terms are
controlled by bounded overlap.  This proves the enlarged version without
identifying values at adjacent endpoints.
\end{proof}

\begin{lemma}[Radial discretization for packet inputs]
\label{lem:allq-radial-discretization}
Fix $0<c_-<c_+<\infty$.  Let $\Theta\in L^\infty(\Sph^{n-1})$,
$\Lambda\ge0$, and let $w$ be a bounded measurable function supported in
$[c_-,c_+]$ with $V^\rho(w)<\infty$.  Let $0<\delta\le1$, let $\{U\}$ be pairwise
disjoint dyadic $\delta$-cubes in $\R^n$, and suppose functions $d_U\in L^1(\R^n)$ satisfy
\[
 \supp d_U\subset U,
 \qquad \|d_U\|_1\le\Lambda|U|.
\]
Then
\begin{align}
 \left\|K_\Theta(\cdot)w(|\cdot|)*\sum_Ud_U\right\|_\infty
 &\le C\Lambda\|\Theta\|_\infty
 \left[
 \|w\|_1+
 \delta^{1/\rho}(\|w\|_\infty+V^\rho(w))
 \right],
 \label{eq:allq-radial-Linf}\\
 \left\|K_\Theta(\cdot)w(|\cdot|)*\sum_Ud_U\right\|_1
 &\le C\|\Theta\|_1\|w\|_1\sum_U\|d_U\|_1.
 \label{eq:allq-radial-L1}
\end{align}
Here the $L^1$ norm of $w$ is one-dimensional; on the fixed annulus it is
equivalent to the natural $dr/r$ norm.
\end{lemma}

\begin{proof}
Only \eqref{eq:allq-radial-Linf} is new.  Fix $x$.  Since $|z|$ stays in a
fixed annulus,
\[
 \left|K_\Theta w*\sum_Ud_U(x)\right|
 \le C\Lambda\|\Theta\|_\infty
 \delta^n\sum_U\sup_{y\in U}|w(|x-y|)|.
\]
Group the cubes according to intervals of length $C\delta$ containing the
distance from $x$ to their centers.  If their centers lie in one such
layer, the union of the corresponding disjoint cubes is contained in a
fixed enlargement of a radial layer of volume $O(\delta)$.  Since every
cube has volume $\delta^n$, the layer contains at most
$C\delta^{1-n}$ cubes.  Hence the last display is at most
\[
 C\Lambda\|\Theta\|_\infty
 \delta\sum_m\sup_{J_m^*}|w|,
\]
where the $J_m^*$ are fixed enlargements of consecutive radial intervals of
length comparable to $\delta$.  Lemma~\ref{lem:allq-upper-riemann} proves
\eqref{eq:allq-radial-Linf}.  Equation \eqref{eq:allq-radial-L1} is Young's
inequality and polar coordinates.
\end{proof}

\subsection{Oscillatory integrals with a finite-variation amplitude}

\begin{lemma}[First derivative estimate with a $V^\rho$ amplitude]
\label{lem:allq-first-derivative}
Let $I$ be a fixed bounded interval, let $\phi\in C^1(I;\mathbb R)$, and
let $a:I\to\mathbb C$ be bounded with $V^\rho(a;I)<\infty$.  Suppose
$\phi'$ is monotone and
\[
 |\phi'(t)|\ge L\ge1\qquad(t\in I).
\]
Then
\begin{equation}\label{eq:allq-first-derivative}
 \left|\int_Ie^{i\phi(t)}a(t)\,dt\right|
 \le C_{I,\rho}L^{-1/\rho}
 \bigl(\|a\|_\infty+V^\rho(a;I)\bigr).
\end{equation}
The same constant is valid when $I$ is replaced by any subinterval of the
fixed ambient interval.
\end{lemma}

\begin{proof}
Partition $I$ into consecutive intervals $I_1,\ldots,I_M$ of length at
most $L^{-1}$ and, except possibly for the last interval, at least
$(2L)^{-1}$.  Thus $M\le C_IL$ because $L\ge1$.  Choose
$t_m\in I_m$ in increasing order.  The error after replacing $a$ by
$a(t_m)$ on $I_m$ is bounded by
\[
 \sum_m|I_m|V^\rho(a;I_m)
 \le C L^{-1}M^{1-1/\rho}
 \left(\sum_mV^\rho(a;I_m)^\rho\right)^{1/\rho}
 \le C L^{-1/\rho}V^\rho(a;I),
\]
using Lemma~\ref{lem:allq-calculus}(i).  Put
$w_m=\int_{I_m}e^{i\phi(t)}dt$.  The ordinary first derivative van der
Corput estimate implies
\[
 \left|\sum_{m=r}^sw_m\right|
 =\left|\int_{\cup_{m=r}^sI_m}e^{i\phi(t)}dt\right|\le CL^{-1}
\]
for every consecutive block.  Put $a_m=a(t_m)$ and
$W_m=\sum_{j=1}^mw_j$.  Discrete Abel summation gives
\[
 \sum_{m=1}^Ma_mw_m
 =a_MW_M+\sum_{m=1}^{M-1}(a_m-a_{m+1})W_m,
 \qquad |W_m|\le CL^{-1}.
\]
Consequently,
\[
 \left|\sum_ma(t_m)w_m\right|
 \le CL^{-1}\left(\|a\|_\infty+
 \sum_{m=1}^{M-1}|a(t_{m+1})-a(t_m)|\right).
\]
Since the samples are ordered,
\[
 \left(\sum_{m=1}^{M-1}|a(t_{m+1})-a(t_m)|^\rho\right)^{1/\rho}
 \le V^\rho(a;I).
\]
H\"older therefore bounds the last sum by
$M^{1-1/\rho}V^\rho(a;I)$, proving
\eqref{eq:allq-first-derivative}.
\end{proof}

\begin{lemma}[Rational van der Corput with a $V^\rho$ amplitude]
\label{lem:allq-rational-vdc}
Let $\psi=N/Q$ on a fixed compact interval, where $N,Q$ are real
polynomials of degree at most $D$, no bound is imposed on the coefficients
of $N$, and
\[
 \|Q\|_{\mathrm{coeff}}\le M,
 \qquad 0<\eta\le |Q|\le M
\]
on a fixed enlargement of the interval. Put
\[
 G:=\|Q^2\psi'\|_{\mathrm{coeff}}
   =\|N'Q-NQ'\|_{\mathrm{coeff}}.
\]
Let $c_0=c_0(D)>0$ be a valid normalized one-dimensional polynomial
sublevel exponent for the derived degree bound of
$N'Q-NQ'$ (at most $2D-1$ when $D\ge1$). Then, with
\begin{equation}\label{eq:allq-vdc-explicit-exponent}
 c_{\mathrm{vdc},\rho,D}
 :=\frac{c_0}{1+c_0\rho}>0,
\end{equation}
every bounded amplitude $a$ supported in the interval with
$V^\rho(a)<\infty$ satisfies
\begin{equation}\label{eq:allq-rational-vdc}
 \left|\int e^{i\psi(v)}a(v)\,dv\right|
 \le C_{\rho,D,\eta,M}
 (\|a\|_\infty+V^\rho(a))
 \min\{1,G^{-c_{\mathrm{vdc},\rho,D}}\}.
\end{equation}
When $G=0$, the minimum on the right is understood to be $1$. The same
assertion holds after a uniformly bounded sectional partition. In
particular, the constants are independent of the coefficients of the
numerator $N$.
\end{lemma}

\begin{proof}
Put $p=Q^2\psi'=N'Q-NQ'$ and choose
\[
 \vartheta=\frac{c_0\rho}{1+c_0\rho}\in(0,1).
\]
The trivial $L^1$ bound proves the assertion for $G\le1$.  If
$1<G<(M^2)^{1/\vartheta}$, the same bound is still dominated by the
right side of \eqref{eq:allq-rational-vdc} after increasing its structural
constant.  We may therefore assume $M^{-2}G^\vartheta\ge1$. The
normalized polynomial $p/G$ has coefficient norm one. Hence the
fixed-degree sublevel estimate gives
\[
 E=\{|p|\le G^\vartheta\},
 \qquad |E|\le C G^{-c_0(1-\vartheta)}.
\]
The contribution of $E$ is at most
$C\|a\|_\infty G^{-c_0(1-\vartheta)}$. On $E^c$ the denominator bound
implies
\[
 |\psi'|=|p|/|Q|^2\ge M^{-2}G^\vartheta.
\]
The boundary of $E$ is contained in the zero sets of the two
bounded-degree polynomials $p\pm G^\vartheta$, and the numerator of
$\psi''$ also has uniformly bounded degree.  If that numerator vanishes
identically, then $\psi'$ is constant and no additional cut is needed;
otherwise it has only $O_D(1)$ zeros. Splitting at all these points
produces $O_D(1)$ intervals on each of which $\psi'$ is monotone and
bounded away from zero. Lemma~\ref{lem:allq-first-derivative}
on every such interval yields
\[
 C_{\rho,D,\eta,M}
 (\|a\|_\infty+V^\rho(a))G^{-\vartheta/\rho}.
\]
The chosen value of $\vartheta$ gives the exact equality
\[
 c_0(1-\vartheta)=\vartheta/\rho
 =\frac{c_0}{1+c_0\rho}
 =c_{\mathrm{vdc},\rho,D}.
\]
Summing the boundedly many intervals proves
\eqref{eq:allq-rational-vdc}.  Indeed, if these intervals are $J_k$, then
\[
 \sum_k V^\rho(a;J_k)
 \le C_D^{1-1/\rho}
 \left(\sum_kV^\rho(a;J_k)^\rho\right)^{1/\rho}
 \le C_{D,\rho}V^\rho(a),
\]
and the finitely many supremum terms cost at most $C_D\|a\|_\infty$.
Equivalently, if the restrictions are zero extended, their endpoint jumps
cost at most two additional supremum terms per interval.  A uniformly
bounded initial partition is handled by the same refinement.
\end{proof}

\begin{corollary}[Coefficient-norm polynomial van der Corput]
\label{cor:allq-polynomial-vdc}
For every fixed degree $D$, every real polynomial $p$ of degree at most
$D$, and every bounded measurable amplitude $a$ supported on a fixed
compact interval with $V^\rho(a)<\infty$,
\begin{equation}\label{eq:allq-polynomial-vdc}
 \left|\int e^{ip(t)}a(t)dt\right|
 \le C_{D,\rho}(\|a\|_\infty+V^\rho(a))
 \min\{1,\|p-p(0)\|_{\mathrm{coeff}}^{-c_{D,\rho}}\}.
\end{equation}
When $p$ is constant, the minimum on the right is understood to be $1$.
\end{corollary}

\begin{proof}
Apply Lemma~\ref{lem:allq-rational-vdc} with denominator identically one.
That lemma imposes no bound on the numerator coefficients, so the resulting
estimate is uniform over all coefficients of $p$.  Finally use
finite-dimensional norm equivalence between the coefficient norm of $p'$
and that of $p$ modulo constants.
\end{proof}

\begin{lemma}[Uniform rational oscillatory operator with sectional $V^\rho$]
\label{lem:allq-rational-operator}
In Lemma~\ref{lem:rational-operator}, retain $B\ge1$ and replace the sectional hypothesis
$\|A(u,\cdot)\|_\infty+\operatorname{Var}(A(u,\cdot))\le B$ by
\[
 \|A(u,\cdot)\|_\infty+V^\rho(A(u,\cdot))\le B
\]
on at most $M_0$ intervals, and retain the same row and column $L^1$
bounds. Let $c_{\mathrm{vdc},\rho}>0$ be the minimum of the finitely many
exponents in \eqref{eq:allq-vdc-explicit-exponent} that occur in the
structural rational families, and put
\begin{equation}\label{eq:allq-rational-operator-exponent}
 \delta_{\mathrm{rat},\rho}
 :=\frac{2c_{\rm sub}c_{\mathrm{vdc},\rho}}
         {c_{\rm sub}+2c_{\mathrm{vdc},\rho}},
 \qquad
 c_{\mathrm{rat},\rho}:=\frac{\delta_{\mathrm{rat},\rho}}4>0.
\end{equation}
Then
\begin{equation}\label{eq:allq-rational-operator}
 \|S\|_{2\to2}
 \le C_\rho B
       \min\{1,L^{-c_{\mathrm{rat},\rho}}\},
\end{equation}
where $L=\|[\Psi]\|_{{\rm rat},Q}$. The constants are independent of all
phase coefficients.
\end{lemma}

\begin{proof}
The Schur bound handles $L\le1$, so assume $L>1$. The kernel of $SS^*$ is
\[
 K(u,u')=\int e^{i[\Psi(u,v)-\Psi(u',v)]}
 A(u,v)\overline{A(u',v)}\,dv.
\]
Refine the two sectional partitions belonging to $u$ and $u'$. On each
refined interval Lemma~\ref{lem:allq-calculus}(iii) gives
\[
 \|A(u,\cdot)\overline{A(u',\cdot)}\|_\infty+
 V^\rho(A(u,\cdot)\overline{A(u',\cdot)})\le CB^2.
\]
With
$G(u,u')=\|\mathfrak C_Q\Psi(u,u',\cdot)\|_{\mathrm{coeff},v}$,
Lemma~\ref{lem:allq-rational-vdc} therefore yields
\begin{equation}\label{eq:allq-rational-kernel-pointwise}
 |K(u,u')|
 \le C_\rho B^2
 \min\{1,G(u,u')^{-c_{\mathrm{vdc},\rho}}\}.
\end{equation}
Let $\mathbf p(u,u')$ be the vector of coefficients in $v$ whose Euclidean
length is $G(u,u')$.  Regrouping coefficients is an isomorphism between two
fixed finite-dimensional coefficient spaces, so
\[
 A_0:=\|\mathbf p\|_{\mathrm{coeff},(u,u')}
 \asymp \|\mathfrak C_Q\Psi\|_{\mathrm{coeff},(u,u',v)}=L.
\]
Apply \eqref{eq:multi-negative} on $I\times I$ to the square of
\eqref{eq:allq-rational-kernel-pointwise}, with
$\sigma=2c_{\mathrm{vdc},\rho}$ and scale parameter $L\ge1$.  Since
$A_0\asymp L$, its normalized factor
$L|\mathbf p(u,u')|/A_0$ is comparable to $G(u,u')$, so its integrand is
comparable, with structural constants, to
$\min\{1,G(u,u')^{-2c_{\mathrm{vdc},\rho}}\}$.  The definition
\eqref{eq:allq-rational-operator-exponent} gives
\[
 \iint |K(u,u')|^2\,du\,du'
 \le C_\rho B^4 L^{-\delta_{\mathrm{rat},\rho}}.
\]
Consequently
\[
 \|S\|_{2\to2}^2=\|SS^*\|_{2\to2}
 \le\|K\|_{L^2(I\times I)}
 \le C_\rho B^2L^{-\delta_{\mathrm{rat},\rho}/2}.
\]
Thus the negative-moment estimate has first been square-rooted to pass
from $\iint|K|^2$ to $\|K\|_2$, and the $SS^*$ identity supplies the
second square root.  Taking that final square root proves
\eqref{eq:allq-rational-operator}. Algebraically degenerate points and
the diagonal are already included in the negative-moment integral, so no
coefficient-dependent deletion is made.
\end{proof}

\begin{lemma}[Sectional pullback calculus]
\label{lem:allq-sectional-calculus}
Let $I,J$ be fixed compact intervals and let
\[
 A(u,v)=w(u,v)\prod_{\ell=1}^m
 a_\ell(\gamma_\ell(u,v)).
\]
Assume that $w$ and every $\gamma_\ell$ are jointly Borel measurable on
$I\times J$; in particular, after choosing the fixed regulated
representatives of the $a_\ell$, the function $A$ is jointly measurable.
Assume that every $a_\ell$ is extended by zero from one fixed compact
radial interval and
\[
 \|a_\ell\|_\infty+V^\rho(a_\ell;\mathbb R)\le B,
 \qquad B\ge1.
\]
Suppose that, in either variable with the other fixed, each
$\gamma_\ell$ has at most $J_0$ monotonicity intervals, uniformly in the
fixed variable.  Suppose also that $w$ has uniformly bounded
$L^\infty$ norm and sectional $W^{1,\infty}$ norms on a fixed partition
into at most $J_0$ rectangles.  Fixed half-open cutoffs with a uniformly
bounded number of sectional jumps may be included in $w$.  Then, for
each fixed $u\in I$, the interval $J$ has a partition into at most
$C_{m,J_0}$ pieces $J_\alpha(u)$ on each of which
\begin{equation}\label{eq:allq-sectional-calculus}
 \|A(u,\cdot)\|_{L^\infty(J_\alpha(u))}
 +V^\rho(A(u,\cdot);J_\alpha(u))
 \le C_{m,J_0,\rho}B^m.
\end{equation}
Separately, for each fixed $v\in J$, the interval $I$ has a partition
into at most $C_{m,J_0}$ pieces $I_\beta(v)$ on each of which
\[
 \|A(\cdot,v)\|_{L^\infty(I_\beta(v))}
 +V^\rho(A(\cdot,v);I_\beta(v))
 \le C_{m,J_0,\rho}B^m.
\]
Moreover
\begin{equation}\label{eq:allq-sectional-schur}
 \sup_u\int_J|A(u,v)|\,dv+
 \sup_v\int_I|A(u,v)|\,du
 \le C_{m,J_0}B^m.
\end{equation}
\end{lemma}

\begin{proof}
For a fixed row take the common refinement of the structural partition and
all monotonicity partitions of the functions
$v\mapsto\gamma_\ell(u,v)$.  Its cardinality is bounded solely in terms
of $m,J_0$.  On every resulting interval, monotone pullback gives
\[
 V^\rho(a_\ell\circ\gamma_\ell)
 \le V^\rho(a_\ell;\mathbb R)\le B.
\]
The zero extension is essential here: crossings of the radial support are
already among the increments counted by the variation of $a_\ell$.
The sectional $W^{1,\infty}$ bound gives
$V^\rho(w)\le V^1(w)\le C$.  Repeated application of
Lemma~\ref{lem:allq-calculus}(iii), including the boundedly many cutoff
jumps, proves the row part of \eqref{eq:allq-sectional-calculus}.  The
column proof is identical.  Finally the rectangles have fixed length and
$|A|\le C B^m$, which proves the two independent Schur bounds
\eqref{eq:allq-sectional-schur}.  No two-variable or mixed variation norm
is used.
\end{proof}

\subsection{Stability of the correlation and high-shell estimates}

\begin{proposition}[Rough correlation with finite radial variation]
\label{prop:allq-rough-correlation}
Fix $1\le\rho<\infty$ and put $d_*=\max(d,1)$. Uniformly for
every polynomial degree $1\le D\le d_*$, there are positive exponents
$c_{\mathrm{eq},\rho}$ and $c_{\mathrm{uneq},\rho}$ and a finite power
$C_{\mathrm{corr},\rho}$, depending only on $n,d,\rho$, such that the
following two statements hold.

First, the conclusion of Lemma~\ref{lem:rough-correlation} remains valid
for phases of degree at most $D$ and profiles satisfying
$\|a\|_\infty+V^\rho(a)\le B$ and
$\|b\|_\infty+V^\rho(b)\le B$, where the profiles are zero extended and
$B\ge1$, with
\begin{equation}\label{eq:allq-equal-correlation}
 \left|\int e^{iS(z)}\overline{k_a(z)}k_b(z-h)\,dz\right|
 \le C_\rho B^{C_{\mathrm{corr},\rho}}\|\Theta\|_\infty^2
 \min\{1,(\Lambda|h|^D)^{-c_{\mathrm{eq},\rho}}\}
\end{equation}
when $n\ge2$, and with
$\min\{1,\Lambda^{-c_{\mathrm{eq},\rho}}\}$ in dimension one.

Second, Corollary~\ref{cor:unequal-correlation} remains valid for
$V^\rho$ profiles, with
\begin{equation}\label{eq:allq-unequal-correlation}
 \left|\int e^{i[P(y+z,a)-P(y+z,y)]}
 \overline{k_R(z)}k_S(z+y-a)\,dz\right|
 \le C_\rho B^{C_{\mathrm{corr},\rho}}\|\Theta\|_\infty^2S^{-n}
 \min\{1,(L_R|\mathbf p(W)|)^{-c_{\mathrm{uneq},\rho}}\}
\end{equation}
for $n\ge2$, with the stated $S^{-1}$ analogue in dimension one.
\end{proposition}

\begin{proof}
We follow the coordinate calculations in the two cited proofs and record
all places where radial regularity enters.

For the equal-shell estimate, on each normalized $(X,Y)$ block the two
profile arguments are, with the notation of the proof of
Lemma~\ref{lem:rough-correlation},
\[
 \gamma_1(X,Y)=\frac{\sqrt{U^2+\alpha^2}}{|q(X,Y)|},
 \qquad
 \gamma_2(X,Y)=\frac{\sqrt{V^2+\alpha^2}}{|q(X,Y)|},
\]
where $U=U_0+HX$, $V=V_0+HY$, and
$q=\Delta+X-Y$ has fixed sign and is bounded away from zero on a fixed
enlargement.  More precisely, if $K=[r_-,r_+]\Subset(0,\infty)$ is the
common fixed support from Lemma~\ref{lem:rough-correlation}, then the
zero-extended amplitude is supported where both pullbacks lie in $K$.
Choose $\chi\in C_c^\infty((0,\infty))$ equal to one on $K$ and supported
in a slightly larger fixed interval $K^+\Subset(0,\infty)$.  Because the
profiles are zero outside $K$, inserting
$\chi(\gamma_1)\chi(\gamma_2)$ does not change the amplitude; it does,
however, give a smooth support factor on whose transition region the same
lower bounds remain valid.  The fixed block cutoffs and this enlarged
support condition give, uniformly on a fixed enlargement,
\[
 |U|+|V|\le C_K,\qquad c_K\le |q|\le C_K,
 \qquad c_K\le s_U,s_V\le C_K,
 \quad s_U=(U^2+\alpha^2)^{1/2},\ s_V=(V^2+\alpha^2)^{1/2}.
\]
Thus every denominator occurring in the normalized Jacobian and in the
fixed rational cutoffs is separated from zero on the actual support.
Equations displayed immediately before
\eqref{eq:rough-fixed-angular-operator} give
\[
 \partial_X\frac{s_U}{q}
  =\frac{HUq-s_U^2}{s_Uq^2},\qquad
 \partial_Y\frac{s_U}{q}=\frac{s_U}{q^2},
\]
and the analogous two formulas for $s_V/q$.  Thus the only possible
critical-point equations are polynomials of degree at most two.  Each row
and column therefore splits into at most $C_{n,d}$ intervals on which both
arguments are monotone. On every such interval
Lemma~\ref{lem:allq-calculus}(ii) gives
\[
 V^\rho(a\circ\gamma)\le V^\rho(a),
 \qquad V^\rho(b\circ\widetilde\gamma)\le V^\rho(b).
\]
Consequently the quotient rule, the preceding lower bounds, and the fixed
size of the normalized rectangles give a uniform sectional
$W^{1,\infty}$ bound for every smooth algebraic Jacobian and fixed cutoff
(including cutoffs of the form $\beta(1/q)$).  These smooth factors are
differentiated only inside their fixed support; a sharp half-open cutoff
contributes instead only a bounded number of sectional jumps.
Lemma~\ref{lem:allq-sectional-calculus} therefore
gives both the required sectional $L^\infty+V^\rho$ bound and the separate
row and column $L^1$ bounds, with size
$C_\rho B^{C_{\mathrm{corr},\rho}}$.  The two rough angular factors remain
the input and output $L^2$ functions in
\eqref{eq:rough-band-direct-integral}; they are not part of this amplitude
and are never differentiated.  Thus
\eqref{eq:rough-fixed-angular-operator} is replaced by
\[
 \|T_{\alpha,\omega}\|_{2\to2}
 \le C_\rho B^{C_{\mathrm{corr},\rho}}
 \min\{1,[\Gamma(\omega)(\alpha H)^D]^{-c_{\mathrm{rat},\rho}}\}.
\]
When $n=2$, the transverse sphere is $\mathbb S^0$ and the two
restrictions have coefficient norm comparable to the full coefficient
norm, exactly as in the proof of Lemma~\ref{lem:rough-correlation}. When
$n\ge3$, the chart coefficient lower bound
\eqref{eq:rough-chart-coefficient-lower} and
\eqref{eq:multi-negative}, with
$\sigma=c_{\mathrm{rat},\rho}$, give an angular exponent
\begin{equation}\label{eq:allq-equal-angular-exponent}
 c_{\mathrm{ang},\rho}
 :=\frac{c_{\rm sub}c_{\mathrm{rat},\rho}}
         {c_{\rm sub}+c_{\mathrm{rat},\rho}}>0.
\end{equation}
The $n=2$ estimate is stronger, so after decreasing its exponent the same
$c_{\mathrm{ang},\rho}$ is valid in every dimension $n\ge2$. Hence one
transverse band is bounded by
\[
 C_\rho B^{C_{\mathrm{corr},\rho}}\|\Theta\|_\infty^2
 \alpha^{n-1}
 \min\{1,[\Lambda(\alpha H)^D]^{-c_{\mathrm{ang},\rho}}\}.
\]
For $n\ge2$ choose
\begin{equation}\label{eq:allq-equal-final-exponent}
 c_{\mathrm{eq},\rho}^{(n\ge2)}
 :=\frac12\min\left\{c_{\mathrm{ang},\rho},
                     \frac{n-1}{d_*}\right\}>0.
\end{equation}
The two geometric sums in the proof of
Lemma~\ref{lem:rough-correlation} then give
\eqref{eq:allq-equal-correlation}. In dimension one,
Corollary~\ref{cor:allq-polynomial-vdc} applies on the finitely many signed
overlap intervals and gives an exponent
$c_{\mathrm{eq},\rho}^{(n=1)}>0$. Define, for the fixed ambient dimension,
\[
 c_{\mathrm{eq},\rho}
 :=\begin{cases}
 c_{\mathrm{eq},\rho}^{(n=1)},&n=1,\\
 c_{\mathrm{eq},\rho}^{(n\ge2)},&n\ge2.
 \end{cases}
\]
This number is strictly positive.

For unequal shells, the two zero-extended profile arguments are exactly
\eqref{eq:unequal-profile-pullback-new}.  In the formulas cited in this
paragraph, write the geometric parameter $|W|$ as $r_W$ (so the earlier
$\rho$, $S\rho$, and $\delta=R/(S\rho)$ are read as $r_W$, $Sr_W$, and
$R/(Sr_W)$); this keeps it distinct from the fixed radial-variation
exponent $\rho$ of the present section.  On the support of a transverse
band and both profiles, the calculation following
\eqref{eq:unequal-jacobian-new} gives
\[
 \begin{gathered}
 |u|+|V|\le C,\qquad 0<c\le D=V-\delta u\le C,
 \qquad 0<\delta\le C,\\
 (u^2+\alpha^2)^{1/2}\ge c,
 \qquad (V^2+\alpha^2\delta^2)^{1/2}\ge c.
 \end{gathered}
\]
As in the equal-shell case, insert fixed smooth cutoffs of the two
pullback variables which equal one on the radial profile support and are
supported in a slightly larger compact subinterval of $(0,\infty)$.
Zero extension makes this insertion exact, and all the displayed lower
bounds persist on the cutoff transition regions.
All constants here are independent of the shell radii and of
$\alpha,\delta$.  It follows directly from the product and quotient rules
that the factor
\[
 D^{n-1}(u^2+\alpha^2)^{-n/2}
 (V^2+\alpha^2\delta^2)^{-n/2}
\]
and every smooth rational/chart cutoff have uniformly bounded sectional
$W^{1,\infty}$ norm on each fixed normalized rectangle.  A half-open
rectangle or band boundary produces only one of a bounded number of
sectional jumps.  The two pullbacks have only the nontrivial sectional
critical equation $uV+\delta\alpha^2=0$, so each section requires at most
one additional cut.  Take the common refinement for the two arguments,
the fixed rational rectangles, and the transverse cutoff.
Lemma~\ref{lem:allq-sectional-calculus} then supplies both sectional
$V^\rho$ and the independent Schur row/column bounds. The zero-extended transverse
cutoffs satisfy
$V^\rho(\beta_k)\le V^1(\beta_k)\le C$. Thus the normalized rational
operator contributes the same exponent $c_{\mathrm{rat},\rho}$.
For $n=2$ the transverse sphere again consists of two points; for
$n\ge3$ use the finite rational atlas and
\eqref{eq:multi-negative}. The oriented coefficient lower bound and this
angular argument give, in either case after decreasing the exponent, a
band estimate
\[
 C_\rho B^{C_{\mathrm{corr},\rho}}\|\Theta\|_\infty^2S^{-n}
 \alpha^{n-1}
 \min\{1,(L_R\alpha^d|\mathbf p(W)|)^{-c_{\mathrm{ang},\rho}}\}.
\]
For $n\ge2$ set
\begin{equation}\label{eq:allq-unequal-final-exponent}
 c_{\mathrm{uneq},\rho}^{(n\ge2)}
 :=\frac12\min\left\{c_{\mathrm{ang},\rho},
                     \frac{n-1}{d_*}\right\}>0.
\end{equation}
Splitting the dyadic $\alpha$-sum at
$\alpha\simeq(L_R|\mathbf p(W)|)^{-1/d_*}$ proves
\eqref{eq:allq-unequal-correlation}. The one-dimensional proof uses
Corollary~\ref{cor:allq-polynomial-vdc} and gives an exponent
$c_{\mathrm{uneq},\rho}^{(n=1)}>0$. Define, for the fixed ambient dimension,
\[
 c_{\mathrm{uneq},\rho}
 :=\begin{cases}
 c_{\mathrm{uneq},\rho}^{(n=1)},&n=1,\\
 c_{\mathrm{uneq},\rho}^{(n\ge2)},&n\ge2.
 \end{cases}
\]
This number is strictly positive and completes the proof.
No rough angular factor has been differentiated.
\end{proof}

\begin{corollary}[Quantitative flat high block and collective high shells]
\label{cor:allq-high-shell-stability}
There are $c_{\mathrm{flat},\rho}>0$ and a finite power $C_{B,\rho}$ such
that Proposition~\ref{prop:flat-high-block}, with radial profiles
controlled in $L^\infty+V^\rho$ by $B\ge1$, has the quantitative form
\begin{equation}\label{eq:allq-flat-high-block}
 \left\|A^{P,\Theta}_a\sum_Qe_Q\nu_Q\right\|_2^2
 \le C_\rho B^{C_{B,\rho}}\|\Theta\|_\infty^2\Lambda M
       (L^{-c_{\mathrm{flat},\rho}}+2^{-c_{\mathrm{flat},\rho}t}).
\end{equation}
For every fixed $\kappa>0$ there is
$c_{\kappa,\rho}>0$ such that Proposition~\ref{prop:collective-high} has
the form
\begin{equation}\label{eq:allq-collective-high}
 \left\|\sum_{R\in\mathcal R}T_R
       \sum_{Q\in\mathcal Q_R}e_Q\nu_Q\right\|_2^2
 \le C_{\rho,\kappa}B^{C_{B,\rho}}\|\Theta\|_\infty^2
       \Lambda M2^{-c_{\kappa,\rho}s}.
\end{equation}
Both estimates are uniform under packetwise coefficients $|e_Q|\le1$.
\end{corollary}

\begin{proof}
Let $c_{\mathrm{pair}}>0$ be the exponent in
\eqref{eq:weighted-pair-sublevel}, reduced if necessary so that it also
works in the finitely many rational charts used in the unequal-shell
incidence estimate. In the proof of Proposition~\ref{prop:flat-high-block},
the bad coefficient sublevel has weighted mass
\[
 C\Lambda M(L^{-c_{\mathrm{pair}}/3}+2^{-c_{\mathrm{pair}}t}).
\]
On its complement the actual oscillatory coefficient is at least
$cL^{2/3}$. For $n\ge2$, the near-pair choice
$r=L^{-1/(6d_*)}$ contributes
$C\Lambda M(L^{-n/(6d_*)}+2^{-nt})$, while on the remaining pairs
\eqref{eq:allq-equal-correlation} gives
$C_\rho B^{C_{\mathrm{corr},\rho}}\|\Theta\|_\infty^2
 L^{-c_{\mathrm{eq},\rho}/2}$. The one-dimensional branch omits the
near-pair split. Therefore one may take
\begin{equation}\label{eq:allq-flat-explicit-choice}
 c_{\mathrm{flat},\rho}
 :=\begin{cases}
 \displaystyle\frac12\min\left\{
   \frac{c_{\mathrm{pair}}}{3},
   c_{\mathrm{pair}},
   \frac{c_{\mathrm{eq},\rho}}2,
   \frac{n}{6d_*}\right\},& n\ge2,\\[2ex]
 \displaystyle\frac12\min\left\{
   \frac{c_{\mathrm{pair}}}{3},
   c_{\mathrm{pair}},
   \frac{c_{\mathrm{eq},\rho}}2\right\},& n=1.
 \end{cases}
\end{equation}
Both values are strictly positive. Taking
$C_{B,\rho}$ larger than the finite radial powers in the correlation
estimate proves \eqref{eq:allq-flat-high-block}.

For the collective estimate, the diagonal and comparable-scale terms are
controlled by \eqref{eq:allq-flat-high-block}. For $R\le S/C_0$, retain
the original choice $\tau_R=L_R^{-1/2}$. The one-sided incidence lemma
bounds the bad row by
\[
 C\Lambda\{L_R^{-c_{\mathrm{pair}}/2}
             +(2^{-s}R/S)^{c_{\mathrm{pair}}}\},
\]
whereas \eqref{eq:allq-unequal-correlation} bounds the good row by
$C_\rho B^{C_{B,\rho}}\Lambda L_R^{-c_{\mathrm{uneq},\rho}/2}$.
Put
\[
 c_{\mathrm{row},\rho}
 :=\frac12\min\left\{
 c_{\mathrm{flat},\rho},
 \frac{c_{\mathrm{pair}}}{2},
 \frac{c_{\mathrm{uneq},\rho}}{2}
 \right\}>0.
\]
Since $L_R\ge2^{\kappa s}$ on every high shell, the geometric row sums in
Proposition~\ref{prop:collective-high} are bounded by a constant times
\[
 2^{-\kappa c_{\mathrm{row},\rho}s}
 +2^{-c_{\mathrm{pair}}s}
 +2^{-c_{\mathrm{flat},\rho}s}.
\]
Thus the explicit choice
\begin{equation}\label{eq:allq-collective-explicit-choice}
 c_{\kappa,\rho}
 :=\frac12\min\{\kappa c_{\mathrm{row},\rho},
                  c_{\mathrm{pair}},c_{\mathrm{flat},\rho}\}>0
\end{equation}
absorbs all diagonal, comparable, incidence, and separated-scale terms.
The original Jensen reduction and row orientation use absolute values and
are unchanged, so packetwise signs remain admissible. This proves
\eqref{eq:allq-collective-high}.
\end{proof}

\subsection{Terminal packets with a finite-variation radial profile}

\begin{lemma}[Smooth fixed-sign packet kernel]
\label{lem:allq-smooth-fixed-sign-packet}
Put $N_0=4n+1$ and $\eta=(16N_0)^{-1}$.  In the unit-scale packet
setting of Proposition~\ref{prop:terminalpacket}, let $c_t$ be supported in
one fixed annular enlargement and satisfy
\[
 |c_t|\le1,
 \qquad
 \|\partial_r^\ell c_t\|_\infty\le C_\ell2^{\eta\ell t}
 \quad(0\le\ell\le N_0).
\]
For arbitrary fixed packet coefficients of modulus at most one, let $S_t$
be the convolution output with kernel $K_\Theta(z)c_t(|z|)$.  Then
\begin{equation}\label{eq:allq-smooth-fixed-sign-packet}
 \|S_t\|_2^2
 \le C_n\|\Theta\|_\infty^2
       2^{-3t/32}\Lambda M.
\end{equation}
The assertion is unchanged by the input and output modulations allowed in
Proposition~\ref{prop:terminalpacket}.
\end{lemma}

\begin{proof}
This is the smooth part of the proof of
Proposition~\ref{prop:terminalpacket}, but we record precisely which
hypotheses are used.  Put
$H_t(z)=K_\Theta(z)c_t(|z|)$.  If $n=1$, choose the center $c_U$ of each
packet interval.  Packet cancellation gives
\[
 H_t*d_U(x)=\int_U[H_t(x-y)-H_t(x-c_U)]d_U(y)\,dy.
\]
The annular support stays a fixed distance from the origin, so the stated
radial derivative bounds imply
\[
 \|H_t'\|_1\le C\|\Theta\|_\infty2^{\eta t}.
\]
Since $\operatorname{diam}U\lesssim2^{-t}$, translation continuity in
$W^{1,1}$ and summation over the disjoint packets yield
\[
 \left\|H_t*\sum_Ua_Ud_U\right\|_1
 \le C\|\Theta\|_\infty2^{-(1-\eta)t}M.
\]
The positive thin-layer majorant gives simultaneously
$\|S_t\|_\infty\le C\Lambda\|\Theta\|_\infty$.  Hence
\[
 \|S_t\|_2^2\le\|S_t\|_1\|S_t\|_\infty
 \le C\|\Theta\|_\infty^2\Lambda M
       2^{-(1-\eta)t}.
\]
Here $\eta=1/80$, so this is stronger than
\eqref{eq:allq-smooth-fixed-sign-packet}.  We may therefore assume
$n\ge2$.  The one-shell kernel has normalized radial
$C^{N_0}$ size at most
$C\|\Theta\|_\infty2^{\eta N_0t}$.  Seeger's decomposition with
$\kappa=1/2$ and $\epsilon=1/4$ therefore gives a microlocal part with
squared $L^2$ norm
\[
 C\|\Theta\|_\infty^2 2^{-t/2}\Lambda M
\]
and a remainder with $L^1$ norm
\[
 C\|\Theta\|_\infty
 2^{-(1/4-\eta N_0)t}M
 =C\|\Theta\|_\infty2^{-3t/16}M.
\]
The support and the bound $|c_t|\le1$ give, by the positive thin-layer
majorant, the simultaneous $L^\infty$ bound
$C\Lambda\|\Theta\|_\infty$.  Splitting the remainder at height
$\Lambda\|\Theta\|_\infty2^{-3t/32}$ and interpolating exactly as in
\eqref{eq:fixedsign} proves \eqref{eq:allq-smooth-fixed-sign-packet}.
No property of a step function, or of the set on which an unsmoothing error
is supported, enters these three estimates.  The proof uses absolute
packet coefficients and packet cancellation, so the modulation statements
are preserved.
\end{proof}

\begin{proposition}[Radially weighted selector-free terminal packet]
\label{prop:allq-terminal-packet}
Fix $1\le\rho<\infty$ and put
\begin{equation}\label{eq:allq-terminal-parameters}
 N_0=4n+1,
 \qquad \eta=\frac1{16N_0},
 \qquad \gamma=\frac\eta2,
 \qquad c_{\rm sm}=\frac3{32}.
\end{equation}
There is $t_{0,\rho}\in\mathbb N$ such that
Proposition~\ref{prop:terminalpacket} remains valid for $t\ge t_{0,\rho}$
after inserting a radial factor $a(|x-y|/R)$ satisfying
$\|a\|_\infty+V^\rho(a;[1,2])\le1$. More precisely, one may take
\begin{equation}\label{eq:allq-terminal-explicit-decay}
 \delta_{T,\rho}
 :=\frac14\min\left\{\gamma,c_{\rm sm},\frac\eta\rho\right\}>0
\end{equation}
so that, with the same packet hypotheses and $R\le\varepsilon\le2R$,
\begin{equation}\label{eq:allq-terminal-packet}
 \|V^2(F_\varepsilon:R\le\varepsilon\le2R)\|_2^2
 \le C_{n,\rho}\|\Theta\|_\infty^2
 2^{-\delta_{T,\rho}t}\Lambda M.
\end{equation}
The modulation invariances in Proposition~\ref{prop:terminalpacket}
remain valid.
\end{proposition}

\begin{proof}
Scale to $R=1$ and partition $[1,2]$ into
$K_t\simeq2^{\gamma t}$ consecutive cells $I_\ell$.  If $G_\ell$ denotes
the full shell output over $I_\ell$, inserting all cell endpoints into an
arbitrary variation partition gives the deterministic decomposition
\[
 V^2(F_\varepsilon:1\le\varepsilon\le2)
 \le C V^2\left(\sum_{\ell\le m}G_\ell:0\le m\le K_t\right)
 +C\left(\sum_\ell V^2(F_\varepsilon:\varepsilon\in I_\ell)^2\right)^{1/2}.
\]
The within-cell argument uses
$V^2\le V^1$, the positive thin-layer majorant, and only
$\|a\|_\infty\le1$. It therefore gives
\begin{equation}\label{eq:allq-terminal-within-cell}
 \left\|\left(\sum_\ell
 V^2(F_\varepsilon:\varepsilon\in I_\ell)^2\right)^{1/2}\right\|_2^2
 \le C\|\Theta\|_\infty^2\Lambda M2^{-\gamma t}.
\end{equation}

For the grid partial sums fix signs
$\epsilon_\ell\in\{-1,0,1\}$ and put
\[
 c(r)=a(r)\sum_{\ell=1}^{K_t}\epsilon_\ell\mathbf1_{I_\ell}(r).
\]
Lemma~\ref{lem:allq-calculus} gives
\begin{equation}\label{eq:allq-signed-profile}
 \|c\|_\infty+V^\rho(c)
 \le C_\rho(1+K_t^{1/\rho}).
\end{equation}
Extend $c$ by zero and mollify at scale $h=2^{-\eta t}$.
Lemma~\ref{lem:allq-translation} yields
\begin{equation}\label{eq:allq-terminal-smooth-error-L1}
 \|c-c_h\|_1
 \le C_\rho h^{1/\rho}(1+K_t^{1/\rho})
 \le C_\rho2^{-(\eta-\gamma)t/\rho}
 =C_\rho2^{-\eta t/(2\rho)}.
\end{equation}
Moreover
$\|c-c_h\|_\infty\le2$ and
$V^\rho(c-c_h)\le C_\rho(1+2^{\gamma t/\rho})$.
With packet side $2^{-t}$,
Lemma~\ref{lem:allq-radial-discretization} gives for the unsmoothing
output $D$
\begin{align}
 \|D\|_\infty
 &\le C_\rho\Lambda\|\Theta\|_\infty
 2^{-\eta t/(2\rho)},
 \label{eq:allq-terminal-unsmooth-Linf}\\
 \|D\|_1
 &\le C_\rho\|\Theta\|_1M
 2^{-\eta t/(2\rho)}.
 \label{eq:allq-terminal-unsmooth-L1}
\end{align}
Indeed the additional discretization term is
$2^{-t/\rho}(1+2^{\gamma t/\rho})$, whose exponent is no smaller than
$\eta/(2\rho)$ because $\gamma=\eta/2<1$. Since $\|\Theta\|_1\le C_n\|\Theta\|_\infty$, it follows that
\begin{equation}\label{eq:allq-terminal-unsmooth-L2}
 \|D\|_2^2
 \le C_{n,\rho}\|\Theta\|_\infty^2\Lambda M
 2^{-\eta t/\rho}.
\end{equation}

It remains to estimate the smoothed signed sum. Since $|c|\le1$,
$|c_h|\le1$, and
$\|\partial_r^\ell c_h\|_\infty\le C_\ell h^{-\ell}$, the smoothed kernel
satisfies Lemma~\ref{lem:allq-smooth-fixed-sign-packet}.  Hence
\begin{equation}\label{eq:allq-terminal-smoothed-L2}
 \|S\|_2^2
 \le C\|\Theta\|_\infty^2\Lambda M2^{-c_{\rm sm}t}.
\end{equation}
For $n=1$, the elementary cancellation argument in the proof of
Proposition~\ref{prop:terminalpacket} gives the stronger exponent
$1-\eta$; thus the same common choice $c_{\rm sm}=3/32$ is valid in every
dimension.

Lemma~\ref{lem:RM} costs at most $C(1+t)^2$ after squaring. The elementary
bound
$(1+t)^2 2^{-at}\le C_a2^{-at/2}$ for $a>0$, together with
\eqref{eq:allq-terminal-within-cell},
\eqref{eq:allq-terminal-unsmooth-L2}, and
\eqref{eq:allq-terminal-smoothed-L2}, proves
\eqref{eq:allq-terminal-packet} with the smaller exponent in
\eqref{eq:allq-terminal-explicit-decay}. Increase $t_{0,\rho}$ once so that the support enlargements caused by
mollification remain in the fixed radial annulus and so that the finitely
many polynomial factors in $t$ are absorbed by the displayed geometric
decays. The proof uses only absolute packet coefficients and packetwise
cancellation, so the stated input and output modulation invariances are
preserved.
\end{proof}

\subsection{The unphased strong theorem for every finite radial exponent}

We next prove the terminal strong input in a form with the exact angular
conjugate-exponent loss needed for the $L\log L$ endpoint.

For $j\in\mathbb Z$ define the annular measure
\begin{equation}\label{eq:allq-shell-measure}
 d\nu_j(x)=
 \frac{\Omega(x/|x|)}{|x|^n}b(|x|)
 \mathbf1_{\{2^j<|x|\le2^{j+1}\}}\,dx.
\end{equation}
The mean-zero condition on $\Omega$ implies $\nu_j(\R^n)=0$.

\begin{lemma}[Shell Fourier envelope]
\label{lem:allq-shell-Fourier}
Assume \eqref{eq:allq-normalize-b} and
$\int_{\Sph^{n-1}}\Omega\,d\sigma=0$.  There exists
$\beta_\rho>0$ such that $\widehat\nu_j(0)=0$ and, for
$t=2^j|\xi|>0$,
\begin{align}
 |\widehat\nu_j(\xi)|
 &\le C t\|\Omega\|_1,
 &&0<t\le1,
 \label{eq:allq-shell-low}\\
 |\widehat\nu_j(\xi)|
 &\le C_\rho t^{-\beta_\rho}\|\Omega\|_2,
 &&t\ge1.
 \label{eq:allq-shell-high-L2}
\end{align}
One may take
\[
 \beta_\rho=\min\left\{\frac14,\frac1{2\rho}\right\}.
\]
Consequently, if $1<s\le2$, $s'=s/(s-1)$, and $\Omega$ has mean zero,
then, for $t>0$,
\begin{equation}\label{eq:allq-shell-Ls}
 |\widehat\nu_j(\xi)|
 \le C_\rho\|\Omega\|_s\,a_{s,\rho}(\log_2t),
\end{equation}
where $a_{s,\rho}:\mathbb R\to(0,\infty)$ is defined by
\begin{equation}\label{eq:allq-envelope}
 a_{s,\rho}(r)=
 \begin{cases}
 C2^r,&r\le0,\\
 C_\rho2^{-2\beta_\rho r/s'},&r>0.
 \end{cases}
\end{equation}
For every $\delta>0$,
\begin{equation}\label{eq:allq-envelope-sum}
 \sum_{r\in\mathbb Z}a_{s,\rho}(r)^\delta
 \le C_{\rho,\delta}s'.
\end{equation}
\end{lemma}

\begin{proof}
For $\xi\ne0$, set $e=\xi/|\xi|$ and make the change of variables
$r=2^ju$.  Then the multiplier is
\[
 \widehat\nu_j(\xi)=
 \int_{\Sph^{n-1}}\Omega(\theta)
 \int_1^2b(2^ju)e^{-itu e\cdot\theta}\,\frac{du}{u}\,d\sigma(\theta).
\]
At frequency zero the same formula before introducing $e$ has a radial
integral independent of $\theta$, so mean zero of $\Omega$ gives
$\widehat\nu_j(0)=0$.  For $\xi\ne0$,
subtracting that zero-frequency value and using
$|e^{-iz}-1|\le|z|$ proves \eqref{eq:allq-shell-low}.  Thus no value of
$\log_2t$ is invoked when $t=0$.

The amplitude $u\mapsto b(2^ju)/u$ has uniformly bounded
$L^\infty+V^\rho$ norm by Lemma~\ref{lem:allq-calculus}.  Applying
Lemma~\ref{lem:allq-first-derivative} to the linear radial phase gives
\[
 \left|\int_1^2b(2^ju)e^{-itu\tau}\frac{du}{u}\right|
 \le C_\rho\min\{1,(t|\tau|)^{-1/\rho}\}.
\]
For $n\ge2$, Cauchy--Schwarz in $\theta$ gives the square root of
\[
 \int_{\Sph^{n-1}}
 \min\{1,(t|e\cdot\theta|)^{-2/\rho}\}\,d\sigma(\theta).
\]
The distribution of $|e\cdot\theta|$ near zero has bounded density.  Thus
this integral is $O(t^{-1})$ for $\rho<2$, is
$O(t^{-1}\log(2+t))$ for $\rho=2$, and is $O_\rho(t^{-2/\rho})$ for
$\rho>2$.  The weaker uniform bound
$C_\rho t^{-2\beta_\rho}$ follows.  When $n=1$, there is no equatorial
integration and the direct bound is even stronger.  This proves
\eqref{eq:allq-shell-high-L2}.

For high frequencies interpolate the linear functional of $\Omega$ between
the trivial $L^1$ bound and \eqref{eq:allq-shell-high-L2}.  Since the
$L^2$ endpoint has interpolation weight $2/s'$, one obtains the second line
of \eqref{eq:allq-envelope}; the low-frequency line follows from
\eqref{eq:allq-shell-low} and $L^s\hookrightarrow L^1$ on the sphere.
Finally \eqref{eq:allq-envelope-sum} is the sum of two geometric series; the
positive-frequency series has ratio
$2^{-2\beta_\rho\delta/s'}$ and hence costs $O_{\rho,\delta}(s')$.
\end{proof}

We shall also use the following standard vector-valued shell estimate.  A
proof is included because it is uniform down to angular $L^1$ and therefore
is exactly what allows spatial interpolation below.

\begin{lemma}[Uniform vector-valued annular estimate]
\label{lem:allq-vector-shell}
For $1<p<\infty$ and arbitrary functions $f_j$,
\begin{equation}\label{eq:allq-vector-shell}
 \left\|\left(\sum_j|\nu_j*f_j|^2\right)^{1/2}\right\|_p
 \le C_p\|b\|_\infty\|\Omega\|_1
 \left\|\left(\sum_j|f_j|^2\right)^{1/2}\right\|_p.
\end{equation}
\end{lemma}

\begin{proof}
For a direction $\theta$, let $M_\theta$ denote the one-dimensional
Hardy--Littlewood maximal operator along lines parallel to $\theta$.  Polar
coordinates give pointwise
\[
 |\nu_j*f_j(x)|
 \le C\|b\|_\infty\int_{\Sph^{n-1}}|\Omega(\theta)|
 M_\theta f_j(x)\,d\sigma(\theta).
\]
Minkowski first in $\ell^2$ and then in $\theta$, followed by the
vector-valued one-dimensional maximal inequality on almost every line,
gives
\[
 \left\|\left(\sum_j|\nu_j*f_j|^2\right)^{1/2}\right\|_p
 \le C_p\|b\|_\infty\|\Omega\|_1
 \left\|\left(\sum_j|f_j|^2\right)^{1/2}\right\|_p.
\]
The constant is uniform in $\theta$ by rotation.
\end{proof}

Fix a smooth Littlewood--Paley function $\psi$ supported in
$\{1/2\le|\xi|\le2\}$ and satisfying
$\sum_{\ell\in\mathbb Z}\psi(2^{-\ell}\xi)^2=1$ for $\xi\ne0$.
Let $\Delta_\ell$ be the corresponding multiplier.  For a relative frequency
$r\in\mathbb Z$ put
\[
 g_{j,r}:=\nu_j*\Delta_{r-j}^2f.
\]
On the Fourier support of $g_{j,r}$ one has $2^j|\xi|\simeq2^r$.

\begin{lemma}[Frequency-separated tail representation]
\label{lem:allq-separated-tail}
For every $1<p<\infty$ and every $f\in L^p(\R^n)$, there is a fixed
integer $M\ge5$ such that, after splitting the shell
indices $j$ into their $M$ residue classes, the tail
\[
 \sum_{\substack{j\ge k\\j\equiv c\ ({\rm mod}\ M)}}g_{j,r}
\]
is, up to repetitions as $k$ varies, exactly a subsequence of the smooth
low-pass family $\{\phi_\ell*h_{r,c}\}_{\ell\in\mathbb Z}$ used in
Proposition~2.1 of Chen--Ding--Hong--Liu, applied to
\[
 h_{r,c}:=\sum_{j\equiv c\ ({\rm mod}\ M)}g_{j,r}.
\]
Here $\widehat\phi=1$ on $|\xi|\le2$ and
$\widehat\phi=0$ on $|\xi|>4$, and
$\widehat{\phi_\ell}(\xi)=\widehat\phi(2^\ell\xi)$.
Consequently, Proposition~2.1 of \cite{ChenDingHongLiu} gives
\begin{equation}\label{eq:allq-tail-jump-by-h}
 \left\|\mathbf J\left(
 \sum_{\substack{j\ge k\\j\equiv c\ ({\rm mod}\ M)}}g_{j,r}:k\in\mathbb Z
 \right)\right\|_p
 \le C_p\|h_{r,c}\|_p.
\end{equation}
\end{lemma}

\begin{proof}
The Fourier support of $g_{j,r}$ lies in
\[
 2^{r-j-1}\le|\xi|\le2^{r-j+1}.
\]
Fix $M\ge5$.  If $j_0\equiv c\pmod M$, then
$\phi_{j_0-r}$ equals one on every support with $j\ge j_0$: indeed the
largest such frequency is $2^{r-j_0+1}$, precisely the inner plateau of
$\widehat\phi(2^{j_0-r}\xi)$.  On the other hand, the preceding index in
the same residue class is $j_0-M$, whose support begins at
$2^{r-j_0+M-1}>2^{r-j_0+2}$; hence
$\phi_{j_0-r}$ vanishes on it and on all earlier supports.  Therefore
\[
 \phi_{j_0-r}*h_{r,c}
 =\sum_{\substack{j\ge j_0\\j\equiv c\ ({\rm mod}\ M)}}g_{j,r}
\]
exactly.  When $k$ moves between two consecutive indices in the residue
class, the tail is merely repeated.  Thus the residue-class tail is a
subsequence, with repetitions, of the very family covered by
Proposition~2.1 of \cite{ChenDingHongLiu}; repetitions cannot create jumps.
The jump count in that proposition is the paired count used in the
definition of $\mathbf J$, so it proves \eqref{eq:allq-tail-jump-by-h}
directly.  The first inequality in
\eqref{eq:chain-pair-functional} also gives the corresponding chain bound
with only an absolute loss; this latter form is used below solely to enter
the normed quotient of Lemma~\ref{lem:jumpnorm}.
\end{proof}

\begin{lemma}[$s'$-tracked dyadic terminal jump]
\label{lem:allq-terminal-dyadic}
Let $1<s\le2$, assume $\Omega\in L^s(\Sph^{n-1})$ has mean zero, and fix
$1<p<\infty$.  Then, for every $f\in L^p(\R^n)$, under
\eqref{eq:allq-normalize-b},
\begin{equation}\label{eq:allq-terminal-dyadic}
 \left\|\mathbf J(T^{0,b}_{\Omega,2^k}f:k\in\mathbb Z)\right\|_p
 \le C_{n,p,\rho}s'\|\Omega\|_s\|f\|_p.
\end{equation}
\end{lemma}

\begin{proof}
First impose finite shell and frequency cutoffs.  For one residue class,
Littlewood--Paley theory, Lemma~\ref{lem:allq-vector-shell}, and bounded
frequency overlap give, for every $1<p_0<\infty$,
\begin{equation}\label{eq:allq-h-uniform}
 \|h_{r,c}\|_{p_0}\le C_{p_0}\|\Omega\|_1\|f\|_{p_0}
 \le C_{p_0}\|\Omega\|_s\|f\|_{p_0}.
\end{equation}
On $L^2$, the separated Fourier supports and
Lemma~\ref{lem:allq-shell-Fourier} give the following bound.  Indeed, on
$\supp\widehat{\Delta_{r-j}f}$ one has
$|\log_2(2^j|\xi|)-r|\le1$.  The two exponential slopes in
\eqref{eq:allq-envelope} are uniformly bounded for $1<s\le2$; including
the possible crossing of $r=0$, this gives
$a_{s,\rho}(\log_2(2^j|\xi|))\le C_\rho a_{s,\rho}(r)$ with a constant
independent of $s$.  Hence
\begin{equation}\label{eq:allq-h-L2}
 \|h_{r,c}\|_2
 \le C_\rho a_{s,\rho}(r)\|\Omega\|_s\|f\|_2.
\end{equation}
If $p<2$, interpolate \eqref{eq:allq-h-L2} with
\eqref{eq:allq-h-uniform} at one fixed $p_0<p$; if $p>2$, use one fixed
$p_0>p$.  For $p=2$ use \eqref{eq:allq-h-L2}.  Thus for some
$\theta_p>0$ independent of $s,r$,
\begin{equation}\label{eq:allq-h-decay}
 \|h_{r,c}\|_p
 \le C_{p,\rho}a_{s,\rho}(r)^{\theta_p}
 \|\Omega\|_s\|f\|_p.
\end{equation}
Lemma~\ref{lem:allq-separated-tail} gives the same estimate for the paired
jump of the corresponding residue-class tail.  By
$\mathbf J^{\rm ch}\le\sqrt2\,\mathbf J$ from
\eqref{eq:chain-pair-functional}, and then the uniform equivalence in
Lemma~\ref{lem:jumpnorm}, the same bound holds in the norm
$\|\cdot\|_{\mathfrak J_D}$ for every finite ordered parameter set $D$.
Summing the finitely many residue classes changes only the constant.

We now remove the cutoffs rather than merely invoking lower
semicontinuity.  Fix a finite ordered set $D\subset\mathbb Z$ of dyadic
truncation indices and use the norm $\|\cdot\|_{\mathfrak J_D}$ supplied by
Lemma~\ref{lem:jumpnorm}.  The preceding paired-to-chain conversion is
what permits this normed summation.  If $R_2>R_1$, its triangle
inequality and \eqref{eq:allq-h-decay} give, uniformly in the finite shell
cutoff,
\[
 \left\|\sum_{R_1<|r|\le R_2}
  \biggl[\sum_{j\ge k}g_{j,r}\biggr]_{k\in D}
 \right\|_{L^p(\mathfrak J_D)}
 \le C_{p,\rho}\|\Omega\|_s\|f\|_p
 \sum_{R_1<|r|\le R_2}a_{s,\rho}(r)^{\theta_p}.
\]
By \eqref{eq:allq-envelope-sum} the right side tends to zero as
$R_1\to\infty$, so the frequency sums are Cauchy in
$L^p(\mathfrak J_D)$.  For a fixed shell $j$, the Littlewood--Paley
identity $\sum_r\Delta_{r-j}^2f=f$ holds in $L^p$; convolution with the
finite-annulus measure $\nu_j$ identifies every finite-shell coordinate of
the limit.

Next remove the shell cutoff with the order of the two tails made
explicit.  Let $\mathcal J_J\uparrow\mathbb Z$ be the growing finite shell
sets and put $E_J=\mathbb Z\setminus\mathcal J_J$.  Given $\epsilon>0$,
first choose $R$ so large that
\[
 C_{p,\rho}\|\Omega\|_s\|f\|_p
 \sum_{|r|>R}a_{s,\rho}(r)^{\theta_p}<\frac\epsilon2;
\]
this estimate is uniform in the shell cutoff.  For the finitely many
$|r|\le R$ and residue classes $c$, put
\[
 h_{r,c}^{E_J}=\sum_{\substack{j\in E_J\\j\equiv c\ ({\rm mod}\ M)}}
      \nu_j*\Delta_{r-j}^2f.
\]
The separated-tail estimate applies to an arbitrary subset
$E\subset\{j:j\equiv c\ ({\rm mod}\ M)\}$, not only to the full residue
class: the low-pass multiplier used in its proof still gives, term by term,
\[
 \phi_{j_0-r}*h_{r,c}^{E}
 =\sum_{\substack{j\in E,\ j\ge j_0\\j\equiv c\ ({\rm mod}\ M)}}g_{j,r}.
\]
Thus the same CDHL bound, with the same constant, holds for every such
$E$.
Interpret this sum first over finite subsets of $E_J$.  The following
uniform estimate makes those partial sums Cauchy in $L^p$ and hence
defines $h_{r,c}^{E_J}$; the finite-parameter tail estimate then passes to
that limit through the norm $\mathfrak J_D$.
The proof of \eqref{eq:allq-h-uniform}, restricted to $E_J$, gives
\[
 \|h_{r,c}^{E_J}\|_p
 \le C_p\|\Omega\|_1
 \left\|\left(\sum_{j\in E_J}
       |\Delta_{r-j}^2f|^2\right)^{1/2}\right\|_p
 \longrightarrow0.
\]
Indeed, for fixed $r$ the map $j\mapsto r-j$ is a bijection of
$\mathbb Z$, and
\[
 \left(\sum_{j\in E_J}|\Delta_{r-j}^2f|^2\right)^{1/2}
 \longrightarrow0\quad\hbox{pointwise},\qquad
 \left(\sum_{j\in E_J}|\Delta_{r-j}^2f|^2\right)^{1/2}
 \le \left(\sum_{\ell\in\mathbb Z}|\Delta_\ell^2f|^2\right)^{1/2}.
\]
The dominating full Littlewood--Paley square function belongs to $L^p$;
dominated convergence proves the asserted norm limit.  This also shows
that the limit is independent of the particular exhaustion
$\mathcal J_J$.
Choose $J$ so that the sum of these finitely many norms is less than
$\epsilon/2$.  Lemma~\ref{lem:allq-separated-tail}, the triangle
inequality in $L^p(\mathfrak J_D)$, and the preceding frequency-tail
choice then show that the full omitted contribution has
$L^p(\mathfrak J_D)$ norm below $\epsilon$.  This proves joint Cauchy
convergence, independently of the order in which the two cutoffs are
removed.

To identify the quotient limit without losing its constant coordinate,
fix $k_0\in D$ and choose the representative whose $k_0$-coordinate is
the ordinary $L^p$ limit supplied by the finite-annulus shell series and
the Littlewood--Paley identity.  Finite-dimensional quotient convergence
controls every difference between a $k$-coordinate and the
$k_0$-coordinate; the same coordinatewise Littlewood--Paley identity
therefore identifies all coordinates with
$T^{0,b}_{\Omega,2^k}f$.  Convergence in $L^p(\mathfrak J_D)$ gives, after
a subsequence, convergence in the finite-dimensional quotient almost
everywhere.  The equivalence in Lemma~\ref{lem:jumpnorm} first returns the
resulting bound to $\mathbf J^{\rm ch}_D$; the inequality
$\mathbf J_D\le2\mathbf J^{\rm ch}_D$ in
\eqref{eq:chain-pair-functional}, or equivalently finite strict-witness
lower semicontinuity, then gives the paired estimate
\[
 \|\mathbf J(T^{0,b}_{\Omega,2^k}f:k\in D)\|_p
 \le C_{p,\rho}\sum_r a_{s,\rho}(r)^{\theta_p}
 \|\Omega\|_s\|f\|_p
 \le C_{p,\rho}s'\|\Omega\|_s\|f\|_p.
\]
Finally exhaust $D$.  To reiterate the interface: Proposition~2.1 of
\cite{ChenDingHongLiu} is stated for the paired jump count and therefore
enters without conversion.  The paired-to-chain comparison is used only
to enter the MSZK normed quotient for summation, and the chain-to-paired
comparison is used on exit.  Both losses are absolute and independent of
$D$.
\end{proof}

\begin{lemma}[$s'$-tracked short terminal variation]
\label{lem:allq-terminal-short}
Under the hypotheses of Lemma~\ref{lem:allq-terminal-dyadic}, and for
every $f\in L^p(\R^n)$, let
$S_2(T^{0,b}_\Omega f)$ denote the square sum of the within-dyadic-shell
$2$-variations. Then
\begin{equation}\label{eq:allq-terminal-short}
 \|S_2(T^{0,b}_\Omega f)\|_p
 \le C_{n,p,\rho}s'\|\Omega\|_s\|f\|_p.
\end{equation}
\end{lemma}

\begin{proof}
For $1\le u\le2$, let
\[
 A_{j,u}f(x)=
 \int_{u2^j<|z|\le2^{j+1}}
 b(|z|)K_\Omega(z)f(x-z)\,dz
\]
and put
\[
 S_r f=
 \left(\sum_jV^2(A_{j,u}\Delta_{r-j}^2f:1\le u\le2)^2\right)^{1/2}.
\]
The path $u\mapsto A_{j,u}f(x)$ is locally absolutely continuous for
almost every $x$. Since $V^2\le V^1$, differentiation of the moving
boundary and integration in $u$ give
\begin{align*}
 V^2(A_{j,u}\Delta_{r-j}^2f:u)
 &\le \int_1^2\left|\frac d{du}
       A_{j,u}\Delta_{r-j}^2f\right|\,du\\
 &\le C\|b\|_\infty\int_{\Sph^{n-1}}|\Omega(\theta)|
    \left[2^{-j}\int_{2^j}^{2^{j+1}}
       |\Delta_{r-j}^2f(x-t\theta)|\,dt\right]d\sigma(\theta)\\
 &\le C\|b\|_\infty\int_{\Sph^{n-1}}|\Omega(\theta)|
       M_\theta(\Delta_{r-j}^2f)(x)\,d\sigma(\theta).
\end{align*}
The vector-valued directional maximal inequality used in
Lemma~\ref{lem:allq-vector-shell} therefore gives, uniformly in $r$ and
for every $1<p_0<\infty$,
\begin{equation}\label{eq:allq-short-uniform}
 \|S_rf\|_{p_0}\le C_{p_0}\|\Omega\|_1\|f\|_{p_0}
 \le C_{p_0}\|\Omega\|_s\|f\|_{p_0}.
\end{equation}

We now justify the $L^2$ frequency envelope uniformly in the moving
endpoint $u$. On the relative frequency $t=2^j|\xi|\simeq2^r$, the
partial-shell multiplier is
\[
 m_{j,u}(\xi)=
 \int_{\Sph^{n-1}}\Omega(\theta)
 \int_u^2 b(2^jv)e^{-itv e\cdot\theta}\,\frac{dv}{v}\,d\sigma(\theta),
 \qquad e=\xi/|\xi|.
\]
The zero-frequency value is independent of $\theta$, so mean zero of
$\Omega$ gives the low-frequency bound $C t\|\Omega\|_1$, uniformly in
$u$. For high frequency, the zero-extended amplitude
\[
 v\longmapsto b(2^jv)v^{-1}\mathbf1_{[u,2]}(v)
\]
has uniformly bounded $L^\infty+V^\rho$ norm by
Lemma~\ref{lem:allq-calculus}(iii)--(iv). Lemma~\ref{lem:allq-first-derivative}, followed by the same equatorial
integration and $L^1$--$L^2$ interpolation as in
Lemma~\ref{lem:allq-shell-Fourier}, therefore yields
\begin{equation}\label{eq:allq-partial-shell-envelope}
 \sup_{1\le u\le2}|m_{j,u}(\xi)|
 \le C_\rho\|\Omega\|_s a_{s,\rho}(r).
\end{equation}
Moreover,
\[
 \frac d{du}m_{j,u}(\xi)
 =-u^{-1}\int_{\Sph^{n-1}}\Omega(\theta)b(2^ju)
 e^{-itu e\cdot\theta}\,d\sigma(\theta),
\]
so the derivative operator has $L^2$ norm at most
$C\|b\|_\infty\|\Omega\|_1$. Apply the anchored
Sobolev--variation argument of Lemma~\ref{lem:shortannular}, with the
operator norm in \eqref{eq:allq-partial-shell-envelope} playing the role
of $a$ and the derivative bound playing the role of $b$. This gives
\[
 \|V^2(A_{j,u}\Delta_{r-j}^2f:u)\|_2
 \le C_\rho a_{s,\rho}(r)^{1/4}\|\Omega\|_s
 \|\Delta_{r-j}^2f\|_2.
\]
After squaring and summing in $j$,
\begin{equation}\label{eq:allq-short-L2}
 \|S_rf\|_2
 \le C_\rho a_{s,\rho}(r)^{1/4}\|\Omega\|_s\|f\|_2.
\end{equation}
Interpolate the sublinear operator $S_r$ between
\eqref{eq:allq-short-uniform} and \eqref{eq:allq-short-L2}. For the target
$p$ this gives an exponent $\vartheta_p>0$, independent of $s,r$:
\[
 \|S_rf\|_p
 \le C_{p,\rho}a_{s,\rho}(r)^{\vartheta_p}
 \|\Omega\|_s\|f\|_p.
\]
For finite $j$- and $r$-sets, the variation triangle inequality and
Minkowski in the outer $\ell^2$ sum give
\[
 \left(\sum_jV^2\left(\sum_{r\in E}
 A_{j,u}\Delta_{r-j}^2f:u\right)^2\right)^{1/2}
 \le\sum_{r\in E}S_rf.
\]
Consequently \eqref{eq:allq-envelope-sum}, with
$\delta=\vartheta_p$, makes the $r$-sum Cauchy in $L^p$, uniformly over
finite shell sets and finite $u$-partitions.  We again record the
two-tail argument.  Given $\epsilon>0$, choose $R$ so that
\[
 \sum_{|r|>R}\|S_rf\|_p<\frac\epsilon2.
\]
For the finitely many $|r|\le R$, let $\mathcal J_J\uparrow\mathbb Z$ be
the retained shell sets, put $E_J=\mathbb Z\setminus\mathcal J_J$, and
restrict the outer square sum defining $S_r$ to $E_J$.  The
vector-valued estimates above, with the same restriction, bound this
quantity by the corresponding Littlewood--Paley square-function tail;
hence
\[
 \left\|\left(\sum_{j\in E_J}
 V^2(A_{j,u}\Delta_{r-j}^2f:u\in\mathcal U)^2
 \right)^{1/2}\right\|_p\longrightarrow0
\]
for every finite $\mathcal U\subset[1,2]$, uniformly in the choice of its
ordered partition.  Choose $J$ so that the sum of these finitely many
norms is below $\epsilon/2$.  Minkowski in the outer $\ell^2$ sum now
shows that the combined frequency--shell approximation is Cauchy in the
finite-path variation norm.
Here the asserted tail convergence is not formal: for each fixed $r$,
$j\mapsto r-j$ is a bijection and the restricted square function is
pointwise dominated by
$\bigl(\sum_\ell|\Delta_\ell^2f|^2\bigr)^{1/2}\in L^p$, while it tends
pointwise to zero.  Dominated convergence therefore removes the shell
cutoff uniformly over the already fixed finite set of $r$'s and finite
$u$-partitions.

For each fixed $j,u$, the Littlewood--Paley identity and the
fixed-annulus $L^1$ kernel identify its coordinate limit with $A_{j,u}f$.
First exhaust the shell sets and then the finite rational
$u$-partitions.  Finite-partition lower semicontinuity and Fatou, followed
by local absolute continuity of $u\mapsto A_{j,u}f(x)$, give
$S_2(T^{0,b}_\Omega f)\le\sum_rS_rf$ in the limiting sense and complete
the proof.
\end{proof}

\begin{proposition}[Terminal $s'$-tracked critical jump]
\label{prop:allq-terminal-sp}
For $1<s\le2$, mean-zero $\Omega\in L^s(\Sph^{n-1})$, and
$1<p<\infty$, the following holds for every $f\in L^p(\R^n)$:
\begin{equation}\label{eq:allq-terminal-sp}
 \|\mathbf J(T^{0,b}_{\Omega,\varepsilon}f:\varepsilon>0)\|_p
 \le C_{n,p,\rho}s'\|\Omega\|_s\|f\|_p
\end{equation}
under normalization \eqref{eq:allq-normalize-b}.
\end{proposition}

\begin{proof}
The deterministic long--short inequality, equivalently (1.13) in
\cite{ChenDingHongLiu}, gives pointwise
\[
 \mathbf J(T^{0,b}_{\Omega,\varepsilon}f)
 \le C\left[
 S_2(T^{0,b}_\Omega f)+
 \mathbf J(T^{0,b}_{\Omega,2^k}f:k\in\mathbb Z)
 \right].
\]
Apply Lemmas~\ref{lem:allq-terminal-dyadic}
and~\ref{lem:allq-terminal-short}.
\end{proof}

\begin{proposition}[Unphased terminal $L\log L$ theorem]
\label{prop:allq-terminal-LlogL}
For mean-zero $\Omega\in L\log L(\Sph^{n-1})$, $1<p<\infty$, and every
$f\in L^p(\R^n)$,
\begin{equation}\label{eq:allq-terminal-LlogL}
 \|\mathbf J(T^{0,b}_{\Omega,\varepsilon}f:\varepsilon>0)\|_p
 \le C_{n,p,\rho}\mathfrak L(\Omega)\|f\|_p
\end{equation}
under \eqref{eq:allq-normalize-b}.
\end{proposition}

\begin{proof}
We use exactly the finite-trajectory Orlicz extrapolation used in the proof
of Proposition~\ref{prop:strong-far-endpoint}, but record the cancellation
issue.
Let
\[
 \Pi\Theta=\Theta-\sigma(\Sph^{n-1})^{-1}\int_{\Sph^{n-1}}\Theta\,d\sigma.
\]
For a finite ordered truncation set $D$, let $E_D$ be the scalar
finite-path quotient endowed with a norm uniformly equivalent to the
pointwise critical jump functional.  Define the linear map
\[
 \mathcal A_{f,D}\Theta
 =\bigl[T^{0,b}_{\Pi\Theta,\varepsilon}f\bigr]_{\varepsilon\in D}
 \in L^p(\R^n;E_D).
\]
Since $\Pi$ is bounded on every $L^s$ on the finite-measure sphere,
Proposition~\ref{prop:allq-terminal-sp} gives uniformly in $D$
\[
 \|\mathcal A_{f,D}\|_{L^s(\Sph^{n-1})\to L^p(E_D)}
 \le C_{p,\rho}s'\|f\|_p,
 \qquad1<s\le2.
\]
Pair with a norm-one functional in $[L^p(E_D)]^*$ and take
$s=k/(k-1)$ for integers $k\ge2$.  For each $k$ the resulting scalar
functional on the sphere is represented by $h_k\in L^k$ satisfying
$\|h_k\|_k\le C_{p,\rho}k\|f\|_p$.  These representatives agree almost
everywhere: two of them represent the same functional on bounded angular
functions, their difference lies in $L^1$ on the finite-measure sphere,
and testing against its bounded complex sign makes that difference zero.
Denote the common representative by $h$.  If $\|f\|_p>0$, Tonelli and the
exponential series give, for a sufficiently large structural $L$,
\[
 \int_{\Sph^{n-1}}\Psi_0\left(\frac{|h|}{L\|f\|_p}\right)d\sigma
 \le\sum_{k=2}^\infty\frac{(C_{p,\rho}k/L)^k}{k!}\le1.
\]
The case $f=0$ is immediate.  Hence
\[
 \|h\|_{\exp L}\le C_{p,\rho}\|f\|_p.
\]
Orlicz--H\"older therefore yields
\[
 \|\mathcal A_{f,D}\Omega\|_{L^p(E_D)}
 \le C_{p,\rho}\mathfrak L(\Omega)\|f\|_p
\]
for mean-zero $\Omega$, uniformly in $D$.  Increase finite rational sets and
use finite-annulus local absolute continuity to recover all positive
parameters.  This proves \eqref{eq:allq-terminal-LlogL}.
\end{proof}

\subsection{Strong polynomial phases}

\begin{proposition}[Radially weighted annular decay]
\label{prop:allq-annular-decay}
Let $P$ be a real polynomial on $\R^n\times\R^n$ of total degree at most
$d$, and let $\Theta\in L^\infty(\Sph^{n-1})$.  For $R>0$ and
$1/2\le u\le1$ put
\[
 A^{P,\Theta,b}_{R,u}f(x)=
 \int_{uR<|x-y|<R}e^{iP(x,y)}b(|x-y|)
 K_\Theta(x-y)f(y)\,dy.
\]
Under \eqref{eq:allq-normalize-b}, for every nonterminal flag position
$\mu$ there are $a_{\mu,\rho}>0$ and $C_{n,d,\rho}$ such that, if the
first nonzero flag coordinate of $P$ is $\pm1$ at position $\mu$, then
\begin{equation}\label{eq:allq-annular-decay}
 \sup_{1/2\le u\le1}\|A^{P,\Theta,b}_{R,u}\|_{2\to2}
 \le C_{n,d,\rho}R^{-a_{\mu,\rho}}\|\Theta\|_\infty,
 \qquad R\ge1.
\end{equation}
Every later phase coefficient remains unrestricted.
\end{proposition}

\begin{proof}
We revisit the three branches in Proposition~\ref{prop:annulardecay}.

For a transverse pivot, the kernel entering $A_Q^*A_Q$ is obtained from
the original $g_h(v)$ by multiplication by
$\overline{b(|v|)}b(|v-h|)$. The compact polynomial operator is applied
to this function only through its $L^2$ norm. Since
$\|b\|_\infty\le1$, the bound \eqref{eq:gh-L2}, the polynomial negative
moment in $h$, and both Hilbert--Schmidt square roots are unchanged.
Thus the original transverse exponent remains valid.

For a skew pivot, the twisted-convolution argument uses only the $L^2$
size of the annular kernel. Inserting $b$ decreases that size by at most
$\|b\|_\infty$; the original $R^{-1}$ decay is unchanged.

For a convolution pivot, write the first nonzero homogeneous convolution
term as $Q_s$, with
$\|Q_s\|_{\mathrm{coeff}}\ge c_{n,d}$. After the change $r=Rv$, the
one-dimensional fibre amplitude is
\[
 v\longmapsto b(Rv)v^{-1}\mathbf1_{[u,1]}(v).
\]
Its zero extension has uniformly bounded $L^\infty+V^\rho$ norm by
Lemma~\ref{lem:allq-calculus}. Let
$c_{\mathrm{poly},\rho}>0$ be the exponent in
Corollary~\ref{cor:allq-polynomial-vdc} for degrees at most $d$. Plancherel
on each line therefore gives the fibre bound
\[
 C_\rho\min\{1,(R^s|Q_s(\theta)|)^{-c_{\mathrm{poly},\rho}}\}.
\]
Choose
\begin{equation}\label{eq:allq-convolution-annular-beta}
 \beta_{s,\rho}
 :=\frac12\min\left\{c_{\mathrm{poly},\rho},\frac1s\right\}>0.
\end{equation}
The last display is bounded by
$C_\rho R^{-s\beta_{s,\rho}}|Q_s(\theta)|^{-\beta_{s,\rho}}$, and
\eqref{eq:negsphere} makes the angular integral finite uniformly in all
lower coefficients. Thus the convolution branch has exponent
$s\beta_{s,\rho}>0$.

Take $a_{\mu,\rho}$ to be the corresponding branch exponent. There are
only finitely many flag positions, so all constants are structural and
coefficient-uniform.
\end{proof}

For test functions, polar coordinates show that
$u\mapsto A^{P,\Theta,b}_{R,u}f$ is an absolutely continuous
$L^2$-valued path even when $b$ has jumps.  For almost every $u$,
\[
 \partial_uA^{P,\Theta,b}_{R,u}f(x)
 =-u^{-1}b(uR)\int_{\Sph^{n-1}}\Theta(\theta)
 e^{iP(x,x-uR\theta)}f(x-uR\theta)\,d\sigma(\theta).
\]
Consequently
\[
 \sup_{1/2\le u\le1}
 \|\partial_uA^{P,\Theta,b}_{R,u}f\|_q
 \le C\|b\|_\infty\|\Theta\|_1\|f\|_q,
 \qquad1\le q\le\infty.
\]
The anchored Sobolev--variation argument of
Lemma~\ref{lem:shortannular} therefore gives
\begin{equation}\label{eq:allq-short-annular-decay}
 \|V^2(A^{P,\Theta,b}_{R,u}f:1/2\le u\le1)\|_2
 \le C_{n,d,\rho}R^{-a_{\mu,\rho}/4}
 \|\Theta\|_\infty\|f\|_2.
\end{equation}
Put
\begin{equation}\label{eq:allq-global-annular-exponent}
 a_{0,\rho}
 :=\min\left(\left\{\frac{a_{\mu,\rho}}4:
                 \mu\ {\rm nonterminal}\right\}\cup\{1\}\right)>0.
\end{equation}
The added value $1$ is used only when there is no nonterminal position.

\begin{lemma}[Weighted $q'$-tracked partial-annulus estimates]
\label{lem:allq-qtracked-annulus}
Let $P$ be normalized and nonterminal, $R\ge1$, $1<q\le2$, and
$q'=q/(q-1)$.  Let $\Theta\in L^q(\Sph^{n-1})$.  Under
\eqref{eq:allq-normalize-b}, for every $1<p<\infty$ set
\[
 \gamma_{p,\rho}:=\frac{a_{0,\rho}}{2\max(p,p')}>0.
\]
Then, for every $f\in L^p(\R^n)$,
\begin{align}
 \sup_{1/2\le u\le1}
 \|A^{P,\Theta,b}_{R,u}f\|_p
 &\le C_{n,d,p,\rho}R^{-\gamma_{p,\rho}/q'}
       \|\Theta\|_q\|f\|_p,
 \label{eq:allq-qtracked-annulus}\\
 \|V^2(A^{P,\Theta,b}_{R,u}f:1/2\le u\le1)\|_p
 &\le C_{n,d,p,\rho}R^{-\gamma_{p,\rho}/q'}
       \|\Theta\|_q\|f\|_p.
 \label{eq:allq-qtracked-short}
\end{align}
All constants are uniform in the later phase coefficients and in the
location of the variation of $b$.
\end{lemma}

\begin{proof}
For fixed $P,R,u,b$, the map
$\Theta\mapsto A^{P,\Theta,b}_{R,u}$ is linear into
$\mathcal B(L^2)$.  Its angular $L^1$ bound follows from the annular
$L^1$ kernel and $\|b\|_\infty\le1$; its angular $L^\infty$ bound is
Proposition~\ref{prop:allq-annular-decay}.  Complex interpolation, whose
$L^\infty$ weight is $1/q'$, gives, after weakening the finite minimum of
the nonterminal exponents to $a_{0,\rho}$,
\[
 \sup_u\|A^{P,\Theta,b}_{R,u}\|_{2\to2}
 \le C_{n,d,\rho}R^{-a_{0,\rho}/q'}\|\Theta\|_q.
\]
The preceding a.e. derivative estimate and the anchor
$A^{P,\Theta,b}_{R,1}=0$ give, exactly as in
\eqref{eq:strong-far-a14b34},
\[
 \|V^2(A^{P,\Theta,b}_{R,u}f:u)\|_2
 \le C_{n,d,\rho}R^{-a_{0,\rho}/(4q')}
       \|\Theta\|_q\|f\|_2.
\]
On the other hand $V^2\le V^1$ and the same derivative formula give
nondecaying strong $(1,1)$ and $(\infty,\infty)$ bounds with norm
$C\|\Theta\|_1$.  Marcinkiewicz interpolation puts weight $2/p'$ on the
$L^2$ endpoint when $p<2$ and weight $2/p$ when $p>2$.  This proves
\eqref{eq:allq-qtracked-short} with the displayed, slightly weakened,
$\gamma_{p,\rho}$.  The same spatial interpolation starting from the
fixed-$u$ $L^2$ estimate proves the stronger fixed-$u$ bound and hence
\eqref{eq:allq-qtracked-annulus}.  Approximation in $f$ and $\Theta$ is
justified by the nondecaying $V^1$ estimate.
\end{proof}

Apply this lemma to the dyadic far path.  The annular pieces and their short
variations have $L^p$ norm
$C2^{-\gamma_{p,\rho}j/q'}\|\Theta\|_q\|f\|_p$; hence their geometric sum
is at most $Cq'\|\Theta\|_q\|f\|_p$.  The deterministic dyadic
long--short decomposition therefore proves the weighted analogue of
Proposition~\ref{prop:strong-far-q}.  On every finite truncation set, pair
the resulting linear angular map with a norm-one dual functional.  Taking
$q=k/(k-1)$ gives the compatible bounds $\|h\|_k\le Ck\|f\|_p$; the
common-representative and exponential-series argument used in
Proposition~\ref{prop:allq-terminal-LlogL} yields the weighted analogue of
Proposition~\ref{prop:strong-far-endpoint}.

Finally, the local triangular reduction of
Proposition~\ref{prop:strong-local-reduction} acquires the factor
$b(|x-y|)$ in its error kernel; its absolute value is bounded by one, and
the residual later-phase path carries the same $b$. Backward induction
over the finite flag, with the unphased terminal estimate
Proposition~\ref{prop:allq-terminal-LlogL}, therefore proves
\begin{equation}\label{eq:allq-strong-core}
 \|\mathbf J(T^{P,b}_{\Omega,\varepsilon}f)\|_p
 \le C_{n,d,p,\rho}\mathfrak L(\Omega)\|f\|_p
\end{equation}
under \eqref{eq:allq-normalize-b}. The normalizing dilation sends $b$ to
$b(R\,\cdot)$, and \eqref{eq:allq-dilation-b} preserves its radial norm;
hence the estimate is uniform in every coefficient of $P$.

\subsection{Phase-adapted packets with finite radial variation}

\begin{proposition}[Radially weighted phase-adapted one-shell packet]
\label{prop:allq-direct-packet}
For every fixed $1\le\rho<\infty$,
Proposition~\ref{prop:direct-packet} remains valid with the radial factor
$b(|x-y|)$ under \eqref{eq:allq-normalize-b}. Namely, there are
$c_{P,\rho},c_{E,\rho}>0$ and $t_{0,\rho}$ such that, whenever the
packet separation parameter satisfies $t\ge t_{0,\rho}$,
\begin{align}
 \|V^2(G_\varepsilon:R\le\varepsilon\le2R)\|_2^2
 &\le C_{n,d,\rho}\|\Theta\|_\infty^2
 2^{-c_{P,\rho}t}\Lambda M,
 \label{eq:allq-direct-G}\\
 \|V^1(E_\varepsilon:R\le\varepsilon\le2R)\|_1
 &\le C_{n,d,\rho}\|\Theta\|_1
 2^{-c_{E,\rho}t}M.
 \label{eq:allq-direct-E}
\end{align}
The estimates are uniform under packetwise scalar coefficients of modulus
at most one.
\end{proposition}

\begin{proof}
We repeat the finite-mode low/high phase-block induction in
Proposition~\ref{prop:direct-packet}, but choose all parameters after the
radial exponents have been fixed. Let
$c_{\mathrm{flat},\rho}$ and $C_{B,\rho}$ be the constants in
Corollary~\ref{cor:allq-high-shell-stability}, let
$\delta_{T,\rho}$ be given by
\eqref{eq:allq-terminal-explicit-decay}, and let $C_{\rm Four}$ dominate
all powers in Lemma~\ref{lem:block-fourier}. Decrease
$c_{\mathrm{flat},\rho}$ if necessary so that it is at most one.
Here and below $C_{B,\rho}$ has first been enlarged once to dominate every
radial-amplitude power in the equal-shell, unequal-shell, flat-block, and
collective-shell estimates; it is not enlarged after the next parameter is
chosen.  Put
\begin{equation}\label{eq:allq-packet-theta}
 \theta_\rho
 :=\frac{c_{\mathrm{flat},\rho}}
          {16(C_{B,\rho}+1)}.
\end{equation}
Choose $B_{*,\rho}>1$ so large that
\begin{equation}\label{eq:allq-packet-ratio}
 \frac{2C_{\rm Four}}{B_{*,\rho}-1}
 \le\min\{\theta_\rho/2,c_{\mathrm{flat},\rho}/8\}.
\end{equation}
Put $N_d=\max(d-1,0)$ and choose
$0<\kappa_{*,\rho}\le1$ so small that, for every possible number
$0\le N\le N_d$ of nonpure homogeneous blocks,
\begin{equation}\label{eq:allq-packet-terminal-budget}
 2C_{\rm Four}\sum_{r=1}^N
 \kappa_{*,\rho}B_{*,\rho}^{\,r-N}
 \le\min\{\delta_{T,\rho}/2,1/2\}.
\end{equation}
Set
\[
 \kappa_r=\kappa_{*,\rho}B_{*,\rho}^{\,r-N},
 \qquad S_r=\sum_{i<r}\kappa_i,
 \qquad \gamma_r=\theta_\rho\kappa_r.
\]
Then $S_r\le\kappa_r/(B_{*,\rho}-1)$ and
\begin{align}
 \gamma_r-2C_{\rm Four}S_r&\ge\gamma_r/2,
 \label{eq:allq-packet-hierarchy-a}\\
 c_{\mathrm{flat},\rho}\kappa_r
 -C_{B,\rho}\gamma_r-2C_{\rm Four}S_r
 &\ge c_{\mathrm{flat},\rho}\kappa_r/2,
 \label{eq:allq-packet-hierarchy-b}\\
 c_{\mathrm{flat},\rho}
 -C_{B,\rho}\gamma_r-2C_{\rm Four}S_r
 &\ge c_{\mathrm{flat},\rho}/2.
 \label{eq:allq-packet-hierarchy-c}
\end{align}
These are the complete compatibility conditions; the choice is
noncircular because all quantities on their right were fixed before
$B_{*,\rho}$ and $\kappa_{*,\rho}$.

Suppose block $r$ is the first high block. The preceding low blocks have
squared Fourier loss at most $2^{2C_{\rm Four}S_rt}$. Partition the
normalized shell into $K_r\simeq2^{\gamma_rt}$ cells. The within-cell
square function uses only $\|b\|_\infty\le1$ and contributes
$C2^{-\gamma_rt}\|\Theta\|_\infty^2\Lambda M$. A fixed
Rademacher--Menshov choice of cell signs produces the radial profile
\[
 b(Ru)\sum_{\ell=1}^{K_r}\epsilon_\ell\mathbf1_{I_\ell}(u).
\]
By \eqref{eq:allq-b-step}, its $L^\infty+V^\rho$ norm is at most
$C_\rho(1+K_r^{1/\rho})$. Applying
\eqref{eq:allq-flat-high-block} and Lemma~\ref{lem:RM} therefore gives,
before the preceding Fourier loss is included,
\begin{equation}\label{eq:allq-direct-high-grid}
 C_\rho t^2\|\Theta\|_\infty^2
 2^{(C_{B,\rho}/\rho)\gamma_rt}\Lambda M
 \{2^{-c_{\mathrm{flat},\rho}\kappa_rt}
   +2^{-c_{\mathrm{flat},\rho}t}\}.
\end{equation}
Since $\rho\ge1$, replace $C_{B,\rho}/\rho$ by $C_{B,\rho}$.
Equations \eqref{eq:allq-packet-hierarchy-a}--
\eqref{eq:allq-packet-hierarchy-c}, after multiplication by the squared
Fourier loss, leave the three exponents
$\gamma_r/2$, $c_{\mathrm{flat},\rho}\kappa_r/2$, and
$c_{\mathrm{flat},\rho}/2$. Thus, after absorbing the polynomial factor
$t^2$, one may take
\begin{equation}\label{eq:allq-direct-high-exponent}
 c_{r,\rho}
 :=\frac14\min\{\gamma_r,
    c_{\mathrm{flat},\rho}\kappa_r,c_{\mathrm{flat},\rho}\}>0.
\end{equation}
Every high branch therefore has $E=0$ and geometric $L^2(V^2)$ decay.

If all blocks are low, expand them all. Freezing the combined input
frequency at each packet centre leaves a mean-zero packet, so
Proposition~\ref{prop:allq-terminal-packet} applies. The squared Fourier
loss is at most
$2^{2C_{\rm Four}\sum_r\kappa_rt}$, and
\eqref{eq:allq-packet-terminal-budget} gives, at the level of the full
squared estimate,
\[
 2^{-\delta_{T,\rho}t}
 \left(2^{C_{\rm Four}\sum_r\kappa_rt}\right)^2
 \le 2^{-\delta_{T,\rho}t/2}.
\]
The input-modulation remainder is estimated
in total truncation $1$-variation exactly as in
Proposition~\ref{prop:direct-packet}; multiplication by $b$ costs only
$\|b\|_\infty\le1$. The weighted Fourier sum and
\eqref{eq:allq-packet-terminal-budget} leave the bound
$C\|\Theta\|_1 2^{-t/2}M$.

All Fourier expansions are first truncated. If $K'>K$ and block $r$
is the first high block, define
\[
 \mathfrak B_{r,t,\rho}
 :=2^{-\gamma_rt}
 +t^2 2^{(C_{B,\rho}/\rho)\gamma_rt}
   \{2^{-c_{\mathrm{flat},\rho}\kappa_rt}
      +2^{-c_{\mathrm{flat},\rho}t}\}.
\]
Applying the within-cell and signed-grid estimates only to Fourier tuples
omitted at level $K$ gives, on every normalized interacting block pair,
\begin{equation}\label{eq:allq-direct-high-fourier-cauchy}
 \|G^{K'}-G^K\|_{L^2(V^2)}
 \le C_\rho\|\Theta\|_\infty(\Lambda\widetilde M)^{1/2}
       \mathfrak B_{r,t,\rho}^{1/2}\omega_K(\{i:i<r\}),
 \qquad E^{K'}-E^K=0.
\end{equation}
If all blocks are low, the radial terminal-packet estimate and the
freezing remainder give the following sharper tail estimates.  Here one
does \emph{not} replace the omitted Fourier coefficients by the global
majorant $2^{C_{\rm Four}\sum_r\kappa_rt}$: the exact common tail
$\omega_K$ is retained.  Consequently the terminal estimate contributes
its raw square-root factor $2^{-\delta_{T,\rho}t/2}$ before the omitted
tuples are summed:
\begin{align}
 \|G^{K'}-G^K\|_{L^2(V^2)}
 &\le C_\rho\|\Theta\|_\infty
       2^{-\delta_{T,\rho}t/2}
       (\Lambda\widetilde M)^{1/2}\omega_K(\{1,\ldots,N\}),
 \label{eq:allq-direct-low-good-cauchy}\\
 \|E^{K'}-E^K\|_{L^1(V^1)}
 &\le C\|\Theta\|_1 2^{-t}\widetilde M
       \omega_K(\{1,\ldots,N\}).
 \label{eq:allq-direct-low-error-cauchy}
\end{align}
Here $\widetilde M$ is the normalized block mass. Since
$\omega_K(S)\to0$ by \eqref{eq:direct-fourier-tail-zero}, the finite-mode
decompositions are Cauchy in the anchored $L^2(V^2)$ and $L^1(V^1)$
spaces. The absolute common Fourier majorants give uniform convergence
of the phase factors on the compact normalized blocks; dominated
convergence therefore identifies the variation-space limit with the
original trajectory. Fatou on finite parameter partitions preserves all
displayed bounds.

For the branch $N=0$ there is no high block and we assign its squared
decay exponent to be $\delta_{T,\rho}/2$.  Thus the final minimum is never
over an empty set: take $c_{P,\rho}$ to be the minimum of
$\delta_{T,\rho}/2$ and all high-branch exponents $c_{r,\rho}$ arising
for the finitely many cases $1\le N\le N_d$, and take
$c_{E,\rho}=1/2$. Increase $t_{0,\rho}$ so that it dominates the
terminal-packet threshold and absorbs the finitely many polynomial
factors used above. Rescaling and the bounded-overlap block summation are
exactly those in the final paragraph of Proposition~\ref{prop:direct-packet};
the $L^2$-square and $L^1$ Jacobians cancel the normalized mass factor.
This proves \eqref{eq:allq-direct-G}--\eqref{eq:allq-direct-E}.
\end{proof}

\subsection{The weak endpoint}

We first treat the unphased long path, where one additional modification is
essential.

\begin{lemma}[Whole-profile smoothing in the unphased weak long path]
\label{lem:allq-whole-profile-smoothing}
Fix a shell $2^j<|z|\le2^{j+1}$ and put
\[
 a_j(u)=b(2^ju)\mathbf1_{(1,2]}(u),
 \qquad 1/2<u<3.
\]
Under \eqref{eq:allq-normalize-b},
$\|a_j\|_\infty+V^\rho(a_j)\le C$ uniformly in $j$.  If
$a_{j,h}=a_j*\phi_h$, then
\begin{align}
 \|a_j-a_{j,h}\|_1&\le C_\rho h^{1/\rho},
 \label{eq:allq-weak-smooth-L1}\\
 \|\partial^Na_{j,h}\|_\infty&\le C_Nh^{-N}.
 \label{eq:allq-weak-smooth-derivative}
\end{align}
Consequently, in the long-variation part of the proof of Proposition~\ref{prop:nonoscweak} one
may replace the smoothing of the sharp shell indicator by smoothing the
entire profile $a_j$.  The sharp--smooth difference has total
$V^1L^1$ norm $O_\rho(h^{1/\rho}\Sigma)$, while the smoothed full-kernel
estimate has exactly the derivative hypothesis used in the proof of Lemma~\ref{lem:fullkernel}.
\end{lemma}

\begin{proof}
The uniform variation bound follows from Lemma~\ref{lem:allq-calculus};
the two estimates are Lemma~\ref{lem:allq-translation}.  Polar coordinates
give the shell-kernel $L^1$ difference
$C\|\omega_s\|_1h^{1/\rho}$.  If $E_j$ denotes this kernel difference,
the sharp and smooth long paths are suffix sums.  Inserting the zero and
total endpoints in an arbitrary variation partition gives the deterministic
bound
\[
 V^1\left(\sum_{j\ge k}E_j*B_{j-s}:k\right)
 \le 2\sum_j|E_j*B_{j-s}|.
\]
Since, for fixed $s$, every bad cube is assigned to at most one shell,
Young's inequality gives
\[
 \left\|V^1(\hbox{sharp long path}-\hbox{smooth long path})\right\|_1
 \le C\|\omega_s\|_1h^{1/\rho}
   \sum_j\|B_{j-s}\|_1
 \le C\|\omega_s\|_1h^{1/\rho}\Sigma.
\]
At physical scale the smoothed kernel is
\[
 H_j^s(z)=K_{\omega_s}(z)a_{j,h}(2^{-j}|z|).
\]
The chain rule and \eqref{eq:allq-weak-smooth-derivative} give
\[
 \left\|\partial_r^N
  [a_{j,h}(2^{-j}r)]\right\|_\infty
 \le C_N2^{-jN}h^{-N},
\]
while polar coordinates give $\|H_j^s\|_1\le C\|\omega_s\|_1$.
Thus its normalized full-kernel constant satisfies
$M_N\le C\|\omega_s\|_\infty h^{-N}$, exactly as required.
\end{proof}


\begin{lemma}[One-dimensional fixed-sign full-kernel estimate]
\label{lem:allq-fullkernel-1d}
Let $H_j$ be odd $C^2$ kernels on $\R$, supported where
$2^{j-2}\le|x|\le2^{j+2}$, and assume
\[
 \|H_j\|_1\le M,
 \qquad
 \|H_j'\|_\infty\le M2^{-2j},
 \qquad
 \|H_j''\|_\infty\le M2^{-3j}.
\]
Let $\mathcal Q$ be a finite family of pairwise disjoint intervals and
suppose
\[
 \supp b_Q\subset Q,\qquad \int b_Q=0,
 \qquad \|b_Q\|_1\le\Lambda|Q|.
\]
For $s\ge5$, pair $H_j$ with the intervals of length $2^{j-s}$.  Then for
arbitrary fixed coefficients $|\epsilon_Q|\le1$,
\begin{equation}\label{eq:allq-fullkernel-1d}
 \left\|\sum_jH_j*\!\sum_{|Q|=2^{j-s}}\epsilon_Qb_Q\right\|_2^2
 \le C M^2 2^{-s}\Lambda\sum_Q\|b_Q\|_1.
\end{equation}
\end{lemma}

\begin{proof}
Write $B_j=\sum_{|Q|=2^{j-s}}\epsilon_Qb_Q$ and
$L_{jk}=\widetilde H_j*H_k$.  Since $H_j$ is odd, $\int H_j=0$.
The family $\mathcal Q$ is finite, so only finitely many $B_j$ are
nonzero and all expansions below are finite.  With
$\widetilde H_j(x)=\overline{H_j(-x)}$, oddness and integrability also give
$\int\widetilde H_j=0$.
For $j\le k$,
\[
 L_{jk}'(x)
 =\int\widetilde H_j(y)
   [H_k'(x-y)-H_k'(x)]\,dy,
\]
and therefore
\begin{equation}\label{eq:allq-Ljk-derivative}
 \|L_{jk}'\|_\infty
 \le C M^2 2^j2^{-3k}.
\end{equation}
Moreover $L_{jk}$ is supported in an interval of length $C2^k$.
Because the intervals supporting $B_k$ are disjoint and have length
$2^{k-s}$, the total $|B_k|$ mass in any interval of length $C2^k$ is at
most $C2^s\Lambda2^{k-s}=C\Lambda2^k$.  Hence
\[
 \|(L_{jk}*B_k)'\|_\infty
 \le C M^2\Lambda 2^{j-2k}.
\]
For each interval $Q$ contributing to $B_j$, cancellation at its center
$x_Q$ gives
\[
 \left|\int_Q\epsilon_Qb_Q(x)(L_{jk}*B_k)(x)dx\right|
 \le \|b_Q\|_1|Q|\,\|(L_{jk}*B_k)'\|_\infty,
\]
where $|\epsilon_Q|\le1$.  Since $|Q|=2^{j-s}$,
\[
 |\langle H_j*B_j,H_k*B_k\rangle|
 \le C M^2\Lambda 2^{-s}2^{-2(k-j)}
 \sum_{|Q|=2^{j-s}}\|b_Q\|_1.
\]
This estimate includes $j=k$.  Sum first in $k\ge j$ and then in $j$;
the geometric factor $2^{-2(k-j)}$ and the identity
$\sum_j\sum_{|Q|=2^{j-s}}\|b_Q\|_1=\sum_Q\|b_Q\|_1$ prove
\eqref{eq:allq-fullkernel-1d} after expanding the squared norm and using
conjugate symmetry for $j>k$.
\end{proof}

\begin{proposition}[Unphased weak endpoint with radial $V^\rho$]
\label{prop:allq-unphased-weak}
For every mean-zero $\Omega\in L\log L(\Sph^{n-1})$ and every bounded
compactly supported $f$, under \eqref{eq:allq-normalize-b},
\begin{equation}\label{eq:allq-unphased-weak}
 \|\mathbf J(T^{0,b}_{\Omega,\varepsilon}f)\|_{L^{1,\infty}}
 \le C_{n,\rho}\mathfrak L(\Omega)\|f\|_1.
\end{equation}
\end{proposition}

\begin{proof}
Impose the finite bad-cube, shell, and rational-parameter cutoffs from
Proposition~\ref{prop:nonoscweak}; all estimates below are uniform in
those cutoffs. If $\mathfrak L(\Omega)=0$, then $\Omega=0$ and there is
nothing to prove. By homogeneity normalize $\mathfrak L(\Omega)=1$ and
perform the ordinary Calder\'on--Zygmund decomposition at height
$\alpha$. Write
\[
 b_{\rm CZ}=\sum_Qb_Q,\qquad
 \Sigma=\sum_Q\|b_Q\|_1\le C\|f\|_1.
\]
The good function is controlled by the $p=2$ case of
Proposition~\ref{prop:allq-terminal-LlogL}.  Thus, for the current finite
parameter set $I$, the set
\[
 E_{\rm good}
 =\{x:\mathbf J_I(T^{0,b}_{\Omega,\varepsilon}g)>\alpha\}
\]
has measure at most $C\alpha^{-1}\|f\|_1$.

Choose
\begin{equation}\label{eq:allq-weak-parameters}
 N=4n+1,\qquad
 \eta_0=\frac1{192N},\qquad
 0<\delta_a<
 \min\left\{\frac{\delta_{T,\rho}}8,\frac1{192}\right\}.
\end{equation}
For the bad function use the mean-zero angular split
$\Omega=\omega_s+\Omega_s^\sharp$ of
\eqref{eq:meanzerosplit} with this value of $\delta_a$. The angular tails
satisfy the original $V^1L^1$ estimate because multiplication by $b$
costs only $\|b\|_\infty\le1$. For the bounded pieces,
\[
 \|\omega_s\|_\infty\le C2^{\delta_as},
 \qquad \|\omega_s\|_1\le C.
\]
Proposition~\ref{prop:allq-terminal-packet} gives the bounded-angular
short square function with squared $L^2$ bound
\begin{equation}\label{eq:allq-weak-short-decay}
 C2^{-[\delta_{T,\rho}-2\delta_a]s}\alpha\Sigma,
\end{equation}
which is geometrically summable by \eqref{eq:allq-weak-parameters}.

For the dyadic long path take $h_s=2^{-\eta_0s}$ and apply
Lemma~\ref{lem:allq-whole-profile-smoothing}. The sharp--smooth
difference has total $V^1L^1$ norm
\begin{equation}\label{eq:allq-weak-unsmooth-sum}
 C_\rho2^{-\eta_0s/\rho}\Sigma.
\end{equation}
The smoothed kernels satisfy
\[
 M_N\le C\|\omega_s\|_\infty h_s^{-N}
 \le C2^{(\delta_a+N\eta_0)s}.
\]
Use exactly the nested adjacent grids, exceptional sets $X_s$, fixed
ranks, and suffix representation from
\eqref{eq:adjacent-grids}--\eqref{eq:rank-shell-pointwise}. A fixed rank
sign is still chosen before the output point is evaluated, so
Lemma~\ref{lem:fullkernel} applies. After the
Rademacher--Menshov factor is squared, the classwise long bound is
\begin{equation}\label{eq:allq-weak-long-decay}
 Cs^2 2^{-[1/24-2\delta_a-2N\eta_0]s}\alpha\Sigma.
\end{equation}
Here
\[
 \frac1{24}-2\delta_a-2N\eta_0
 >\frac1{24}-\frac1{96}-\frac1{96}
 =\frac1{48}.
\]
Thus \eqref{eq:allq-weak-long-decay} is summable. In dimension one,
\eqref{eq:meanzerosplit} gives
$\omega_s(1)+\omega_s(-1)=0$.  Since the smoothed radial profile is even
as a function of the spatial variable, the kernel $H_j^s$ is odd.
Lemma~\ref{lem:allq-fullkernel-1d} therefore applies; the $C^2$ smoothing norm is at most
$C2^{(\delta_a+2\eta_0)s}$, so its exponent
$1-2\delta_a-4\eta_0$ is larger than the common positive exponent
$1/48$. Hence the same summable majorant works in every dimension.

For clarity, here is the exact weighted counterpart of the majorants in
Proposition~\ref{prop:nonoscweak}.  Denote by
$U^b_{s,\varepsilon}$ the bounded-angular level-$s$ path, by
$H_j^{s,b}$ its whole-profile smoothed annular kernel, by $B_{j-s}$ the
same bad-cube input as there, and by
$\operatorname{Tail}^b_{s,\varepsilon}$ the angular-tail path.  Put
\begin{align*}
 \mathcal S_s^b
 &=\left(\sum_j
   V^2(U^b_{s,\varepsilon}:2^j\le\varepsilon\le2^{j+1})^2
   \right)^{1/2},\\
 \widetilde U_s^b(k)&=\sum_{j\ge k}H_j^{s,b}*B_{j-s},\\
 \mathcal L_s^b
 &=\mathbf1_{X_s^c}V^2(\widetilde U_s^b(k):k),\\
 \mathcal E_s^b
 &=V^1(U^b_{s,2^k}-\widetilde U_s^b(k):k),\\
 \mathcal T_s^b
 &=V^1(\operatorname{Tail}^b_{s,\varepsilon}:\varepsilon>0).
\end{align*}
The terminal packet estimate \eqref{eq:allq-weak-short-decay} controls
$\mathcal S_s^b$ in $L^2$, the fixed-rank whole-profile estimate
\eqref{eq:allq-weak-long-decay} controls $\mathcal L_s^b$ in $L^2$,
\eqref{eq:allq-weak-unsmooth-sum} controls $\mathcal E_s^b$ in $L^1$,
and the unchanged angular-tail estimate controls $\mathcal T_s^b$ in
$L^1$.  Thus, exactly as in \eqref{eq:nonosc-global-majorants}, set
\[
 \mathcal G=\sum_{s\ge s_0}(\mathcal S_s^b+\mathcal L_s^b),
 \qquad
 \mathcal W=\sum_{s\ge s_0}(\mathcal E_s^b+\mathcal T_s^b).
\]
Their sums satisfy
\[
 \|\mathcal G\|_2^2\le C_{\rho}\alpha\Sigma,\qquad
 \|\mathcal W\|_1\le C_{\rho}\Sigma,\qquad
 |X|\le C\alpha^{-1}\|f\|_1.
\]
All majorants and exceptional sets are independent of the jump amplitude.
The finite-path long--short inequality and finite-path jump normability
therefore give, off the fixed exceptional set,
\[
 \mathbf J_I(T^{0,b}_{\Omega,\varepsilon}b_{\rm CZ}:
                \varepsilon\in I)
 \le C_\rho(\mathcal G+\mathcal W).
\]
Let $E_{\rm exc}$ be the union of $E_{\rm good}$, the fixed enlarged
Calder\'on--Zygmund set, and $X=\bigcup_sX_s$ (an empty component is simply
omitted).  Its measure is at most
$C\alpha^{-1}\|f\|_1$.  Choose a fixed
$A_\rho>2$ larger than the constants in the pointwise jump
quasi-triangle inequality and in the last display.  Outside
$E_{\rm exc}$, the inclusion
\[
 \{\mathbf J_I(T^{0,b}_{\Omega,\varepsilon}f)>A_\rho\alpha\}
 \subset
 \{\mathcal G>c_\rho\alpha\}\cup
 \{\mathcal W>c_\rho\alpha\}
\]
holds for a fixed $c_\rho>0$.  Consequently Chebyshev, with its quadratic
form on $\mathcal G$ and linear form on $\mathcal W$, gives
\[
 \begin{aligned}
 |\{\mathbf J_I(T^{0,b}_{\Omega,\varepsilon}f)>A_\rho\alpha\}|
 &\le |E_{\rm exc}|+
 C_\rho\alpha^{-2}\|\mathcal G\|_2^2+
 C_\rho\alpha^{-1}\|\mathcal W\|_1\\
 &\le C_\rho\alpha^{-1}\|f\|_1 .
 \end{aligned}
\]
This estimate is uniform in every cutoff.  For an arbitrary requested
threshold $\tau>0$, run the decomposition at
$\alpha=\tau/A_\rho$; no majorant or exceptional set depends on a jump
amplitude chosen afterward.

Finally remove the bad-cube and shell cutoffs, then exhaust the positive
rational parameters, in the order used in
Proposition~\ref{prop:nonoscweak}.  At each cutoff-removal step, fix
$\tau>0$ and a finite parameter set $I$.  The finite-annulus $L^1$ kernel
gives, after a subsequence, coordinatewise almost-everywhere convergence
$F_N\to F$ on $I$.  If $\mathbf J_I(F)>\tau$, a finite strict witness has
positive excess over $\tau$ and persists for all sufficiently large $N$;
hence
\[
 \mathbf1_{\{\mathbf J_I(F)>\tau\}}
 \le\liminf_{N\to\infty}
   \mathbf1_{\{\mathbf J_I(F_N)>\tau\}} .
\]
Fatou preserves the distribution estimate at the same threshold.
Increasing finite sets exhaust $\mathbb Q_+$, and local absolute
continuity in $\varepsilon$ identifies the rational and full-parameter
functionals.  Thus the cutoff estimates pass to the full pointwise
functional. Restoring angular homogeneity proves
\eqref{eq:allq-unphased-weak}.
\end{proof}

\begin{proposition}[Polynomial weak endpoint with radial $V^\rho$]
\label{prop:allq-polynomial-weak}
Under \eqref{eq:allq-normalize-b}, for every polynomial $P$ of degree at
most $d$, every mean-zero $\Omega\in L\log L(\Sph^{n-1})$, and every
bounded compactly supported $f$,
\begin{equation}\label{eq:allq-polynomial-weak}
 \|\mathbf J(T^{P,b}_{\Omega,\varepsilon}f)\|_{L^{1,\infty}}
 \le C_{n,d,\rho}\mathfrak L(\Omega)\|f\|_1.
\end{equation}
\end{proposition}

\begin{proof}
If $\mathfrak L(\Omega)=0$, then $\Omega=0$ and the assertion is
immediate. By angular homogeneity assume throughout the induction that
$\mathfrak L(\Omega)=1$.  At each inductive stage impose finite cube and
shell cutoffs and fix a finite rational truncation set $I$; every bound
below is uniform in these choices.  We induct on the largest total degree $m$
carrying a nonpure homogeneous class. If no such class exists, write
$P(x,y)=p(x)+q(y)$ and use
Proposition~\ref{prop:allq-unphased-weak} after input and output
modulation.

Suppose the result is known through nonpure degree $m-1$. Let $P_m$ be
the highest nonpure homogeneous block and put
\[
 C_m=\|P_m\|_{\rm coeff},\qquad r_0=C_m^{-1/m}.
\]
The exact conjugacy
\begin{equation}\label{eq:allq-weak-dilation-conjugacy}
 T^{P,b}_{\Omega,r_0\varepsilon}f(r_0X)
 =T^{\widetilde P,\widetilde b}_{\Omega,\varepsilon}f_{r_0}(X),
 \qquad
 \widetilde b(u)=b(r_0u),
\end{equation}
and \eqref{eq:allq-dilation-b} reduce the proof to
$\|P_m\|_{\rm coeff}=1$ without changing the normalized radial bound.

We now make the far-field parameter choice explicit. Put
\begin{equation}\label{eq:allq-polynomial-weak-parameters}
 d_*=\max(d,1),\qquad
 \kappa_0=\frac1{100d_*},\qquad
 c_{H,\rho}=c_{\kappa_0,\rho},
\end{equation}
where $c_{\kappa_0,\rho}>0$ is given by
\eqref{eq:allq-collective-explicit-choice}, and choose
\begin{equation}\label{eq:allq-polynomial-angular-choice}
 0<\delta_a<
 \frac18\min\{c_{P,\rho},c_{E,\rho},c_{H,\rho}\}.
\end{equation}
Perform the phase-adapted Calder\'on--Zygmund decomposition and write
\[
 b_{\rm CZ}=\sum_Qb_Q,\qquad
 \Sigma=\sum_Q\|b_Q\|_1\le C\|f\|_1.
\]
The good part is controlled by \eqref{eq:allq-strong-core}: since the
Calder\'on--Zygmund decomposition gives
$\|g\|_2^2\le C\alpha\|f\|_1$, the finite-path set
\[
 E_{\rm good}^P
 =\{x:\mathbf J_I(T^{P,b}_{\Omega,\varepsilon}g)>\alpha\}
\]
has measure at most $C_{n,d,\rho}\alpha^{-1}\|f\|_1$.  Use the
mean-zero angular split \eqref{eq:meanzerosplit} with the value
$\delta_a$ in \eqref{eq:allq-polynomial-angular-choice}. The angular
tails have the same summable $V^1L^1$ bound as in the unweighted proof.

For the bounded angular part, Proposition~\ref{prop:allq-direct-packet}
gives the short square function with squared $L^2$ decay
$2^{-[c_{P,\rho}-2\delta_a]s}$ and the packet error with $L^1$ decay
$2^{-c_{E,\rho}s}$. Split the dyadic shells by
\[
 1\le R^m\le2^{\kappa_0s}
 \quad\hbox{and}\quad
 R^m>2^{\kappa_0s}.
\]
There are $O(1+s)$ low shells. Cauchy--Schwarz and the endpoint
differences supplied by Proposition~\ref{prop:allq-direct-packet} give a
low-shell squared $L^2$ bound
\[
 C(1+s)2^{-[c_{P,\rho}-2\delta_a]s}\alpha\Sigma,
\]
plus a summable $V^1L^1$ error.

On the high shells use the same output localization, fixed adjacent-grid
classes, exceptional sets, and rank suffixes as in
Proposition~\ref{prop:farweak}. The normalized sharp-annulus profiles are
products of $b(R\,\cdot)$ with fixed step cutoffs and hence have uniformly
bounded $L^\infty+V^\rho$ norm. Corollary~\ref{cor:allq-high-shell-stability}, followed by Lemma~\ref{lem:RM},
gives
\[
 Cs^2 2^{-[c_{H,\rho}-2\delta_a]s}\alpha\Sigma
\]
for the high dyadic tails. By
\eqref{eq:allq-polynomial-angular-choice}, every displayed exponent is
strictly positive and all polynomial factors in $s$ are summable.

Here is the precise comparison with the majorants in
Proposition~\ref{prop:farweak}.  Write
$G^b_{j,s}+E^b_{j,s}$ for the weighted one-shell decomposition supplied
by Proposition~\ref{prop:allq-direct-packet}, and
$F^{G,b}_{j,s}+F^{E,b}_{j,s}$ for its endpoint difference.  Define
\[
 \mathcal S_s^b=\left(\sum_jV^2(G^b_{j,s})^2\right)^{1/2},\qquad
 \mathcal L_s^b=2\sum_{j\ {\rm low}}|F^{G,b}_{j,s}|,
\]
\[
 \mathcal H_s^b=\mathbf1_{X_s^c}
 V^2(\text{weighted high dyadic tails}),\qquad
 \mathcal E_s^{*,b}=\sum_jV^1(E^b_{j,s}).
\]
Let $\mathcal W$ be the total $V^1$ majorant of the weighted angular-tail
paths plus $C\sum_s\mathcal E_s^{*,b}$, including the low-shell endpoint
errors, and set
\[
 \mathcal G=\sum_s(\mathcal S_s^b+\mathcal L_s^b+\mathcal H_s^b),
 \qquad X=\bigcup_sX_s.
\]
These are exactly the four objects in
\eqref{eq:far-shell-error-majorant}--\eqref{eq:far-global-GW}, with the
new direct-packet and collective-high estimates substituted; multiplication
by $|b|\le1$ leaves the angular-tail comparison unchanged.  Hence the
resulting nonnegative majorants satisfy
\[
 \|\mathcal G\|_2^2\le C_{n,d,\rho}\alpha\Sigma,\qquad
 \|\mathcal W\|_1\le C_{n,d,\rho}\Sigma,\qquad
 |X|\le C\alpha^{-1}\|f\|_1,
\]
and, off the fixed complement,
\[
 \mathbf J(\mathcal F^{P,b}b_{\rm CZ})
 \le C_{n,d,\rho}(\mathcal G+\mathcal W).
\]
The signs are fixed packetwise before the output point is evaluated, and
none of these objects depends on the jump threshold.

For the local path apply
Proposition~\ref{prop:degree-local-reduction}. Its error kernel is simply
multiplied by $|b(|x-y|)|\le1$, while the residual phase has nonpure
degree at most $m-1$ and carries the same radial function. On each unit
output cube the induction hypothesis applies to the translated,
purely modulated input; the fixed input enlargements have bounded
overlap. The piecewise error satisfies
\[
 \|V^1(E)\|_1
 \le C_{n,d}\|\Omega\|_1\|f\|_1.
\]
Combining the far path, local residual, and error with the pointwise jump
quasi-triangle inequality now gives a literal threshold argument.  Namely,
choose a fixed $A_{n,d,\rho}$ larger than all quasi-triangle constants.
Outside $E_{\rm good}^P$, the enlarged-CZ set, and $X$, a strict
$A_{n,d,\rho}\alpha$ witness for the full path forces either
$\mathcal G>c_{n,d,\rho}\alpha$,
$\mathcal W>c_{n,d,\rho}\alpha$, a strict
$c_{n,d,\rho}\alpha$ witness for the lower-degree residual, or
$V^1(E)>c_{n,d,\rho}\alpha$.  Chebyshev, the two displayed majorant
bounds, the induction hypothesis, and the $V^1(E)$ estimate therefore give
\[
 |\{\mathbf J_I(T^{P,b}_{\Omega,\varepsilon}f:
                    \varepsilon\in I)>
       A_{n,d,\rho}\alpha\}|
 \le C_{n,d,\rho}\alpha^{-1}\|f\|_1.
\]
Putting $\alpha=\tau/A_{n,d,\rho}$ proves the required distribution
estimate for every strict threshold $\tau>0$, uniformly at the
finite-cutoff stage and for the normalized top block.

All estimates above were first proved with finite cube, shell, and
rational-parameter cutoffs.  To remove them, fix $\tau>0$ and a finite
rational parameter set $I$.  Finite-annulus coordinate convergence gives,
after a subsequence, $F_N\to F$ pointwise on $I$.  A strict witness for
$\mathbf J_I(F)>\tau$ has positive excess and survives for all sufficiently
large $N$, so
\[
 \mathbf1_{\{\mathbf J_I(F)>\tau\}}
 \le\liminf_N\mathbf1_{\{\mathbf J_I(F_N)>\tau\}}.
\]
Fatou preserves the finite-parameter distribution estimate without
changing its threshold.  Increasing finite sets exhaust
$\mathbb Q_+$, and local absolute continuity recovers all positive
parameters, exactly as in the final part of
Proposition~\ref{prop:farweak}. Finally the input $L^1$ norm and weak output quasi-norm acquire
the same factor under the dilation $r_0$, so the constant is independent
of $C_m$. This closes the finite degree induction under the normalization
$\mathfrak L(\Omega)=1$; angular homogeneity restores the factor
$\mathfrak L(\Omega)$.
\end{proof}

\begin{remark}[Quantitative closure of the finite-radial proof]
\label{rem:allq-parameter-ledger}
For clarity, the order of every small-parameter choice in this section is
as follows.  Once $n,d,\rho$ are fixed, the polynomial sublevel constants
are structural.  Equation~\eqref{eq:allq-vdc-explicit-exponent} first fixes
$c_{\mathrm{vdc},\rho}>0$, then
\eqref{eq:allq-rational-operator-exponent} fixes
$c_{\mathrm{rat},\rho}>0$.  These determine
$c_{\mathrm{eq},\rho}$ and $c_{\mathrm{uneq},\rho}$ in
\eqref{eq:allq-equal-final-exponent} and
\eqref{eq:allq-unequal-final-exponent}, hence
$c_{\mathrm{flat},\rho}$ and $c_{\kappa,\rho}$ in
\eqref{eq:allq-flat-explicit-choice} and
\eqref{eq:allq-collective-explicit-choice}.  Independently,
\eqref{eq:allq-terminal-explicit-decay} fixes
$\delta_{T,\rho}>0$.  Only after these exponents are known are
$B_{*,\rho}$ and $\kappa_{*,\rho}$ selected in
\eqref{eq:allq-packet-ratio}--\eqref{eq:allq-packet-terminal-budget};
the three inequalities
\eqref{eq:allq-packet-hierarchy-a}--\eqref{eq:allq-packet-hierarchy-c}
then fix $c_{P,\rho},c_{E,\rho}>0$.  Finally the angular splitting
parameters are chosen in \eqref{eq:allq-weak-parameters} and
\eqref{eq:allq-polynomial-angular-choice}.  Thus no decay exponent is
defined using a parameter that is chosen later, and every geometric
series in the strong and weak arguments has a strictly positive exponent
for each fixed finite $\rho$.  The resulting constants may deteriorate
as $\rho\to\infty$, as they must in view of
Proposition~\ref{prop:radial-infinity-counterexample}.
\end{remark}

\subsection{Completion of the proof}

We give the completion because it also explains why point jumps of $b$
do not create jumps in the truncation trajectory.  Assume first
\eqref{eq:allq-normalize-b}.  For bounded compactly supported $f$ set
\[
 \mathcal A_x(r)
 :=\int_{\Sph^{n-1}}e^{iP(x,x-r\theta)}
       \Omega(\theta)f(x-r\theta)\,d\sigma(\theta).
\]
For every $0<a<c<\infty$,
\begin{equation}\label{eq:allq-completion-annulus-kernel}
 \|b(|\cdot|)K_\Omega(\cdot)
       \mathbf1_{\{a<|\cdot|\le c\}}\|_1
 \le \|b\|_\infty\|\Omega\|_1\log(c/a).
\end{equation}
Fubini, first on the countable annuli
$\{N^{-1}<|z|<N\}$, gives one full-measure set on which
\begin{equation}\label{eq:allq-weighted-radial-density}
 \int_a^c|b(r)\mathcal A_x(r)|\,\frac{dr}{r}<\infty,
 \qquad
 T^{P,b}_{\Omega,a}f(x)-T^{P,b}_{\Omega,c}f(x)
 =\int_a^c b(r)\mathcal A_x(r)\,\frac{dr}{r}
\end{equation}
for every $0<a<c<\infty$.  Thus the core path is locally absolutely
continuous.  In particular, a jump of the regulated function $b$ at one
radius changes neither the Lebesgue integral nor the truncation path: no
derivative of $b$ occurs in \eqref{eq:allq-weighted-radial-density}.
For core data the path is also zero for all sufficiently large
truncations at each fixed $x$.

We next construct all coordinates for general data.  In the strong case
choose bounded compactly supported $f_m\to f$ in $L^p$ so that
\[
 \sum_m\|f_{m+1}-f_m\|_p<\infty.
\]
For $g_m=f_{m+1}-f_m$, the zero anchor for core data and
\eqref{eq:anchor-controls-maximal}, followed by
\eqref{eq:allq-strong-core}, give
\begin{equation}\label{eq:allq-strong-uniform-completion}
 \left\|\sup_{q\in\mathbb Q_+}
       |T^{P,b}_{\Omega,q}g_m|\right\|_p
 \le \|\mathbf J(T^{P,b}_{\Omega,\varepsilon}g_m)\|_p
 \le C_{n,d,p,\rho}\mathfrak L(\Omega)\|g_m\|_p.
\end{equation}
Hence the sum of the rational suprema is finite almost everywhere, and
$T^{P,b}_{\Omega,q}f_m$ is uniformly Cauchy in
$q\in\mathbb Q_+$ outside one null set.

In the weak case choose instead bounded compactly supported $f_m\to f$
in $L^1$ with $\|g_m\|_1\le2^{-3m}$.  The core weak estimate and the
same anchor imply
\begin{equation}\label{eq:allq-weak-borel-cantelli}
 \sum_m\left|\left\{x:
   \sup_{q\in\mathbb Q_+}|T^{P,b}_{\Omega,q}g_m(x)|>2^{-m}
 \right\}\right|
 \le C_{n,d,\rho}\mathfrak L(\Omega)
       \sum_m2^m\|g_m\|_1<\infty.
\end{equation}
Borel--Cantelli again gives almost-everywhere uniform Cauchy convergence
on $\mathbb Q_+$.  In either case we define every rational coordinate by
this uniform limit.  Moreover, if
\[
 \mathcal R_m(q,x)
 :=\sum_{k\ge m}T^{P,b}_{\Omega,q}g_k(x),
\]
then $\sup_{q\in\mathbb Q_+}|\mathcal R_m(q,x)|\to0$ almost everywhere.
For fixed $m$ the core path of $f_m$ vanishes for all sufficiently large
$q$; consequently
\begin{equation}\label{eq:allq-uniform-rational-anchor}
 \lim_{R\to\infty}
 \sup_{\substack{q\in\mathbb Q_+\\q\ge R}}
 |T^{P,b}_{\Omega,q}f(x)|=0
\end{equation}
almost everywhere in both completions.

It remains to identify these coordinates with the desired finite-shell
increments.  By \eqref{eq:allq-completion-annulus-kernel}, for every
rational $0<a<c<\infty$ the core finite-annulus expressions converge to
the corresponding expression for $f$ in $L^p$ in the strong case and in
$L^1$ in the weak case.  A diagonal subsequence over the rational pairs
therefore gives, without changing the preceding uniform rational limit,
the identity
\[
 T^{P,b}_{\Omega,a}f(x)-T^{P,b}_{\Omega,c}f(x)
 =\int_a^c b(r)\mathcal A_x(r)\,\frac{dr}{r}
\]
on one full-measure set.  Fubini on the same countable annuli shows that
the right-hand density is locally integrable for general $f$ as well
(Young's inequality applies in the strong case).  Define the real
coordinates from the base point $1$ by these finite-annulus integrals.
The resulting path is locally absolutely continuous, agrees with the
rational coordinates, and \eqref{eq:allq-uniform-rational-anchor}, by
continuity and density of $\mathbb Q_+$, yields its zero-at-infinity
anchor for all real radii.

For every finite rational parameter set, coordinatewise convergence and
\eqref{eq:finite-J-lsc} transfer the strong estimate by Fatou.  For the
weak estimate, fix its threshold $\tau>0$ before taking the limit: if the
limiting finite path has $\mathbf J>\tau$, one of its finite strict
witnesses has positive excess and therefore persists along the
coordinatewise convergent subsequence.  The corresponding level-set
indicator is bounded by the liminf of the approximating indicators, so
Fatou transfers the distribution estimate at the same $\tau$.
Increasing through finite sets exhausting $\mathbb Q_+$ and then using
local continuity and Lemma~\ref{lem:finite-witness} recovers the full
positive-parameter functional.  This also settles measurability: all
coordinates are measurable, and the completed jump functional is a
countable supremum over finite rational witnesses.  Two locally
continuous paths with the same finite-annulus increments differ by a
constant in the parameter, and the zero anchor makes the completion
unique.

We have proved \eqref{eq:finite-radial-strong} and
\eqref{eq:finite-radial-weak} under \eqref{eq:allq-normalize-b}.  If
$\|b\|_{\mathcal V^\rho}=0$, then $b=0$ almost everywhere and the theorem
is immediate.  Otherwise apply the normalized result to
$b/\|b\|_{\mathcal V^\rho}$; linearity of every coordinate and
homogeneity of $\mathbf J$ restore exactly one factor
$\|b\|_{\mathcal V^\rho}$.  Finally the pointwise weak-$2$
jump-to-$r$-variation inequality from Lemma~\ref{lem:jumpcalc} gives every
$r>2$ consequence, while the zero anchor gives the maximal assertion.
This completes the proof of Theorem~\ref{thm:finite-radial}.


\begin{remark}[Dependence on the radial exponent]
\label{rem:finite-radial-dependence}
The proof uses the positive exponents $1/\rho$ in both the translation
modulus and the first-derivative oscillatory estimate.  These exponents tend
to zero as $\rho\to\infty$, and no constant uniform in $\rho$ is asserted.
Proposition~\ref{prop:radial-infinity-counterexample} shows that a merely
bounded radial endpoint is false even for the one-dimensional Hilbert
kernel.
\end{remark}


\subsection{A shellwise quadratic radial-variation refinement for the Hilbert transform}

The abstract transfer in Proposition~\ref{prop:young-radial-transfer} stops
at $\rho=2$.  In the one-dimensional Hilbert model, Fourier orthogonality
allows a sharper shellwise statement.

For bounded $b:(0,\infty)\to\mathbb C$, put
\[
 H^b_\varepsilon f(x)=\int_{|y|>\varepsilon}b(|y|)f(x-y)\,\frac{dy}{y},
\]
where on Schwartz data the integral is understood as the symmetric
improper integral at infinity, and on $L^2$ the compatible strong-$L^2$
representative is constructed in the proof below,
and define
\begin{equation}\label{eq:hilbert-shell-r2}
 q_j=V^2(b;(2^j,2^{j+1})),\qquad
 R_2(b)=\left(\sum_{j\in\mathbb Z}q_j^2\right)^{1/2}.
\end{equation}
Concatenating disjoint shell partitions gives $R_2(b)\le V^2(b)$ whenever
the latter is finite.

\begin{proposition}[Shellwise radial $V^2$ Hilbert estimate]
\label{prop:hilbert-radial-v2}
For every bounded $b$ with $R_2(b)<\infty$ and every $f\in L^2(\R)$,
\begin{equation}\label{eq:hilbert-radial-v2}
 \|\mathbf J(H^b_\varepsilon f:\varepsilon>0)\|_2
 \le C\bigl(\|b\|_\infty+R_2(b)\bigr)\|f\|_2.
\end{equation}
In particular, the right-hand side may be replaced by
$C(\|b\|_\infty+V^2(b))\|f\|_2$.
\end{proposition}

\begin{proof}
Put $B=\|b\|_\infty$.  On $I_j=(2^j,2^{j+1})$, finite $2$-variation
makes $b$ regulated.  Introduce the closed-shell representative
\[
 \overline b_j(2^j)=b(2^j+),\qquad
 \overline b_j(r)=b(r)\quad(r\in I_j),\qquad
 \overline b_j(2^{j+1})=b(2^{j+1}-).
\]
Approximating the two endpoints from inside $I_j$ in every finite
partition gives
\begin{equation}\label{eq:hilbert-closed-shell-variation}
 V^2(\overline b_j;[2^j,2^{j+1}])\le q_j.
\end{equation}
Changing the values at the two endpoints does not change any of the
Lebesgue kernels below.  Set $c_j=\overline b_j(2^j)$ and
$e_j=\overline b_j-c_j$ on the closed shell.  Then
\begin{equation}\label{eq:hilbert-shell-decomposition}
 |c_j|\le B,\qquad e_j(2^j)=0,\qquad
 \|e_j\|_{L^\infty([2^j,2^{j+1}])}\le q_j,\qquad
 V^2(e_j;[2^j,2^{j+1}])\le q_j.
\end{equation}
For $2^j\le\varepsilon\le2^{j+1}$, the within-shell part of the path is
\[
 A^b_{j,\varepsilon}f(x)
 =\int_{\varepsilon<|y|\le2^{j+1}}b(|y|)f(x-y)\,\frac{dy}{y}.
\]
For the constant component, write
\[
 \mathcal H_\varepsilon f(x)
 =\int_{|y|>\varepsilon}f(x-y)\,\frac{dy}{y}.
\]
Then
$A^1_{j,\varepsilon}f=\mathcal H_\varepsilon f-
\mathcal H_{2^{j+1}}f$ for $\varepsilon\in I_j$, and hence the classical
Hilbert short-$2$-variation square-function estimate
\cite[Theorem~1.4]{CJRW} gives
\[
 \left\|\left(\sum_j
 V^2(A^1_{j,\varepsilon}f:\varepsilon\in I_j)^2
 \right)^{1/2}\right\|_2\le C\|f\|_2.
\]
The finite-annulus integral is continuous in $\varepsilon$ for almost
every $x$ on locally integrable data, so the same estimate holds after
adjoining both dyadic endpoints.  Together with $|c_j|\le B$, this
controls the constant components.  For the error component
$E_{j,\varepsilon}$, every finite
partition of $I_j$ gives pointwise
\[
 V^1(E_{j,\varepsilon}f:\varepsilon\in I_j)
 \le q_j\int_{2^j<|y|\le2^{j+1}}|f(x-y)|\,\frac{dy}{|y|}.
\]
Put
\[
 K_j(y)=|y|^{-1}\mathbf1_{\{2^j<|y|\le2^{j+1}\}},
 \qquad \|K_j\|_1=2\log2.
\]
Since $V^2\le V^1$, Tonelli in $L^2(\ell^2)$ and Young's inequality give
\[
\begin{aligned}
 \left\|\left(\sum_j
 V^2(E_{j,\varepsilon}f:\varepsilon\in I_j)^2\right)^{1/2}\right\|_2^2
 &\le\sum_jq_j^2\|K_j*|f|\|_2^2\\
 &\le(2\log2)^2R_2(b)^2\|f\|_2^2.
\end{aligned}
\]
The endpoint continuity noted above and the triangle inequality in
$L^2(\ell^2)$ therefore give
\begin{equation}\label{eq:hilbert-short-bound}
 \left\|\left(\sum_j
 V^2(A^b_{j,\varepsilon}f:
       2^j\le\varepsilon\le2^{j+1})^2\right)^{1/2}\right\|_2
 \le C(B+R_2(b))\|f\|_2.
\end{equation}

For the dyadic path define the odd shell kernels
\begin{equation}\label{eq:hilbert-shell-kernel}
 \nu_j^b(y)=\frac{b(|y|)}{y}
 \mathbf1_{\{2^j<|y|\le2^{j+1}\}}.
\end{equation}
We claim that, with $t=2^j|\xi|$,
\begin{equation}\label{eq:hilbert-shell-multiplier}
 |\widehat{\nu_j^b}(\xi)|
 \le C(B+q_j)\min\{t,t^{-1/2}\}.
\end{equation}
Oddness gives the bound $CBt$ for $t\le1$.  For $t\ge1$, scale the radial
integral to $[1,2]$ and write $e(u)=e_j(2^ju)$.  By
\eqref{eq:hilbert-closed-shell-variation}, $e(1)=0$ and
$V^2(e;[1,2])\le q_j$.  The constant part in
\eqref{eq:hilbert-shell-decomposition} is $O(Bt^{-1})$ by integration by
parts.  It remains to prove
\begin{equation}\label{eq:hilbert-v2-oscillatory}
 \left|\int_1^2 e(u)\frac{\sin(tu)}u\,du\right|
 \le C V^2(e;[1,2])t^{-1/2}
\end{equation}
when $e(1)=0$.  Partition $[1,2]$ at the zeros of $\sin(tu)$.  There are
$M\le C(1+t)$ consecutive intervals $J_m$.  Choose $u_m\in J_m$ and split
$e=e(u_m)+(e-e(u_m))$ on each interval.  The error is at most
\[
 C t^{-1}\sum_mV^2(e;J_m)
 \le C t^{-1}M^{1/2}
       \left(\sum_mV^2(e;J_m)^2\right)^{1/2}
 \le C t^{-1/2}V^2(e).
\]
For the last inequality, choose near-extremizing partitions on the
consecutive $J_m$ and concatenate them; this gives
\[
 \sum_mV^2(e;J_m)^2\le V^2(e;[1,2])^2.
\]
For the frozen part put
$\omega_m=\int_{J_m}\sin(tu)u^{-1}\,du$.  Integration by parts shows that
every consecutive partial sum of the $\omega_m$ is $O(t^{-1})$.
Discrete Abel summation, followed by Cauchy--Schwarz, gives
\[
 \left|\sum_me(u_m)\omega_m\right|
 \le Ct^{-1}\left(|e(u_1)|+
       \sum_m|e(u_{m+1})-e(u_m)|\right)
 \le Ct^{-1/2}V^2(e).
\]
Indeed, take the $u_m$ in increasing order.  Since $e(1)=0$, the
partition $1<u_1<\cdots<u_M$ gives
\[
 |e(u_1)|^2+\sum_{m=1}^{M-1}
 |e(u_{m+1})-e(u_m)|^2\le V^2(e;[1,2])^2;
\]
Cauchy--Schwarz over at most $M$ terms proves the last bound and also
accounts for the endpoint term in Abel summation.
This proves \eqref{eq:hilbert-v2-oscillatory} and hence
\eqref{eq:hilbert-shell-multiplier}.

For $\xi\ne0$ define
\[
 m_b(\xi)=\sum_{j\in\mathbb Z}\widehat{\nu_j^b}(\xi),
 \qquad m_b(0)=0.
\]
This series is absolutely convergent at each nonzero frequency.  More
precisely, for every finite $J\subset\mathbb Z$, the geometric summation
of the $B$ part of \eqref{eq:hilbert-shell-multiplier} and
Cauchy--Schwarz for the $q_j$ part give
\[
\begin{aligned}
 \sum_{j\in J}\left|\widehat{\nu_j^b}(\xi)\right|
 &\le CB\sum_j
  \min\{2^j|\xi|,(2^j|\xi|)^{-1/2}\}\\
 &\quad+C R_2(b)\left(\sum_j
  \min\{2^j|\xi|,(2^j|\xi|)^{-1/2}\}^2\right)^{1/2}\\
 &\le C(B+R_2(b)),
\end{aligned}
\]
uniformly in $\xi$ and $J$.  If $J_N\uparrow\mathbb Z$, the corresponding
multipliers converge pointwise to $m_b$ and remain uniformly bounded by
the displayed estimate.  Plancherel and dominated convergence therefore
show, for every $f\in L^2$, that
\[
 \sum_{j\in J_N}\nu_j^b*f
 \longrightarrow T^bf:=\mathcal F^{-1}(m_b\widehat f)
 \quad\hbox{strongly in }L^2.
\]
In particular,
\begin{equation}\label{eq:hilbert-static-bound}
 \|T^bf\|_2\le C(B+R_2(b))\|f\|_2.
\end{equation}

Choose a smooth radial low-pass multiplier $\widehat\phi$ equal to one on
$[-2,2]$ and supported in $[-4,4]$, and write
$\widehat{\phi_k}(\xi)=\widehat\phi(2^k\xi)$.  The exact identity
\begin{align}
 H^b_{2^k}f
 &=\phi_k*T^bf
 +\sum_{s\ge0}(\delta_0-\phi_k)*\nu_{k+s}^b*f
 -\sum_{s<0}\phi_k*\nu_{k+s}^b*f
 \label{eq:hilbert-dyadic-decomposition}
\end{align}
holds first for finite shell sums.  The $T^b$ term passes to the limit by
the strong convergence just proved, while the two relative-frequency
series pass to the limit by the summable square estimates below.  Thus
the identity holds in $L^2$ and, in particular, defines the dyadic
truncations $H^b_{2^k}f$ compatibly with the strong shell limit.
For $2^k<\varepsilon<2^{k+1}$ we use the compatible representative
\[
 H^b_\varepsilon f=H^b_{2^{k+1}}f+A^b_{k,\varepsilon}f.
\]
On Schwartz data this agrees with the symmetric improper integral.  The
estimates below are first proved with finite shell, relative-frequency,
and parameter cutoffs; the strong $L^2$ limits just constructed, followed
by the finite-parameter lower-semicontinuity argument below, remove those
cutoffs.  Proposition~2.1 of
Chen--Ding--Hong--Liu \cite{ChenDingHongLiu} and
\eqref{eq:hilbert-static-bound} control the critical jump of the first
term.  Here Proposition~2.1 uses the paired jump count, so this estimate
already follows the convention of $\mathbf J$.  The chain functional below
is introduced only to sum the two relative-frequency series in the MSZK
normed quotient, and \eqref{eq:chain-pair-functional} returns that sum to
the paired convention after the cutoffs are removed.

For a scalar sequence $(h_k)$,
\begin{equation}\label{eq:sequence-jump-square}
 \mathbf J(h_k:k\in\mathbb Z)
 \le V^2(h_k:k\in\mathbb Z)
 \le2\left(\sum_k|h_k|^2\right)^{1/2}.
\end{equation}
Using \eqref{eq:hilbert-shell-multiplier}, the support and first-order
vanishing of $1-\widehat\phi$, and dyadic summation, one obtains uniformly
in $\xi$
\begin{align}
 \sum_k |1-\widehat\phi(2^k\xi)|^2
 |\widehat{\nu_{k+s}^b}(\xi)|^2
 &\le C2^{-s}(B^2+R_2(b)^2),&&s\ge0,
 \label{eq:hilbert-high-tail-square}\\
 \sum_k |\widehat\phi(2^k\xi)|^2
 |\widehat{\nu_{k+s}^b}(\xi)|^2
 &\le C2^{2s}(B^2+R_2(b)^2),&&s<0.
 \label{eq:hilbert-low-tail-square}
\end{align}
For completeness, the $B^2$ terms follow by summing the dyadic envelope in
\eqref{eq:hilbert-shell-multiplier}; for the $q_{k+s}^2$ terms one uses the
corresponding pointwise envelope bound, respectively $C2^{-s}$ and
$C2^{2s}$, and then sums $q_{k+s}^2$.  Plancherel and
\eqref{eq:sequence-jump-square} now give geometric bounds
$C2^{-s/2}(B+R_2(b))\|f\|_2$ for $s\ge0$ and
$C2^s(B+R_2(b))\|f\|_2$ for $s<0$.

To sum the relative-frequency pieces in
\eqref{eq:hilbert-dyadic-decomposition}, fix
$K_R=\{-R,\ldots,R\}$ and first restrict $s$ to a finite set.  Use the
chain functional and the normed finite-path quotient from
Lemma~\ref{lem:jumpnorm}; its triangle inequality has a constant
independent of $K_R$, and
\eqref{eq:chain-pair-functional} converts between the chain and paired
functionals.  Write
\[
 U_k^{(s)}=(\delta_0-\phi_k)*\nu_{k+s}^b*f\quad(s\ge0),
 \qquad L_k^{(s)}=\phi_k*\nu_{k+s}^b*f\quad(s<0).
\]
The preceding estimates and the two geometric series give,
uniformly in both cutoffs,
\[
 \left\|\mathbf J^{\rm ch}_{K_R}\left(
   \left(\sum_{s\ge0}U_k^{(s)}\right)_{k\in K_R}\right)\right\|_2
 +\left\|\mathbf J^{\rm ch}_{K_R}\left(
   \left(\sum_{s<0}L_k^{(s)}\right)_{k\in K_R}\right)\right\|_2
 \le C(B+R_2(b))\|f\|_2,
\]
where each sum is finite at this stage.  Moreover, the tails tend to zero
in $L^2_x(\ell^2_k)$: by Minkowski and the geometric estimates their
norms are bounded by the corresponding tails of
$\sum_{s\ge0}2^{-s/2}$ and $\sum_{s<0}2^s$.  Thus the relative-frequency
sums converge in $L^2_x(\ell^2_k)$ and, after passing to a subsequence, at
every $k$ almost everywhere.  Lemma~\ref{lem:finite-witness},
\eqref{eq:finite-chain-J-lsc}, and Fatou remove the $s$ cutoff on
$K_R$.  Now let $R\to\infty$.  The finite-parameter chain functionals
increase to the full dyadic functional, so monotone convergence, followed
by \eqref{eq:chain-pair-functional}, proves
\begin{equation}\label{eq:hilbert-dyadic-jump}
 \|\mathbf J(H^b_{2^k}f:k\in\mathbb Z)\|_2
 \le C(B+R_2(b))\|f\|_2.
\end{equation}
Finally combine \eqref{eq:hilbert-short-bound} and
\eqref{eq:hilbert-dyadic-jump} through the pointwise long--short inequality
\eqref{eq:pointwise-longshort}.  There is a common full-measure set on
which all dyadic coordinates are defined and every finite-shell path is
continuous in $\varepsilon$; the compatible representative above is used
on this set.  This justifies the pointwise long--short step for the full
continuum path.
\end{proof}

\begin{remark}\label{rem:hilbert-shellwise-sharper}
The condition $R_2(b)<\infty$ is strictly weaker than global finite
$V^2(b)$.  A bounded function may be constant on each half-open dyadic
shell, so that $R_2(b)=0$, while its jumps from shell to shell have infinite
global $2$-variation.  Thus even this basic model identifies a radial class
strictly larger than the global $V^2$ class of Theorem~\ref{thm:finite-radial}.
\end{remark}


\section{Obstructions and sharpness}
\label{sec:sharpness-obstructions}

This section records four complementary negative results.  They distinguish
the pointwise critical-jump functional from the weaker mixed quasi-seminorm,
show why the soft radial transfer cannot reach exponent two, prove that the
finite-radial theorem has no $\rho=\infty$ endpoint, and show that the output
critical jump cannot be strengthened to true $2$-variation.

\subsection{The mixed-norm exchange is false}

\begin{proposition}[Finite genuine-jump mixed-norm counterexamples]
\label{prop:mixed-norm-counterexample}
For every $1\le p<\infty$ and every integer $M\ge2$ there exist a finite
measure space $X_M$ and a two-point scalar path $F$ such that
\begin{equation}\label{eq:mixed-outer-bounded}
 \sup_{\lambda>0}\|\lambda N_\lambda(F)^{1/2}\|_{L^p(X_M)}\le C_p,
\end{equation}
whereas
\begin{equation}\label{eq:mixed-inner-growth}
 \left\|\sup_{\lambda>0}\lambda N_\lambda(F)^{1/2}\right\|_{L^p(X_M)}
 \simeq(\log M)^{1/p}.
\end{equation}
In addition, for every $M\ge2$ there is a scalar path $G$ on one common
finite ordered parameter set such that
\begin{equation}\label{eq:mixed-weak-outer-bounded}
 \sup_{\lambda>0}
 \|\lambda N_\lambda(G)^{1/2}\|_{L^{1,\infty}(X_M)}\le\sqrt2,
\end{equation}
but
\begin{equation}\label{eq:mixed-weak-inner-growth}
 \left\|\sup_{\lambda>0}\lambda N_\lambda(G)^{1/2}
 \right\|_{L^{1,\infty}(X_M)}=\sqrt2\,M.
\end{equation}
Thus neither the strong nor the weak spatial norm can in general be
interchanged with the supremum in the jump amplitude.
\end{proposition}

\begin{proof}
For the strong example, take $X_M=\{1,\ldots,M\}$ with counting measure
and set $a_k=k^{-1/p}$.
On the ordered parameter set $\{0,1\}$ define $F_0(k)=0$ and $F_1(k)=a_k$.
Then $N_\lambda(F(k))=1$ for $0<\lambda<a_k$ and vanishes for
$\lambda\ge a_k$.  Thus the pointwise supremum equals $a_k$, and its
$p$th power sums to $\sum_{k\le M}k^{-1}\simeq\log M$.  For fixed
$\lambda$,
\[
 \|\lambda N_\lambda(F)^{1/2}\|_p
 =\lambda\#\{k:a_k>\lambda\}^{1/p}\le C_p,
\]
which proves \eqref{eq:mixed-outer-bounded}--
\eqref{eq:mixed-inner-growth}.

For the weak endpoint use the same counting space and put
\[
 I_M=\{0,1,\ldots,2\cdot4^M\},\qquad
 a_k=2^{-k},\qquad L_k=4^k.
\]
For $j\in I_M$ define
\[
 G_j(k)=a_k\,\mathbf1_{\{j<2L_k\}}\mathbf1_{\{j\ \mathrm{odd}\}}.
\]
Thus, at the point $k$, the path alternates
$0,a_k,0,a_k,\ldots,0$ through the first $2L_k+1$ coordinates and is
zero thereafter.  For $0<\lambda<a_k$, its $2L_k$ consecutive changes
are admissible paired strict jumps (successive pairs may share an
endpoint).  Conversely, every strict jump interval must contain one of
these changes, and ordered jump intervals cannot use the same change
twice.  Hence, exactly,
\[
 N_\lambda(G(k))=
 \begin{cases}
  2L_k,&0<\lambda<a_k,\\
  0,&\lambda\ge a_k.
 \end{cases}
\]
It follows that
\[
 \mathbf J(G)(k)=a_k\sqrt{2L_k}=\sqrt2
\]
for every $1\le k\le M$, proving
\eqref{eq:mixed-weak-inner-growth}.

Fix $\lambda>0$.  If the path is nonzero at no spatial point there is
nothing to prove.  Otherwise let $K$ be the largest integer in
$\{1,\ldots,M\}$ for which $\lambda<2^{-K}$.  The nonzero values of
$\lambda N_\lambda(G)^{1/2}$ have decreasing rearrangement
\[
 A,\frac A2,\ldots,\frac{A}{2^{K-1}},
 \qquad A=\sqrt2\,\lambda2^K<\sqrt2.
\]
Consequently
\[
 \|\lambda N_\lambda(G)^{1/2}\|_{\ell^{1,\infty}}
 =\sup_{1\le m\le K}m\frac{A}{2^{m-1}}
 \le A<\sqrt2,
\]
uniformly in $\lambda$ and $M$.  This proves
\eqref{eq:mixed-weak-outer-bounded}.
\end{proof}

The first construction shows that an outer supremum in the jump amplitude
cannot be upgraded to the pointwise functional by an abstract exchange of
strong spatial norms.  The second construction proves the same strict
failure in $L^{1,\infty}$ with a genuine jump trajectory on a common finite
parameter set.  In particular, the pointwise-inside weak endpoint in
Theorems~\ref{thm:main} and~\ref{thm:finite-radial} is genuinely stronger
than the conventional mixed weak quasi-seminorm.

\subsection{No black-box endpoint radial-variation path transfer}

\begin{proposition}[Abstract endpoint transfer obstruction]
\label{prop:abstract-radial-v2-counterexample}
For every $N\ge2$ there exist a finite scalar path $F$, a scalar coefficient
path $\beta$, and the weighted-tail path
\begin{equation}\label{eq:abstract-weighted-tail}
 G_i=\sum_{k=i}^{2N-1}\beta_k(F_k-F_{k+1}),\qquad G_{2N}=0,
\end{equation}
such that
\begin{equation}\label{eq:abstract-input-bounds}
 \mathbf J(F)\le C,
 \qquad \|\beta\|_\infty+V^2(\beta)\le C,
\end{equation}
but
\begin{equation}\label{eq:abstract-output-growth}
 \mathbf J(G)\ge c\sqrt{\log N}.
\end{equation}
Hence no universal scalar-path estimate of the form
\[
 \mathbf J(G)\le C(\|\beta\|_\infty+V^2(\beta))\mathbf J(F)
\]
can hold for the operation \eqref{eq:abstract-weighted-tail}.
Moreover, with absolute changes in the constants, this discrete example
embeds into the same weighted-tail structural class that appears in
Proposition~\ref{prop:young-radial-transfer}, now at the formal endpoint
$\rho=2$:
there are a locally absolutely continuous path
$\widetilde F:(0,\infty)\to\mathbb C$ with
$\widetilde F_s\to0$ as $s\to\infty$ and a bounded right-continuous
$\widetilde\beta$ of finite $2$-variation such that
\[
 \widetilde G_s=-\int_s^\infty \widetilde\beta(u)\,d\widetilde F_u
\]
has the analogues of \eqref{eq:abstract-input-bounds} and
\eqref{eq:abstract-output-growth}.  Thus local absolute continuity of the
input path does not restore an abstract endpoint transfer at $\rho=2$.
\end{proposition}

\begin{proof}
Let
\[
 H_N=\sum_{j=1}^N\frac1j,\qquad a_j=j^{-1/2},\qquad \sigma_j=(-1)^j,
\]
and define
\begin{equation}\label{eq:abstract-F-definition}
 F_{2j}=0\quad(0\le j\le N),
 \qquad F_{2j-1}=\sigma_ja_j\quad(1\le j\le N).
\end{equation}
If both endpoints of an increment have modulus at most $\lambda/2$, the
increment cannot be a strict $\lambda$-jump.  Each nonzero endpoint can
occur in at most two ordered paired jumps, and therefore
\[
 N_\lambda(F)
 \le2\#\{j:a_j>\lambda/2\}\le C\lambda^{-2}.
\]
Thus $\mathbf J(F)\le C$.

Put $\delta_j=(jH_N)^{-1/2}$, take $\beta_0=0$, and prescribe
\begin{equation}\label{eq:abstract-beta-increments}
 \beta_{2j-1}-\beta_{2j-2}=\sigma_j\delta_j,
 \qquad \beta_{2j}-\beta_{2j-1}=0.
\end{equation}
The nonzero increments alternate in sign and decrease in modulus.  Every
consecutive block sum is therefore bounded by its first nonzero term.
Every increment selected by a variation partition is such a consecutive
block, and these blocks are disjoint and occur in order.  Their first
nonzero indices are consequently distinct, so
\[
 V^2(\beta)^2\le\sum_{j=1}^N\delta_j^2=1,
 \qquad \|\beta\|_\infty\le1.
\]
Discrete summation by parts in \eqref{eq:abstract-weighted-tail}, using
$\beta_0=F_{2N}=0$, gives
\[
 G_0=\sum_{k=1}^{2N-1}(\beta_k-\beta_{k-1})F_k
 =\sum_{j=1}^N\delta_ja_j=\sqrt{H_N}.
\]
Since $G_{2N}=0$, one ordered jump yields
$\mathbf J(G)\ge\sqrt{H_N}\simeq\sqrt{\log N}$.

It remains to realize the same obstruction inside the locally absolutely
continuous path class.  Define $\widetilde F$ by
$\widetilde F(s)=0$ for $0<s\le1$, by affine interpolation from $F_k$ to
$F_{k+1}$ on every interval $[k+1,k+2]$, $0\le k<2N$, and by
$\widetilde F(s)=0$ for $s\ge2N+1$.  Thus $\widetilde F$ is locally
absolutely continuous and has the zero anchor.  Define the right-continuous
step function
\[
 \widetilde\beta(s)=\beta_k\quad(k+1\le s<k+2,\ 0\le k<2N),
\]
take $\widetilde\beta(s)=\beta_0=0$ for $0<s<1$, and take
$\widetilde\beta(s)=\beta_{2N}$ for $s\ge2N+1$.  Since
$\beta_{2N}=\beta_{2N-1}$, this extension adds no terminal jump.  Hence
\[
 \|\widetilde\beta\|_\infty+V^2(\widetilde\beta)
 \le \|\beta\|_\infty+V^2(\beta)\le C.
\]
For every $0\le i\le2N$, absolute continuity on the finitely many affine
pieces gives the exact identity
\begin{equation}\label{eq:abstract-continuous-embedding}
 -\int_{i+1}^\infty\widetilde\beta(s)\,d\widetilde F_s
 =\sum_{k=i}^{2N-1}\beta_k(F_k-F_{k+1})=G_i.
\end{equation}

For completeness, the interpolation does not enlarge the input critical
jump by a factor depending on $N$.  Given $\lambda>0$, let
$\widetilde F^{(\lambda)}$ be obtained by flattening to zero every
triangular spike whose height satisfies $a_j\le\lambda/4$.  Then
\[
 \|\widetilde F-\widetilde F^{(\lambda)}\|_\infty\le\lambda/4,
 \qquad
 V^1(\widetilde F^{(\lambda)})
 \le2\sum_{a_j>\lambda/4}a_j\le C\lambda^{-1}.
\]
Every strict $\lambda$-jump of $\widetilde F$ is therefore a strict
$\lambda/2$-jump of $\widetilde F^{(\lambda)}$.  Disjoint jump intervals
consume disjoint portions of total variation, and consequently
\[
 N_\lambda(\widetilde F)
 \le N_{\lambda/2}(\widetilde F^{(\lambda)})
 \le \frac{2}{\lambda}V^1(\widetilde F^{(\lambda)})
 \le C\lambda^{-2}.
\]
Thus $\mathbf J(\widetilde F)\le C$.  On the integer subsequence
$s=i+1$, \eqref{eq:abstract-continuous-embedding} retains
$\widetilde G_1=G_0=\sqrt{H_N}$ and
$\widetilde G_{2N+1}=G_{2N}=0$, so
$\mathbf J(\widetilde G)\ge c\sqrt{\log N}$.  This proves the asserted
locally absolutely continuous embedding.
\end{proof}

This obstruction is the sequence-space failure
$\ell^{2,\infty}\cdot\ell^2\not\subset\ell^1$.  Proposition~\ref{prop:hilbert-radial-v2}
and the general argument of Section~\ref{sec:finite-radial} escape it by
using Fourier and packet orthogonality absent from an arbitrary scalar path.

\subsection{Failure at the bounded-radial endpoint}

For the limiting seminorm put
$V^\infty(b)=\sup_{s<t}|b(t)-b(s)|$.

\begin{proposition}[No uniform merely bounded radial theory]
\label{prop:radial-infinity-counterexample}
There is no uniform $L^2$ bound for the one-dimensional radial Hilbert
operator in terms of $\|b\|_\infty+V^\infty(b)$ alone.  Consequently,
Theorem~\ref{thm:finite-radial} cannot be extended to $\rho=\infty$.
\end{proposition}

\begin{proof}
Let
\[
 \operatorname{sgn}_+(\sin r)=
 \begin{cases}
  \operatorname{sgn}(\sin r),&\sin r\ne0,\\
  \displaystyle\lim_{h\downarrow0}\operatorname{sgn}(\sin(r+h)),
     &\sin r=0.
 \end{cases}
\]
For $N>2$ take the right-continuous representative
\[
 b_N(r)=\mathbf1_{[1,N)}(r)\operatorname{sgn}_+(\sin r).
\]
The half-open cutoff makes the value at $N$ agree with the right limit,
and the definition at each zero of $\sin r$ makes all interior values
right continuous; the value at $1$ also agrees with its right limit.
Changing these countably many endpoint values does not change any
Lebesgue integral below.  Then $\|b_N\|_\infty\le1$ and
$V^\infty(b_N)\le2$.  The convolution
operator
\[
 T_Nf(x)=\int_{\mathbb R}\frac{b_N(|y|)}y f(x-y)\,dy
\]
has an integrable kernel.  With
$\widehat f(\xi)=\int f(x)e^{-ix\xi}\,dx$, its continuous multiplier at
$\xi=1$ is
\[
 m_N(1)=-2i\int_1^N\frac{|\sin r|}{r}\,dr.
\]
Splitting the integral into intervals of length $\pi$ and comparing $r^{-1}$
with its endpoint values gives
$\int_1^N|\sin r|r^{-1}dr=(2/\pi)\log N+O(1)$.  Hence
$\|T_N\|_{2\to2}\ge|m_N(1)|\ge c\log N$.

For the associated truncation path, $T^{b_N}_\varepsilon=T_N$ when
$0<\varepsilon<1$ and $T^{b_N}_\varepsilon=0$ when $\varepsilon>N$.
Therefore $|T_Nf|\le\mathbf J(T^{b_N}_\varepsilon f:\varepsilon>0)$.
A jump estimate with a constant depending only on
$\|b_N\|_\infty+V^\infty(b_N)$ would contradict the preceding multiplier
lower bound.
\end{proof}

\subsection{Failure of quadratic output variation for sharp Hilbert truncations}

Let
\[
 H_\varepsilon f(x)=\int_{|y|>\varepsilon}f(x-y)\,\frac{dy}{y},
 \qquad
 \widehat{P_sf}(\xi)=e^{-s|\xi|}\widehat f(\xi).
\]
For $u,t>0$ define
\begin{equation}\label{eq:mellin-definitions}
 \Phi(u)=\int_u^\infty\frac{\sin v}{v}\,dv,
 \qquad
 K(t)=\frac1\pi\operatorname{sech}^2(\log t)
 =\frac{4t^2}{\pi(1+t^2)^2}.
\end{equation}

\begin{lemma}[Mellin identity]
\label{lem:mellin-hilbert-poisson}
For every $s>0$,
\begin{equation}\label{eq:mellin-identity}
 e^{-s}=\int_0^\infty K(t)\Phi(s/t)\,\frac{dt}{t},
 \qquad
 \int_0^\infty K(t)\,\frac{dt}{t}=\frac2\pi.
\end{equation}
\end{lemma}

\begin{proof}
Let $I(s)$ denote the first integral.  Since $|\Phi|$ is bounded and
$K(t)dt/t$ is integrable, the integral is absolutely convergent.  Moreover
$\Phi'(u)=-\sin(u)/u$.  To justify differentiation, fix
$0<a<b<\infty$.  Uniformly for $s\in[a,b]$,
\[
 \left|\partial_s\{K(t)\Phi(s/t)\}\right|\frac1t
 =\frac{K(t)}s|\sin(s/t)|\frac1t
 \le C_{a,b}
 \begin{cases}
  t,&0<t\le1,\\
  t^{-4},&t\ge1,
 \end{cases}
\]
where for $t\ge1$ we used $|\sin(s/t)|\le s/t$.  The right-hand side is
integrable, so local dominated convergence gives
\[
 I'(s)=-\frac4{\pi s}\int_0^\infty
 \frac{t}{(1+t^2)^2}\sin(s/t)\,dt
 =-\frac4{\pi s}\int_0^\infty
 \frac{y}{(1+y^2)^2}\sin(sy)\,dy.
\]
The standard contour integral
\[
 \int_0^\infty\frac{\cos(sy)}{(1+y^2)^2}\,dy
 =\frac\pi4(1+s)e^{-s}
\]
and differentiation in $s$ show that the last sine integral equals
$\frac\pi4se^{-s}$.  Hence $I'(s)=-e^{-s}$.  Dominated convergence gives
$I(s)\to0$ as $s\to\infty$, so $I(s)=e^{-s}$.  Finally,
\[
 \int_0^\infty K(t)\frac{dt}{t}
 =\frac4\pi\int_0^\infty\frac{t}{(1+t^2)^2}\,dt=\frac2\pi.
\]
\end{proof}

For $\xi>0$, the multiplier of $H_\varepsilon$ is
$-2i\Phi(\varepsilon\xi)$.  Lemma~\ref{lem:mellin-hilbert-poisson}
therefore gives, first for Schwartz functions with Fourier support compactly
contained in $(0,\infty)$,
\begin{equation}\label{eq:poisson-hilbert-representation}
 P_sf=\frac i2\int_0^\infty K(t)H_{s/t}f\,\frac{dt}{t}.
\end{equation}
Here the integral is initially a Bochner integral in $L^2(\mathbb R)$.
For fixed $s$, the map $t\mapsto H_{s/t}f$ is strongly measurable (on the
Schwartz core this follows from joint measurability, or from continuity of
the multiplier in $L^2$), the Hilbert truncation multipliers are uniformly
bounded, and $K(t)dt/t$ has finite mass.  The multiplier identity
identifies this Bochner integral with $P_sf$ in $L^2$.

\begin{proposition}[Poisson variation dominated by Hilbert variation on the core]
\label{prop:poisson-by-hilbert}
If $f\in\mathcal S(\mathbb R)$ and
$\operatorname{supp}\widehat f\Subset(0,\infty)$, then
\begin{equation}\label{eq:poisson-by-hilbert}
 V^2(P_sf:s>0)
 \le\frac1\pi V^2(H_\varepsilon f:\varepsilon>0)
\end{equation}
pointwise.
\end{proposition}

\begin{proof}
For the present Schwartz data, $(t,x)\mapsto H_{s/t}f(x)$ is jointly
measurable.  Define
\[
 G_s(x)=\int_0^\infty K(t)|H_{s/t}f(x)|\,\frac{dt}{t}.
\]
The uniform $L^2$ multiplier bound and Minkowski give
\[
 \|G_s\|_2
 \le\int_0^\infty K(t)\|H_{s/t}f\|_2\,\frac{dt}{t}
 \le C\|f\|_2\int_0^\infty K(t)\,\frac{dt}{t}<\infty.
\]
Thus $G_s(x)<\infty$ almost everywhere.  The Bochner--Fubini theorem
(equivalently, approximation of
$t\mapsto K(t)H_{s/t}f/t$ in $L^1((0,\infty);L^2_x)$ by simple
functions) says that its Bochner integral is represented by the
pointwise absolutely convergent integral off this null set.  Since the
multiplier computation identifies the same Bochner integral with $P_sf$
in $L^2$, \eqref{eq:poisson-hilbert-representation} holds there after
choosing representatives.  Taking the union of these null sets over any
fixed finite set of $s$-parameters gives one common full-measure set.
For a finite partition $0<s_0<\cdots<s_J$, apply
\eqref{eq:poisson-hilbert-representation} and Minkowski's inequality in
$\ell^2_J$:
\begin{align*}
 \left(\sum_{j=1}^J|P_{s_j}f-P_{s_{j-1}}f|^2\right)^{1/2}
 &\le\frac12\int_0^\infty K(t)
 \left(\sum_{j=1}^J|H_{s_j/t}f-H_{s_{j-1}/t}f|^2\right)^{1/2}\frac{dt}{t}\\
 &\le\frac1\pi V^2(H_\varepsilon f:\varepsilon>0).
\end{align*}
Both the Poisson and sharp Hilbert trajectories of this Schwartz input
are continuous in their positive parameters.  It is therefore enough to
take the supremum over rational finite partitions, for which one common
full-measure set suffices.  This proves the pointwise inequality.  No
infinite-parameter $L^2$ closure is used here; the finite-parameter
$L^2$ closure needed below is carried out explicitly in the proof of
Theorem~\ref{thm:hilbert-v2-failure}.
\end{proof}

We record explicitly the finite-parameter transference used below.

\begin{lemma}[Finite-dimensional de Leeuw restriction]
\label{lem:finite-dimensional-deleeuw}
Let $E$ be a finite-dimensional Banach space and let $M:\mathbb R\to E$ be
continuous.  If
\[
 \left\|\mathcal F^{-1}(M\widehat f)\right\|_{L^2(\mathbb R;E)}
 \le C\|f\|_{L^2(\mathbb R)}
\]
for Schwartz $f$, then for every trigonometric polynomial $g$,
\begin{equation}\label{eq:deleeuw-vector-restriction}
 \left\|\sum_{n\in\mathbb Z}M(n)\widehat g(n)e^{inx}
 \right\|_{L^2(\mathbb T;E)}
 \le C\|g\|_{L^2(\mathbb T)}.
\end{equation}
\end{lemma}

\begin{proof}
This is the restriction argument of de Leeuw \cite{deLeeuw1965}, applied
with a finite-dimensional range.  Let $\Lambda\subset\mathbb Z$ be the
finite Fourier support of $g$.  Choose a nonzero
$\eta\in\mathcal S(\mathbb R)$ with
$\operatorname{supp}\widehat\eta\subset(-1/4,1/4)$, periodically extend
$g$, and put $f_\delta(x)=\eta(\delta x)g(x)$, $0<\delta<1$.
Then
\[
 \widehat f_\delta(\xi)
 =\sum_{n\in\Lambda}\widehat g(n)\delta^{-1}
   \widehat\eta((\xi-n)/\delta),
\]
and these packets are disjoint.  On the packet centered at $n$,
continuity gives
\[
 \sup_{|\zeta|\le1/4}\|M(n+\delta\zeta)-M(n)\|_E\longrightarrow0.
\]
There are only finitely many packets.  To record the normalization of the
error, choose functionals $\ell_1,\ldots,\ell_m\in E^*$ and a constant
$C_E$ such that
\[
 \|v\|_E\le C_E\left(\sum_{a=1}^m|\ell_a(v)|^2\right)^{1/2}.
\]
Set
\[
 \omega_\Lambda(\delta)
 =\max_{n\in\Lambda}\sup_{|\zeta|\le1/4}
   \|M(n+\delta\zeta)-M(n)\|_E,
\]
so $\omega_\Lambda(\delta)\to0$.  Packet disjointness and scalar
Plancherel applied to the $\ell_a$ give, up to the fixed Fourier
normalization,
\[
 \delta\left\|
 \mathcal F^{-1}\!\left(
   \sum_{n\in\Lambda}[M(\xi)-M(n)]\widehat g(n)\delta^{-1}
   \widehat\eta((\xi-n)/\delta)\right)
 \right\|_{L^2(\mathbb R;E)}^2
 \le C_{E,\eta}\omega_\Lambda(\delta)^2
       \sum_{n\in\Lambda}|\widehat g(n)|^2 .
\]
Consequently
\[
 \delta^{1/2}\left\|
 \mathcal F^{-1}(M\widehat f_\delta)
 -\eta(\delta\,\cdot)\sum_{n\in\Lambda}
  M(n)\widehat g(n)e^{in\,\cdot}
 \right\|_{L^2(\mathbb R;E)}
 \longrightarrow0.
\]

We also use the elementary periodic averaging formula
\[
 \delta\int_{\mathbb R}|\eta(\delta x)|^2A(x)\,dx
 \longrightarrow
 \left(\int_{\mathbb R}|\eta(u)|^2\,du\right)
 \frac1{2\pi}\int_0^{2\pi}A(x)\,dx
\]
for every continuous $2\pi$-periodic scalar function $A$.  Apply it first
to $A=|g|^2$ and then to the squared $E$-norm of the finite
$E$-valued trigonometric polynomial above.  With the harmless Fourier
normalization incorporated into $c_\eta>0$, this gives
\begin{align*}
 \delta\|f_\delta\|_2^2
 &\longrightarrow c_\eta\|g\|_{L^2(\mathbb T)}^2,\\
 \delta\|\mathcal F^{-1}(M\widehat f_\delta)\|_{L^2(E)}^2
 &\longrightarrow c_\eta
 \left\|\sum_nM(n)\widehat g(n)e^{inx}\right\|_{L^2(\mathbb T;E)}^2.
\end{align*}
The packet error tends to zero at precisely this normalization, and the
same constant $c_\eta>0$ occurs in both limits.  Applying the assumed
Euclidean inequality before taking the limit proves
\eqref{eq:deleeuw-vector-restriction} without loss in the operator norm.
\end{proof}

\begin{theorem}[Failure of true $V^2$ for sharp Hilbert truncations]
\label{thm:hilbert-v2-failure}
There is no constant $C$ such that
\begin{equation}\label{eq:hilbert-v2-false}
 \|V^2(H_\varepsilon f:\varepsilon>0)\|_2\le C\|f\|_2
\end{equation}
for every $f\in L^2(\mathbb R)$.
\end{theorem}

\begin{proof}
Assume \eqref{eq:hilbert-v2-false}, and denote its constant by $C_H$.
For every
\[
 f\in\mathcal S_+
 :=\{g\in\mathcal S(\mathbb R):
       \operatorname{supp}\widehat g\Subset(0,\infty)\},
\]
Proposition~\ref{prop:poisson-by-hilbert} and the assumed estimate give
\begin{equation}\label{eq:positive-core-poisson-v2}
 \|V^2(P_sf:s>0)\|_2\le\frac{C_H}{\pi}\|f\|_2.
\end{equation}

We now pass only on a fixed finite parameter set.  Let
$D=\{s_0<\cdots<s_J\}\subset(0,\infty)$ and write $V_D^2$ for the
corresponding finite variation.  For every scalar $D$-path $z$,
\begin{equation}\label{eq:finite-D-coordinate-control}
 V_D^2(z)\le2\left(\sum_{s\in D}|z_s|^2\right)^{1/2};
\end{equation}
indeed, in any increasing variation chain each selected coordinate occurs
in at most two consecutive differences.  Since every $P_s$ is an
$L^2$ contraction, \eqref{eq:finite-D-coordinate-control} implies
\begin{equation}\label{eq:finite-D-L2-continuity}
 \|V_D^2(P_sh:s\in D)\|_2
 \le2|D|^{1/2}\|h\|_2.
\end{equation}

The class $\mathcal S_+$ is dense in the positive-frequency subspace
$L^2_+(\mathbb R)$.  Given $f\in L^2_+$, choose $f_m\in\mathcal S_+$
with $f_m\to f$ in $L^2$.  Equation
\eqref{eq:finite-D-L2-continuity} gives convergence in
$L^2(V_D^2)$, while \eqref{eq:positive-core-poisson-v2} has a constant
independent of $D$ and $m$.  Hence
\begin{equation}\label{eq:positive-finite-D-poisson-v2}
 \|V_D^2(P_sf:s\in D)\|_2
 \le\frac{C_H}{\pi}\|f\|_2,
 \qquad f\in L^2_+(\mathbb R).
\end{equation}
Reflection gives the same estimate on $L^2_-(\mathbb R)$.  Decomposing
$f=f_++f_-$, using subadditivity of $V_D^2$, and using orthogonality of
the frequency projections, we obtain
\begin{equation}\label{eq:full-finite-D-poisson-v2}
 \|V_D^2(P_sf:s\in D)\|_2
 \le\frac{\sqrt2C_H}{\pi}\|f\|_2
\end{equation}
for every $f\in L^2(\mathbb R)$, uniformly in the finite set $D$.

Let
$E_D=\mathbb C^{J+1}/\mathbb C\mathbf1$ with norm
\[
 \|[z]\|_{E_D}
 :=\sup_{0\le j_0<\cdots<j_m\le J}
 \left(\sum_{\ell=1}^m
 |z_{j_\ell}-z_{j_{\ell-1}}|^2\right)^{1/2}.
\]
This is the finite $V^2$ norm and is a genuine norm after quotienting by
constant vectors.  Equation \eqref{eq:full-finite-D-poisson-v2} says that
the continuous
$E_D$-valued symbol
\[
 M_D(\xi)=[(e^{-s|\xi|})_{s\in D}]
\]
defines an $L^2(\mathbb R)$ to $L^2(\mathbb R;E_D)$ multiplier with norm
independent of $D$.  Notice that it is continuous also at $\xi=0$ because
the vector $(1)_{s\in D}$ is zero in the quotient.

Lemma~\ref{lem:finite-dimensional-deleeuw} restricts $M_D$ to the
integers with the same uniform bound.  The restricted coordinates are the
periodic Poisson multipliers $e^{-s|n|}$.  If trigonometric polynomials
$g_m$ tend to $g$ in $L^2(\mathbb T)$, then
\eqref{eq:finite-D-coordinate-control} and the $L^2$ contractivity of
each periodic $P_s$ give
\[
 \|V_D^2(P_s(g_m-g):s\in D)\|_{L^2(\mathbb T)}
 \le2|D|^{1/2}\|g_m-g\|_{L^2(\mathbb T)}.
\]
Thus scalar density extends the restricted bound to every
$g\in L^2(\mathbb T)$ and every finite $D$, with the original bound
uniform in $D$.

Enumerate $\mathbb Q_+$ and let $D_m$ be the increasing finite sets
formed by the first $m$ rationals.  Pointwise,
\[
 V_{D_m}^2(P_sg:s\in D_m)\uparrow
 V_{\mathbb Q_+}^2(P_sg:s\in\mathbb Q_+).
\]
For $s$ in a compact subinterval of $(0,\infty)$, Cauchy--Schwarz in the
Fourier series shows that $s\mapsto P_sg(x)$ is continuous (indeed the
series is uniformly absolutely convergent).  Consequently
\[
 V_{\mathbb Q_+}^2(P_sg:s\in\mathbb Q_+)
 =V^2(P_sg:s>0).
\]
Monotone convergence applied to the squares and the uniform finite-$D_m$
bounds now yields
\[
 \|V^2(P_sg:s>0)\|_{L^2(\mathbb T)}
 \le \frac{\sqrt2C_H}{\pi}\|g\|_{L^2(\mathbb T)}.
\]

Jones and Wang
\cite[Theorem~1.13, statement on p.~4495 and proof on p.~4517]
{JonesWang2004} give a bounded function on $\mathbb T$ whose periodic
Poisson trajectory has infinite $2$-variation almost everywhere.  This
function lies in $L^2(\mathbb T)$, so it is in the domain of the estimate
just obtained.  Their $V_2$ is the supremum of the same quadratic
increment sum over finite ordered parameter chains; adding or omitting a
single base-point term does not affect the assertion that it is infinite.
Their periodic Poisson operator has multiplier $r^{|n|}$ for $0<r<1$.
Under $r=e^{-s}$ this is exactly $e^{-s|n|}$.  Although this change of
variables reverses order, reversing each finite chain leaves the set of
increments, and hence its quadratic variation, unchanged.  Thus their
counterexample contradicts the preceding periodic estimate.
\end{proof}

\begin{remark}[Sharpness of the output fluctuation exponent]
The failure already occurs for $n=1$, $P=0$, $b\equiv1$, and the odd kernel
on $\mathbb S^0$.  Thus the pointwise critical jump $\mathbf J(F)$---the
paired/chain weak-$\ell^2$ increment functional, up to the absolute
comparison constants in \eqref{eq:chain-pair-functional}---cannot be
replaced in the main theorems by the true strong variation $V^2(F)$.
\end{remark}

\begin{thebibliography}{99}

\bibitem{ACP}
H.~M. Al-Qassem, L.~Cheng, and Y.~Pan,
\emph{Rough oscillatory singular integrals on $\mathbb R^n$},
Studia Math. \textbf{221} (2014), no.~3, 249--267.
\href{https://doi.org/10.4064/sm221-3-4}
{doi:10.4064/sm221-3-4}.

\bibitem{AlSalman}
A.~Al-Salman,
\emph{Rough oscillatory singular integral operators of nonconvolution type},
J. Math. Anal. Appl. \textbf{299} (2004), 72--88.
\href{https://doi.org/10.1016/j.jmaa.2004.06.006}
{doi:10.1016/j.jmaa.2004.06.006}.

\bibitem{AlSalmanGrafakos}
A.~Al-Salman and L.~Grafakos,
\emph{On rough oscillatory singular integral operators},
Forum Math. \textbf{38} (2026), no.~3, 955--975.
\href{https://doi.org/10.1515/forum-2025-0067}
{doi:10.1515/forum-2025-0067}.

\bibitem{BJT}
L.~Becker, A.~Jamneshan, and C.~Thiele,
\emph{Quantitative polynomial Wiener--Wintner theorems},
to appear in Israel J. Math.; arXiv:2601.03999v3 (2026).
\url{https://arxiv.org/abs/2601.03999v3}.

\bibitem{BhojakShrivastava}
A.~Bhojak and S.~Shrivastava,
\emph{Endpoint variation and jump inequalities for rough singular integrals},
arXiv:2602.21888v2 (2026).
\url{https://arxiv.org/abs/2602.21888v2}.

\bibitem{CZ56}
A.~P. Calder\'on and A.~Zygmund,
\emph{On singular integrals},
Amer. J. Math. \textbf{78} (1956), 289--309.
\href{https://doi.org/10.2307/2372517}{doi:10.2307/2372517}.

\bibitem{CJRW}
J.~T. Campbell, R.~L. Jones, K.~Reinhold, and M.~Wierdl,
\emph{Oscillation and variation for the Hilbert transform},
Duke Math. J. \textbf{105} (2000), no.~1, 59--83.
\href{https://doi.org/10.1215/S0012-7094-00-10513-3}
{doi:10.1215/S0012-7094-00-10513-3}.

\bibitem{CJRW03}
J.~T. Campbell, R.~L. Jones, K.~Reinhold, and M.~Wierdl,
\emph{Oscillation and variation for singular integrals in higher dimensions},
Trans. Amer. Math. Soc. \textbf{355} (2003), no.~5, 2115--2137.
\href{https://doi.org/10.1090/S0002-9947-02-03189-6}
{doi:10.1090/S0002-9947-02-03189-6}.

\bibitem{ChenDingHongLiu}
Y.~Chen, Y.~Ding, G.~Hong, and H.~Liu,
\emph{Weighted jump and variational inequalities for rough operators},
J. Funct. Anal. \textbf{274} (2018), no.~8, 2446--2475.
\href{https://doi.org/10.1016/j.jfa.2018.01.009}
{doi:10.1016/j.jfa.2018.01.009}.

\bibitem{ChanilloChrist}
S.~Chanillo and M.~Christ,
\emph{Weak $(1,1)$ bounds for oscillatory singular integrals},
Duke Math. J. \textbf{55} (1987), no.~1, 141--155.
\href{https://doi.org/10.1215/S0012-7094-87-05508-6}
{doi:10.1215/S0012-7094-87-05508-6}.

\bibitem{CSS}
S.~S. Choudhary, S.~Shrivastava, and K.~Shuin,
\emph{Sparse bounds for maximal oscillatory rough singular integral
operators},
Bull. Sci. Math. \textbf{201} (2025), Paper No.~103612, 20~pp.
\href{https://doi.org/10.1016/j.bulsci.2025.103612}
{doi:10.1016/j.bulsci.2025.103612}.

\bibitem{CRdF88}
M.~Christ and J.~L. Rubio de Francia,
\emph{Weak type $(1,1)$ bounds for rough operators. II},
Invent. Math. \textbf{93} (1988), no.~1, 225--237.
\href{https://doi.org/10.1007/BF01393693}{doi:10.1007/BF01393693}.


\bibitem{deLeeuw1965}
K.~de Leeuw,
\emph{On $L^p$ multipliers},
Ann.\ of Math. (2) \textbf{81} (1965), 364--379.

\bibitem{DHL}
Y.~Ding, G.~Hong, and H.~Liu,
\emph{Jump and variational inequalities for rough operators},
J. Fourier Anal. Appl. \textbf{23} (2017), no.~3, 679--711.
\href{https://doi.org/10.1007/s00041-016-9484-8}
{doi:10.1007/s00041-016-9484-8}.

\bibitem{DingLai}
Y.~Ding and X.~Lai,
\emph{Weak type $(1,1)$ bound criterion for singular integrals with rough
kernel and its applications},
Trans. Amer. Math. Soc. \textbf{371} (2019), no.~3, 1649--1675.
\href{https://doi.org/10.1090/tran/7346}
{doi:10.1090/tran/7346}.

\bibitem{DRdF86}
J.~Duoandikoetxea and J.~L. Rubio de Francia,
\emph{Maximal and singular integral operators via Fourier transform estimates},
Invent. Math. \textbf{84} (1986), no.~3, 541--561.
\href{https://doi.org/10.1007/BF01388746}{doi:10.1007/BF01388746}.

\bibitem{FS72}
C.~Fefferman and E.~M.~Stein,
\emph{$H^p$ spaces of several variables},
Acta Math. \textbf{129} (1972), 137--193.
\href{https://doi.org/10.1007/BF02392215}{doi:10.1007/BF02392215}.

\bibitem{JiangLu}
Y.-S.~Jiang and S.-Z.~Lu,
\emph{Oscillatory singular integrals with rough kernel},
in \emph{Harmonic Analysis in China}
(M.~Cheng, D.-G.~Deng, S.~Gong, and C.-C.~Yang, eds.),
Mathematics and Its Applications, vol.~327,
Kluwer Academic Publishers, Dordrecht, 1995, 135--145.
\href{https://doi.org/10.1007/978-94-011-0141-7_7}
{doi:10.1007/978-94-011-0141-7\_7}.

\bibitem{JSW}
R.~L. Jones, A.~Seeger, and J.~Wright,
\emph{Strong variational and jump inequalities in harmonic analysis},
Trans. Amer. Math. Soc. \textbf{360} (2008), no.~12, 6711--6742.
\href{https://doi.org/10.1090/S0002-9947-08-04538-8}
{doi:10.1090/S0002-9947-08-04538-8}.


\bibitem{JonesWang2004}
R.~L. Jones and G.~Wang,
\emph{Variation inequalities for the Fej\'er and Poisson kernels},
Trans. Amer. Math. Soc. \textbf{356} (2004), no.~11, 4493--4518.
\href{https://doi.org/10.1090/S0002-9947-04-03397-5}
{doi:10.1090/S0002-9947-04-03397-5}.

\bibitem{KrauseLacey}
B.~Krause and M.~T. Lacey,
\emph{Sparse bounds for maximally truncated oscillatory singular integrals},
Ann. Sc. Norm. Super. Pisa Cl. Sci. (5) \textbf{20} (2020), no.~2,
415--435.
\href{https://doi.org/10.2422/2036-2145.201706_023}
{doi:10.2422/2036-2145.201706\_023}.

\bibitem{Lai}
X.~Lai,
\emph{Weak $(1,1)$ estimate for maximal truncated rough singular integral
operator},
arXiv:2508.17737v3 (2025).
\url{https://arxiv.org/abs/2508.17737v3}.

\bibitem{LuWu03}
S.~Lu and H.~Wu,
\emph{Oscillatory singular integrals and commutators with rough kernels},
Ann. Sci. Math. Qu\'ebec \textbf{27} (2003), no.~1, 47--66.

\bibitem{LuWu04}
S.~Lu and H.~Wu,
\emph{A class of multilinear oscillatory singular integrals related to
block spaces},
Tohoku Math. J. (2) \textbf{56} (2004), 299--315.

\bibitem{LuZhang}
S.~Lu and Y.~Zhang,
\emph{Criterion on $L^p$-boundedness for a class of oscillatory singular
integrals with rough kernels},
Rev. Mat. Iberoam. \textbf{8} (1992), no.~2, 201--219.
\href{https://doi.org/10.4171/RMI/122}{doi:10.4171/RMI/122}.

\bibitem{MSZK}
M.~Mirek, E.~M.~Stein, and P.~Zorin-Kranich,
\emph{Jump inequalities via real interpolation},
Math. Ann. \textbf{376} (2020), 797--819.
\href{https://doi.org/10.1007/s00208-019-01889-2}
{doi:10.1007/s00208-019-01889-2}.

\bibitem{Ojanen}
H.~Ojanen,
\emph{Weighted estimates for rough oscillatory singular integrals},
J. Fourier Anal. Appl. \textbf{6} (2000), no.~4, 427--436.
\href{https://doi.org/10.1007/BF02510147}
{doi:10.1007/BF02510147}.

\bibitem{RicciStein}
F.~Ricci and E.~M.~Stein,
\emph{Harmonic analysis on nilpotent groups and singular integrals. I.
Oscillatory integrals},
J. Funct. Anal. \textbf{73} (1987), no.~1, 179--194.
\href{https://doi.org/10.1016/0022-1236(87)90064-4}
{doi:10.1016/0022-1236(87)90064-4}.

\bibitem{Seeger96}
A.~Seeger,
\emph{Singular integral operators with rough convolution kernels},
J. Amer. Math. Soc. \textbf{9} (1996), 95--105.
\href{https://doi.org/10.1090/S0894-0347-96-00185-3}
{doi:10.1090/S0894-0347-96-00185-3}.

\bibitem{SteinSI}
E.~M.~Stein,
\emph{Singular Integrals and Differentiability Properties of Functions},
Princeton Mathematical Series, vol.~30,
Princeton University Press, Princeton, NJ, 1970.


\bibitem{Young1936}
L.~C. Young,
\emph{An inequality of the H\"older type, connected with Stieltjes
integration},
Acta Math. \textbf{67} (1936), 251--282.
\href{https://doi.org/10.1007/BF02401743}
{doi:10.1007/BF02401743}.


\end{thebibliography}

\end{document}
